% 本文件由 LaTeX 排版 Agent 根据有效 translation.json 自动生成。
% author=John Chrysostom; work_id=2202; title=格林多后书讲道集

\chapter{格林多后书讲道集}

% structure: p0001 source=h1 target=omitted reason=page-title-already-rendered-by-parent

\Emph{讲道1} \BibleRef{格后 1:1-5}\newline{}\Emph{讲道2} \BibleRef{格后 1:6-11}\newline{}\Emph{讲道3} \BibleRef{格后 1:12-22}\newline{}\Emph{讲道4} \BibleRef{格后 1:23-2:11}\newline{}\Emph{讲道5} \BibleRef{格后 2:12-17}\newline{}\Emph{讲道6} \BibleRef{格后 3:1-6}\newline{}\Emph{讲道7} \BibleRef{格后 3:7-18}\newline{}\Emph{讲道8} \BibleRef{格后 4:1-7}\newline{}\Emph{讲道9} \BibleRef{格后 4:8-18}\newline{}\Emph{讲道10} \BibleRef{格后 5:1-10}\newline{}\Emph{讲道11} \BibleRef{格后 5:11-21}\newline{}\Emph{讲道12} \BibleRef{格后 6:1-10}\newline{}\Emph{讲道13} \BibleRef{格后 6:11-7:1}\newline{}\Emph{讲道14} \BibleRef{格后 7:2-7}\newline{}\Emph{讲道15} \BibleRef{格后 7:8-13}\newline{}\Emph{讲道16} \BibleRef{格后 7:13-8:6}\newline{}\Emph{讲道17} \BibleRef{格后 8:7-15}\newline{}\Emph{讲道18} \BibleRef{格后 8:16-24}\newline{}\Emph{讲道19} \BibleRef{格后 9:1-9}\newline{}\Emph{讲道20} \BibleRef{格后 9:10-15}\newline{}\Emph{讲道21} \BibleRef{格后 10:1-6}\newline{}\Emph{讲道22} \BibleRef{格后 10:7-18}\newline{}\Emph{讲道23} \BibleRef{格后 11:1-12}\newline{}\Emph{讲道24} \BibleRef{格后 11:13-21}\newline{}\Emph{讲道25} \BibleRef{格后 11:21-33}\newline{}\Emph{讲道26} \BibleRef{格后 12:1-10}\newline{}\Emph{讲道27} \BibleRef{格后 12:11-15}\newline{}\Emph{讲道28} \BibleRef{格后 12:16-21}\newline{}\Emph{讲道29} \BibleRef{格后 13:1-9}\newline{}\Emph{讲道30} \BibleRef{格后 13:10-14}\par

\SourceRecord{Translated by Talbot W. Chambers.；Nicene and Post-Nicene Fathers, First Series；Vol. 12.；Edited by Philip Schaff.；Buffalo, NY: Christian Literature Publishing Co.,；1889.；<http://www.newadvent.org/fathers/2202.htm>.}

% structure: p0001 source=h1 target=section reason=page-title

\section{\WorkTitle{格林多后书讲道第一讲}}

% structure: p0002 source=h2 target=subsection reason=relative-heading-rank; base=h2

\subsection{\BibleRef{格后 1:1-4}}

\BibleQuote{因天主的旨意，做基督耶稣宗徒的\Person{保禄}和\Person{弟茂德}兄弟，写信给在\Place{格林多}的天主教会和全\Place{阿哈雅}的众位圣徒：愿恩宠与平安由我们的父天主和主耶稣基督赐给你们。愿我们的主耶稣基督的天主和父，仁慈的父和施与各种安慰的天主受赞扬，他在我们的各种磨难中，常安慰我们，为使我们能以自己由天主所亲受的安慰，去安慰那些在各种磨难中的人。}\par

首先应当探究，为何他在前书之后又加写后书？又为何要以天主的怜悯和安慰作为开头？\par

为何他又加写后书？在前书中他曾说过：\BibleQuote{我却要到你们那里去，并且要知道的不是那些自高自大之人的言语，而是他们的能力；}\BibleRef{格前 4:19}又近末尾处以较温和的措辞作了同样的许诺，即：\BibleQuote{我要从\Place{马其顿}经过；因为我必须经过\Place{马其顿}；并且也许我将在你们那里住下，甚至过冬；}\BibleRef{格前 16:5}然而过了许久，他仍未前来；反而仍然拖延，即使约定的时间已过；是圣神将他留在其他更为必要的事上。为此，他需要写第二封书信；若他只是稍微延迟，原本是不需要的。\par

% structure: p0006 source=h2 target=subsection reason=relative-heading-rank; base=h2

\subsection{\BibleRef{格后 1:5}}

但不仅为此原因，还因为他们因前书而改正了；因为那个犯奸淫的人，从前他们赞赏并自高自大的，他们已将其除去并完全隔绝了。这一点他在以下的话中表明：\BibleQuote{但若有人使人忧苦，他所忧苦的不是我，而是你们众人，至少一部分（我免得说得太重）。对这样的人，由许多人施加的处罚已经够了。}\BibleRef{格后 2:5}随着论述的进行，他再次暗示同一件事，说：\BibleQuote{因为看，你们依着天主的标准感到忧苦，在你们身上产生了何等热诚的关心，何等清楚的表白，何等义怒，何等恐惧，何等仰慕，何等热忱，何等报应！在一切事上，你们证明了自己在这事上是纯洁的。}\BibleRef{格后 7:11}此外，他所吩咐的捐款，他们也十分热心地收集了。因此他也说：\BibleQuote{因为我知道你们慷慨的心志，我曾以此向\Place{马其顿}人夸耀过你们，说\Place{阿哈雅}从去年就预备好了。}\BibleRef{格后 9:2}还有他所派遣的\Person{弟铎}，他们也以极友善的态度接待了他，正如他再次所说：\BibleQuote{他的深情越发向你们流露，因为他回忆你们众人的服从，怎样以恐惧战兢接待了他。}\BibleRef{格后 7:15}为了这些原因，他写了后书。因为正当如此：正如他们在有过错时他责备了他们，如今他们改正了，他理应认可并称赞他们。因此，这封书信并非全篇严厉，只有末尾几处。因为在他们中间，有些犹太人自视甚高，指责\Person{保禄}是个夸口的人，不值得尊重；因此有他们这样的话：\BibleQuote{他的书信严厉而有力，但他本人却软弱无力，言语又粗俗不堪。}\BibleRef{格后 10:10}意思就是，他在场时显得无足轻重（因为这就是\BibleQuote{他的身体软弱}的意思），但当他不在时，他就在所写的书信中大大夸口（因为这就是\BibleQuote{他的书信严厉而有力}的意思）。此外，为了抬高自己的声望，这些人假装不收任何东西，他也暗示了这一点，说：\BibleQuote{使他们在所夸耀的事上，他们也可以被找到如同我们一样。}\BibleRef{格后 11:12}而且，他们还拥有语言的能力，因而立即大为得意。因此他也称自己\BibleQuote{在言语上粗俗}\BibleRef{格后 11:6}，表明他并不以此为耻；也不认为相反的是多大的成就。既然看来这些人可能会迷惑一些人，在称赞了他们正确行为之后，并驳斥他们在犹太主义事物上的愚昧骄傲，因为他们不合时宜地争执要遵守那些事物，他也对此方面给予温和的责备。\par

2. 这样说，总的来说，依我看来，这封书信的宗旨便是如此。现在要思考引言，并说明为何他在惯常的问候之后，就以天主的怜悯开始。但首先，有必要从开头说起，探究为何他在这里将\Person{弟茂德}与自己联合。因为他说：\BibleQuote{\Person{保禄}，因天主的旨意做\Person{基督耶稣}宗徒的，和\Person{弟茂德}兄弟。}在前书中，他许诺要打发他去；并吩咐他们说：\BibleQuote{若是\Person{弟茂德}到了你们那里，你们要使他安心在你们那里，不要使他惧怕。}\BibleRef{格前 16:10}那么，他怎能在开头就将他与自己联合呢？在他到了他们那里之后，遵照他老师的那个应许：\BibleQuote{我打发\Person{弟茂德}到你们那里去，他将在主内提醒你们，我在基督内的道路。}\BibleRef{格前 4:17}并处理好一切之后，他便回到了\Person{保禄}那里；\Person{保禄}在打发他去时曾说过：\BibleQuote{你们要送他平安前行，使他回到我这里，因为我正等待他同弟兄们一起来。}\BibleRef{格前 16:11}\par

既然\Person{弟茂德}已回到他的老师那里，并与他一起安排好了\Place{亚细亚}的事\BibleRef{格前 16:8}，又再次渡海到了\Place{马其顿}；\Person{保禄}很自然地从此将他视为与自己同住的同伴。因为那时他是从\Place{亚细亚}写信，但现在是从\Place{马其顿}。如此与他联合，立时使\Person{弟茂德}更受尊重，也显明\Person{保禄}自己极大的谦逊：因为\Person{弟茂德}远不及他，然而爱德使一切融合为一。因此，他在各处都将\Person{弟茂德}与己相等：有时说\BibleQuote{如同儿子服事父亲一样，他同我一起服事了}\BibleRef{斐 2:22}；有时说\BibleQuote{因为他做主的工，正如我也做一样}\BibleRef{格前 16:10}；而在此处，他甚至称他为\Quote{兄弟}；借此使他在\Place{格林多}人面前成为可敬的对象，这些\Place{格林多}人，如我所说，曾与他同在，并证明了他的价值。\par

\Quote{致天主在\Place{格林多}的教会。}他又称他们为\Quote{教会}，为要将他们全部合而为一。因为若其中的成员分裂、彼此分离，就不能成为一体。\Quote{并与在全\Place{阿哈雅}的所有圣人。}通过给\Place{格林多}人所写的书信向全体致意，他既尊重这些人，又将整个民族聚集起来。但他称他们为\Quote{圣人}，暗示若有人不洁，便无份于此问候。但为何他写给母城，却通过她向所有人致辞？因为他并非处处如此。例如，在致\Place{得撒洛尼}人的书信中，他没有同时向\Place{马其顿}人致辞；同样，在致\Place{厄弗所}人的书信中，他没有包括整个\Place{亚细亚}；致\Place{罗马}人的信也没有写给居住在\Place{意大利}的所有人。但在这封书信中，他如此做了；在致\Place{迦拉达}人的书信中也是如此。因为在那里，他也不是写给一个或两个或三个城市，而是写给所有散布各处的人，说：\BibleQuote{保禄宗徒（不是由于人，也不是借着人，而是借着耶稣基督和使祂从死者中复活的天主父），以及同我一起的众弟兄，致\Place{迦拉达}众教会。愿恩宠与平安归于你们。}\BibleRef{迦 1:1}致\Place{希伯来}人也写了一封书信给全体，没有按各城分别。这其中的原因是什么？我以为，因为在此情况下，所有的人都染上了同一种混乱，因此他共同致函，因为他们需要同一疗法。所有\Place{迦拉达}人都受感染。\Place{希伯来}人也是如此，我认为这些\Place{阿哈雅}人同样。\par

3. 因此，当他将整个民族聚集一处，并以惯常的问候向他们致意，他说：\BibleQuote{愿恩宠与平安由天主我们的父和主耶稣基督赐予你们。}\BibleRef{格后 1:2}请看他如何恰当地开始：\BibleQuote{愿我们的主耶稣基督的天主和父，仁慈的父和施与各种安慰的天主，受赞颂。}\BibleRef{格后 1:3}你问，这与主题有何关系？我回答：关系极大；因为请看，他们非常烦恼不安，因为宗徒没有到他们那里去，尽管他许诺过，却把全部时间花在\Place{马其顿}，似乎偏爱他人而轻视他们。于是，为应对他们对他的这种不满，他说明自己缺席的原因；然而，他并不直接陈述，如说\Quote{我确实答应要来，但因受苦难阻扰，请原谅我，不要认为我有什么轻视或疏忽你们之处}；而是以另一种方式，使主题更显尊严和可信，并以安慰的性质赋予其重要性，以致此后他们甚至不会问起他延迟的原因。就像一个人答应去见他渴望的人，经历无数危险后终于到达，应说：\Quote{光荣归于你，天主，因为你让我见到了如此渴望的景象——他亲爱的面容！愿受赞颂，天主，你从何等危险中拯救了我！}因为这样的赞颂之词就是对他那准备责备的人的回答，不再容许他抱怨延迟；因为一个正为从如此大难中得救而感谢天主的人，不能羞惭地将他拖到审判台前，让他为自己拖延道歉。因此\Person{保禄}这样开始：\Quote{愿仁慈的天主受赞颂，}用这些话暗示他曾陷入并脱离了大危险。因为正如\Person{达味}并非处处以同一方式或同一称号称呼天主；但当论及战争和胜利时，他说：\BibleQuote{上主，我的力量，我爱慕你；上主是我的盾牌；}而当论及从苦难和笼罩他的黑暗中获得解救时，他说：\BibleQuote{上主是我的光明，我的救援；}\BibleRef{咏 26:1}并且随着眼前情境，他时而从慈爱、时而从公义、时而从公义审判来称呼祂——同样，\Person{保禄}在此处开头也从祂的慈爱来描述祂，称祂为\Quote{仁慈的天主}，即是\Quote{曾向我显示如此大仁慈，以至于将我从死门之中提拔上来的那一位。}\par

因此，仁慈是天主独有的卓越属性，也是祂本性中最内在的；故此祂被称为\Quote{仁慈的天主}。\par

请你们也留意保禄在此处的谦卑。因为他虽然因他所传的福音而身处危险，但他并没有说，他是因自己的功劳而得救，而是因天主的仁慈。但他后来更清楚地宣告了这一点，现在他接着说，\BibleQuote{祂在一切磨难中安慰我们。}\BibleRef{格后 1:4}他没有说：\Quote{祂不让我们进入磨难，}而是\Quote{祂在磨难中安慰。}因为这一方面彰显了天主的德能，另一方面增加了受磨难者的忍耐。因为，他说，\BibleQuote{患难生忍耐。}\BibleRef{罗 5:3}同样，先知也说：\BibleQuote{我在困苦中，你使我宽广。}\BibleRef{咏 4:1}他没有说：\Quote{你没有让我陷入患难，}也没有说：\Quote{你迅速挪去了我的患难，}而是在患难持续之时，说：\BibleQuote{你使我宽广，}\BibleRef{达 3:21}即\Quote{你赐给了我极大的自由和舒畅。}这在三个青年身上也确实发生了，因为天主既没有阻止他们被投入火焰，当投入时也没有熄灭火焰，而是在火窑燃烧时赐给了他们自由。天主的行事方式永远如此，正如保禄所说的：\BibleQuote{祂在一切磨难中安慰我们。}\par

然而，他在这些话中还教导了更多：你问是什么？即天主这样做不是一次、两次，而是不间断地。因为祂不是一时安慰，一时不安慰，而是永远持续地。因此他说：\Quote{安慰的，}而不是\Quote{安慰过的，}并且\Quote{在一切磨难中，}而不是\Quote{在这个或那个磨难中，}而是\Quote{在一切中。}\par

\Quote{使我们能借着天主所安慰我们的安慰，去安慰那些在一切患难中的人。}你们没看见他如何预先为自己的辩护铺垫，使听者想到某种大患难吗？在此也再次显出了他的谦逊：他说这恩慈并不是因他们自己的功德而赐下，而是为了那些需要他们帮助的人；因为，他说，\Quote{祂安慰我们，正是为了这个目的：使我们可以彼此安慰。}他也借此彰显了宗徒们的卓越，表明他们在得到安慰和喘息之后，并不像我们那样软绵绵地躺下，而是继续前往，去傅油、去激励、去唤醒他人。然而，有些人认为这是宗徒的意思：\Quote{我们的安慰也是他人的安慰。}但我的看法是，在这开篇中，他也暗中责备了假宗徒，那些徒然自夸、安居家中、生活奢侈的人；但他只是隐晦地、顺便地提及，主要目的仍是为自己的延迟辩解。（他可能会说）\Quote{如果我们蒙安慰是为了也能安慰他人，那么就不要责怪我们没有来；因为我们所有的时间都花在应对围攻我们的阴谋、暴力和恐怖上了。}\par

4. \Quote{因为我们多受基督的苦难，也借着基督多得安慰。}他并非要以夸张的受苦描述来压垮门徒，反而宣告安慰也是伟大而超绝的，不仅以此，还通过提醒他们想起基督，并称这些苦难为\Quote{祂的，}并且在安慰之前，从苦难本身获得安慰。因为我有什么喜乐能比与基督同受苦难、并为祂的缘故忍受这些事更大呢？什么安慰能与这相比？但不仅从这一源头他振奋了受苦者的心情，还有另一源头。你问另一源头是什么？在于他说：\Quote{充满，}因为他并不说：\Quote{基督的苦难}是\Quote{在我们内，}而是它们\Quote{充满，}从而宣告他们所受的不只是祂的苦难，甚至超过这些。因为他说，我们不是受了祂所遭受的一切；\Quote{甚至更多，}因为，想一想，\Quote{基督被驱逐、迫害、鞭打、死亡，}但他说，我们\Quote{超过这一切，}这本身就足够了安慰。现在不要让任何人指责这大胆的言辞；因为他在别处说：\BibleQuote{如今我在苦难中喜乐，并在我身上补满基督苦难所欠缺的。}\BibleRef{哥 1:24}但那既不是出于胆大，也不是出于任何狂妄。因为正如他们行了比祂更大的奇迹，按祂说的：\BibleQuote{信我的人要做比这更大的事，}\BibleRef{若 14:12}但一切都是那在他们内行事的祂所做；同样，他们也受了比祂更多的苦，但一切又都是那位安慰他们、并使他们能承受临到身上的邪恶的祂所做。\par

考虑到这一点，保禄意识到他说了多么伟大的一件事，就再次显著地加以限制，加上\Quote{我们的安慰也借着基督充满；}从而立刻将一切归于祂，并在此也宣告祂的仁慈；因为他不是说：\Quote{我们的安慰等于我们的苦难，}而是\Quote{远超过；}因此，奋斗的季节也是新桂冠的季节。因为，你说，有什么能与为基督的缘故受鞭打、与天主交谈、超乎万有、胜过驱逐我们的人、不被整个世界所征服、并期待如此美好的事物\BibleQuote{眼睛未曾看见，耳朵未曾听见，人心也未曾想到的}\BibleRef{格前 2:9}相比呢？有什么能与为虔敬的缘故忍受苦难、从天主领受无限的安慰、从如此大罪中被救赎、被视配得圣神、被圣化和成义、不畏惧任何人、并且在危险中出类拔萃相比呢？\par

5. 那么，我们不要在受诱惑时消沉。因为没有一个放纵私欲的人与基督有份，也没有沉睡的人、怠惰的人，或任何这些松散放荡的人。但凡在困苦和试探中的人，这人亲近祂，凡行走在窄路上的人。因为祂亲自行了这条路；因此祂也说：\Quote{人子没有枕首之处。}所以，当你处于困苦时，不要忧伤；要思想你与谁有份，以及你如何被试炼所净化，你的收益何等巨大。因为没有比得罪天主更可悲的事了；但除此以外，无论是困苦、阴谋，还是任何其他事，都不能使一个心志正直的灵魂忧伤：犹如一点火星，你若把它投入深水中，立刻就扑灭了；同样，一种完全过度的忧愁，若落在良心上，也容易消散而消失。\par

这样的乃是保禄持续喜乐的泉源：因为凡属于天主的，他都满怀希望；他甚至不计较如此大的苦难，虽作为人也忧伤，却不沉沦。同样，那位圣祖在诸多痛苦患难中也充满喜乐；试想，他离弃了故乡，经历漫长艰辛的旅途；到了异乡，\BibleQuote{连立足之地也没有}\BibleRef{宗 7:5}。接着又有饥荒等待他，使他再次漂泊；饥荒之后，妻子被夺，然后是死亡的恐惧、无子、战争、危险、阴谋，最后是那最终的考验——杀害他的独生真爱之子，那是惨重而无法挽回的祭献。请你们不要以为，因他欣然服从，就没有感受到他所经历的一切。因为尽管他的义德，如同事实上那样，是无法估量的，但他仍是一个人，感受着本性所要求的一切。然而，这些事没有一件能击倒他；他像一位高贵的运动员一样站立，为每一件事被宣告胜利并加冕。同样，蒙福的保禄，虽然每天看到试探如雪片般袭来，却喜乐欢腾，如同置身于乐园的欢乐之中。正如被这种喜乐所充满的人不会成为绝望的猎物；同样，不以此为己有的人容易被一切所压倒；他就像一个穿着不坚固铠甲的人，经不起轻轻一击便受伤：但那位四面严加防护、能抵挡一切来袭之箭的人则不然。确实，比任何铠甲更坚固的，是天主内的喜乐；谁拥有了它，就没有任何事能使他垂头丧气或面露愁容，反而能高贵地承受一切。因为有什么比火更难忍受？有什么比持续的折磨更痛苦？确实，它在痛苦上超过失去无数财富、子女或任何东西；因为经上说：\BibleQuote{以皮换皮，人为了自己的性命，情愿舍弃所有的一切。}\BibleRef{约伯 2:4} 所以没有什么比身体的痛苦更难忍受；然而，由于这天主内的喜乐，连听都难以忍受的事，也变得可忍受且令人渴望：你若将刚断气的殉道者从十字架或烤架上取下，你会发现他心中有一种无法言喻的喜乐珍宝。\par

6. 也许有人说：我该怎么办？因为现在不是殉道的时候？你怎么说？现在不是殉道的时候？从来没有不是的时候；它常在我们眼前，只要我们肯睁眼。因为并非只有悬挂在十字架上才算殉道者，若是如此，约伯就被排除在这冠冕之外了；因为他既没有站在法庭上，也没有听见法官的声音，也没有看见刽子手；不，也没有高悬在木架上被撕裂肋旁；然而他遭受的比许多殉道者更痛苦；比任何打击更尖锐的是那些接踵而来的使者对他的报告，从四面八方刺痛他；比千名刽子手更折磨人的是啃噬他全身的虫子的咬啮。\par

然而，他能与哪位殉道者相提并论而不失其当呢？诚然可与千万人相比。因为他在每一种（磨难）中都奋战并获得荣冠：在财产、子女、身体、妻子、朋友、仇敌和仆人身上（连这些人都曾唾他的脸），在饥饿、异象、痛苦和恶臭中。正因如此，我说他可当之无愧地与一位、两位、三位，而是与千万殉道者并论。因为除了我所提及的，时间也大大增添了他的荣冠：那是在律法之前，在恩宠之前，他如此受苦，而且长达数月，且每种苦难都以最恶劣的形式临到；这一切灾祸同时攻击他。然而，每个单独的灾祸本身已是无法忍受的，甚至连那看似最可忍受的——财产的丧失——也是如此。因为许多人能忍耐鞭打，却不能忍受财产的损失；他们宁愿不放弃任何部分，甚至甘愿为此受鞭笞并遭受无数苦难；这一打击——财产的损失——在他们眼中比一切都沉重。因此，对于崇高承受损失的人，这是另一种殉道方式。有人会问：我们如何能崇高承受？当你懂得一句感恩的话便能赢得比你所失去的一切更多时。因为若在我们得知损失时不为所动，反而说：\Quote{愿天主受赞颂，}我们就找到了更丰盈的财富。因为真的，即使你将全部财富施舍给穷人，四处寻找贫困之人，将你的财物分给饥饿者，你所获得的果实也不及因这句话所得的。因此，我并非如此钦佩约伯慷慨敞庭接待穷人，而是更受震撼和赞扬他以感恩之心承受财产的掠夺。在失去子女的事上也是如此。因为在这里，你也会得到不少于那献上自己儿子作为祭献的人所得的赏报，只要你看见自己的子女死去时感谢慈爱的天主。因为这样的人怎能比亚巴郎逊色呢？他没有看到自己的儿子伸展成为尸体，只是有意这样做。所以，若他因决心杀害并伸手拿刀\BibleRef{创 22:10}而在比较中得胜，但在这里孩子却已躺卧死去。此外，他还因完成善事的远景以及想到这卓越成就是自己刚毅的成果，以及他所听见的声音来自天上而感到安慰，因此更容易承受。但这里却没有这样的事。所以，必须有钢铁般的灵魂，才能平静地看见一个孩子——独生子——在富裕中长大，正值美好前景，挺直地躺在灵床上成为一具尸体。而如果这样的人，平息了本性的激动，得坚强地不流一滴泪说出约伯的话：\Quote{上主赐予的，上主收回；}\BibleRef{约 1:21}仅凭这句话，他就将与亚巴郎本人并列，并与约伯一同被宣告为胜利者。而且，如果他阻止妇女的哀哭，驱散哀悼的人群，并唤醒他们一同颂扬光荣（归于天主），他将在天上地下获得数不清的赏报：世人赞叹，天使鼓掌，天主为他加冕。\par

7. 你或许会说：人怎能不悲伤呢？我的回答是：若你思量圣祖和约伯，他们都是人，却未表现出任何此类情绪；而且这都是在律法、恩宠以及我们所有法律卓越智慧之前发生的：若你想到亡者已迁往更好的国度，飞逝至更幸福的继承之地，你并未失去儿子，而是今后将他安置于一个不可侵犯之地。那么，我恳求你，不要再说：我不再被称为\Quote{父亲}了，因为当你的儿子尚在时，你为何不再被如此称呼呢？难道你真的与孩子分离或失去了儿子吗？相反，你得到了他，而且拥有得更安全。因此，你不仅在此世被称为\Quote{父亲}，在天堂也将如此；所以你并未失去\Quote{父亲}这一称谓，反而获得了更高尚的意义；因为今后你不再被称为必死之子的父亲，而是不死之子的父亲；一位高贵士兵的父亲；常在天宫内服役。因为不要以为他不在眼前就失去了；因为即使他去了异国他乡，你亲缘的称谓也不会随他的躯体消失。所以，不要凝视躺在那里的面容，因为那只会重新点燃你的悲伤；要将你的思想从躺在那里的他提升到天堂。躺在那里并不是你的孩子，而是已经飞走并腾升到无限高空中的人。当你看到眼睛闭合，嘴唇紧闭，身体静止，哦，不要有这样的想法：\Quote{这些嘴唇不再说话，这些眼睛不再看见，这些脚不再行走，而是都走向腐朽！}哦，不要这样说：应反过来说：\Quote{这些嘴唇将说出更好的话，眼睛将看见更伟大的事物，脚将腾云驾雾；而这个现在腐烂的身体将穿戴上不朽，我将接回更荣耀的儿子。}但若你所见的令你痛苦，可以这样对自己说：这只是衣服，他脱下是为了更珍贵地收回；这是一所房子，被拆毁是为了更辉煌地重建。因为正如我们打算拆毁房屋时，不允许住户留在里面，以免他们沾到灰尘和噪音；而是让他们暂时搬出，等我们牢固地建好房舍后，再自由地让他们进入；同样，天主也是如此：祂拆毁这正在腐朽的帐棚，暂时将他接到祂的父家和祂自己那里，以便当它被拆毁并重建后，再更荣耀地归还给他。\par

那么，不要说：\Quote{他灭亡了，不再存在。}因为这是不信者的话；而要说：\Quote{他睡了，必要复活。}\Quote{他出外旅行，要与君王同回。}谁说这话？是有基督在他内说话的那位。因为他说：\Quote{我们若信耶稣死了，也复活了，}\Quote{同样，那些在耶稣内睡了的人，天主也必把他们与耶稣一同带来。}\BibleRef{得前 4:14}你若寻找你的儿子，就该到君王那里去寻找，到天使的军队那里去；不要到坟墓里，不要到地上；免得他高高受到荣耀，你自己却仍匍匐在地。\par

如果我们有这样真正的智慧，就能轻易抵挡这一切痛苦；愿\Quote{慈悲的天主和一切安慰的圣父}安慰我们众人的心，既安慰那些受此忧愁压迫的人，也安慰那些被其他任何忧愁所困的人；愿祂赐我们脱离一切失望，增加属神的喜乐，并获得将来的美善。愿我们众人赖我们的主耶稣基督的恩宠和慈爱，达到这一切；愿光荣、权能、尊崇归于祂和圣父及圣神，自今至永远，世世无穷。阿们。\par

\SourceRecord{Translated by Talbot W. Chambers.；Nicene and Post-Nicene Fathers, First Series；Vol. 12.；Edited by Philip Schaff.；Buffalo, NY: Christian Literature Publishing Co.,；1889.；<http://www.newadvent.org/fathers/220201.htm>.}

% structure: p0001 source=h1 target=section reason=page-title

\section{讲道第二篇：格林多后书}

% structure: p0002 source=h2 target=subsection reason=relative-heading-rank; base=h2

\subsection{格林多后书 1:6-7}

\BibleQuote{我们或受磨难，是为你们的安慰和救恩，这救恩在忍耐同样我们所受的苦难中运行；并且我们为你们所怀的希望是坚定的。}\par

说完这一点——也就是安慰和劝慰的首要根基，即（藉受苦）与基督有分——他如今提出第二点，即门徒自身的救恩由此而得。\Quote{所以，不要灰心，也不要因我们受磨难而惊慌害怕；因为这反倒成为你们欢欣的理由：因为我们若不遭受磨难，那将是你们全体的丧亡。}如何并为何？因为若是我们因缺乏勇气和惧怕危险，而没有向你们宣讲那使你们学到真知的话语，你们的处境将无可救药。你们是否再次看见\Person{保禄}的激烈和恳切争战？正是那些困扰他们的事，他用来安慰他们。因为，他说，我们的迫害越猛烈，你们的好望就越应增长；因为你们的救恩和安慰也相应地更加丰盛。因为还有什么比藉宣讲获得了如此美物，具有同等安慰的力量呢？然后，为了不显得他把赞美独揽于己身，请看他是如何使他们也分享这些赞美。因为对于\BibleQuote{我们或受磨难，是为你们的安慰和救恩}这句话，他加上\BibleQuote{这救恩在忍耐同样我们所受的苦难中运行} \BibleRef{格后 1:7}。随后他确实更清楚地说明这一点，说道：\BibleQuote{正如你们分担了受苦，也同样分担了安慰}；但与此同时，他也在这句话\BibleQuote{同样苦难}中暗示了这一点，从而使他的话包括他们。因为他所说的是：\Quote{你们的救恩不仅是我们单独的工作，也是你们自己的；因为我们宣讲话语给你们受磨难，你们领受它也同样受磨难；我们将所领受的传授给你们，你们领受所传授的并不放弃。}如今有什么谦卑能与这相比呢？他竟把远远不及他的人提升到相同的忍耐尊严。因为他说：\BibleQuote{在同样苦难的忍耐中运行}；因为你们的救恩不仅通过相信而来，也通过与我们一同受苦和忍耐同样的事而来。因为正如一个拳击手，当他只是露面、训练有素且技能在身时，固然令人钦佩；但当他上场行动，忍受打击并击打对手时，他最光芒四射，因为那时他的良好训练最付诸行动，其技术的证明也显而易见；同样，你们的救恩也正是在拥有忍耐，在受苦并高尚地承受一切时，才格外付诸行动，即被显示、增长、提高。所以，救恩的工作不在于行恶，而在于受苦。况且，他说的不是\BibleQuote{那运行}，而是\BibleQuote{那成就}，为要表明，与他们自己的心志愿意一同，那在他们内运行的恩宠也贡献了许多。\par

% structure: p0005 source=h3 target=subsubsection reason=relative-heading-rank; base=h2

\subsubsection{格林多后书 1:7}

\BibleQuote{并且我们为你们所怀的希望是坚定的。}这就是说，即使你们将遭受无数灾祸，我们深信你们不会转身，无论是因为你们自己的考验，还是因为我们的迫害。因为我们远未怀疑你们因我们的受苦而惊慌，反而当你们自己处于危险时，我们那时对你们满怀信心。\par

2. 你们是否看见自从前一封书信以来，他们的进步有多大？因为这里他为他们作的见证，远比在整封书信中赞扬和称颂的马其顿人更大。因为他为马其顿人担忧，说：\BibleQuote{我们派遣弟茂德到你们那里，为坚固你们，并为你们的信德安慰你们，使没有人被这些苦难动摇，因为你们自己知道，我们是被派定为此的。} \BibleRef{得前 3:2}又说：\BibleQuote{为此，当我不能再忍耐时，就派人去探知你们的信德，恐怕那诱惑者诱惑了你们，以致我们的劳苦归于徒然。} \BibleRef{格后 1:5}但对这些格林多人，他完全没有说类似的话，而是恰恰相反：\Quote{我们为你们所怀的希望是坚定的。}\par

% structure: p0008 source=h3 target=subsubsection reason=relative-heading-rank; base=h2

\subsubsection{格林多后书 1:6-7}

\BibleQuote{我们或受安慰，是为你们的安慰和救恩。因为知道你们怎样分担了受苦，也怎样分担了安慰。}\par

他藉着说\Quote{我们若受磨难，是为你们的安慰与救恩}，表明了宗徒们是为他们而受磨难；他也愿意表明他们也是为了他们的缘故而得安慰。他其实在之前一小段也说了这话，虽然有些笼统，他说：\Quote{愿天主受讚颂，祂在我们的一切磨难中安慰我们，为使我们能去安慰那些在任何磨难中的人}。他在这里也用其他话语重述了它，更清楚、更切合他们的需要。他说：\Quote{因为我们或受安慰，是为你们的安慰}。他的意思是：我们的安慰成为你们的振作，即使我们不是用言语安慰你们。我们若稍得振作，这就有助于鼓励你们；我们若自己受安慰，这就成为你们的安慰。因为正如你们将我们的苦难视为自己的，你们也必使我们的安慰成为自己的。因为当你们与我们同遭厄运时，岂有不分享好运之理？因此，若你们在一切事上都有分，既在磨难中，也在安慰中，你们就绝不会因我们这次来迟而责备我们，因为我们既为你们受磨难，也为你们得安慰。为免有人以为这话难听，说\Quote{我们如此受苦是为你们}，他又加上\Quote{我们得安慰也是为你们}，并且说\Quote{不只是我们单独涉险；因为你们也}，\Quote{是同受一样苦难的人}。这样，他承认他们是有分于危难的人，并把自身安慰的原因归于他们，就缓和了自己的话。因此，若我们因诡计而受困，你们要鼓起勇气；我们忍受这事，是为使你们的信德得以增长。若我们得安慰，也要以此夸耀；因为我们也为你们而享用此安慰，好使你们因分享我们的喜乐而得着一些鼓励。他所说的安慰，不仅来自他们自己得安慰，也来自知道他们（宗徒们）在安歇，听他接着宣告：\Quote{知道你们怎样在苦难上有分，也照样在安慰上有分}。因为就如我们受迫害时，你们也忧伤，仿佛是自己在受苦；同样我们也确信，当我们得安慰时，你们也认为这享受是你们的。还有什么比这种精神更谦卑？他在危难中远超过众人，却称他们为\Quote{有分的人}，他们并未经历任何部分；至于安慰，他却将全部原因归于他们，而非自己的劳苦。\par

3. 随后，先前只是笼统地提到磨难，现在他也提到了他们（Ben. 他）忍受磨难的地方。\par

% structure: p0012 source=h2 target=subsection reason=relative-heading-rank; base=h2

\subsection{格林多后书 1:8}

\Quote{弟兄们，我们不愿意你们不知道，我们在亚细亚所遭遇的磨难。}\par

他说：\Quote{我们讲这些话，是叫你们不至于不知道我们所遭遇的事；因为我们愿意，实在尽心竭力，要你们知道我们的事。}这是爱的一个很高证明。这事他在前书已预先提过，在那里他说：\Quote{因为我在厄弗所开了广大有效之门，且有许多敌人}，见：\BibleRef{格前 16:8}。他于是提醒他们这件事，并诉说他所受的苦，他说：\Quote{我不愿意你们不知道我们在亚细亚所遭遇的磨难。}他在致厄弗所人的书信中也说了同样的话。因为他派提希苛到他们那里去，就把这作为他旅行的理由；因此他说：\Quote{但为叫你们也知道我的事，并我如何，主内亲爱的兄弟和忠信的仆役提希苛，会告诉你们一切；我打发他到你们那里去，正是为这个缘故，叫你们知道我们的景况，并叫你们的心得安慰。}，见：\BibleRef{弗 6:21}。他在其他书信中也这样做。这并非多余，甚至极为必要：一方面是由于他对门徒的极大深情，另一方面是由于他们持续的试炼；在这种试炼中，互知对方境况是一种极大的安慰；这样，若这些境况是悲惨的，他们就可以预备好，既振作精神，又能更安全地不跌倒；若这些境况是好的，他们就可以与他们同乐。然而，他在这里既谈到从试炼中得救，也谈到受试炼攻击，说：\Quote{我们被压得太重，超过我们的力量。}如同一艘船在巨重之下沉没。他这里似乎只说了两件事：\Quote{非常}和\Quote{超过我们的力量}；然而，这不是一件，而是两件；为免有人反驳：\Quote{那么怎样？就算危险非常，但对你并不大；}他补充说，它既大又超过我们的力量，是的，如此超过，\Quote{以致我们连活命的指望都绝了。}\par

这就是说，我们再也没有活命的指望了。达味所说的\Quote{地狱的门，剧痛}和\Quote{死亡的阴影}，这正是他用\Quote{我们忍受了必死无疑的危险}来表达的。\par

% structure: p0016 source=h2 target=subsection reason=relative-heading-rank; base=h2

\subsection{格林多后书 1:9}

\Quote{但我们自己有死亡的回答，使我们不依靠自己，而依靠那使死人复活的天主。}\par

这是什么？\Quote{死亡的回答}？就是表决、判决、期望。因为我们的处境就是这样说的；我们的命运给了这个回答：\Quote{我们必定要死。}\par

诚然，这并未成为事实，只停留在我们的预期中，便止步了：因为我们的处境确实如此宣告，但天主的权能未让那宣告生效，只容许它发生在我们思想和期望中。因此他说：\Quote{我们自身已领受了死亡的判决，}并非实际如此。为何祂容许如此大的危险，以致夺去我们的希望，并使我们绝望？\Quote{为叫我们不依靠自己，}他说，\Quote{而只依靠天主。}保禄说这些话，并非这是他本人的性情——切莫如此想——而是为了以他所言调节其余的人，并出于极大的谨慎以谦逊说话。因此后来他又说：\Quote{有一根刺加在了我身上}（意指他的试探），\Quote{免得我过于高傲。}\BibleRef{格后 12:7}然而天主并未说祂容许这些是为了这个，而是为了另一个理由。什么理由？为使祂的能力更显明；\Quote{因为，}他说，\Quote{我的恩宠为你足够了，因为我的能力在软弱中才得以完全。}\BibleRef{格后 1:9}但是，如我所言，他从未忘记自己的独特品格，将自己与那些极为欠缺、急需许多管教和纠正的人同列。因为如果一两次试探就足以使普通人清醒，那么他——一生中最培养谦卑之心、受苦无人能及的人，在这么多年和堪配诸天的智慧操练之后——怎会还需要这种训诫呢？由此可见，在这里也是出于谦逊，并为了安抚那些自视过高、自夸的人，他才如此说：\Quote{为叫我们不依靠自己，而只依靠天主。}\par

4. 请看他在此处又是如何温柔地对待他们。他说，这些试探之所以临到我们，是为了你们的缘故；你们在天主眼中何等宝贵；因为\Quote{我们若受磨难，那是为你们的安慰和救恩；}但它们\Quote{超过限度}是为了我们，免得我们心高气傲。\Quote{因为我们被压得太重，过于我们所能承受，为叫我们不依靠自己，而只依靠那使死人复活的天主。}他再次提醒他们关于复活的道理，这一点他在前书已详述，并借当前的情形加以确认；因此他接着说：\par

% structure: p0021 source=h2 target=subsection reason=relative-heading-rank; base=h2

\subsection{格林多后书 1:10-11}

\Quote{祂从如此大的死亡中救了我们。}\par

他没有说\Quote{从如此大的危险中}，这样既显出试探那不堪忍受的严酷，又证实了我所提到的道理。因为虽然复活是未来的事，他却表明它每日都在发生：因为当天主使一个已绝望、已到了阴府门口的人重新站起时，祂无非是显明了复活，从死亡的獠牙中夺回那落入其中的人。因此，对于那些已绝望却从重病或不胜任的试探中恢复的人，我们常说：我们在他身上看见了死人复活。\par

\Quote{并且我们指望祂还要继续救我们；你们也一同以祈祷帮助我们，好使许多人因着那藉着多人所赐予我们的恩惠，为我们献上感谢。}\par

既然\Quote{为叫我们不依靠自己}这话可能被视为一种共同的指责和对某些人的控告，他便再次缓和语气，将他们的祈祷称为极大的保护，同时表明我们的生命必须是一场持续的争战。因为在那话中，\Quote{并且我们指望祂还要继续救我们，}他预言了未来如冰雹般众多的试探：但仍无一处提到被遗弃，反而提到再次的帮助与支持。然后，免得他们因听到将不断处于危险而灰心，他先指明了危险的益处，例如：\Quote{为叫我们不依靠自己；}即，为使祂使我们保持持续的谦卑，并使他们的救恩得以成就；此外还有许多益处：与基督同苦（因为他说：\Quote{基督的苦难在我们身上洋溢）；}为信友受苦（因为他说：\Quote{我们若受磨难，那是为你们的安慰和救恩）；}这后者（即他们的救恩）应放射出更耀眼的光辉；他说：\Quote{这要在你们忍耐同样苦难时，得以运行；}使他们变得坚毅；此外还有，在他们眼前生动描绘复活的情景：因为\Quote{祂从如此大的死亡中救了我们；}他们怀着热切的心，不断仰望祂，\Quote{因为，}他说，\Quote{我们指望祂要救我们；}这又使他们紧系于祈祷，因为他说：\Quote{你们也一同以祈祷帮助我们。}如此，他既表明了苦难的益处，又使他们充满活力；随后再次膏抹他们的精神（为作战），并通过见证他们祈祷的伟大之处来激励他们行善，因为天主正是因这些祈祷而赐予了保禄；正如他所说：\Quote{你们藉着祈祷一同帮助我们。}然而\Quote{好使许多人因着那藉着多人所赐予我们的恩惠，为我们献上感谢}是什么意思？他说：\Quote{祂从那些死亡中救了我们，你们也藉着祈祷一同帮助；}即，你们众人都为我们祈祷。因为\Quote{那所赐予我们的恩惠，}即我们的救恩，祂乐意赐予你们众人，好使多人向祂献上感谢，因为多人也领受了这恩赐。\par

5. 他说这话，一面是为了激励他们为他人祈祷，一面是为了使他们习惯于为他人所遭遇的一切事常感谢天主，表明祂也非常愿意这样做。因为那些小心为他人做这两件事的人，更会为自己树立两者的榜样。除此之外，他也教导他们谦卑，并引导他们走向更热烈的爱德。因为如果这位远超于他们之上的人承认自己曾因他们的祈祷而得救，并且自己作为天主的恩赐被赐予他们的祈祷，那么试想他们应有的谦逊和态度。再者，我请你们注意这一点：即使天主以慈悲行事，祈祷也对此大有贡献。因为起初他将自己的救恩归于天主的慈悲；因为他说：\Quote{慈悲的天主}自己\Quote{拯救了我们}，但在此也归于祈祷。因为那个欠一万塔冷通的人也是在俯伏在祂脚前之后才蒙了慈悲；\BibleRef{玛 18:24}尽管经上记载：祂\Quote{动了慈心，就释放了他}。再如，对那个\Quote{客纳罕妇人}，也是在她长期等候和坚持之后，\BibleRef{玛 15:22}祂才最终赐给她女儿痊愈，虽然祂是出于慈悲医治了她。由此我们学到，即使我们将要蒙受慈悲，也必须先使自己配得慈悲；因为虽有慈悲，它也只寻找那些配得的人。它不会不加区别地降临在所有人身上，甚至那些没有感觉的人身上；因为祂说：\BibleQuote{我要怜悯谁，就怜悯谁；我要慈悲谁，就慈悲谁。}\BibleRef{罗 9:15}至少请看他在这里所说的：\Quote{你们也一同用祈祷帮助我们。}他既没有将全部善工归给他们，免得高举他们，也没有剥夺他们所有的份，为鼓励他们，激发他们的热忱，并使他们彼此联合。因此他也说：\Quote{他也将我的安全赐给了你们。}因为常常天主也会因众人同心合意的祈祷而感到羞愧。因此祂对先知说：\BibleQuote{难道我不该怜惜这座有十二万多人的大城吗？}\BibleRef{纳 4:11}然后，免得你以为祂只尊重人多，祂又说：\BibleQuote{以色列的人数虽多如海沙，但只有残存的子民要得救。}\BibleRef{依 10:22}那么，祂是如何拯救尼尼微人的呢？因为在他们的例子中，不仅是人数众多，而且是人数众多且德行兼备。因为每一个人\Quote{离开}了自己的\Quote{恶道}。\BibleRef{纳 3:10}此外，在拯救他们时，祂说他们不能分辨\Quote{右手和左手}：由此明显看出，甚至以前，他们犯罪更多是由于单纯而非邪恶；从他们听了几句话就归正也可见一斑。但如果他们有十二万之众本身就足以拯救他们，那么以前是什么阻止他们得救呢？为什么祂不对先知说：难道我不该怜惜这正归正的城市吗？却提出那十二万之数。祂也将此作为一个额外的理由提出。因为他们的归正先知是知道的，但先知不知道他们的人数或单纯。所以祂用一切可能的考虑来软化他们。因为人数的众多是有力量的，当有德行伴随的时候。确实，圣经在别处也清楚地表明了这一点，它说：\BibleQuote{教会为他一心热切地向天主祈祷。}\BibleRef{宗 12:5}它具有如此大的力量，即使当门被锁、锁链加身、守卫左右睡卧时，它却使宗徒走出并脱离他们所有人。但是，哪里有德行，人数众多就有强大的力量；哪里邪恶，它就毫无益处。因为以色列人，祂说他们的人数如海沙，却全部灭亡了；诺厄时代的众人也是人数众多，甚至无数，但这对他们毫无益处。因为人数的众多本身没有力量，只有当它是附属时。\par

6. 那么，让我们勤勉地聚在一起恳求；让我们彼此代祷，如同他们为宗徒们所做的一样。因为这样做，我们既履行了诫命，又被傅油于爱：（当我说爱时，我是指一切美善之事：）也学习以更大的热诚感恩：因为那些为他人之事感恩的人，更加会为自己的事感恩。这也是达味常做的，他说：\Quote{你们同我一起赞扬主，让我们一同颂扬祂的名号；}\BibleRef{咏 33:3} 宗徒也处处要求这样做。让我们也努力这样做；让我们向众人显明天主的恩惠，以便在赞美中获得同伴：因为如果我们从人那里得到好处，通过宣扬出来，我们就使他们更愿意服事我们；那么，我们通过传扬天主的恩惠，更能吸引祂施予更多的善意。如果我们从人那里得到恩惠，尚且激励他人也加入我们感恩的行列，那么，我们更应当带领许多人到天主面前，请他们为我们感恩。因为如果保禄在天主面前有那么大的信赖尚且这样做，我们更有必要这样做。那么，让我们劝勉圣徒为我们感恩；我们自己也彼此这样做。这善工特别属于司铎，因为这是极大的特恩。因为走近时，我们首先为普世和众人共享的美善感恩。因为即使天主的恩赐是共同的，但共同的保存也包括了你自己的；所以你既为自己的独特祝福欠下共同的感恩，也应当为共同的祝福献上自己独特的赞美：因为祂点燃太阳不是为了你一人，而是为了所有人共享；但尽管如此，你个人却拥有完整的太阳。因为它被造得如此之大是为了公共益处；然而你个人看到它如同所有人看到的一样大；因此你欠下的感恩与所有人加起来一样多；你应当为所有人共有的东西感恩，同样也为别人的德行感恩：因为因着别人，我们也领受许多祝福：因为如果在索多玛只找到十个义人，他们就不会遭受那样的灾祸。因此，让我们也为别人在天主前的信赖感恩。因为这种习俗是古老的，从起初就植根于教会。这样，保禄也为罗马人感恩，\BibleRef{罗 1:8} 为格林多人感恩，\BibleRef{格前 1:4} 为全体世人感恩，\BibleRef{弟前 2:1} 不要对我说：\Quote{这善工不是我的；} 因为即使不是你的，你仍应当感恩，因为你的肢体是这样一位。此外，通过你的欢呼，你使它成为你自己的，你有份于冠冕，并且你本人也将领受恩赐。正因如此，教会的法律也命令这样祈祷，不仅为信友，也为慕道者。因为法律激励信友为未领洗者恳求。因为当执事说：\Quote{我们要恳切地为慕道者祈祷，} 他无非是激发全体信友为他们祈祷；虽然慕道者仍是局外人。因为他们还不是基督身体的肢体，还没有领受奥迹，仍然与属灵的羊群分开。但如果我们该为这些人代祷，那么更该为我们自己的肢体代祷。因此他甚至说：\Quote{让我们恳切地祈祷，} 叫你们不要不认他们为局外人，不要忽视他们如同陌生人。因为他们还没有基督引入的指定祈祷；他们还没有信赖，却需要已领洗者的帮助。因为他们站在君王的庭院之外，远离圣所。因此，当那些可敬畏的祈祷献上时，他们甚至被赶走。因此他也劝勉你为他们祈祷，使他们成为你们的一份子，不再作陌生人或外人。因为\Quote{让我们祈祷}这话不仅是对司铎说的，也是对组成会众的人说的：因为当他说：\Quote{让我们列队：让我们祈祷；} 他劝勉所有人祈祷。\par

7. 然后他开始祈祷，说：\Quote{愿全怜悯而仁慈的天主垂听他们的祈祷。} 为了叫你不至于说：我们祈祷什么？他们是外人，还没有与身体联合。我怎能强制天主的垂顾？我怎能说服祂赐予他们怜悯和宽恕？为了叫你不被这些问题困惑，请看他如何解开你的困惑，说：\Quote{全怜悯而仁慈的天主。} 你听见了吗？\Quote{全怜悯的天主。} 不要再困惑了。因为全怜悯者怜悯一切，包括罪人和朋友。所以不要说：\Quote{我该如何为他们接近祂？} 祂自己会垂听他们的祈祷。慕道者的祈祷，除了使他们不再做慕道者，还能是什么呢？接着，他也提示了祈祷的方式。这是什么？\Quote{愿祂开启他们心灵的耳朵；} 因为它们至今仍闭塞受阻。\Quote{耳朵，}他说，不是外在的耳朵，而是悟性的耳朵，\Quote{以便听到\InnerQuote{眼所未见，耳所未闻，人心所未想到的事。}}\BibleRef{格前 2:9} 因为他们还没有听到那说不尽的奥迹；他们站在离奥迹很远的地方；即使他们听到，也不知道所说的是什么；因为那些奥迹需要很多理解，不仅是听闻而已：他们还没有内在的耳朵；因此他接着为他们祈求先知的恩赐，因为先知这样说话：\Quote{天主赐给我训诲的舌头，使我懂得适时说话；因为祂开启了我的口；祂清晨就赐给我；祂赐给我听话的耳朵。}\BibleRef{依 1:4} 因为先知们听到的方式与众人大不相同，同样，信友也与慕道者不相同。借此，慕道者也受教导：不要向人学习听这些事，（因为祂说：不要称地上的人为师父，而要来自上面，来自天上，\Quote{因为他们都要受天主的教导。}\BibleRef{依 54:13}）\par

因此他说：\Quote{将真理之言灌输给他们，}使之能被内心学习，因为他们尚不按所应当的方式认识真理之言。\Quote{祂愿将祂的敬畏播撒在他们心中。}但这还不够；因为\Quote{有的落在路旁，有的落在石头上。}但我们并非如此祈求；而是如同在肥沃的土地上，犁翻开沟畦，我们也祈祷在这里也能如此，使他们心田的休闲地深耕之后，能接受撒在他们身上的种子，并准确记住所听见的一切。因此他又说：\Quote{并在他们心中坚定祂的信德；}就是说，使信德不浮于表面，而是向下深深扎根。\Quote{祂愿向他们揭示正义之福音。}他指出这幔子是双重的：一方面他们悟性的眼睛被闭住，另一方面福音向他们隐藏。因此他稍前说：\Quote{祂愿开启他们心中的耳朵，}而这里说：\Quote{祂愿向他们揭示正义之福音；}也就是说，既使他们有智慧并适合接受种子，又教导他们将种子撒入；因为即使他们适合，但若天主不启示，也无益处；若天主揭示而他们不接受，同样无益。因此我们为两者祈求：祂既开启他们的心，又揭示福音。因为即使王室的饰品放在幔子下面，眼睛看着也无济于事；同样，即使饰品显露，眼睛若不醒着也无益。但两者都会赐予，如果他们自己首先渴望。那么，\Quote{正义之福音}是什么呢？就是使人成为正义的福音。藉着这些话，他引导他们渴望圣洗，表明福音不仅为罪过的赦免，也为正义而工作。\par

8. \Quote{祂愿赐给他们一个属神的心智、健全的判断和有德行的生活方式。}让那些依附于现世事物的信友们留意。因为如果我们被命令为未领洗者祈求这些，试想我们这些为他人祈求的人应当从事什么样的事情。因为生活方式应与福音保持一致。因此，祈祷的顺序也从（福音的）道理转向品行：因为对于\Quote{祂愿向他们揭示正义之福音}这句话，它加上了\Quote{祂愿赐给他们一个属神的心智。}这\Quote{属神的}是什么意思？就是天主能居住在其中。因为祂说：\BibleQuote{我要在他们内居住，在他们中行走；}\BibleRef{肋 26:12}因为当心智成为正义的、脱去罪过时，它就成为了天主的居所。\BibleRef{罗 6:16}但当天主内住时，人的任何东西都不会留下。这样，心智就成了属神的，每一句话都出于祂，就如同天主的居所真实地居住在它内。那么，言语污秽的人当然没有属神的心智，喜欢戏谑和嬉笑的人也没有。\par

\Quote{健全的判断。}拥有\Quote{一个健全的判断}能是什么呢？就是享受灵魂的健康；因为被邪恶的私欲压制、被现世事物迷惑的人，永远不能健全，即健康。但如同患病的人贪求不适合他的东西，他也这样。\Quote{和有德行的生活方式，}因为道理需要相应的生活方式。请注意这事，你们这些在生命临终时来领受圣洗的人，因为我们确实祈祷你们在圣洗后也能拥有这种品行，但你们却在寻求并尽力在无此品行时离去。因为，即使你被成义，那也只是藉着信德。但我们祈祷你们也拥有来自善行的信心。\par

\Quote{不断地思念那些属于祂的事，留意那些属于祂的事，实行那些属于祂的事：}因为我们祈求不仅拥有一天的健全判断和有德行的品行，或两三天，而是贯穿我们整个生活的脉络和时期；并作为一切善事的基础，\Quote{留意那些属于祂的事。}因为许多人\BibleQuote{寻求自己的事，而非耶稣基督的事。}\BibleRef{斐 2:21}那么这怎能成就呢？（因为除了祈祷，我们也需要贡献自己的努力。）如果我们日夜专心于祂的法律。因此他继续祈求这一点：\Quote{专心于祂的法律；}正如他上面说\Quote{不断地，}现在这里说\Quote{日夜。}因此我甚至为那些一年几乎只在教堂出现一次的人脸红。因为那些奉命不仅\Quote{日夜}与法律交流，而且\Quote{专心于}法律，即永远与它交谈，却几乎在生命中最微小的一部分也不这样做的人，还有什么借口呢？\par

\Quote{记念祂的诫命，遵守祂的典章。}你看见这里何等卓越的链条吗？每一个链环如何彼此衔接，比任何金链更坚固、更美丽？因为祈求了一个属神的心智后，他告诉我们这如何能够产生。如何？通过不断地实践它。这又如何实现？通过不断注意法律。如何能说服人们做到这一点？如果他们遵守祂的诫命；不仅如此，注意法律也产生对祂诫命的遵守；同样，留意那些属于祂的事、拥有属神的心智，也产生实行那些属于祂的事。因为所提及的每一件事都共同产生并被下一件事产生，既连接它，又被它连接。\par

9. \Quote{让我们更恳切地为他们祈求。}因为长时间说话，灵魂常会变得昏沉，所以他再次唤醒它，因为他打算再次祈求某些伟大崇高之事。因此他说：\Quote{让我们更恳切地为他们祈求。}这是什么？\Quote{求祂拯救他们脱离一切邪恶和混乱无序之事。}这里我们为他们祈求，使他们不陷于诱惑，而是从每个陷阱中被救出，一种身体和精神的拯救。因此他接着说：\Quote{脱离一切魔鬼的罪恶和仇敌的围攻，}意指诱惑和罪恶。因为罪恶容易围攻，从四面八方、前后左右站定，从而打倒。因为，在告诉我们应当做什么之后——即：沉浸于祂的法律，记住祂的诫命，遵守祂的典章——他接着保证，即使这样也不够，除非祂亲自站在旁边并施以援手。因为，\BibleQuote{除非主建造房屋，建造的人就枉然劳力；}\BibleRef{咏 126:1}尤其是在那些仍暴露于魔鬼并受其统治的人身上。你们这些已领洗者深知此理。比如，回想那些话语，你们在其中弃绝了他的篡权统治，屈膝并投奔君王，说出了那些令我们学会绝不顺服他的可怕言辞。但他称他为仇敌和控告者，因为他既向人控告天主，又向天主控告我们，也使我们互相控告。因为有一次他向天主控告约伯说：\BibleQuote{约伯无故敬畏天主吗？}\BibleRef{约 1:16}另一次，天主对约伯：\BibleQuote{火从天降下。}又再一次，天主对亚当，\BibleRef{创 3:5}当祂说他们的眼睛要被打开。并且对今日的许多人说，天主不关心可见的事物秩序，而是将你们的事务委托给了邪魔。并且对许多犹太人，他控告基督，称祂为骗子或巫师。但也许有人想听他以何种方式运作。当他找不到虔敬的心灵，找不到健全的理解时，那么，就像进入一个空洞的灵魂，他就在那里狂欢；当一个人不记得天主的诫命，不遵守祂的典章时，他就掳掠他并离去。例如，假如亚当记得那条诫命：\BibleQuote{你可以吃各种树上的果子；}\BibleRef{创 2:16}假如他遵守了那条判断：\BibleQuote{你在吃的那天，必定要死；}他就不会落到那样的下场。\par

\Quote{求祂使他们在适当的时候配得圣洗的重生、罪过的赦免。}因为我们祈求有些事情现在来临，有些事情将来来临；并且我们阐述圣洗的道理，通过祈求教导他们认识其能力。因为此后所说的话使他们已经知道那里所发生的是重生，并且我们由水重生，正如由母腹重生；免得他们像尼苛德摩一样说：\BibleQuote{人已经老了，如何能重生呢？难道他还能进入母亲的母腹而重生吗？}然后，因为他已经提到了\Quote{罪过的赦免}，他通过接下来的话确认这一点：\Quote{不朽的衣裳；}因为那穿上儿子名分的人显然成为不朽的。但什么是\Quote{在适当的时候？}当任何人心怀善意，当任何人以热忱和信德前来；因为这就是信友的\Quote{适当时候}。\par

10. \Quote{祂要祝福他们的出入，即他们一生的全部路程。} 这里他们被指示甚至祈求一些身体上的好处，因为他们还相当软弱。\Quote{他们的房屋和他们的家庭，} 即，如果他们拥有仆人或亲属或任何其他属于他们的人。因为这是旧盟约的赏报；当时没有比寡居、无子女、不期而至的哀悼、遭受饥荒、事情不顺遂更令人恐惧的。因此，他允许他们怀着喜爱之情停留在较为物质的祈求上，使他们逐步攀升到更高的事物。因为基督也是如此，保禄也是如此，他们提及古代的祝福：基督说：\BibleQuote{温良的人是有福的，因为他们要承受土地；} 保禄说：\BibleQuote{应孝敬你的父亲和你的母亲……并在世上延年益寿。} \Quote{祂要增加他们的子女，祝福他们，使他们长大成人，并教导他们智慧。} 这里又同时是身体和灵性的事，如同对待在性情上仍太过孩子气的人。随后的话全然是灵性的：\Quote{祂要把他们面前的一切导向益处；} 因为他并非简单地说\Quote{他们面前的一切，} 而是\Quote{他们面前的一切导向益处。} 因为人面前常有旅程，但并非有益；或其他类似的事，没有益处。他们由此被教导在一切事上感谢天主，因为一切都为好处而发生。在这之后，他吩咐他们在接下来的部分站立。因为他先前将他们俯伏在地，当他们求问了所当求的并充满了信心之后，现在所宣示的话使他们站起，并吩咐他们在接下来的部分也为他们自己向天主恳求。因为我们自己说一部分，又允许他们说一部分，现在为他们打开祈祷之门（正如我们先教儿童说什么，然后让他们自己说一样），说：\Quote{请你们祈祷，慕道者，为平安的天使；} 因为有惩罚的天使，如经上说：\BibleQuote{一群邪恶的天使} \BibleRef{咏 77:49}，有毁灭的天使。因此我们吩咐他们求平安的天使，教导他们寻求那万善之链的平安；好使他们能摆脱一切争斗、一切战争、一切纷争。\Quote{愿你们面前的一切都平安；} 因为即使一件事是沉重的，如果人有平安，它也是轻省的。因此基督也说：\BibleQuote{我把我的平安赐给你们} \BibleRef{若 14:27}。因为魔鬼没有比争斗、敌意和战争更强大的武器。\Quote{祈求这一天和你们一生的所有日子都充满平安。} 你看到他如何再次坚持整个生命要在德行中度过吗？\Quote{愿你们的结局是基督徒的；} 你们最高的善，即荣耀与有益；因为不荣耀的也不是有益的。因为我们对益处的本质看法与多数人不同。\Quote{把你们自己交托给永生的天主和祂的基督；} 因为此时我们还不放心他们为别人祈祷，只需他们能为自己祈祷就够了。\par

你看到这祈祷无论在教义还是行为上的完整性吗？因为当我们提及福音、不朽的衣袍和重生之洗时，我们提及了全部的教义；当我们再说虔敬的心、健全的理智以及我们所说的其余一切时，我们提示了生活方式。然后我们吩咐他们低头；视之为天主祝福他们、他们的祈祷蒙垂听的证明。因为这祝福的确实不是人；而是借着他的手和他的舌，我们将在场者的头带到君王本人面前。所有人都同声高呼\Quote{阿们。}\par

现在我说这一切是为什么？为教导你们，我们应当寻求别人的事，好使信徒不认为当这些事被说出时与他们无关。因为执事说\Quote{让我们为慕道者祈祷}，难道是对着墙壁说的吗？但有些人竟如此无知、如此愚蠢、如此败坏，不仅在慕道者时间站立交谈，甚至也在信徒时间如此。因此一切都颠倒了；因此一切都完全丧失了：因为正当我们最应该平息天主的时候，我们却离去而激怒了祂。同样，在信徒的祈祷中，我们被吩咐接近那爱人的天主，为主教、司铎、君王、掌权者、地、海、季节和整个世界祈祷。那么，当我们这些本应如此大胆以至为他人祈祷的人，甚至为自己祈祷时也几乎不醒觉，我们怎能原谅自己？怎能获得宽恕？因此我恳求你们，把这一切放在心上，认识祈祷的时刻，被提升起，脱离尘世，触及天穹本身；好使我们有能力平息天主，获得所应许的美物。愿我们众人借我们的主耶稣基督的恩宠和爱人之心，偕同祂，及圣神，归于圣父，享受荣耀、权能、尊崇，自今至永远，及世之世。阿们。\par

\SourceRecord{Translated by Talbot W. Chambers.；Nicene and Post-Nicene Fathers, First Series；Vol. 12.；Edited by Philip Schaff.；Buffalo, NY: Christian Literature Publishing Co.,；1889.；<http://www.newadvent.org/fathers/220202.htm>.}

% structure: p0001 source=h1 target=section reason=page-title

\section{论格林多后书第三篇讲道}

% structure: p0002 source=h2 target=subsection reason=relative-heading-rank; base=h2

\subsection{格林多后书 1:12}

\BibleQuote{我们的夸耀就是这良心作证，即我们以正直和真诚，不是凭属世的智慧，而是凭天主的恩宠，在这世上为人处世。}\par

他又为我们开启了另一种安慰的依据，而且不是小小的，而是极大的，很适合振作在危险中沉沦的心灵。因为他曾说过：天主安慰了我们，天主拯救了我们，并将一切都归于祂的怜悯和他们的祈祷；免得这样使听者变得懒惰，只依赖天主的怜悯和他人的祈祷，他指出他们自己也贡献了不少。实际上，他先前也表明了这一点，当他说：\BibleQuote{因为基督的苦难怎样在我们身上越分，我们所得的安慰也怎样越分。}\BibleRef{格后 1:5}但这里他谈论的是另一种善工，确实是他们自己的。这是什么？他说：以纯洁无伪的良心，我们在各处为人处世；这大大有助于我们的鼓励和安慰；甚至不只是安慰，而是比安慰更伟大的东西，甚至我们的夸耀。他说这话，也是教导他们在苦难中不要消沉，反而如果他们有纯洁的良心，甚至要以此为荣；同时，他也温和而柔和地击打假宗徒。正如在前一封书信中他说：\BibleQuote{基督派遣我传福音，不是凭言辞的智慧，免得基督的十字架失去效力。}\BibleRef{格前 1:17}以及\BibleQuote{你们的信德不在于人的智慧，而在于天主的德能。}\BibleRef{格前 2:5}所以这里也说：\Quote{不是凭智慧，而是凭基督的恩宠。}\par

他还暗示了另外一些事，用了\Quote{不是凭智慧}这些词，也就是说，\Quote{不是出于诡诈}，这里也打击了异教的教学。\Quote{我们的夸耀}，他说，\Quote{就是这良心作证}；也就是说，我们的良心没有可以定我们罪的，好像我们因作恶而受迫害。因为我们虽然遭受无数恐怖，虽然四面受敌、处于危险，但足以安慰我们，甚至不仅安慰，还给我们加冕，就是我们的良心纯洁，见证我们不是因作恶，而是因中悦天主的事而受苦；为德行的缘故，为属天的智慧，为许多人的救恩。之前的安慰来自天主，但这是他们自己贡献的，来自他们生活的纯洁。所以他也称之为他们的夸耀，因为这是他们自己德行的成就。那么，这夸耀是什么，我们的良心见证什么？\Quote{那就是以真诚}，也就是说，没有任何诡诈、虚伪、假装、谄媚、埋伏或奸诈，以及其他类似之事，而是全然坦率、朴实、真实，以纯洁无恶意的心灵，无狡诈的思维，没有隐藏，没有溃烂的疮。\Quote{不是凭属世的智慧}，也就是说，不是凭邪恶的巧计，也不是凭邪恶，不是凭言辞的巧妙，不是凭诡辩的罗网，因为他以此意指\Quote{属世的智慧}；而他们所大大自夸的，他否认并推开：充分表明这不是可夸的根据；不仅他不寻求，反而拒绝并以此为耻。\par

\BibleQuote{而是凭天主的恩宠，我们在世上为人处世。}\par

什么是\Quote{凭天主的恩宠}？显示那来自祂的智慧，祂赐予我们的德能，通过所行的奇迹，征服了智者、修辞学家、哲学家、君王、民众，我们本身没有学问，也没有带来任何外来的智慧。然而，这并不是普通的安慰和夸耀，就是意识到自己不是使用了人的能力，而是藉着天主的恩宠取得了全部成功。\par

\Quote{在这世上。}所以不只在格林多，也在世界的每一部分。\par

\BibleQuote{向你们更加丰富。}什么是向你们更加丰富？\BibleQuote{凭天主的恩宠，我们为人处世。}因为我们在你们中间行了神迹和奇事，以及更大的严格，和无可指摘的生活；因为他称这些也是天主的恩宠，也将自己的善工归于它。因为在格林多，他甚至超越了目标，白白传福音，因为他宽容他们的软弱。\par

% structure: p0010 source=h2 target=subsection reason=relative-heading-rank; base=h2

\subsection{格林多后书 1:13}

\BibleQuote{因为我们写给你们的，不是别的，正是你们所读的，甚至所承认的。}\par

因为他谈到自己的大事，似乎为自己作证，这是可憎的事，他再次呼吁他们作为他所说的见证。因为，他说，不要让任何人认为我所说的是写作上的夸耀；因为我们向你们宣告你们自己所知的事；我们并没有撒谎，你们比任何人都更能为我们作证。因为当你读的时候，你承认你们所知我们在行动中实践的事，我们也在写作中说，你们的见证与我们的书信并不矛盾；但你们先前对我们的认识与你们的阅读是一致的。\par

% structure: p0013 source=h2 target=subsection reason=relative-heading-rank; base=h2

\subsection{格林多后书 1:14}

\BibleQuote{正如你们也部分地承认了我们。}\par

因为他对你们的认识，他说，不是来自传闻，而是来自实际经验。\Quote{部分地}这个词是他出于谦卑而加的。因为这是他的习惯，当有必要说出任何高调的话时（因为他从不无故这样说），他希望迅速压制因他所说的话而产生的得意。\par

\BibleQuote{并且我希望你们将一直承认到底。}\par

2. 你再次看到，他如何从过去为将来取得保证；不仅从过去，也从天主的德能？因为他并非绝对肯定，而是将一切交托给天主和他在祂内的希望。\par

\BibleQuote{我们是你们的夸耀，正如你们也是我们的夸耀，在我们的主耶稣基督的日子。}\par

他在此处断绝了他言语可能引发的嫉妒的根源，使众人分享并参与他善工的光荣。\Quote{因为这一切并不停留于我们，而是也传到你们那里，再从你们传给你们。}因为他既已赞扬了自己，又提供了过去的证据，并为将来作了保证；唯恐听众因他说话傲慢而责备他，或者，如我所说，仓促地陷入嫉妒，他便使这喜乐成为共同的，并宣告这赞美的冠冕是属于他们的。因为，他说，如果我们显明自己是如此，我们的赞美便是你们的光荣；正如当你们也蒙悦纳时，我们便喜乐、欢跃，并获得冠冕。在此，他又再次借所说的话显明他的极大谦卑。因为他如此谦逊地表达，不是像师傅对门徒讲话，而是像门徒对同等地位的同门门徒讲话。请看，他如何高举他们，以智慧充满他们，将他们送往那一日。因为，他说，不要对我提及现在的事，即众人的指责、辱骂、讥笑，因为此处的事，无论是善是苦，都无关紧要；也无关于人的讥笑或赞美；但请记起，我恳求你们，那充满恐惧和战栗的日子，那时一切都要显露。因为那时，我们将在你们身上光荣，你们也在我们身上光荣；当你们被看见有这样一些师傅，他们不教导人的道理，也不生活在邪恶中，也不给（人）任何把柄；而我们则有这样的门徒，既不按人的方式受影响，也不动摇，而是以爽快的心接受一切，不被来自任何方面的诡辩所迷惑。因为这事即使现在，对于有智慧的人也是明白的，但那时则对所有人显明。所以，即使我们现在受苦，我们仍有这安慰，这安慰并非轻微，即良心此刻所提供的，以及那时彰显本身。因为现在，我们的良心知道我们凭天主的恩宠做一切事，正如你们也知道并将知道；但那时，所有的人也都会知道我们和你们的行为，并要看见我们彼此因对方而受光荣。免得他显得只从这夸耀中独自获得光芒，他便也给予他们一个夸耀的理由，并引领他们离开当前的痛苦。正如他在安慰方面所说：\Quote{我们因你们而受安慰，} \BibleRef{格后 1:6} 他在这里也如此行，说：\Quote{我们因你们而夸耀，正如你们也因我们而夸耀，}处处使他们参与一切：他的安慰、他的苦难、他的保全。因为他将这保全归于他们的祈祷。\Quote{天主教拯救了我们，}他说，\Quote{你们藉祈祷一同帮助。}同样地，他也使夸耀成为共同的。因为正如他在那处所说：\Quote{知道你们既分担了苦难，也必分担安慰；}同样，此处：\Quote{我们是你们的夸耀，正如你们也是我们的夸耀。}\par

% structure: p0020 source=h2 target=subsection reason=relative-heading-rank; base=h2

\subsection{格后 1:15}

\Quote{我既然这样深信，就打算先到你们那里去。}\par

什么深信？\Quote{就是极其信赖你们，以你们为夸耀，是你们的夸耀，极其爱你们，自知心中毫无邪恶，确信我们中间一切都是属灵的，并且你们为此作见证。}\par

\Quote{我打算到你们那里去，然后由你们那里往马其顿去。}\par

然而，他在前书中所应许的却相反，这样说：\Quote{我要从马其顿经过时到你们那里去，因为我正要经过马其顿。} \BibleRef{格前 16:5} 那么他在这里怎么说相反的话呢？他并没有说相反的话：决不会。因为所说的话的确与他所写的相反，但并不与他所愿意的相反。\par

因此他在这里没有说：\Quote{我曾写道，我要经过你们那里往马其顿去；}而是说：\Quote{我打算。}\Quote{因为我虽然未曾那样写过，}他说，\Quote{我却极其渴望，并且\Quote{打算}在此之前到你们那里去；我远非愿意迟于我的许诺，而是乐意提前来到。}\Quote{使你们获得第二次的恩惠。} 什么是第二次的恩惠？\Quote{使你们获得双重的恩惠，既从我写作中，又从我亲临中。}这里他用\Quote{恩惠} 意指喜乐。\par

% structure: p0026 source=h2 target=subsection reason=relative-heading-rank; base=h2

\subsection{格后 1:16-17}

\Quote{我要由你们那里往马其顿去，再从马其顿回到你们那里，然后由你们送我到犹太去。我既有这样的心意，难道我是反复无常吗？}\par

3. 在此，在接下来的部分，他直接消除了因他的迟延和缺席而引起的指控。他所说的话性质如下：\Quote{我打算到你们那里去。}\Quote{那么我为什么没有来呢？难道我是轻浮和善变的吗？}因为这就是\Quote{反复无常}？决不是。但是为什么？\Quote{因为我所决定的，不是凭肉性决定的。} 什么是\Quote{不是凭肉性}？我决定\Quote{不是属血肉的。}\par

格后 1:17。\Quote{以致在我身上有是是和非非吗？} \BibleRef{格后 1:17}\par

但即使是这一点也仍晦涩不明。那么他说的是什么？属血肉的人，就是那紧附于现有事物并不断沉溺其中、且远离圣神感化范围的人，有能力随意游走，四处漂泊。然而，那作圣神之仆、被祂引领且处处受其引导的人，并不能完全主宰自己的意愿，因为他已使自己的意愿依附于从那里所得的权柄；他就像一位忠实的仆人，一举一动皆受主人命令的支配，自身毫无自主之力，甚至不能稍作歇息，若他向同僚许下诺言，却因主人另有安排而无法履行。这就是他所谓\Quote{我不是按照血肉来决意}的意思。我并未超越圣神的治理，也没有自由随意前往。因为我受制于护慰者的主权和命令，被祂的谕令引领且处处引导。正因如此，我未能前来，因为圣神不允。这正如\WorkTitle{宗徒大事录}中屡次发生的事；当他决意去某处时，圣神却命他往别处去。因此，我的未至并非出于轻浮，即反复无常，而是因为我服从圣神。你是否再次注意到他惯用的逻辑？他们想借以证明他\Quote{按照血肉决意}的理由——即他诺言的未履行——反而被他用作论据，特别证明他是按照圣神决意，而相反的决意倒是按照血肉。那么，有人会说：难道他当初许下的诺言不是出于圣神吗？绝非如此。因为我已经说过，保禄并非预知一切将要发生或适宜的事。正因如此，他在前书说：\BibleQuote{你们可在我无论往哪里去时，送我上路。}\BibleRef{格前 16:6}正是出于这种担忧：他若说了\BibleQuote{往犹太去}，可能被迫前往别处；而现在他的意图受挫，他却说：\BibleQuote{并由你们送我上路往犹太去}。凡是由于爱德而行的，他明说，即来到他们那里；但与她们无关的事——即从他们那里往犹太去——他没有明确吩咐。然而当他被证明有误后，他后来在此大胆地说：\BibleQuote{往犹太去}。这事也出于好意，免得他们中间有人对自己存有\BibleRef{宗 14:13}过高的想法。因为如果面对这些事，他们还想向他们献公牛，那么他们若不表现出许多人性的软弱，将会陷入何等不敬？何况他在许多事上尚且不知道将要发生的事，又何必惊奇？因为有时他在祈祷中也不知道何为宜。\par

因为他说：\BibleQuote{我们不知道该如何祈求才合乎正理。}而且，为了不显得他只是在谦逊地说此话，他不仅这样讲，还举出实例：他在祈祷中不知道什么为宜。那么是何实例？当他恳求脱离考验，说：\BibleQuote{有一根刺加在我身上，就是撒殚的使者来击打我。关于这事，我三次求主，祂对我说：\InnerQuote{我的恩宠够你用，因为我的德能是在软弱中得以成全。}}\BibleRef{格后 12:7}你看见吗？他不知道祈求什么为宜，因此虽然他多次祈求，却未获得。\par

% structure: p0032 source=h2 target=subsection reason=relative-heading-rank; base=h2

\subsection{格林多后书 1:18}

\Quote{但天主是忠实的，我们对你们所说的话，并不是\InnerQuote{是}而又\InnerQuote{非}。}\par

他巧妙地推翻了一个可能提出的异议。因为有人或许会说：如果你允诺之后又推迟前来，并且你的\Quote{是}不是\Quote{是}，\Quote{非}不是\Quote{非}，你先前说的后来又被否定——正如这次旅程中所行的那样——那么我们有祸了，如果这同样发生在宣讲中。为避免他们产生这些想法并因此困扰，他说：\BibleQuote{但天主是忠实的，我们对你们所说的话，并不是\InnerQuote{是}而又\InnerQuote{非}。}他说，这并非发生在宣讲中，而只发生在我们的旅行和行程中；而我们在宣讲中所说的任何话，都坚定不移、不可动摇（因为他在此称其宣讲为\Quote{话}）。接着，他提出一个无可辩驳的证明，将一切归于天主。他所说的是：\Quote{我前来的承诺是我的，我凭自己作了那个承诺；但宣讲不是我的，也不是出于人，而是出于天主；出于天主的，不可能说谎。}因此他也说：\Quote{天主是忠实的}，即\Quote{真实的}。\Quote{所以，不要怀疑那来自祂的，因为其中毫无人的成分。}\par

4. 既然他说了\Quote{话}，他就接着解释他指的是哪种话。那么是哪种呢？\par

% structure: p0036 source=h2 target=subsection reason=relative-heading-rank; base=h2

\subsection{格林多后书 1:19}

\Quote{因为天主之子，即由我们——我、息耳瓦诺和弟茂德——在你们中间所宣讲的，并不是\InnerQuote{是}而又\InnerQuote{非}。}\par

为此他列出教师们的同行，也以此证明那见证的可信性，因它来自教导者，而非仅听闻者。但他们原是门徒；然而他出于谦逊，视他们与教师同列。但所谓\Quote{并不是\InnerQuote{是}而又\InnerQuote{非}}是什么意思？他说，我从未否定过先前在宣讲中所说的话。我对你们所讲的不是此时这样说，彼时那样说。因为这不是出于信德，而是出自错谬的心意。\par

\Quote{在祂内是\InnerQuote{是}。}这就是说，正如我所言，这话坚定不移、稳固可靠。\par

% structure: p0040 source=h2 target=subsection reason=relative-heading-rank; base=h2

\subsection{格林多后书 1:20}

\Quote{因为天主的恩许不论有多少，在祂内都是\InnerQuote{是}；为此也藉着祂而说\InnerQuote{阿们}，以光荣天主，经由我们。}\par

这是什麼？\Quote{天主的恩许，无论有多少}？宣道曾应许了许多事；他们提出并宣讲了这些许多事。因为他们谈论复活、被提、不朽坏，以及那些伟大的赏报和不可言喻的美善。因此，关于这些恩许，他说它们坚定不移，其中没有是而又非；也就是说，所说的这些事并非时而真实、时而虚假，如同我在你们中间那样，而是永远真实。首先，他确实为信仰的条款和关于基督的圣言辩护，说：\Quote{我的话}和我的宣讲\Quote{不是是而又非}；其次，为恩许辩护：\Quote{因为天主的恩许，无论有多少，在祂内都是\InnerQuote{是}。}但如果祂所应许的事是确定的，祂必定会赐予，那么祂本身和关于祂的圣言就更是确定的，不能说祂时而存在、时而不存在，而是祂\Quote{永远}存在，并且是同一的。但\Quote{在祂内是\InnerQuote{是}和阿们}是什么意思？祂指的是那必定会成就的事。因为恩许是在祂内，而非在人内，有其存在和成就。因此，不要害怕；因为这不是人，使你应当怀疑；而是天主，祂既说了又成就。\Quote{借着我们，归于祂的光荣。}\Quote{借着我们归于祂的光荣}是什麼？祂借着我们成就它们，即借着祂对我们的恩惠，归光荣于祂；因为这就是\Quote{为光荣天主}。如果它们是为光荣天主，就必定会成就。因为祂不会轻视自己的光荣，即使祂曾轻视我们的救恩。但事实上，祂也不轻视我们的救恩，既因祂极爱人，也因我们的救恩与从这些事而来的祂的光荣相连。所以，如果恩许是为祂的光荣，我们的救恩也必定随之而来；在致厄弗所人书中，他不断回到这一点，说：\BibleQuote{为颂扬祂的光荣；}\BibleRef{弗 1:14}并且处处提出这一点，表明这一结果的必然性。为此，他在这里说，祂的恩许不撒谎：因为它们不仅拯救我们，也光荣祂。因此，不要因这些恩许是由我们而发出的而忧虑怀疑。因为它们不是由我们成就，而是由祂成就。是的，恩许也是由祂而来；因为我们向你们所说的，不是我们自己的话，而是祂的话。\par

% structure: p0043 source=h2 target=subsection reason=relative-heading-rank; base=h2

\subsection{格林多后书 1:21-22}

\Quote{然而那与你们同在基督内坚固我们，并给我们傅油的天主，就是那在我们心上盖印，并赐下圣神作为抵押的天主。}\par

祂又从过去的事确立将来的事。因为如果是祂在基督内坚固我们（即不允许我们从在基督内的信仰动摇），并且是祂给我们傅油，将圣神赐在我们心里，祂岂能不赐给我们将来的事呢？\par

因为若祂赐予了原则和基础、根和泉源（即对祂的真知识、对圣神的分享），祂岂不赐予由这些而来的一切？因为若为了这些，那些得了赐予，那么如今成为朋友后，祂岂不更要\Quote{白白赐予}他们那些？因此祂没有简单地说\Quote{圣神}，而称之为\Quote{抵押}，好叫你由此对全部之事怀有美好的希望。因为如果祂不打算赐予全体，祂绝不会选择赐予\Quote{抵押}而毫无目的和结果地浪费它。请留意保禄的坦率。因为何必说呢？他说，恩许的真实不在于我们？你们坚定不移的事实不在于我们，这也是出于天主；因为他说：\Quote{坚固我们的是天主。}不是我们坚固你们；因为我们也需要那坚固者。所以，不要有人认为宣道寄托在我们身上是危险的。祂担当了一切，祂关心了一切。\par

而\Quote{傅油}和\Quote{盖印}是什麼？祂赐予圣神，藉着祂同时做了这两件事，立我们为先知、司祭和君王，因为在古时，这三种人是要受傅的。但如今，我们没有其中一种尊位，却拥有所有三种卓越的尊位。因为我们既将享用王国，又通过献上我们的身体作为祭献而成为司祭（因为他说：\Quote{将你们的肢体献给天主，作为生活的祭品}）；同时，我们也被立为先知：因为\BibleQuote{眼所未见，耳所未闻}\BibleRef{格前 2:9}的事，已启示给我们了。\par

5. 此外，我们也在另一种方式上成为君王：如果我们有心思去支配我们不受控制的思念，因为这样的人就是君王，且比那戴冠冕者更为高贵，我现在要向你们说明这一点。他有许多军队，但我们也有思念，数量超过它们；因为我们无法数清我们内心无穷的思念。不仅要注意它们的众多，而且在这众多的思念中，还有许多将军、上校、队长、弓箭手、投石手。还有什么使一个人成为君王呢？他的衣裳？但这个人也穿着一件更好更勇敢的长袍，蛀虫不能蛀蚀，岁月不能损坏。他也有一个精巧的冠冕，那是光荣的冠冕，是天主怜悯的冠冕。因为\BibleBook{}{圣咏}作者说：\BibleQuote{我的灵魂，请赞美上主，那以怜悯和柔情给你加冕的。}\BibleRef{咏 103:2} 再次，光荣的冠冕：\BibleQuote{因为你以光荣和尊严给他加冕。}\BibleRef{咏 8:6} 以及\BibleQuote{你以恩惠给我们加冕，如同盾牌。}\BibleRef{咏 5:12} 再次，恩宠的冠冕：\BibleQuote{因为你必领受恩宠的冠冕戴在头上。}\BibleRef{箴 1:9} 你们看见这个由许多花环组成的冠冕吗？它在恩宠上超越了另一个。但让我们从起初重新更严格地探讨这些君王的状况。那个君王统治他的卫兵，向所有人发令，众人都服从并侍奉他；但在这里，我向你们显示更大的权威。因为这里的数目同样大甚至更大：还需要查问他们的服从。不要向我提出那些治理不当的，因为我也提出那些被自己的卫兵赶出王国并杀害的。那么让我们提出这些例子，但寻求那些妥善治理自己王国的人。你可以提出任何你愿意的人。我向你提出圣祖来对抗所有人。因为当他奉命祭献自己的儿子时，想想有多少思念起来反对他。尽管如此，他使所有思念屈服，它们在他面前颤抖，胜过君王在他的卫兵面前；只用一瞥，他就使它们全都静默，没有一个敢喃喃低语；反而都俯伏，如同向君王让位，一个也不落，虽然它们极其恼怒且毫不留情。因为许多士兵竖起的矛头，也不如当时那些可怕的思念那样令人恐惧，它们没有装备矛，但比许多矛更难对付——是自然的情感！因此它们能够刺透他的灵魂，胜过锋利的长矛。因为从来没有长矛能像那些思念的刺激那样锋利，这些刺激从下面，从他的情感中磨利并升起，穿透了那义人的心灵。因为在这里，需要时间、意图、一击和痛苦，然后死亡随之而来；但在那里，这些都不需要，伤口更迅速更剧烈。尽管如此，虽然如此多的思念当时武装起来反对他，却是一片深沉的平静，它们都整齐排列；更多是装饰他而不是恐吓他。至少看他伸出刀来，任凭你提出多少君王、皇帝、凯撒，你也不能说出像这样的事，没有类似的神态可以指出，如此高贵，如此配得上天。因为那位义人在那一刻建起了一座战胜最专制的暴政的纪念碑。因为没有比自然更暴虐的了；找到一万个诛杀暴君者，你绝不能向我们展示一个像这样的人。因为那一刻是天使的胜利，不是人的。试想。自然连同其一切武器和一切军队被击倒在地；而他伸出手，握着的不是冠冕，而是一把比任何冠冕更光荣的刀，天使的群众喝彩，天主从天上宣布他为得胜者。\par

因为他的国籍是在天上，从那里他也领受了那宣告。\BibleRef{斐 3:20} 还有什么比这更光荣的呢？或者，有什么纪念碑能与之相比呢？因为如果在一次摔跤手得胜的场合，不是地上的传令官，而是天上的君王站起来，亲自宣布他为奥林匹克的胜利者，这岂不在他看来比冠冕更光荣，并吸引了整个剧场的人注视他吗？那么，当不是人间的君王，而是天主自己，不是在这个剧场，而是在宇宙的剧场，在天使和总领天使的集会中，以高声呼喊从天上传扬他的名字时，请告诉我，我们应当给这位圣人安排什么地位？\par

6. 但如果你们愿意，让我们也听一听那声音本身。那声音是什么呢？\BibleQuote{亚巴郎，亚巴郎，不要下手害依撒格，也不要对他做什么。现在我知道你敬畏天主，且为了我的缘故，你没有吝惜你的儿子，你的爱子。}\BibleRef{创 22:11} 这是什么意思？那在万物存在之前就知悉一切者，难道他现在才知道吗？然而，即使对人来说，圣祖对天主的敬畏也是显而易见的：因为他给出了许多证据，证明他的心向天主是正直的，比如当时主对他说：\BibleQuote{离开你的故乡，你的亲族。}\BibleRef{创 12:1} 当他为了主的缘故和应该给主的尊敬，而将优先权让给他姊妹的儿子；当他从那么多危险中拯救了他；当他吩咐他去埃及，而他的妻子被夺去时，他没有抱怨，以及其他更多事例；正如我所说的，从这些事中，即使是人也能了解到圣祖对天主的敬畏，更何况那无需等行动就知道结局的天主呢？而且，如果他不知道，又怎能称他为义呢？因为经上记载：\BibleQuote{亚巴郎相信了，这就算作他的正义。}\BibleRef{创 15:6}\par

那么，\Quote{现在我知道？}是什么意思？叙利亚文是：\Quote{现在祢已使人知道；}也就是说，向人显明。因为我早已知道，甚至在这一切诫命之前。那么，为什么对人甚至说\Quote{现在？}难道那些行为不足以证明他的心对天主是正直的吗？它们固然足够，但这一件如此超越它们全体，以致在它旁边它们显得微不足道。所以，祂以推崇这件善工并表明它超越一切的方式如此说话。因为对于超过并超越先前一切的事，大多数人常这样说：例如，若有人从另一个人那里得到一件比先前任何礼物更大的礼物，他常说：\Quote{现在我知道某人爱我}，这句话并非表示他在过去不知道，而是意图表明现在给予的比一切都大。同样，天主以人的方式说话，说：\Quote{现在我知道}，只意图标明这功绩的无比伟大；并非祂\Quote{那时}才知道他的敬畏或敬畏的伟大。因为当祂说：\BibleQuote{来，我们下去，在那里看看}\BibleRef{创 11:7}，祂这样说并非因为需要下去（因为祂充满万有，确切知道一切），而是为教导我们不要轻易下判断。当祂说：\BibleQuote{主从天上垂看}\BibleRef{咏 13:2}，这是以人的比喻来描述祂的完全知识。同样，祂在这里说：\Quote{现在我知道}，是为了宣告这一件比先前所有都更伟大。祂自身也为此提供了证据，补充说：\Quote{因为你为了我的缘故，没有怜惜你的儿子，你的独生子。}祂不只是说\Quote{你的儿子}，更说\Quote{你的独生子}。因为不仅是本性，还有父母的疼爱，他既有天性的倾向，又有儿子的极好品性，却仍然胆敢在儿子身上蔑视这一切。如果父母对无价值的儿女也不容易无动于衷，甚至为他们悲哀；那么当这儿子是他的儿子、独生子、心爱的依撒格，而父亲本人正要祭杀他时，谁能描述这种超凡之情的极致呢？这一功绩胜过千千万万冠冕和无数的荣冠。因为佩戴那冠冕的人，死亡常常袭击并困扰他，死亡之前还有无数环境的冲击；但这冠冕，没人能有力量从佩戴者身上夺去，甚至死后也不能，无论是家人还是外人。让我指给你看这冠冕中最珍贵的宝石。因为就像一块珍贵的宝石，它出现在末尾并将冠冕扣紧。那么这宝石是什么？就是\Quote{为了我的缘故}这句话；因为此处令人惊叹的不是他不怜惜，而是\Quote{为了祂的缘故}。\par

哦！有福的右手，它被配得怎样的刀？哦！奇妙的刀，它被配得怎样的右手？哦！奇妙的刀，它被预备作何用途？它服事于何种职分？它服务于何种预像？它如何被染红？它如何未被染红？因为我不知该说什么，那奥秘如此可畏。它没有触及孩童的颈项，也没有穿过那圣者的喉咙；也没有被义人的血染红；不，它既触及了，穿过了，被染红了，也浸透了，却又未被浸透。或许我在你们看来是疯癫的，说出这样矛盾的话。因为，确实，我因思念那义人的奇妙行为而疯癫，但我并没有说出矛盾。因为那义人的手确实将刀刺入少年的喉咙，但天主的手却没有容许这样刺入的刀被少年的血玷污。因为不单是亚巴郎收回了手，天主也收回了：他凭借意愿作出了击刺，天主凭借话语制止了它。因为同一句话既武装了那只右手，也解除了它的武装，那只右手在天主麾下，如同在统帅之下，依祂的示意行事一切，一切都听从祂的声音。请看：祂说：\Quote{杀}，立时它武装起来；祂说：\Quote{不可杀}，立时它解除了武装；因为一切〔之前〕都已完全预备好了。\par

如今，天主向普世显示了这位战士和统帅；这位戴着冠冕的得胜者登上众天使的剧场；这位司祭、这位君王，用那刀戴上了比冠冕更美的冠冕，这位凯旋者、这位斗士、这位不战而胜的征服者。就如一位统帅若有一名最勇猛的士兵，便以他的武艺、姿态、整齐的动作来震慑敌人；同样，天主仅以那义人的意愿、姿态、举止，就震慑并击溃了我们共同的敌人——魔鬼。因为我认为连魔鬼当时也惊恐退缩。但若有人说：\Quote{为何不容那只手染血，然后立刻使他在被祭杀后复活？}因为天主不会接受如此血腥的祭品；那样的宴席是属于复仇的魔鬼的。但在这里，两件事同时显明了：主人的慈爱与仆人的忠诚。之前，他离开了自己的家乡；但那时，他甚至放弃了本性。因此，他也连本带利地收回了他的本金：这是非常合理的。因为他选择失去父亲的名分，以表明自己是一位忠信的仆人。因此，他不仅成为父亲，也成为司祭；并且因为为了天主的缘故放弃了自己的，所以天主也将祂自己的连同这些一并赐给了他。当敌人图谋恶事时，祂容许事情甚至发展到考验的地步，然后行奇迹，就像在火窑和狮子坑的情况（\EditorialNote{达内尔书第三章和第六章}）；但如果是祂亲自命令，一旦准备就绪，祂便停止命令。那么，我问，在这崇高的事迹中还有什么缺失呢？亚巴郎预知了将要发生的事吗？他讨价还价以求天主的怜悯吗？即使他是先知，先知也并非知道一切。所以，实际的祭杀后来是多余的，也不配天主。如果他认为应该知道天主有能力从死者中复活，那么通过母胎他更奇妙地学到了这一点，或者不如说他在那证明之前就已经学到了，因为他有了信德。\par

那么，不要只钦佩这位义人，也要效法他；当你看见他在如此巨大的喧嚣和浪涛中航行如同风平浪静时，你也要以同样的方式拿起顺从和刚毅的舵。因为，请看，求你，不仅看他筑起祭台和木柴，也记住那童子的声音，并思考有多少如雪暴般的队伍袭击他，使他沮丧，当他听到童子说：\Quote{我的父亲，羊羔在哪里？} 你要思量，那时有多少念头被激动起来，它们不是以铁装备，而是以火焰的标枪，从四面八方刺入并切割他。即使现在，许多人，甚至不是父母，也会心碎，并且会哭泣，如果他们不知道结局；而许多人，我看见，尽管知道结局，仍然哭泣；那么，你认为那生养他、在老年拥有他、只有他一个、如此样的一个儿子、看见他、听见他、并且即将杀他的人，会感到什么呢？话语中有何等的智慧！问题中有何等的温良！那么，此刻谁在工作？魔鬼，要使自然燃烧？天主不许！而是天主，为更多地考验那义人的金质灵魂。因为当约伯的妻子说话时，是魔鬼在工作。因为那种劝告就是如此。但这一位没有说出任何亵渎的话，而是非常虔诚和有思想的；有极大的恩宠笼罩着这些话语，有许多蜂蜜从中滴下，流自一个平静温和的灵魂。这些话语足以软化石心。但它们没有使他转向，不，没有摇动那块金刚石。他也没有说：\Quote{你为什么称他为父亲，他很快就不再是你的父亲，甚至已经失去了那尊荣的头衔？} 那童子为什么问这个问题？不是仅仅由于无礼，也不是出于好奇，而是因对所准备之事感到焦虑。因为他想，如果他的父亲无意使他成为所行之事的伙伴，就不会把仆人留在下面，而只带他一人。也正是出于这个原因，当他们独处时，他才问他，那时无人听到所说的话。这童子的判断力如此之大。你们众人，无论男女，不都对他心生温暖吗？你们每一位不是在心里拥抱并亲吻那孩子，惊叹于他的判断力，并尊敬那虔敬吗？当他被捆绑并放在木柴上时，那虔敬使他不惊慌、不挣扎、也不指责他的父亲是疯子；相反，他甚至被捆绑、举起、放在上面，并沉默地忍受一切，如同一只羔羊，是的，更像万有的共同主。因为他既效法了祂的温良，也保持了预像。因为\BibleQuote{祂被牵到宰杀之地，又像羊在剪毛的人前无声。} \BibleRef{依 53:7} 然而依撒格说话了；因为他的主也说话了。那么如何说是无声呢？这意味着他没有说出任何任性或苛刻的话，而是所有的话都甜美温和，并且言语比沉默更显明了他的温良。因为基督也说：\BibleQuote{我若说得不对，你指证哪里不对；若说得对，你为什么打我？} \BibleRef{若 18:23} 并且比祂保持沉默更显明了祂的温良。正如这一位从祭台上对他父亲说话，同样，祂也从十字架上说话，说：\Quote{父啊，宽恕他们，因为他们不知道他们所做的是什么。} 那么，族长说了什么？\BibleRef{创 22:8} \BibleQuote{天主自会预备作全燔祭的羔羊，我的儿子。} 两者都用了本性的名称：前者称父亲，后者称儿子；两边都激起了艰巨的战争，强大的风暴，然而没有任何翻船：因为宗教战胜了一切。之后，他听到有关天主的话，就不再言语，也没有无礼地好奇。这孩子甚至在青春盛年就有这样的判断力。你看见那君王吗？在多少军队、多少围攻他的战役中，他取得了胜利？因为野蛮人多次进攻耶路撒冷城时，并不像四面围攻他的思想那样令此人恐惧；但他仍然战胜了一切。你愿意也看看司祭吗？例子就在手边。因为当你看见他拿着火和刀，站在祭台之上，之后你还怀疑他的司祭职吗？但如果你愿意也看见祭献，看，这里是一个双重的祭献。因为他献上了一个儿子，他也献上了一只公羊，是的，超越一切，他献上了自己的意志。他用羔羊的血祝圣了他的右手，用他儿子的祭献祝圣了他的灵魂。这样，他藉着他独生子的血，藉着一只羔羊的祭献，被祝圣为司祭；因为司祭们也是藉着献给天主的牺牲的血被祝圣的。你愿意也看见先知吗？经上记载：\BibleQuote{你们的父亲亚巴郎曾欢欣喜乐地期望看见我的日子，他看见了，就欢喜。} \BibleRef{若 8:56}\par

同样，你也在洗礼池中被立为君王、司祭和先知；君王，因为你将一切恶行摔在地上，歼灭了你的罪；司祭，因为你将自己献给天主，牺牲了你的身体，自己也同被牺牲，如经上说：\BibleQuote{如果我们与祂同死，也必与祂同生；} \BibleRef{弟后 2:11} 先知，因为你预知将来之事，并受天主默感，且被印封。因为正如士兵身上有印，圣神也盖在信徒身上。你若背弃，众人便因此认出你。因为犹太人以割损为印记，但我们以圣神为抵押。既知道这一切，并思念我们崇高的地位，让我们活出堪当恩宠的生活，好使我们也能获得未来的天国；愿我们众人藉着我们的主耶稣基督的恩宠和爱人，与祂一起，偕同圣神，归于圣父，享有光荣、权能、尊崇，从今世直到永远，阿们。\par

\SourceRecord{Translated by Talbot W. Chambers.；Nicene and Post-Nicene Fathers, First Series；Vol. 12.；Edited by Philip Schaff.；Buffalo, NY: Christian Literature Publishing Co.,；1889.；<http://www.newadvent.org/fathers/220203.htm>.}

% structure: p0001 source=h1 target=section reason=page-title

\section{《格林多后书》第四篇讲道词}

% structure: p0002 source=h2 target=subsection reason=relative-heading-rank; base=h2

\subsection{格林多后书 1:23}

但我呼求天主为我的灵魂作证：我因为要宽免你们，就没有再到格林多去。\par

蒙福的保禄，你说什么？为了宽免他们，你就没有到格林多去？你确实给我们提出了一个似乎矛盾的说法。因为刚才你曾说，你之所以没有来，是因为你不是按肉性行事，也不是凭自己作主，而是处处服从圣神的引导，并展示你的苦难。但在这里，你说你不来是你自己的行动，并非出于圣神的权威；因为他说：\BibleQuote{我因为要宽免你们，就没有再到格林多去。}那么该怎么说呢？要么，这也是出于圣神，他本人愿意来，但圣神建议他不要这样做，并提出了宽免他们的动机；要么，他指的是另一次来临，并要表明在写前一封书信前，他本有意前来，但因爱而克制自己，免得发现他们仍未改正。也许在第二封书信之后，尽管圣神不再禁止他去，他却因这个原因不由自主地留在远处。这一推测更有可能：起初是圣神禁止他；但后来他自己也相信这样更可取，就留在了远处。\par

请你注意，他再次沿用了他一贯的作风（我将永不停止地观察这一点）：把看似不利于他的事转化为有利于他。因为他们自然会尊重这一点并说：\Quote{你是因为恨我们才不肯到我们这里来}，他却相反地表明，他不来的原因正是他爱他们。\par

\Quote{宽免你们}这话是什么意思？他说，我听说你们中间有人犯了奸淫；因此我不愿前来使你们忧愁；因为如果我来了，就必须调查此事，究办处罚，对许多人施行公道。我认为最好还是离开，给予悔改的机会，而不是与你们同在，追究查办，越发激怒。因为在这封书信的结尾，他已明确宣告说：\BibleQuote{我怕当我再来的时候，我的天主在你们面前再使我谦卑，并且我要为许多从前犯罪，至今没有悔改自己所行的淫乱、污秽和邪荡的人哀哭。}\BibleRef{格后 12:20}因此，他在这里也暗示了这一点；他虽是在为自己辩护，却严厉地责备他们，使他们恐惧；因为他暗示他们应受惩罚，除非他们迅速改正，否则还将遭受某些事。他在书信的结尾也这样说：\BibleQuote{若我再来，必不宽容。}\BibleRef{格后 13:2}只不过那里说得更明白；但这里因为是序言，他并不这样说，而是语气克制；他甚至不满足于此，还加以缓和，进行修正。因为这句话显示了很大的权威（因为一个人宽免那些他有能力惩罚的人），为了减轻效果，掩盖似乎严厉之处，他说：\par

% structure: p0007 source=h2 target=subsection reason=relative-heading-rank; base=h2

\subsection{格林多后书 1:24}

\BibleQuote{这并不是因为我们主管了你们的信德。}\par

就是说，我并不是因为做了你们的主人，才说\Quote{我为了宽免你们而没来}。他也不是说\Quote{你们}，而是说\Quote{你们的信德}，这样既更温和，也更真实。因为对于一个不愿意相信的人，谁能强迫呢？\par

\BibleQuote{而是我们协助你们的喜乐。}\par

因为他说，你们的喜乐就是我们的喜乐，所以我没有来，免得使你们陷入忧愁，增加我的沮丧；我留在远处，好使你们因受到警告而改正，从而喜乐。因为我们所做的一切都是为了你们的喜乐，并为此尽心竭力，因为我们自己也分享这喜乐。\BibleQuote{因为你们在信德上站立的稳。}\par

看，他再次克制地说话。因为他害怕再次责备他们；因为他在前一封书信中严厉地对待了他们，而他们也有所改善。如今他们既已改善，若再受到同样的责备，这很可能使他们倒退。因此，这封书信比前一封温和得多。\par

% structure: p0013 source=h2 target=subsection reason=relative-heading-rank; base=h2

\subsection{格林多后书 1:25}

\BibleQuote{但我为自己决定，不再带着忧愁到你们那里去。}\par

\Quote{再}这个词证明他早已在那里感到忧愁；他看似在为自己辩护，却暗中责备他们。若他们已使他忧愁，又将再次使他忧愁，试想这不满会何等之大？但他没有说\Quote{你们使我很忧愁}，而是改变措辞，却同样暗示了这一点：\Quote{因此我没有来，免得使你们忧愁。}这与我所说的效力相同，却更容易接受。\par

% structure: p0016 source=h2 target=subsection reason=relative-heading-rank; base=h2

\subsection{格林多后书 2:1-2}

2.\BibleQuote{因为如果我叫你们忧愁，那么，除了那个受我忧愁的人以外，谁能叫我喜乐呢？}\par

这是什么结论？确实非常恰当。他说，你们看，我不愿到你们那里去，免得增加你们的忧愁，责备、发怒、厌恶。随后，他见这话仍很严厉，暗示了他们的生活使保禄忧愁，便用以下的话加以缓和：\BibleQuote{因为如果我叫你们忧愁，那么，除了那个受我忧愁的人以外，谁能叫我喜乐呢？}\par

他所说的意思是这样：\Quote{即使我必须忧愁，被迫责备你们，并看见你们忧愁，尽管如此，这件事本身也会使我喜乐。因为这是至大之爱的证明：你们如此尊重我，以至于因我对你们不满而感到伤心。}\par

也请看他的智慧。他们把众门徒所做的事——即听到责备时感到刺痛——当作取悦他的事例；因为他说：\Quote{没有一个人像那听从我话、见我发怒就伤心的人那样使我喜乐。}\par

然而，他自然会说的是：\Quote{因为如果我叫你们忧愁，那么谁能叫你们喜乐呢？}但他并没有这样说，而是把话转回来，温和地对待他们，说：\Quote{即使我叫你们忧愁，你们仍在这事上给了我极大的恩惠，因为你们因我的话而感到伤心。}\par

% structure: p0022 source=h2 target=subsection reason=relative-heading-rank; base=h2

\subsection{格林多后书 2:3}

\BibleQuote{为此，我也曾写了那件事给你们。}\par

什么？为此缘故我没有来，是为宽容你们。他何时写了？在前一封书信中，他说：\Quote{我不愿现在路过时见你们}？\BibleRef{格前 16:7} 我想不是；而是在这封信中，他说：\Quote{免得我再来时，我的天主在你们面前降卑我。}\BibleRef{格后 12:21} 那么，他写道，我就在结尾写了这件事，说：\Quote{免得我再来时，我的天主降卑我，并且我要为许多以前犯罪的人哀哭。}\par

但你为什么写信？\Quote{免得我来时，本应从他们得喜乐，反倒忧愁；我信赖你们众人，知道我的喜乐就是你们众人的喜乐。}因为他说他因他们的忧愁而喜乐，这显得过于傲慢和严厉，于是他再次改变说法，通过下文使之缓和。因为他说，我先前这样写信给你们，为的是不至于忧愁地发现你们仍未悔改；我说这话，\Quote{免得我忧愁}，是出于关心，不是为我自己的利益，而是为你们的。因为我知道，你们若见我喜乐，你们也喜乐；若见我忧愁，你们也忧愁。因此再看他说这些话的连贯性；这样他的话就更容易理解了。他说，我没有来，免得我见你们仍未悔改而使你们忧愁。我这样做，不是求自己的益处，而是求你们的。因为就我自己而言，你们忧愁时，我得到不少快乐，因为看到你们如此在乎我，以至于因我不悦而忧愁难过。\Quote{因为那个使我喜乐的人是谁呢？岂不就是那个被我弄忧愁的人吗？}然而，虽然我如此，但我因求你们的益处，就写了这事给你们，免得我忧愁，在这里也是不求自己的益处，而求你们的；因为我知道，你们若见我忧愁，你们也会忧愁；正如你们见我喜乐，你们也喜乐一样。现在注意他的智慧。他说，我没有来，免得我使你们忧愁；虽然，他说，这使我喜乐。然后，免得他显得以他们的痛苦为乐，他说，就这一点而言，我是喜乐的，因为我使你们有感觉；因为就另一方面而言，我是忧愁的，因为我不得不使那些如此爱我的人忧愁，不仅藉此责备，也因我自己忧愁，从而给你们带来新的忧愁。\par

但请看他是如何说这话以混合赞美的；他说，\Quote{本应从他们得喜乐}，因为这是一个人见证亲情和极大温柔之情的话；如同一个人谈论儿子，他曾赐予他们许多恩惠，并为他们劳苦。如果因此我写信而不来；那么我不来是有重要意义的，不是因为仇恨或厌恶，而是出于极大的爱。\par

3. 接下来，既然他说，那使我忧愁的人使我喜乐；免得他们说：\Quote{那么这就是你所追求的，为要得喜乐，并向众人显示你的权能有多大；}他接着说，\par

% structure: p0028 source=h2 target=subsection reason=relative-heading-rank; base=h2

\subsection{格林多后书 2:4}

\Quote{因为我从极大的患难和心中的痛苦，流了许多眼泪写信给你们，不是要使你们忧愁，而是叫你们知道，我对你们倍加的爱。}\par

还有什么比这个人的心灵更温柔的呢？因为他显示自己不比那些犯罪的人更痛苦，甚至远超过他们。因为他不仅说\Quote{出于患难}，而是\Quote{出于极大的}；不仅说\Quote{流眼泪}，而是说\Quote{流许多眼泪}和\Quote{心中的痛苦}，也就是说，我窒息了，被沮丧憋住了；当我再不能忍受沮丧的阴云时，\Quote{我写信给你们：不是要使你们忧愁，而是叫你们知道}，他说，\Quote{我对你们倍加的爱。}然而自然接下去该说的是：不是要使你们忧愁，而是要使你们改正；（因为他写信正是为此目的。）但他没有这样说，而是（为了更使他的话变得甜美，并赢得他们更大的爱）他以此代之，显明他一切都是出于爱。他不简单地说\Quote{爱}，而是\Quote{我对你们倍加的爱}。因为借此他也想赢得他们，显明他爱他们胜过一切，并将他们视为特选的门徒。因此他说：\Quote{即使我对别人不是宗徒，对你们至少是}\BibleRef{格前 9:2}；又说：\Quote{你们虽有千万导师，却没有许多父亲}\BibleRef{格前 4:15}；又说：\Quote{凭着天主的恩宠，我们在世上为人，尤其对你们更是如此}\BibleRef{格后 1:12}；以及后面说：\Quote{即使我越发热爱你们，却越少得人的爱}；以及这里：\Quote{我对你们倍加的爱}\BibleRef{格后 12:15}。所以，即使我的话充满愤怒，但这愤怒是出于极大的爱和悲伤；并且写这封书信时，我受苦，我疼痛，不仅因为你们犯了罪，也因为我被迫使你们忧愁。而这件事本身也是出于爱。正如一位父亲，他的合法儿子患了坏疽，被迫使用刀和烙铁，他因两方面而痛苦：儿子生病了，以及他被迫对儿子动刀。所以，你们认为我恨你们的标志，实际上是过分爱的标志。如果使你们忧愁是出于爱，那么我为那忧愁而喜乐更是如此。\par

4. 他这样为自己辩护后，（因为他经常为自己辩护，并不以为耻；因为如果天主也这样做，说：\Quote{我的百姓啊，我对你做了什么？}\BibleRef{米 6:3}，更何况保禄呢），我说，他为自己做了这番辩护后，现在将要转到为那犯奸淫的人说情，免得他们因接受相互矛盾的命令而分心，也不致吹毛求疵，因为正是他当时发怒，现在又命令宽恕那人，请看他是如何通过已说的和将说的来预先安排的。他说了什么？\par

% structure: p0032 source=h2 target=subsection reason=relative-heading-rank; base=h2

\subsection{格林多后书 2:5}

\Quote{但如果有人使人忧愁，他所忧愁的不是我。}\par

他先称赞他们，说他们与他同感喜乐和忧伤，然后切入这个人的话题，先说：\Quote{我的喜乐就是你们众人的喜乐。}但如果我的喜乐是你们众人的喜乐，那么现在你们也应该与我同乐，正如那时你们曾与我同忧；因为你们曾忧伤，使我喜悦；现在你们若喜乐（如果我料想你们将感到快乐），你们也会这样做。他没有说\Quote{我的忧伤是你们众人的忧伤}，而是在其余的话中确立了这一点，如今只提出了他最渴望的，即喜乐，说：\Quote{我的喜乐就是你们众人的喜乐。}然后，他也提到了前一件事，说：\par

\Quote{但若有人使人忧愁，他不是使我忧愁，而是部分地（我不说得太重）使你们众人忧愁。}\par

他说：我知道你们与我一同对那犯奸淫的人感到愤怒和憎恶，所发生的事使你们众人部分地忧伤。因此我说\Quote{部分地}，并非像你们所受的伤比我轻，而是为了不压垮那犯奸淫的人。那时他不仅使我忧愁，也同样使你们忧愁，即使为了宽容他，我说了\Quote{部分地}。你看到他是如何立即缓和了他们的怒气，宣称他们也分享了他的愤慨。\par

% structure: p0037 source=h2 target=subsection reason=relative-heading-rank; base=h2

\subsection{格林多后书 2:6}

\Quote{这样的人，受了大多数人的责罚，已经够了。}\par

他不说\Quote{那犯奸淫的人}，而是再次说\Quote{这样的人}，如同在前一封信中一样。但理由不同：那里是出于羞愧，这里是出于怜悯。因此他后来再也没有提及那罪行，因为现在是该宽容的时候了。\par

% structure: p0040 source=h2 target=subsection reason=relative-heading-rank; base=h2

\subsection{格林多后书 2:7}

\Quote{你们倒不如饶恕他，安慰他，免得这样的人被过度的忧愁吞没了。}\par

他吩咐他们不仅要撤销责罚，还要将他恢复原状；因为如果放走受鞭打的人而不医治他，那等于什么也没做。你看他也如何抑制他，免得他因宽恕而变得更糟。尽管他已经认罪悔改，但他表明他得到赦免主要不是由于他的忏悔，而是由于这白白的恩赐。因此他说\Quote{饶恕他，安慰他}，后面的话也使这事更明白。他说：\Quote{并非因为他配受，也不是因为他显示了足够的忏悔，而是因为他软弱，我才请求这事。}为此他又补充说：\Quote{免得这样的人被过度的忧愁吞没了。}这既证明了他的深切悔改，又不让他陷入绝望。\par

但\Quote{吞没}是什么意思？或是像犹大那样行事，或是甚至在生活中变得更糟。他说：因为如果他因无法忍受这长久谴责的痛苦而逃避，也许绝望之下他会去上吊，或者日后犯更大的罪。那么，我们应当预先采取措施，免得伤口变得难以处理；免得因缺乏节制而失去我们已做好的事。\par

他说这话（如我已指出的），既是为使他谦卑，也是为教训他在这次恢复后不要过于漫不经心。因为他说：我接纳他，并非因为他完全洗净了，而是因为担心他做出更深的祸害。由此我们学到：我们必须根据罪的性质以及犯罪者的习性和状态来决定补赎，正如宗徒在那实例中所做的。因为他害怕他软弱，所以说\Quote{免得他被吞没}，如同被野兽、风暴、巨浪吞没一样。\par

% structure: p0045 source=h2 target=subsection reason=relative-heading-rank; base=h2

\subsection{格林多后书 2:8}

\Quote{所以我恳求你们。}\par

他不再发命令，而是恳求，不是以教师的身份，而是以同等的身份；他让他们坐在审判席上，而将自己置于辩护人的地位；因为他的目的已达到，他出于喜悦而毫无拘束地采用了恳求的语气。那么你恳求什么呢？请告诉我。\par

\Quote{向他显出你们的爱德。}\par

就是说，\Quote{坚定你们的爱德}，不仅是与他交往，也不是随便的。这里他又一次证明他们的德行非常伟大；因为那些曾如此友善、如此称赞他甚至自高自大的人，如今竟如此疏远，以致保禄费力要他们坚定对他的爱德。这是门徒的卓越，这是教师的卓越；使他们如此服从缰绳，他如此引导他们的情感。如果现在也是如此，那么犯罪的人就不会愚昧地过犯。因为既不应轻率地相爱，也不应无故地疏远。\par

% structure: p0050 source=h2 target=subsection reason=relative-heading-rank; base=h2

\subsection{格林多后书 2:9}

5.  \Quote{为此，我也曾写信给你们，要知道你们是否在一切事上服从。}不仅在断绝上，也在重新联合上。你看到他又如何把危险带到他们门前。因为当他犯罪时，除非他们将他隔绝，他警告他们的心说：\Quote{一点面酵能使全团发起来}，\BibleRef{格前 5:6}以及其他几件事；同样，在这里他又用不顺服的恐惧面对他们，几乎在说：\Quote{正如那时你们不仅为他，也为自己顾虑；现在你们也当不仅为自己，也为他顾虑，免得你们被看作好争论、缺乏人情，且不能在一切事上服从。}因此他说：\Quote{为此我也曾写信给你们，要知道你们是否在一切事上服从。}\par

因为前者似乎可能出于嫉妒和恶意，但这显示出服从是多么纯洁，以及你是否倾向于慈爱。因为这是正直门徒的试金石：他们不仅在被命令做某些事时服从，而且在被命令做相反的事时也服从。因此他说，\Quote{在一切事上}，表明如果他们不服从，他们羞辱的不是他，而是他们自己，赢得喜好争辩的名声；他这样做也是为了由此驱使他们服从。因此他也说，\Quote{为此我写信给你们}；然而他写信并非为此目的，但他这样说为了赢得他们。因为主要目的是那个人的救恩。但在没有损害的地方，他也满足他们。通过说，\Quote{在一切事上}，他再次称赞他们，回忆并展示他们先前的服从。\par

% structure: p0053 source=h2 target=subsection reason=relative-heading-rank; base=h2

\subsection{\BibleRef{格后 2:10}}

\BibleQuote{你们宽恕谁什么，我也宽恕。}\par

你看见他如何再次将第二部分归于自己，显示他们为起始，他自己跟随。这是软化愤怒、平息争辩之灵的方法。然后，免得他使他们粗心，好像他们是仲裁者，而他们拒绝宽恕；他再次迫使他们这样做，说他自己也宽恕了他。\par

\Quote{因为我所宽恕的——如果我宽恕了什么——是为你们的缘故而宽恕的。} 因为，他说，我正是为你们的缘故做了这事。正如当他命令他们把他开除时，他并没有把宽恕的权力留给他们，说，\Quote{我已经判断要把这样的人交给撒殚，} \BibleRef{格前 5:3} 然后又使他们成为他决定的伙伴，说，\Quote{你们聚集在一起，把他交给}（同上4-5节）（由此确保了两件最重要的事：即判决得以通过；然而并非没有他们的同意，免得在此事上他似乎伤害他们；）他既不是单独宣布判决，免得他们认为他固执己见，而自己则被忽视；也不是将一切留给他们，免得当掌握权力时，他们不合时宜地宽恕而背叛了犯罪者：他在这里也如此行，说：\Quote{我已经宽恕了，我在前书的信中已经判断了。} 然后，免得他们因感到被忽视而受到伤害，他补充说，\Quote{为你们的缘故。} 那么怎样？他为了人的缘故而宽恕吗？不；因为为此他补充说，\Quote{在基督的位格内。}\par

什么是\Quote{在基督的位格内？} 他要么指按照天主的[旨意]，要么指向基督的光荣。\par

% structure: p0058 source=h2 target=subsection reason=relative-heading-rank; base=h2

\subsection{\BibleRef{格后 2:11}}

\BibleQuote{免得撒殚占我们的便宜，因为我们不是不知道他的诡计。}\par

你看见他如何既把权力交给他们，又再次收回，以便通过前者软化他们，通过后者根除他们的私意。但他这样做所预备的还不止这些，而且还表明如果他们不顺从，损害将波及众人，正如他起初所做的那样。因为那时他也说，\Quote{一点面酵能使整个面团发酵。} \BibleRef{格前 5:6} 这里再次，\Quote{免得撒殚占我们的便宜。} 并且贯穿始终，他使这次宽恕成为他自己和他们的共同行动。从起初考虑。\Quote{但如果有人，}他说，\Quote{使人忧愁，他忧愁的不是我，而是部分地（我不说得太重）是你们众人。} 然后再次，\Quote{这样的人，受了\Quote{大多数人的}责罚，已经够了。} 这是他自己的决定和意见。然而他并不止于此决定，而是再次使他们成为伙伴，说，\Quote{所以你们倒要\Quote{宽恕}他\Quote{并安慰}他。} \Quote{为此我恳求你们，对他再表示你们的爱德。} 如此再次将整个事情变成他们的行动后，他转到自己的权威，说，\Quote{为此我也曾写了那封信，为要知道你们是否在一切事上服从到底。} 然后，他又将恩惠归于他们，说，\Quote{你们宽恕谁什么，} 然后是他自己的，\Quote{我也宽恕：}说，\Quote{如果我宽恕了什么，是为你们的缘故。} 然后是他们的和他的，\Quote{因为，}他说，\Quote{如果我宽恕了什么，是为你们的缘故，在基督的位格内宽恕的，} 要么那是为了基督的光荣，要么好像是基督也命令了这事，这对说服他们最为有效。因为此后，他们必不敢拒绝那归于他光荣和他所愿意的事。然后，他又指出如果他们不顺从，共同的损害，当他说，\Quote{免得撒殚占我们的便宜；} 他很好地称之为占便宜。因为他不再取他自己的，而是暴力夺取我们的，因为他已改过。不要告诉我只有这一个成了野兽的猎物，也要考虑这一点：羊群的数目减少了，而现在尤其能够恢复它失去的。\Quote{因为我们不是不知道他的诡计，} 他甚至在虔诚的外表下毁灭。因为不仅通过引人行淫他能毁灭，而且通过相反的事，即悔改后无节制的忧愁。那么，当他除了自己的之外还夺取我们的，当他既通过命令犯罪而毁灭，又当我们命令悔改时暴力夺取，这难道不是占便宜吗？因为他不仅满足于通过犯罪打击，甚至通过悔改他也这样做，除非我们警惕。因此，他也合理地称之为占便宜，当他甚至征服我们自己的武器时。因为通过犯罪夺取是他的本职；然而通过悔改夺取不再是他本职；因为那是我们的武器，不是他的。那么，当他甚至能通过这个夺取时，想想这失败多么可耻，他将如何嘲笑我们，认为我们软弱可怜，如果他要以我们自己的武器征服我们。因为这是极大的嘲笑和最后的耻辱，他竟能通过我们自己的补救措施给我们带来创伤。因此他说，\Quote{因为我们不是不知道他的诡计，} 揭露了他的多变性、他的狡猾、他的恶毒计谋、他的恶意、他在虔诚外表下的伤害能力。\par

6. 既然如此，愿我们也不轻视任何人；即使我们陷入罪恶，也总不绝望；另一方面，又不可事后掉以轻心，而应在犯罪时折磨自己的心灵，而不只是口头发泄。因为我知道许多人，他们确实说自己哀哭罪恶，却毫无实际作为。他们禁食、穿粗衣；但追逐钱财比小贩更急切，比野兽更易发怒，比他人喜爱称赞更喜爱毁谤。这些不是悔改，这些只是悔改的幻影和阴影，而非悔改本身。因此，对这些人的情形，最好说：提防\BibleQuote{免得撒殚占了我们的便宜，因为我们并非不晓得他的诡计}；因为有些人他藉罪恶毁灭，有些人藉悔改毁灭；但这些人是藉另一种方式，即容许他们从悔改中得不到任何果实。因为当他找不到以直接攻击毁灭他们的方法时，他便走了另一条路，加重他们的劳苦，同时夺去果实，并说服他们，仿佛他们已成功完成了当做的一切，因而忽略其余的事。\par

为使我们不至徒然折磨自己，我们向这品性的妇女略进数言；因为这错乱尤属妇女。你们现在所做的禁食、躺卧于地、蒙灰，诚然可嘉；但若不加添其余，则毫无效益。天主已显示了祂如何赦罪。为何离弃那条道路，却为自己另辟蹊径呢？古时尼尼微人犯了罪，他们所做的正是你们现在所做的。然而，让我们看看什么使他们得益。正如在病人身上，医生用许多药方；但有识之人不看重病人试过这个那个，而是看重什么对他有益；我们在此也当如此查考。那么，使那些外邦人得益的是什么呢？他们以禁食敷于伤口，是的，极端禁食，也躺卧于地，披上苦衣，蒙灰，悲哀哭泣；他们也改变了生活。让我们看看哪一样治好了他们。有人会说：我们怎能知道？若到医生那里，若问祂，祂必不向我们隐瞒，甚至急切地揭开。更确切地说，为使无人无知，也无需询问，祂甚至将那治愈他们的药方记载于书。这药方是什么呢？祂说：\BibleQuote{天主看见他们都离开了自己的邪道，就后悔了，不再降灾。}\BibleRef{纳 3:10}祂没有说祂看见了他们的禁食、苦衣和灰。我说这些并非推翻禁食（绝无此意！），而是劝勉你们在禁食的同时，做那比禁食更好的事——戒绝一切邪恶。达味也犯了罪。\BibleRef{撒下 12:17}让我们看看他是如何悔改的。他坐在灰中三天。但这并非为罪，而是为那孩子，当时仍被那苦难所麻木。然而，他以别的方式洗除了那罪：藉谦卑、痛悔的心、良心的刺痛、不再犯同样的罪、常常记念它、甘受一切遭遇、宽恕那些苦待他的人、容忍那些谋害他的人；是的，甚至领先于那些想要这样做的人。例如，当\Person{史米}对他百般辱骂\BibleRef{撒下 16:5}，与他同在的将领极其愤慨时，他说：\BibleQuote{由他骂吧，因为上主吩咐了他。}因为他有痛悔和谦卑的心，这正是洗除他罪的。因为这便是忏悔，这便是悔改。但若我们在禁食时骄傲，则非但无益，反而受损。\par

7. 那么，你也当谦卑你的心，好使你能吸引天主到你身边。\BibleQuote{主亲近于那些心灵痛悔的人。} \BibleRef{咏 32:19} 你难道没有看见那些失宠于华丽宫室中的人吗？他们如何不还口，甚至当卑微的仆人侮辱他们时，也默默忍受，只因他们的过错所招致的耻辱？你也当如此：若有人辱骂你，不要恼怒，反而要哀叹，但不是为那侮辱，而是为那使你陷入耻辱的罪。当你犯罪时，要哀叹，不是因为你要受罚（这不算什么），而是因为你得罪了你的主，一位如此温柔、如此仁慈、如此爱你、渴望你救恩，甚至为你赐下祂独生子的一位。为此哀叹，并要不断如此行：因为这就是告解。不要今天喜乐，明天忧愁，然后又喜乐；而要在哀恸和痛悔中恒久不变。因为，\BibleQuote{哀恸的人是有福的，} 他说，\BibleQuote{因为他们必得安慰。} 意思是那些恒久如此的人。那么，你要恒久如此行，并省察自己，痛心自责，如同一个失去爱子的人哀哭。\BibleQuote{要撕裂你们的心，} 他说，\BibleQuote{而不是你们的衣服。} \BibleRef{岳 2:13} 撕裂的东西不会高举自己；破碎的不能再升起。因此一位说：\BibleQuote{撕裂，} 另一位说：\BibleQuote{天主不会轻视一颗破碎和痛悔的心。} \BibleRef{咏 50:17} 是的，即使你聪明、富有或有权势，也要撕裂你的心。不容它有高傲的念头或被膨胀。因为撕裂的东西不会膨胀，即使有什么使它升起，因着撕裂，它也不能保持膨胀。同样，你也要谦卑自下。想想那个税吏，只凭一句话就成了义，虽然那并非自谦，而是真实的告解。如果这有如此大的力量，那么谦卑就更甚了。宽恕那些得罪你的人，因为这也宽恕罪过。关于前者，祂说：\BibleQuote{我看见他忧愁地行走，我就医治了他的道路。} \BibleRef{依 57:17} 在阿哈布的例子中，这也平息了天主的怒气。\BibleRef{列上 21:29} 关于后者，祂说：\Quote{宽恕，你们就必被宽恕。} 还有另一条道路给我们带来这药方：谴责我们所做的错事；因为，\BibleQuote{先声明你的过犯，好使你成义。} \BibleRef{依 43:26} 并且在困苦中感恩能解除罪过；而施舍更是胜过一切。\par

因此，要数算那些医治你创伤的药方，并坚持不懈地使用所有的：谦卑、告解、忘怨、在困苦中感恩、在施舍和行为上显示仁慈、恒心祈祷。那个寡妇就是这样软化了残忍无情的法官。如果她能软化不义的法官，那么你更能软化那温柔的一位。还有另一条路与此并行：保护受压迫的人；因为，\BibleQuote{要审判孤儿，为寡妇辩护；} 祂说，\BibleQuote{然后来，让我们一同辩论；即使你们的罪如朱红，我也要使它们白如雪。} \BibleRef{依 1:17} 如果藉着这么多通向天堂的道路，和这么多医治创伤的药方，我们甚至在圣洗之后依然停留在原处，那么我们能有什么借口呢？我们不仅要如此继续，而且那些从未跌倒的人应当持守他们原有的可爱；是的，他们更要不断培养它（因为这些善工，在找不到罪的地方，使美丽更加增大）；而我们这些在许多事上犯了过错的人，为了纠正我们的罪，应使用上述的方法，好使我们在基督的审判台前怀有极大的信心。愿我们众人借着我们的主耶稣基督的恩宠和祂对人的爱，得以到达那里；愿光荣、权能和尊崇归于祂，与圣父及圣神，从今直到永远，世世无尽。阿们。\par

\SourceRecord{Translated by Talbot W. Chambers.；Nicene and Post-Nicene Fathers, First Series；Vol. 12.；Edited by Philip Schaff.；Buffalo, NY: Christian Literature Publishing Co.,；1889.；<http://www.newadvent.org/fathers/220204.htm>.}

% structure: p0001 source=h1 target=section reason=page-title

\section{\WorkTitle{格林多后书第五篇讲道}}

% structure: p0002 source=h2 target=subsection reason=relative-heading-rank; base=h2

\subsection{\BibleRef{格后 2:12-13}}

\BibleQuote{当我为基督的福音来到\Place{特洛阿}时，主为我开了门，但因我没有找到我的弟兄\Person{弟铎}，我的心灵不得安宁。}\par

这些话一方面似乎不配\Person{保禄}，若因弟兄不在而放弃如此大的拯救机会；另一方面，又似乎与上下文脱节。那么你们愿意我们先证明它们与上下文相连，还是证明他没有说任何不配自己的话？依我看，后者，因为这样另一个点也会更容易更清楚。\par

那么这些话如何与前面的话相连呢？让我们回想前面的话，这样我们就会明白。前面的话是什么？就是他在开头说的：\BibleQuote{弟兄们，我不愿意你们不知道，我们在亚细亚所受的苦难，我们被压得太重，力不能胜。}\BibleRef{格后 1:8}。在显示了他获救的方式，并插入中间内容后，他必然再次教导他们，他在另一方面也曾受苦。如何受苦？在什么方面？就是没有找到\Person{弟铎}。\EditorialNote{参格后7:6; 8:6,16,22,23; 12:18}。的确，忍受试炼已是可怕并足以使灵魂沮丧；但当没有人安慰和帮助承担重担时，风暴就变得更大了。\Person{弟铎}就是他在后面提到从他们那里来到他那里的人，他对他有许多赞扬，并说他曾打发他去。为了表明他在这一点上也是为他们而受苦，他说了这些话。\par

那么这些话与前面相连，从这一切可以清楚看出。我也要努力证明它们并不配\Person{保禄}。因为他并没有说\Person{弟铎}的缺席阻碍了那些将要归信的人的得救，也没有说他因此忽略了那些相信的人，而是说\BibleQuote{我不得安宁}，即\BibleQuote{我受苦，我因弟兄不在而忧伤}，表明弟兄缺席是何等大事；因此他离开了那里。但\BibleQuote{当我为福音来到\Place{特洛阿}}是什么意思？他不仅说\BibleQuote{我到了}，而是\BibleQuote{为了宣讲}。然而，尽管我为此而来并发现有很多工作可做（因为\BibleQuote{主为我开了门}），他说\BibleQuote{我不得安宁}，并不是为此阻碍了工作。那么他怎么说呢？\par

\BibleRef{格后 2:13} \BibleQuote{我辞别了他们，就去了那里？}\par

也就是说，\BibleQuote{我没有多停留，因受困和忧伤。}也许工作甚至因他的缺席而受阻。这对他们也是一种不小的安慰。因为如果在那里开了门，并且他为此而来，却因没有找到弟兄而迅速离开，那么你们更应该体谅那些牵动我们到处奔走的事务的必然性，它们不容我们随意旅行或在我们愿意停留的人中久留。因此他在这里又将自己的行程归于天主，如同上面归于圣神一样，说：\par

% structure: p0009 source=h2 target=subsection reason=relative-heading-rank; base=h2

\subsection{\BibleRef{格后 2:14}}

\BibleQuote{感谢天主，祂常常使我们在基督内凯旋，并藉着我们在各处显扬认识祂的香气。}\par

为使他不会显得在忧伤中哀叹这些事，他向天主献上感谢。他所说的是：\BibleQuote{处处是困难，处处是窘迫。我来到\Place{亚细亚}，被压得太重。我来到\Place{特洛阿}，没有找到弟兄。我没有到你们那里，这也给我带来不小的沮丧，甚至极大的沮丧，因为你们中许多人犯了罪，也因为我因此没有见到你们。因为他说：\InnerQuote{为了宽免你们，我还没有到\Place{格林多}去。}}因此，为使他这样说时不显得在抱怨，他加上说：\BibleQuote{我们不仅在这些苦难中不忧伤，反而欢喜；而且，更大的事是，不仅是因将来的赏报，也因现时的赏报。因为即使在这里，我们也因这些事而荣耀和显赫。因此，我们远非哀叹，反而称这为凯旋，并为所发生的事而夸耀。}为此他也说：\BibleQuote{感谢天主，祂常常使我们在基督内凯旋}，即\BibleQuote{祂使我们在众人中出名。因为那似乎是羞辱的事，即处处被迫害，在我们看来却是极大的光荣。}因此他没有说\BibleQuote{祂使我们被众人看见}，而是\BibleQuote{祂使我们凯旋}，表明这些迫害在世界各地为魔鬼树立了一系列的凯旋。然后，在提到凯旋的作者的同时，也提到凯旋的主题，他以此激励听众。\BibleQuote{我们不仅被天主所凯旋，而且\InnerQuote{在基督内}}，即因基督和福音的缘故。\BibleQuote{既然必须凯旋，那么我们必须带着凯旋的旗帜被众人看见，因为我们承载着祂。为此我们变得被看见和显眼。}\par

2. \BibleRef{格后 2:14} \BibleQuote{并藉着我们在各处显扬认识祂的香气。}\par

他在上文说：\Quote{那常常率领我们获胜的。}这里他说\Quote{在各处}，表明一切地方、一切时候都充满了宗徒们的劳苦。他又用了另一个比喻，就是\OriginalTerm{馨香}{sweet savor}的比喻。因为\Quote{我们好像那些携带香膏的人一样}，他说，\Quote{向众人显明}；他把这知识称为极贵重的香膏。不仅如此，他说的不是\Quote{知识}，而是\Quote{\OriginalTerm{知识之馨香}{savor of the knowledge}}；因为当前的知识正是这样，不太清晰，也非显露。所以他在前书也说：\BibleQuote{因为现在我们对着镜子观看，模糊不清。}\BibleRef{格前 13:12}在这里他把这种情形称为\Quote{馨香}。那闻到馨香的人知道有香膏藏在某处，但还不知道它的本质是什么，除非他先前见过。\Quote{我们也一样。我们知道天主存在，但还不知道祂的本体是什么。我们就像是一尊皇室的香炉，无论走到哪里，都发出天上的香膏和属神的馨香。}他说这话，一方面是为彰显宣讲的大能，因为即使面对针对他们的阴谋，他们比那些迫害他们并使全世界都知道他们的凯旋和馨香的人更加闪耀；另一方面是为劝勉他们在苦难和试炼中勇敢地承受一切，因为甚至在获得赏报之前，他们就已经得到了这种无法言说的光荣。\par

% structure: p0014 source=h2 target=subsection reason=relative-heading-rank; base=h2

\subsection{格林多后书 2:15}

\BibleQuote{因为我们就是基督在天主前所发出的馨香，在得救的人身上，也在丧亡的人身上。}\par

他说，无论是得救还是丧亡，福音始终保持其固有的德能。正如光虽然使虚弱者目眩，但仍是光，尽管造成失明；正如蜜虽然对患病者是苦的，但本质上仍是甜的。同样，福音也是馨香的，即使有些人不信而丧亡。因为并非福音，而是他们自己的乖戾导致丧亡。而正是藉此，福音的\OriginalTerm{馨香}{sweet savor}最大程度地显明出来，使败坏和邪恶的人丧亡。因此，不仅藉善人的得救，也藉恶人的丧亡，福音的卓越得以显明。因为太阳正因其极其明亮而伤及虚弱者的眼目；救主也是\BibleQuote{使许多人跌倒和复兴的}\BibleRef{路 2:34}，但即使一万人跌倒，祂仍是救主；祂的来临为不信者带来更重的惩罚，但祂仍充满医治。因此他说：\Quote{我们在天主前是馨香}；意思是\Quote{即使有些人丧亡，我们仍然是原来的我们。}他说的不仅是\Quote{馨香}，而且是\Quote{在天主前}。当我们是在天主前是馨香，而天主如此定断时，谁还能反驳呢？\par

那个\Quote{\OriginalTerm{基督的馨香}{sweet savor of Christ}}的说法，按我看有双重解释：要么是指在死亡中他们将自己献为祭品；要么是指他们是基督死亡的馨香，如同有人说，这乳香是这祭品的\OriginalTerm{牺牲之馨香}{sweet savor of this victim}。那么，馨香这个说法要么指此，要么如我最初所说，他们每日为基督而牺牲。\par

3. 你看到他把试炼提升到何等高度了吗？他称之为凯旋、馨香和向天主献上的祭品。然后，他说\Quote{我们乃是馨香，甚至在丧亡的人中}，免得你以为这些人也被悦纳，他补充道：\par

% structure: p0019 source=h2 target=subsection reason=relative-heading-rank; base=h2

\subsection{格林多后书 2:16}

\BibleQuote{对一班人，是由死亡归于死亡的馨香；对另一班人，是由生命归于生命的馨香。}\par

因为这\OriginalTerm{馨香}{sweet savor}，有些人领受了而得救，有些人领受了而丧亡。所以，如果有人丧亡，错误在于他自己：因为香膏据说能呛死猪，光（如我刚才所说）能使目眩者失明。优秀之物的本性正是如此：它们不仅纠正相近之物，也毁灭相反之物，并且以此方式最大程度地显示其能力。因为火也是这样，不仅当它发光、炼净金子时，甚至当它烧毁荆棘时，它大大彰显其本有的能力，显明它是火；基督也是这样，当祂\BibleQuote{以口中的气息消灭那恶者，以祂来临的显现使之化为乌有}时，也显示了祂的威严。\BibleRef{得后 2:8}\par

\BibleQuote{谁能担得起这些事呢？}\par

他既说出了伟大的事，说\Quote{我们是基督的祭品和馨香，并随处被率领得胜}，便又节制地引用一切归功于天主。因此他说：\BibleQuote{谁能担得起这些事呢？}他说：\Quote{因为一切都是基督的，没有一样是我们的。}你看见他的言语与假宗徒们何等相反吗？因为他们夸耀的似乎是他们自己为这信息贡献了什么；他则相反，他说他之所以夸耀，是因为他说没有一样是自己的。\Quote{因为我们夸耀的正是这事：良心的见证，就是我们在世上行事，不是凭着\OriginalTerm{属人的智慧}{fleshly wisdom}，而是凭着\OriginalTerm{天主的恩宠}{grace of God}。}他们认为要获取的光荣，我指的是外在的智慧，他却以除去它为荣。因此他在这里说：\BibleQuote{谁能担得起这些事呢？}但如果没有人能担得起，那么所成就的就是出于恩宠。\par

% structure: p0024 source=h2 target=subsection reason=relative-heading-rank; base=h2

\subsection{格林多后书 2:17}

\BibleQuote{因为我们不像那许多人，他们\OriginalTerm{败坏天主的圣言}{corrupt the word of God}。}\par

\Quote{因為即使我們使用響亮的言詞，我們卻宣稱自己所成就的沒有一樣是出於自己，而全是基督的。因為我們不會效法那些假宗徒；他們說大部分是出於自己。}因為這就是\Quote{攙雜}，當人攙雜了酒；當人把本應白白施予的東西拿來賣錢。因為在我看來，他這裡既是在金錢方面嘲諷他們，又再次暗示我所說的那件事，即他們將自己的東西與天主的東西混雜；這就是依撒意亞所指責的，當他說：\Quote{你的調酒師用水攙酒。} \BibleRef{依 1:22} 因為即使這話是關於酒的，但用來解釋道理也不會錯。\Quote{但是我們，}他說，\Quote{卻不這樣：我們如何被託付，就如何提供給你們，傾注出未經稀釋的聖言。}因此他加上：\Quote{但如同出於真誠，如同出於天主，在天主面前，我們在基督內說話。}\par

\Quote{我們不}，他說，\Quote{欺哄你們而這樣宣講，如同將禮物賜給你們，或如同從自己帶來並攙雜些什麼，\Quote{而是出於天主}；}這就是說，我們不說我們貢獻了任何自己的東西，而是說天主賜予了一切。因為\Quote{出於天主}的意思就是：不以任何事物為榮，好像我們自己擁有，而是將一切歸於祂。\Quote{我們在基督內說話。}\par

不是靠我們自己的智慧，而是受自祂而來的能力教導。那些自誇的人不這樣說話，而是如同從自己帶來什麼。因此他在別處也嘲笑他們，說：\Quote{你有什麼不是領受的呢？既然是領受的，為什麼你還誇耀，好像不是領受的呢？} \BibleRef{格前 4:7} 這是最高的德行：將一切歸於天主，不以為任何東西是自己的，做事不顧人的看法，而只隨從天主的旨意。因為祂是那要求交賬的。但如今這秩序顛倒了：對那將坐在審判座上並要求交賬的祂，我們沒有極大的敬畏，卻在那些站著與我們一同受審的人面前戰慄。\par

4. 那麼，這毛病從何而來？它何時在我們的靈魂中爆發？來自於不持續默想那世界的事，卻被眼前的事物所束縛。因此我們既容易陷入惡行，即使做了什麼好事，也是為了炫耀，從而使我們也遭受損失。例如，有人經常以放縱的目光看著一個人，沒被那女子或與她同行的人看見，卻沒有逃過那從不睡覺的眼睛。因為甚至在犯罪之前，它已看見那放縱的靈魂，內心的瘋狂，以及如風暴和波濤般翻騰的思想；因為祂不需要見證人和證據，祂知道一切。不要看你的同僕：因為雖然人讚揚，若天主不接受，也無益；雖然人責備，若天主不責備，也無害。唉！不要這樣激怒你的審判者：對你的同僕十分看重，而當祂發怒時，卻不敬畏戰慄。那麼，讓我們輕視來自人的讚揚。我們要到幾時才心志卑微、匍匐在地？幾時，當天主將我們高舉到天上，我們卻努力被拖在地上？若瑟的兄弟們，若他們將天主的敬畏放在眼前，如同人應該有的，就不會在偏僻的地方捉住他們的兄弟並殺死他。\EditorialNote{創世紀37} 加音也曾經，若他畏懼那判決如同應該畏懼的，就不會說：\Quote{來，我們到田間去。} \BibleRef{創 4:8} 因為，可悲又可憐的人啊！你為何將他從生養他的人那裡帶走，領他到偏僻的地方？難道天主在田野裡看不見這大膽的行為嗎？你還沒有從你父親的遭遇學到祂知道一切，臨在於一切所發生的事嗎？為什麼當他否認時，天主沒有對他說：\Quote{你躲避我嗎？我無所不在，知道隱密的事。}因為那時他還不正確地理解這些高深的真理。但祂說了什麼？\Quote{你弟弟的血的聲音從地上向我哀號。} 不是說血有聲音，而是如同我們說事情明顯清楚時，\Quote{事情本身說話。}\par

因此，確實應該將天主的判決放在眼前，那麼所有的恐懼都熄滅了。同樣，在祈禱中我們也能保持警醒，如果我們記住我們與誰交談，如果我們反省我們是在獻祭，手中握有刀、火和柴；如果我們在思想上敞開天堂的門，如果我們將自己轉移到那裡，拿起聖神的劍，插入犧牲的喉嚨：以警醒為祭，以眼淚為奠祭獻給祂。因為這犧牲的血就是如此。這屠殺使那祭壇變紅。那時不要讓任何世俗的思想佔據你的靈魂。你要想想亞巴郎獻祭時，也不讓妻子、僕人或任何其他人同在。那麼，你也不要容許任何奴役性和卑劣的情慾在你那裡，而要單獨上到他所上的那座山，那裡不許第二個人上去。如果有任何這樣的思想試圖與你一同上去，你要以權威命令它們，說：\Quote{你們在這裡等候，我和孩子要去朝拜，然後回到你們這裡來。} \BibleRef{創 22:5} 將驢和僕人留在下面，以及一切沒有理性和知覺的東西，你帶著有理性的東西上去，如同他帶著依撒格。並像他一樣建造你的祭壇，好像沒有人性，而是超越了本性。因為他若不是超越了本性，就不會殺死自己的孩子。那時不要讓任何事情攪擾你，而要超升到諸天之上。痛切哀號，獻上認罪，\BibleRef{依 43:26}，獻上痛悔的心。這些犧牲不會化為灰燼，不會消散成煙，也不會融化在空氣中；它們不需要柴和火，只需要一顆深深刺痛的心。這就是柴，這就是火，燃燒卻不燒毀它們。因為熱心祈禱的人被焚燒，卻不被燒毀；反而像經試煉的金子一樣，變得更明亮。\par

5. 此外，当谨慎注意一事：祈祷时，切勿说出任何激怒你主的话语；也不要走近[祈祷]反对仇敌。因为如果拥有仇敌是一种羞辱，那么想一想祈祷反对他们是多大的恶。因为你需要为自己辩护，说明你为何有仇敌；但你反而控告他们。你又能获得什么宽恕呢？当你既咒骂，又在你自身需要极大怜悯之时。因为你走近是为了恳求赦免你自己的罪：那么就不要提别人的罪，免得你唤起对自己罪的记忆。因为如果你说：\Quote{击打我的仇敌}，你就封住了自己的口，断绝了舌头的胆量；首先，确实因为你一开始就激怒了审判者；其次，因为你所求的与祈祷的性质相悖。因为如果你来是为求罪的赦免，你怎么谈论惩罚呢？实在需要做相反的事，并为他们祈祷，好使我们也能放胆为自己求这赦免。但现在你以自己的判断抢先于审判者的判决，要求他惩罚那些犯罪的人：因为这剥夺了所有宽恕。但如果你为他们祈祷，即使你为自己罪的事什么也不说，你也成就了一切。试想法律中有多少祭献：赞美祭、感恩祭、和平祭、洁净祭，以及无数其他祭献，却没有一个是反对仇敌的，而是全部为着自身的罪或自身的得胜。难道你来到另一位天主面前？你来到那位曾说过\BibleQuote{为你们的仇敌祈祷。}\BibleRef{路 6:27}的天主面前。你怎么还能呼喊反对他们？你怎么恳求天主违背他自己的法律？这不是恳求者的样子。没有人恳求别人的毁灭，而是自己的安全。那么，你为什么披着恳求者的外衣，却说出控告者的话语？然而，当我们为自己祈祷时，我们抓挠自己、打哈欠，陷入万千思绪；但当我们反对仇敌时，我们却如此清醒。因为魔鬼知道我们正把剑刺向自己，他那时就不来分心或阻止我们，好给我们造成更大的伤害。但是，有人说：\Quote{我受了冤屈，遭受痛苦。} 那么，为什么不祈祷反对那加害我们最甚的魔鬼呢？这也曾命令你们说：\Quote{救我们免于凶恶。} 他是你不可和好的仇敌，但人，无论他做什么，都是朋友和兄弟。那么，让我们都对他发怒；让我们恳求天主反对他，说：\BibleQuote{将撒殚踏在我们脚下。}\BibleRef{罗 16:20}，因为正是他生出了我们所有的仇敌。但如果你祈祷反对仇敌，你的祈祷正合乎他的心意，就如你为仇敌祈祷，就是反对他。那么，为什么放过那真正是你的仇敌，却撕裂你自己的肢体，在此事上比野兽更残忍？但有人说：\Quote{他侮辱了我，抢了我的钱；} 那么，谁更需要悲伤？是受伤害的人，还是施加伤害的人？显然是施加伤害的人，因为他在得钱的同时，将自己抛出天主的恩宠，失去的比得到的更多：所以他是受害的一方。那么，实在不需要祈祷反对他，而要为他祈祷，求天主怜悯他。请看那三位少年受了多少苦，尽管他们并无过错。他们失去了祖国、自由，被掳为奴；被带到外邦和野蛮之地时，甚至因那梦而无缘无故面临被杀的危险。\BibleRef{达 2:13} 结果如何？当他们与达尼尔一同进去时，他们祈祷什么？他们说了什么？打倒拿步高，摘下他的王冠，将他从宝座上抛下？没有这样的事；他们却求了\BibleQuote{天主的怜悯。}\BibleRef{达 2:18}。当他们在火窑中时，也是如此。但你们却不然：当你们所受的苦远比他们轻，且时常是应得的时，你们却不停地发出万千咒骂。有人说：\Quote{打倒我的仇敌，如同你淹没法郎的战车；} 另一个说：\Quote{毁坏他的肉体；} 又一个说：\Quote{报应在他子女身上。} 你们认不出这些话吗？那么，你们的笑声从何而来？当这些话不带激情地说出时，你们难道看不出多么可笑吗？同样，一切罪在除去行凶者的心境后，才显露出其卑劣。如果你提醒一个发怒的人他在盛怒中说的话，他会羞愧难当，鄙视自己，宁愿受千般惩罚也不愿那些话是他说的。如果你在拥抱过后，将不洁者带到与他一同犯罪的女人面前，他也会厌恶地避开她。因此，你们现在因为不受情绪影响而发笑。因为他们和醉酒老妇的话语是值得嘲笑的，且出自女性般狭隘的灵魂。然而若瑟，虽曾被卖为奴，又身陷囹圄，那时也未曾对使他受苦的人说过一句苦毒的话。但他怎么说？\BibleQuote{我实在是从希伯来人地方被拐来的；}\BibleRef{创 40:15}，却没有说是被谁拐的。因为他为弟兄的邪恶感到的羞耻，比那些行恶的人更深。这也应该是我们的态度：为那亏负我们的人感到悲伤，比他们自己更甚。因为伤害归到他们身上。正如那些踢钉子却以此自傲的人，因这疯狂而成为可怜哀叹的对象；同样，那些伤害无恶之人的人，因伤害了自己的灵魂，适宜于许多哀恸和悲叹，而非咒骂。因为没有比咒骂的灵魂更污秽的，也没有比供奉这样祭献的舌头更不洁的。你是人，不要吐出毒蛇的毒液。你是人，不要变成野兽。因为你的口被造，不是为了咬人，而是为了医治别人的伤口。\Quote{记住我给你的命令，}天主说，\Quote{要宽恕和赦免。但你还恳求我成为颠覆我自己诫命的同谋，吞吃你的兄弟，染红你的舌头，如同疯子用牙齿咬自己的肢体。} 你认为，魔鬼听见这样的祈祷，如何欢喜并嗤笑？而天主，当你恳求这样的事时，如何被激怒，转离你并憎恶你？还有比这更危险的吗？因为如果有仇敌的人都不应接近奥迹，那么那不仅拥有仇敌，还祈祷反对他们的人，岂不连外院也被排除？那么，想到这些，并思虑祭献的主题——祂曾为仇敌而牺牲，让我们不要有仇敌；如果我们有，就让我们为他祈祷，好使我们自己也获得所犯罪过的赦免，得以在基督的审判台前坦然无惧地站立；愿光荣归于祂，直到永远。阿们。\par

\SourceRecord{Translated by Talbot W. Chambers.；Nicene and Post-Nicene Fathers, First Series；Vol. 12.；Edited by Philip Schaff.；Buffalo, NY: Christian Literature Publishing Co.,；1889.；<http://www.newadvent.org/fathers/220205.htm>.}

% structure: p0001 source=h1 target=section reason=page-title

\section{\WorkTitle{格林多后书第六篇讲道}}

% structure: p0002 source=h2 target=subsection reason=relative-heading-rank; base=h2

\subsection{格林多后书 3:1}

\BibleQuote{我们又开始推荐自己了么？或者，我们岂不是像某些人那样，需要向你们或从你们那里要推荐信么？}\par

他预料到并预先提出了别人会用来反对他的异议：\Quote{你夸耀你自己。}尽管他先前已经用\BibleQuote{谁足够胜任这些事呢？}和\BibleQuote{我们本着真诚……说话}这些表达施加了强烈的纠正。\BibleRef{格后 2:16}然而他并不满足于此，因为这就是他的性格。他极力避免说任何自夸的话，甚至到了过分的程度。请注意，我请你们看这个例子，也显出他智慧的丰盛。一件令人痛苦的事，我说的是磨难，他竟如此高举并使之光明灿烂，以至于因他所说的，这个异议就起来反对他。他在结尾处也是如此。在列举了无数危险、凌辱、困境、急需以及诸如此类的一切之后，他加上说：\BibleQuote{我们不是推荐自己，而是说话给你们一个夸耀的缘由。}\BibleRef{格后 5:12}他在那里再次激烈地表达，且带着更多的鼓励。因为这里的话是出于爱：\BibleQuote{我们需要像某些人那样，有推荐信么？}但那里的话却充满了一种骄傲，这是必要且恰当的，我说是骄傲和愤怒。\BibleQuote{因为我们不是再次推荐自己，}他说，\BibleQuote{而是说话给你们一个夸耀的缘由}\BibleRef{格后 5:12}；以及\BibleQuote{你们又以为我们是向你们分诉么？其实我们在\Person{基督}里是在天主面前说话。因为我怕我来的时候，找不到你们如我所愿意的，而我自己也被你们发现为你们所不愿意的。}\BibleRef{格后 12:19,20}为了避免任何奉承的嫌疑，好像他渴望从他们那里得到尊荣，他这样说：\BibleQuote{我怕我来的时候，找不到你们如我所愿意的，而我自己也被你们发现为你们所不愿意的。}这却是在许多指责之后；但在开始时他说话不像这样，而是更温和。他说什么呢？他讲到了他的试炼和危险，以及他处处被天主在\Person{基督}里凯旋游行，全世界都知道这些胜利。既然他说了关于自己的大事，他就为自己提出这个异议：\BibleQuote{我们又开始推荐自己了么？}他所说的是：也许有人会反对说：\Quote{\Person{保禄}啊，这是什么事？你自己说这些，夸奖自己么？}为消除这猜疑，他说：我们并不渴望这个，就是说，夸耀并高举自己；是的，我们如此不需要向你们有推荐信，以至于你们就是我们的信。\BibleQuote{因为，}他说，\par

% structure: p0005 source=h2 target=subsection reason=relative-heading-rank; base=h2

\subsection{格林多后书 3:2}

\BibleQuote{你们就是我们的信。}\par

\BibleQuote{你们就是我们的信。}\BibleQuote{你们就是？}是什么意思？\Quote{我们需要向别人推荐自己么？我们会在他们面前呈上你们，代替一封推荐信。}他在前书也这样说。\BibleQuote{因为你们就是我宗徒职务的印证。}\BibleRef{格前 9:2}但他在这里没有这样说，而是以讥讽的方式使他的问题\BibleQuote{我们需要推荐信么？}更尖锐。并影射假宗徒，他加上了\BibleQuote{像某些人那样，向你们有推荐信，或从你们那里有推荐信}给别人。然后因为他所说的很严厉，他用加上\BibleQuote{你们就是我们的信，写在我们的心上，被众人所认识}来缓和。\par

% structure: p0008 source=h2 target=subsection reason=relative-heading-rank; base=h2

\subsection{格林多后书 3:3}

\BibleQuote{你们显明是\Person{基督}的信。}\par

这里他不仅见证他们的爱，也见证他们的善行：因为他们以自身的德行向众人显示他们老师的崇高价值，这就是\Quote{你们就是我们的信}的意思。\par

书信所能做的，就是推荐我们并为我们赢得尊重，而你们无论被看见或被听见，都做到了这一点；因为门徒的德行通常比任何书信更能装饰和推荐老师。\par

\BibleQuote{写在我们的心上。}\par

就是说，尽人皆知；我们处处带着你们，心中有你们。好像他说：你们就是我们对别人的推荐，因为我们不断把你们放在心上，并向众人宣扬你们的善行。因为既然你们自己就是我们对别人的推荐，我们就不需要从你们那里得到推荐信；此外，因为我们极其爱你们，我们就无需向你们推荐自己。因为对于陌生人，需要书信，但你们在我们的心目中。然而他没有仅仅说\Quote{你们是}，而是\BibleQuote{写在心上}，就是说，你们不能从那里滑脱。因为正如通过阅读书信，同样通过感知我们的心，众人都知道我们对你们的爱。那么，如果书信的目的是证明\Quote{某某是我的朋友，让他与你们自由交往}，你们的爱就足以保证这一切。因为如果我们到你们那里去，我们不需要别人推荐我们，因为你们的爱已预先做到了；如果我们到别人那里去，我们同样不需要书信，因为同样的爱又足以代替书信，因为我们心中带着这封信。\par

2. 然后，为更升高他们，他甚至称他们为\Person{基督}的书信，说：\par

\BibleRef{格后 3:3}。\BibleQuote{你们显明是\Person{基督}的信。}\par

说了这话之后，他由此取得立场和机会来讨论法律。他在这里称他们为他的信还有另一个目的。因为前面为推荐他，他称他们为一封信；但这里却是\Person{基督}的信，因为他们有写在他们心中的天主的法律。因为天主愿意向众人并向你们启示的事，都写在你们心中。但却是我们预备了你们接受这书写。就如\Person{梅瑟}凿了石版，我们则预备了你们的灵魂。因此他说：\par

\BibleQuote{由我们服务的。}\par

然而在这件事上他们是平等的；因为前者是由天主写上，而这些是由圣神写上。那么区别在哪里呢？\par

\BibleQuote{不是用墨写的，而是用永生天主的圣神写的；不是在石版上，而是在肉心的版上。}\par

正如圣神与墨水，石版与肉心之间的差别那样大，这些与那些之间的差别也是如此；因此，他们这些服役的人与那服事他们的人之间也是如此。然而，因为他所说的是大事，他便立即抑制自己，说道：\par

% structure: p0021 source=h2 target=subsection reason=relative-heading-rank; base=h2

\subsection{\BibleRef{格后 3:4}}

\BibleQuote{我们藉着基督对天主有这样的信赖。}\par

他又将一切归功于天主：因为他说，基督是这些事在我们身上的创始者。\par

% structure: p0024 source=h2 target=subsection reason=relative-heading-rank; base=h2

\subsection{\BibleRef{格后 3:5-6}}

\BibleQuote{这不是说我们凭自己能承担什么事，好像出于自己似的。}\par

请看，又一个修正。因为他拥有这种美德，我指的是谦卑，达到了非凡的完美。因此，每当他谈及自己的任何大事时，他都竭尽全力、千方百计地极力柔化自己所说的话。他在这里也是如此，说道：\BibleQuote{这不是说我们凭自己能承担什么事，好像出于自己似的。}意思是说，我并没有说\InnerQuote{我们有信赖}，仿佛一部分是我们的，一部分是天主的；而是将全部归于祂，归功于祂。\par

\BibleQuote{因为我们所能承担的，乃是出于天主；祂也使我们能够作成新约的仆役。}\par

\Quote{使我们能够}是什么意思？就是使我们有能力且合宜。将这比从前更大得多的版和信件带给世界，并非小事。因此他又加上：\par

\BibleQuote{不是凭着字句，乃是凭着圣神。}再看另一个区别。那么怎样？那法律不是属神的吗？那么他怎么说：\BibleQuote{我们知道法律是属神的？}\BibleRef{罗 7:14} 它确实是属神的，但它并没有赐予圣神。因为梅瑟带来的不是圣神，而是字句；但我们受托要赐予圣神。因此，在进一步完成这一对比后，他说：\par

\BibleQuote{因为字句叫人死，圣神却叫人活。}\par

然而他说这些话并非绝对如此，而是针对那些以犹太教事物自夸的人。他在这里所说的\Quote{字句}是指惩罚犯法者的法律；而\Quote{圣神}是指那通过圣洗使因罪而死的人得生命的恩宠。因为，在提及了版的性质所产生的差异后，他并未过多停留，而是迅速越过，更加着力于这一点，这使他能够从利益和容易的角度抓住听众；因为他说，这并不费力，而且所提供的恩赐更大。因为，如果当论及基督时，他特别提出那些属于祂慈爱的事，多于我们的功德，并且这些事是相互关联的，那么当他论及盟约时，更必须如此。那么\Quote{字句叫人死}是什么意思？他提到了石版和肉心：到此他似乎没有提到很大的区别。他又加上前一个盟约是用文字或墨水写的，而这一个是用圣神写的。但这仍未完全激发他们，他最后说了一句足以使他们展翅高飞的话：一个\Quote{叫人死}，另一个\Quote{叫人活}。这是什么意思？在法律中，有罪的人受惩罚；在这里，有罪的人来受洗，成为义人，成为义人后就活着，脱离了罪的死亡。法律若抓住一个杀人犯，就处死他；福音若抓住一个杀人犯，就光照他，赐他生命。我为什么举杀人犯为例？法律抓住一个在安息日捡柴的人，就用石头打死了他。\BibleRef{户 15:32}这就是\Quote{字句叫人死}的意思。福音抓住了成千上万的杀人犯和强盗，用圣洗拯救他们脱离以前的恶习。这就是\Quote{圣神叫人活}的意思。前者使被抓住的人从活变死，后者使被定罪的人从死复活。因为他说：\BibleQuote{凡劳苦和负重担的，你们都到我这里来，}\BibleRef{玛 11:28} 他没有说\InnerQuote{我要惩罚你们}，而是说\Quote{我要使你们安息。}因为在圣洗中，罪被埋葬，过去的事被涂抹，人得以活着，整个恩宠像写在版上一样写在他心里。试想圣神的尊荣何其崇高，因为祂的版比前一个更好；因为祂所显示的事甚至比复活本身更伟大。的确，祂所拯救的死亡状态比之前那个更不可救药：正如灵魂比身体更宝贵一样：而这生命是由那一位，由圣神所赐的。但如果祂能赐予这个，那么更小的事就更不用说了。因为先知们曾行过奇事，但这他们却不能：因为除了天主，没有人能赦罪；先知们若不藉着圣神，也不能赐予那生命。但这不仅是奇妙之处，即祂赐予生命，而且祂也使他人也能这样做。因为祂说：\BibleQuote{你们领受圣神吧。}\BibleRef{若 20:22} 为什么？因为没有圣神就不能吗？是的，但天主为显示圣神具有至高的权威和那君王的本质，并与祂有同样的权能，也说了这话。因此祂又加上：\BibleQuote{你们赦免谁的罪，谁的罪就得赦免；你们保留谁的罪，谁的罪就被保留。}\BibleRef{若 20:23}\par

3. 既然祂赐给了我们生命，就让我们保持活着，不再回到从前的死亡中：因为\BibleQuote{基督不再死；因为祂死，是死于罪一次} \BibleRef{罗 6:9}，而祂不愿我们总是藉恩宠得救：否则我们将空无一物。因此祂要我们也从自己这边贡献些什么。那么，让我们贡献，并保全灵魂的生命。灵魂的生命是什么，要从身体学习。因为我们也说身体活着，是当它以健康的方式运动时；但当它伏卧无力，或运动杂乱时，虽然它保留着生命或运动的表象，这样的生命比任何死亡更痛苦；而且如果它说不出任何清醒的话，只说疯言乱语，看此物为彼物，这样的人又比死人更可怜。同样，灵魂没有健康时，虽然保留着生命的表象，也是死的：当它不看金子为金子，却看作伟大贵重的；当它不想未来，却爬在地上；当事物错位。因为从哪里清楚我们有灵魂？不就是从它的运作吗？那么当它不做属于它的事时，不是死的吗？例如，当它不关心德行，却贪婪违法；我从哪里能说你有灵魂？因为你行走？但这也属于无理性的受造物。因为你吃喝？这也属于野兽。那么，因为你用两脚直立？这反而让我相信你是人形的野兽。因为当你在所有其他方面与兽相似，只在直立方式上不同时，你更使我困扰和恐惧；我更认为我所看见的是一个怪物。因为如果我看见一头野兽以人的声音说话，我不会因此说它是人，反而正是因此它是比野兽更怪异的野兽。那么，我从哪里能知道你有人的灵魂，当你像驴一样踢，像骆驼一样怀恨，像熊一样咬，像狼一样吞，像狐一样偷，像蛇一样狡猾，像狗一样无耻？我从哪里能知道你有人的灵魂？你们愿意我给你们看一个死的灵魂和一个活的吗？让我们把讨论转回到那些古人；如果你们愿意，让我们把\Person{拉匝禄}故事中的富人摆在我们面前，我们就知道灵魂中的死亡是什么；因为他有一个死的灵魂，从他所作的就可以看出。因为灵魂的工作，他没有做一件，只知吃喝享乐。即使现在，那些无怜悯心肠和残忍的人也是如此，因为这些人的灵魂也像他一样是死的。因为那从爱近人而流出的所有温暖都已耗尽，它比无生命的身体更死。但是穷人不这样，而是站在至高天上智慧的高峰闪耀；虽然与持续的饥饿搏斗，连所需的食物也得不到，他却没有说出任何亵渎天主的话，而是高尚地忍受一切。现在这不是灵魂的小工作；而是它被扎实且健康的高等证明。当没有这些品质时，显然是因为灵魂死了，它们才消失了。或者，告诉我，我们不该说那灵魂是死的，当魔鬼攻击、打击、咬伤、踢它，它却毫无知觉，如同死了一样，在被掠夺财富时不悲伤；但魔鬼甚至跳上它，它仍不动，像灵魂离开的身体，甚至不感觉？因为当对天主的敬畏不严格存在时，灵魂必然如此，然后死者更悲惨。因为灵魂不是腐坏成灰烬和尘土，而是变成比这些更臭的东西，变成醉酒、愤怒和贪得无厌，变成不当的爱好和不合时宜的欲望。但如果你要更确切地知道它有多臭，给我一个纯净的灵魂，然后你就会清楚看到这个肮脏不洁的灵魂有多臭。因为目前你无法察觉。因为我们惯常接触臭味时，就不会察觉。但当我们以属灵的话喂养时，就会意识到那邪恶。然而对许多人来说，这似乎无关紧要。我还不提地狱；但让我们，如果你们愿意，考察目前的情况，以及那人是多么可笑，不是那行污秽的，而是那说污秽的；他首先如何给自己加上辱骂；就像一个人从口中喷出污秽，他这样玷污了自己。因为如果水流如此不洁，想想这污秽的源头该是什么！\BibleQuote{因为心里充满什么，口里就说什么。} \BibleRef{玛 12:34} 然而不仅为这事我悲伤，而是因为对有些人来说，这甚至不被视为不当之事。因此所有的罪恶变得更糟，当我们犯罪，却不以为我们做错了。\par

4. 你要学习污秽言语的害处有多大吗？请看听者如何因你的无礼而脸红。有什么比好说污秽话的人更卑劣？有什么比他们更无耻？因为这种人把自己降格为小丑和娼妓之流，甚至这些女人比你还知耻。你用这样的言语训练妻子走向淫荡，又怎能教她端庄呢？从口中吐出腐烂之物，也比说一句污秽的话好。如果你的口有恶臭，你甚至不能享用普通的食物；那么，当你灵魂中有如此恶臭时，请问，你竟敢领受奥迹吗？如果有人拿一个肮脏的器皿放在桌子上，你必会用棍棒打他，把他赶出去；然而，在天主自己的桌子上（因为我们充满感恩的口就是祂的桌子），你倒出比任何不洁器皿更令人作呕的话语，请问，你以为你不会激怒祂吗？这怎么可能呢？因为没有什么比这些话更惹怒圣洁和纯洁的人；没有什么比说和听这些话更使人变得厚颜无耻；没有什么比它们所点燃的火焰更能松懈端庄的筋腱。天主在你的口中放置了香料，你却储存了比尸体更臭的话语，摧毁灵魂本身，使它不能动弹。因为当你侮辱人时，这不是灵魂的声音，而是愤怒的声音；当你讲污秽的话时，这是淫荡，而不是灵魂在说话；当你毁谤时，这是嫉妒；当你图谋时，这是贪吝。这些不是灵魂的工作，而是属于灵魂的情欲和疾病的作为。正如腐败不仅来自身体，也来自死亡和身体内的情欲；同样，事实上，这些事来自于生长在灵魂上的情欲。因为如果你愿意听一个活灵魂的声音，请听\Person{保禄}说：\Quote{只要有衣有食，我们就当知足。}\BibleRef{弟前 6:8}并说：\Quote{虔敬是大利。}\BibleRef{弟前 6:6}又说：\Quote{世界于我已被钉在十字架上，我于世界也被钉。}\BibleRef{迦 6:14}听\Person{伯多禄}说：\Quote{金银我没有，但把我所有的给你。}\BibleRef{宗 3:6}听\Person{约伯}感恩说：\Quote{上主赐予，上主收回。}\BibleRef{约 1:21}这些是一个活灵魂的话，是一个灵魂尽其实当职能的话。同样，\Person{雅各伯}也说：\Quote{若上主赐我食物吃和衣服穿。}\BibleRef{创 28:20}\Person{若瑟}也说：\Quote{我怎能行这大恶，得罪天主呢？}\BibleRef{创 39:9}但那外邦女子却不然，她如醉酒和疯狂的人那样说话，说：\Quote{与我同寝吧。}\BibleRef{创 39:7}既然知道这些，让我们切切渴求活灵魂，逃离死灵魂，好使我们也能获得来生；愿我们众人都有分于那来生，赖我们的主耶稣基督的恩宠和爱人之心，藉着祂并偕同祂，归于圣父，偕同圣神，获享光荣、权能、尊崇，从今世直到永远，世世无尽。阿们。\par

\SourceRecord{Translated by Talbot W. Chambers.；Nicene and Post-Nicene Fathers, First Series；Vol. 12.；Edited by Philip Schaff.；Buffalo, NY: Christian Literature Publishing Co.,；1889.；<http://www.newadvent.org/fathers/220206.htm>.}

% structure: p0001 source=h1 target=section reason=page-title

\section{第七篇讲道辞\WorkTitle{格林多后书}}

% structure: p0002 source=h2 target=subsection reason=relative-heading-rank; base=h2

\subsection{\BibleRef{格后 3:7-8}}

\BibleRef{格后 3:7,8}\par

\BibleQuote{但如果那以文字刻在石头上属于死亡的职事，尚且带着荣耀，以致以色列子民不能注视梅瑟的面貌，因为他脸上的荣耀原是暂时的；那么，属圣神的职事岂不更该带有荣耀吗？}\par

他说，梅瑟的石版是石头的，也是用字母写成的；而新约的石版是血肉的，即宗徒们的心，是由圣神写成的。的确，文字叫人死，圣神却叫人活。但这个比较中还缺少一个并非无关紧要的补充，就是梅瑟的荣耀；而在新约中，没有人用肉身之眼见过这样的荣耀。正因如此，旧约这事显得伟大，因为荣耀是可感知的（它肉眼可见，即使不能靠近）；但新约的荣耀是藉理智领会的。因为对于较弱的人，这种优越性的理解并不清晰；而旧约的荣耀更能抓住他们，引导他们归向自己。保禄既已涉足这个比较，又因听众迟钝而极难显出新约的优越性，请看他是怎样行事、怎样入手：先用论证将差异摆在面前，而这些论证是从之前所说的构造出来的。\par

因为如果那职事属于死亡，而此职事属于生命，那么无疑，他说，后者的荣耀也大于前者。由于他无法向肉眼呈现，他便用这种逻辑推理确立了其优越性，说道：\par

\BibleRef{格后 3:8}。\BibleQuote{但如果那属于死亡的职事带着荣耀，那么属圣神的职事岂不更该带有荣耀吗？}\par

他说的\Quote{死亡的职事}是指法律。请注意他在比较中多么谨慎，以免给异端者留下把柄；因为他没有说\Quote{致死的}，而是说\Quote{死亡的职事}；因为法律服务于死，而非死的原因；导致死亡的是罪；法律却带来惩罚，显示罪，而非导致罪。因为它更清楚地揭示邪恶并惩罚它；它没有推动邪念；它不是为罪或死亡的存在服务，而是为罪人承受报应服务。这样说来，它甚至消灭了罪。因为既然它显明罪恶如此可怕，显然也使人避免它。就如有人持刀砍杀被判刑者，他是为判案的法官服务，并非他毁灭了那人，虽然是他砍的；同样，判案和定罪的人也不是毁灭者，而是被罚者的邪恶。此处也是如此：不是法律毁灭，而是罪。罪造成毁灭和定罪，而法律藉惩罚削弱了罪的力量，以惩罚的恐惧遏制它。但他并不满足于这一个考虑来确立优越性；还加上另一个，说道：\Quote{以文字刻在石头上}。请看他是怎样再次打击犹太人的骄傲根源。因为法律不过是文字：从文字中没有涌出一种帮助来激励与战斗的人，像圣洗中那样；而是柱子和文字给违反文字的人带来死亡。你们看到吗？他在纠正犹太人的好辩时，甚至用措辞削弱了法律的权威，提到石头、文字和死亡的职事，又加上\Quote{刻在石头上}。因为借此他无非是说：法律固定在某个地方；不像圣神那样无所不在，向众人吹入极大的力量；或者说文字充满威胁，而且是无法磨灭、永远留存的威胁，因为是刻在石头上。然后，即使是在看似赞美旧约时，他又掺入了对犹太人的指控。因为他说了\Quote{以文字刻在石头上，带着荣耀}之后，又补充说\Quote{以致以色列子民不能注视梅瑟的面貌}；这证明他们的极大软弱和卑下心态。他也没有说\Quote{因石版的荣耀}，而是\Quote{因他面貌的荣耀，这荣耀原是暂时的}；因为他显示：承载体有荣耀，而非石版本身。因为他没有说\Quote{因为他们不能注视石版}，而是\Quote{梅瑟的面貌}；又再次没有说\Quote{因石版的荣耀}，而是\Quote{因他面貌的荣耀}。然后，在他赞扬之后，请看他又怎样降低它，说道：\Quote{这原是暂时的}。但这并非指责，而是贬低；因为他没有说\Quote{那是朽坏的、邪恶的}，而是\Quote{停止、有终结的}。\par

\Quote{那么属圣神的职事岂不更该带有荣耀吗？}因为从此他自信地称赞新约之事，无可争议。注意他做了什么：他将\Quote{石头}与\Quote{心}对立，\Quote{文字}与\Quote{圣神}对立。然后，在显明了各自的结果后，他并没有列出各自的结果；而是列出了法律的结果即死亡和定罪，却没有列出圣神的结果即生命和义德；而是列了圣神本身；这给论证增添了伟大。因为新约不仅赐予生命，还提供了\Quote{圣神}，而圣神赐予生命，远大于生命本身。因此他说\Quote{属圣神的职事}。然后他又回到同一件事，说：\par

% structure: p0010 source=h2 target=subsection reason=relative-heading-rank; base=h2

\subsection{\BibleRef{格后 3:9}}

\BibleQuote{因为如果定罪的职事尚有荣耀}\par

他又更清楚地解释了\Quote{文字叫人死}这句话的意思，表明它正是我们上面所说过的：即法律显示罪，而非导致罪。\par

\BibleQuote{那称义的职事更远超过荣耀}。因为那些石版确实显示罪人并惩罚他们，而这职事不但不惩罚罪人，反而使他们成为义人：这就是圣洗所赋予的。\par

% structure: p0014 source=h2 target=subsection reason=relative-heading-rank; base=h2

\subsection{\BibleRef{格后 3:10}}

\Quote{诚然，那曾成为荣耀的，因这超越的荣耀，在这点上并不算荣耀。}\par

\Quote{在前面的话中，他确实已表明这恩宠也具有荣耀；而且不仅是具有荣耀，更是超越于它：因为他没有说：\Quote{圣神的职事怎会不更在荣耀中？}而是说：\Quote{在荣耀上超越；}他从先前所陈述的论据中得出证明。在这里，他也表明这优越性是何等巨大，说道：\Quote{如果我将此与那相比，旧约的荣耀根本不是荣耀；}他并非绝对地断言没有荣耀，而是就比较而言。因此他又加上：\Quote{在这点上，}即就比较而言。这并不是贬低旧约，反而高度赞扬了它：因为比较通常是在同类事物之间进行的。接着，他又从另一个角度提出另一论据来证明优越性。这论据是什么呢？是基于持久性的论据，他说：}\par

% structure: p0017 source=h2 target=subsection reason=relative-heading-rank; base=h2

\subsection{\BibleRef{格后 3:11}}

\Quote{因为如果那消逝的尚且有荣耀，那长存的更在荣耀中。}\par

因为前者停止了，而后者持续常存。\par

% structure: p0020 source=h2 target=subsection reason=relative-heading-rank; base=h2

\subsection{\BibleRef{格后 3:12}}

\Quote{我们既怀有这样的希望，就大大地放胆宣讲。}\par

因为当听者听到关于新约如此众多而伟大的事时，他会渴望亲眼目睹这荣耀，请注意他将听者引向何处——到来世。因此他也提出了\Quote{希望，}说道：\Quote{我们既怀有这样的希望。}这样的？什么样的？就是我们被算为配得比梅瑟更伟大的事；不仅我们宗徒，而且所有信徒。\Quote{我们大大地放胆宣讲。}向着谁？请告诉我。向着天主，还是向着门徒？向着你们这些受教导的人，他说；也就是说，我们随处自由地讲话，毫不隐瞒，毫不保留，毫不怀疑，而是公开宣讲；我们并不担心会损伤你们的眼目，如同梅瑟损伤犹太人的眼目那样。因为他暗示了这一点，请听下文；或者更确切地说，有必要先叙述这段历史，因为他自己不断提及它。那么这段历史是什么呢？当梅瑟第二次领受约版下来时，从他脸上发出的某种荣耀如此耀眼，以致犹太人无法靠近与他交谈，直到他用帕子蒙上脸。在出谷纪中这样记载，\BibleRef{出 34:29} \BibleQuote{当梅瑟从山上下来时，手里拿着两块约版。梅瑟不知道自己的脸皮因同上主说话而发光。他们害怕走近他。梅瑟便叫他们来，对他们说话。梅瑟向他们说完了话，就用帕子蒙上自己的脸。但当他进到上主面前与他说话时，就揭去帕子，直到出来。}\par

然后，提醒他们这段历史，他说：\par

% structure: p0024 source=h2 target=subsection reason=relative-heading-rank; base=h2

\subsection{\BibleRef{格后 3:13}}

\Quote{不像梅瑟，他曾在脸上蒙上帕子，免得以色列子民定睛注视那消逝者的结局。}\par

现在他所说的是这样的意思：我们不需要像梅瑟那样遮盖自己；因为你们能够看见我们周围所环绕的这荣耀，尽管它比那另一个更大更明亮。你看到进步了吗？因为他在前封书信中说：\Quote{我曾用奶哺养你们，没有用饭；}而在这里却说：\Quote{我们大大地放胆宣讲。}他在他们面前提出梅瑟，通过比较来推进论述，从而引导听者向上。\par

目前，他将他们置于犹太人之上，说道：\Quote{我们不需要帕子，像他治理那些人时那样；}但在随后的话中，他将他们提升到立法者本身的尊严，甚至更高。\par

然而在此期间，让我们听听接下来是什么。\par

% structure: p0029 source=h2 target=subsection reason=relative-heading-rank; base=h2

\subsection{\BibleRef{格后 3:14}}

\Quote{但他们的心灵变得迟钝；因为直到今日，在诵读旧约时，同样的帕子仍然存在，没有向他们启示出它已在基督内被废除。}\par

请注意他由此确立了什么。因为在梅瑟身上那时发生了一次的事，在法律身上却不断发生。因此，所说的并非对法律的控告，也并非对梅瑟那时蒙上自己的脸的控告，而只是对愚昧的犹太人的控告。因为法律有它应有的荣耀，但他们却无法看见。他说：\Quote{那么你们为什么困惑，如果他们不能看见这恩宠的荣耀，既然他们连梅瑟那较小的荣耀也看不见，也不能定睛注视他的面容？为什么你们因犹太人不信基督而烦恼呢，既然他们连法律也不信？因为他们也因此不认识恩宠，因为他们既不知道旧约，也不知道其中的荣耀。因为法律的荣耀在于引领人归向基督。}\par

3. 你看见吗？从这一考虑，他也打击了犹太人的骄傲？借着他们认为自己占优势的那一点，即梅瑟的脸发光，他证明了他们的粗俗和卑下本性。因此，不要让他们因此自夸，因为对那些没有享受它的犹太人来说，那算什么？因此他也反复强调，一时说：\Quote{在旧盟约诵读时，同一帕子仍然存在，}它\Quote{没有被启示出来，就是它在基督内已被废除：}另一次，说：\Quote{直到今日，当诵读梅瑟时，}\BibleRef{格后 3:15}同样的帕子仍在他们心上；显示出帕子既在诵读上，也在他们心上；而上面说：\Quote{以致以色列子民不能定睛注视梅瑟的面，因为他的面容的光荣；}这光荣\BibleRef{格后 3:7}\Quote{是正在消逝的。}还有什么比这更能标志他们价值之低呢？既然甚至对于一个将要消逝的光荣，或者更确切地说，相比之下根本不是光荣，他们都不能观看，反而被遮盖了，\Quote{以致他们不能定睛注视那消逝之事的结局；}也就是说，法律的结局，因为它有一个终结；\Quote{但他们的心意变得顽固。}有人说：\InnerQuote{那么，这与帕子有什么关系呢？}因为它预表了将要发生的事。因为他们不仅当时没有察觉，而且直到现在也不认识法律。过错在于他们自己，因为这种顽固是一种不易感动和乖谬判断的顽固。所以正是我们认识法律；但对他们而言，不仅恩宠，就连这点也被阴影遮盖；他说：\Quote{因为直到今日，在旧盟约诵读时，同一帕子仍然存在，}它\Quote{没有被启示出来，就是它在基督内已被废除：}现在他所说的是这个。他们不能看见的正是这一点，即法律已被终结，因为他们不相信基督。因为如果法律被基督终结，正如它确已终结，而且法律早已预言了这事，那么那些不接受基督（祂已废除法律）的人，如何能看见法律已被废除呢？既然不能看见这一点，很明显，就连法律本身（它宣称了这些事）的能力和完全光荣，他们也不认识。有人说：\InnerQuote{它在哪里说过这事，即法律在基督内已被废除？}它不仅是说了，而且还借着所做的事显示了。首先，确实是将祭献和全部礼仪限制在一个地方，即圣殿，然后摧毁了它。因为如果祂无意终结这些事和整个相关法律，祂就会做两件事中的一件：要么不摧毁圣殿，要么摧毁后不禁止在其他地方献祭。但是，事实是，整个世界，甚至耶路撒冷本身，祂都禁止进行这类宗教仪式；只允许并指定他们在圣殿进行。然后，在摧毁圣殿本身之后，祂完全借着所做的事显示了法律的事在基督内已被终结；因为基督也摧毁了圣殿。但如果你要也在言语中看见法律如何在基督内被废除，就听立法者亲自这样说：\Quote{上主要从你们弟兄中给你们兴起一位先知，像我一样；}\BibleRef{申 17:15}\Quote{你们要听从祂，凡祂所吩咐你们的一切，都要遵行。凡不听从那先知的，必要从民中完全灭绝。}\BibleRef{宗 3:22}你看见法律如何显示它在基督内被废除了吗？因为这位先知，即按肉身的基督，梅瑟命令他们听从的，使安息日、割损礼和所有其他事情都停止了。而达味也显示同样的事，论到基督说：\Quote{你照默基瑟德的品位永为司祭，}\BibleRef{咏 109:4}而不是照亚郎的品位。因此，保禄清楚地解释这点说：\Quote{司祭职位既然变更，法律也必须变更。}\BibleRef{希 7:12}在另一处他又说：\Quote{牺牲和素祭你都不喜欢，全燔祭和赎罪祭你也不喜悦；于是我说：看，我来。}\BibleRef{希 10:5}此外，从旧约中还可引证远多于这些的证言，显示法律如何被基督废除。因此，当你离弃法律时，你就会清晰地看见法律；但只要你紧守法律而不相信基督，你连法律本身也不认识。因此他也加上这句，以更清楚地确立此事；\par

% structure: p0033 source=h2 target=subsection reason=relative-heading-rank; base=h2

\subsection{格林多后书 3:15}

\Quote{但是直到今日，每当诵读梅瑟时，帕子还留在他们心上。}\par

因为既然他说在旧约诵读时帕子还在，免得有人以为这话是出于法律的隐晦，他既借其他事预先显示了他的意思（因为借说\Quote{他们的心意变得顽固，}他显示过错在于他们自己），在这里也同样。因为他说不是\Quote{帕子留在书写上}，而是\Quote{在诵读中}（现在诵读是读者之行动）；又说\Quote{当诵读梅瑟时。}然而他在接下来的话中更清楚地显示了这点，毫无保留地说：\Quote{帕子留在他们心上。}因为甚至帕子在梅瑟的脸上，不是因为梅瑟，而是因为他们粗俗和属肉体的心。\par

4. 在适当地控告他们之后，他也指出了他们修正的方式。这方式是什么？\par

% structure: p0037 source=h2 target=subsection reason=relative-heading-rank; base=h2

\subsection{格林多后书 3:16}

\Quote{然而当一人转向主时，}即是离弃法律，\Quote{帕子就被除去了。}\par

你看见那帕子不是放在梅瑟的脸上，而是放在犹太人的眼目上吗？因为这事的发生，不是要隐藏梅瑟的光荣，而是不让犹太人看见。因为他们不能承受。所以缺陷是在他们那里，因为这不使他（梅瑟）对任何事无知，而是使他们。他确实没有说：\Quote{当你们放弃法律时，}但他暗示了，因为\Quote{当你们转向主时，帕子就被揭去了。}他始终紧扣历史。因为当梅瑟与犹太人说话时，他蒙着脸；但当他转向天主时，脸是露着的。这正是将来要发生之事的预表：当我们转向主时，我们就会看见法律的光荣，和法律颁布者的脸露出来；是的，不仅如此，我们甚至要站在与梅瑟同等的地位上。你看他是如何邀请犹太人归信，通过表明，藉着来到恩宠，他们不仅能看见梅瑟，还能与法律颁布者站在完全同等的地位上。他说：\Quote{因为不但你们要看见那时你们所看不见的光荣，而且你们自己也要被包括在这同一光荣中；是的，甚至更大的光荣，大到那一个与这个相比似乎不再是光荣了。}怎样并以何种方式？\Quote{因为当你们转向主并被纳入恩宠时，你们将享受那光荣，与之相比，梅瑟的光荣是如此之小，以至于毫无光荣可言。但尽管它很小，且远远比不上那另一个，但当你们还是犹太人时，连这个也不赐给你们；但成为信徒后，那时就会赐给你们看见远比它更大的。}当他向信徒说话时，他说：\Quote{那有光荣的竟没有光荣；}但这里他并不这么说；而是怎样？\Quote{当一个人转向主时，帕子就被揭去了。}他循序渐进地引导他，先把他放在梅瑟的地位上，然后使他分享更大的事。因为当你们看见梅瑟在光荣中，此后你们也要转向天主，享受这更大的光荣。\par

5. 那么请看，从一开始他陈述了多少事，作为构成差异并显示新约相对于旧约的优越性，而非敌对或矛盾。他说，那是文字、石头、死亡的服务，且被废除；然而犹太人甚至未被赐予这光荣（或：这光荣）。这桌子是属血肉的、属神的、正义的，且长存；它被赐予我们所有人，而非仅赐予一人，像当时那较小的赐予梅瑟那样。\BibleRef{格后 3:18} 因为，他说，\BibleQuote{我们众人以揭开的脸面反映主的光荣}，不是梅瑟的光荣。但既然有些人坚持\Quote{当一个人转向主}这句话是指圣子而言，这与公认的相反；让我们更准确地查考这一点，首先说明他们认为可以确立此点的理由。这是什么理由呢？他们说：正如经上说：\BibleQuote{天主是神}\BibleRef{若 4:24}，同样这里说：\Quote{那圣神就是主}。但他没有说\Quote{主是神}，而是说\Quote{那圣神就是主}。并且这种结构与那种有很大差别。因为当他想按照你说的方式说话时，他就不将冠词放在谓语上。此外，让我们从头回顾他全部的话：他说的是谁？例如，当他说：\BibleQuote{文字叫人死，神却叫人活}\BibleRef{格后 3:6}，以及\BibleQuote{不是用墨写的，而是用永生活天主的圣神写的}\BibleRef{格后 3:3}，他是在说天主，还是说圣神？很明显是在说圣神；因为他正呼召他们从文字转向圣神。为了免得有人听见\Quote{圣神}后，又想到梅瑟转向了主，而他自己转向了圣神，就以为自己更差，为纠正这种猜疑，他说，\par

% structure: p0041 source=h2 target=subsection reason=relative-heading-rank; base=h2

\subsection{格林多后书 3:17}

\Quote{如今那圣神就是主。}这位也是主，他说。为使你们知道他在说护慰者，他补充说，\par

\Quote{主的圣神在哪里，那里就有自由。}\par

因为你们必定不会断言，他说：\InnerQuote{主的主在哪里。}他说：\Quote{自由，}指以前的奴役而言。然后，免得你们以为他在说将来的事，他说，\par

% structure: p0045 source=h2 target=subsection reason=relative-heading-rank; base=h2

\subsection{格林多后书 3:18}

\Quote{但我们众人，以揭开的脸面，像镜子一样反映主的光荣。}\par

不是那被废去的，而是那长存的。\par

\Quote{被变成同样的形像，从光荣到光荣，正如从主的圣神而来的。}\par

你看见他如何再次将圣神置于天主的地位（见下文），又将他们提升至宗徒的地位。因为他先前说过：\Quote{你们是基督的书信；}而这里说：\Quote{我们众人以揭开的面容。}可是他们像梅瑟一样，带来了一条法律。但他说，正如我们不需要帕子，你们这些接受它的人也不需要。然而，这荣耀远更大，因为这不属于我们的面貌，而是属于圣神；但尽管如此，你们也能像我们一样坚固地注视它。因为他们确实连藉着中保也不能，但你们甚至不用中保就能坚固地注视一个更大的。他们不能注视梅瑟的荣耀，你们却能注视圣神的荣耀。倘若圣神有任何低下，祂就不会将这些事定为比那些更大。但\Quote{我们犹如镜子反映主的荣耀，就变成同样的形象}是什么意思？这在奇迹神恩运行之时确实更清楚地显示出来；然而，即使现在，对有信德眼睛的人也不难看见它。因为我们一受洗，灵魂被圣神洁净，就比太阳更闪耀；我们不仅看见天主的荣耀，而且也从它领受一种光辉。正如纯银转向太阳光线时，它自身也会射出光线，不仅出于它自身的自然属性，也出于太阳的光辉；同样，灵魂被洁净并变得比银子更亮时，从圣神的荣耀领受一道光线，并将其反射回去。因此他也说：\Quote{我们犹如镜子反映，从荣耀变成同样的形象，}即圣神的荣耀，\Quote{到荣耀，}即我们自己的、在我们内产生的荣耀；并且这种荣耀，正如人们可期待于主那圣神的那样。请看这里他也称圣神为主。在别处也可看到祂的这种主权。因为他说：\Quote{当他们敬拜主并禁食时，圣神说：你们为我分开保禄和巴尔纳伯。}\BibleRef{宗 13:2}因此他说：\Quote{当他们敬拜主时，你们为我分开，}是为了显示圣神的尊荣平等。再者，基督说：\Quote{仆人不知道他主人所做的事；}但正如人知道自己内心的事，同样圣神也知道天主的事；不是藉着受教导，因为这样类比就不成立了。祂随己意行事也显示了祂的权威和主权。这使我们变化；这不让我们被同化于此世；因为这就是那创造者所创造的。正如他所说：\Quote{在基督耶稣内受造的}\BibleRef{弗 2:10}，同样他也说：\Quote{天主，求祢在我内再造一颗纯洁的心，在我五内更新一个正直的灵。}\BibleRef{咏 50:10}\par

6. 你愿意我亦从宗徒们身上，更显明地以感官示你这事吗？请看保禄，他的衣物行奇事；伯多禄，他的影子亦有大能。\BibleRef{宗 19:12}因为倘若他们身上没有君王的肖像，他们的光辉不可接近，那么他们的衣物和影子岂能有如此大能？君王的衣物，连强盗也畏惧。你愿见这光辉甚至透过身体照射出来吗？经上说：\Quote{定睛注视，}他说，\Quote{斯德望的面容，他们见它好像天使的面容。}\BibleRef{宗 6:15}但这相较于内在闪耀的荣耀，算不得什么。因为梅瑟面容上的，这些人却携带着在灵魂上，是的，甚至远超过它。梅瑟的荣耀确实更显于感官，而此荣耀是无形的。正如火亮的物体从发光体流向邻近物体，也将自身的光辉分予它们一些，同样，在信徒身上也确实如此。因此，凡是这样的人，必脱离尘世，只梦想天上之事。我有祸了！我们实在应当在此哀号痛哭，因为我们虽享有如此高贵的出生，竟不知所言为何，因我们迅速失去实相，被感官之物所眩惑。因为这不可言喻、可敬畏的荣耀，在我们内只存留一两天，然后我们便将它熄灭，以世俗之忧的冬天覆盖它，以那些云层的厚度驱退它的光芒。因为世俗之事乃是冬天，且比冬天更阴沉。它不生寒霜，亦不降雨，也不产生淤泥和深沼，但它制造比这一切更可怕的事，它形成地狱和地狱的苦楚。正如严霜使四肢僵硬而死，同样，灵魂在罪的冬天中战栗，也无法履行它应有的职能，如同被霜冻僵了良心。因为寒冷之于身体，正是邪恶良心之于灵魂，由此也生出胆怯。因为没有什么比固着于世俗事物的人更胆怯；这样的人过着加音的生活，日日战栗。我何必提死亡、损失、冒犯、谄媚和服役？即使没有这些，他也恐惧万千变故。他的宝库满盈黄金，但他的灵魂并未摆脱贫穷的恐惧。这很合理。因为他如同系在腐朽而迅速变动的事物上，即使自己未遭遇逆运，看到别人遭遇也使他败坏；他的胆怯巨大，他的懦弱巨大。因为这样的人不仅在危险面前无勇气，在其他一切事上也是如此。若贪财欲袭击他，他不能像自由人一样击退攻击，却如买来的奴隶，行一切它所命之事，侍奉爱财之心如同凶残的主母。若他又看见美丽的少女，他便立刻俯伏被掳，像一条狂吠的狗紧随其后，虽然本应做相反的事。因为当你看见美丽的女人时，不应考虑如何满足情欲，而应考虑如何从情欲中得释放。\Quote{这怎么可能？}有人说，\Quote{因为爱欲并非我自己的行为。}那么是谁的？请告诉我。是来自魔鬼的恶意。你深信那谋划害你的是魔鬼；那么，与之角力，与这疾病搏斗吧。\Quote{但我不能，}他说。那么来，我们先教你这一点：所发生的事来自你自己的懈怠，是你起初给魔鬼开了门，现在你若愿意，很容易就能驱逐它。那些犯奸淫的人，是因情欲而犯，还是单纯出于对危险的爱慕？显然是由于情欲。那么他们因此得到宽恕吗？当然不。为什么不？因为罪是他们自己的。\Quote{但是，}有人说，\Quote{请问，为何要堆砌三段论？因为我的良心为我作证，我愿意克制这情欲，却不能；它紧贴我，压迫我，使我痛苦。}人啊，你愿意克制它，但你并未做克制者应做的事；你就像一个发烧的人，畅饮冷泉，却说：\Quote{我用了多少方法想要退烧，却不行；它们反而使我的火焰更炽。}那么我们要看看，你是否也在做那些助长情欲的事，却以为自己在做抑制的事。\Quote{我没有，}他说。那么请告诉我，为了扑灭情欲，你曾试过做什么？最后，什么会助长情欲？即使我们这些人都不受这些指控的影响（因为被钱财之爱俘虏的人可能多于被美貌俘虏的人），但将提出的疗法对所有人和这两类人都共通。因为两者都是不理性的情欲，且后者比前者更强烈、更凶猛。当我们战胜了较大的，显然我们也会轻易制服较小的。\Quote{又怎样呢？}有人说，\Quote{如果这情欲更强烈，为何并非所有人都被这恶习所俘，而更多的人是为金钱疯狂？}因为首先，这最后的欲望似乎没有危险；其次，即使对美貌的情欲更猛烈，它却更快熄灭；因为如果它像对金钱的情欲那样持续，它会完全毁灭它的俘虏。\par

7. 那么，让我们与你们讨论这个问题，即对美丽的爱慕，并且让我们看看灾祸是如何增加的；这样我们就会知道过错是否在于我们。如果在于我们，就让我们竭尽全力去战胜它；如果不在于我们，我们为何无故折磨自己？为何我们不宽恕那些被它俘虏的人，反而责备他们？那么，这种爱是从何产生的？\Quote{从容貌的美丽，}有人说，\Quote{当伤人的女子美丽且面容姣好时。}这种说法是徒然无益的。因为如果美貌吸引恋人的话，那么这样的女子就会使所有男子都成为她的恋人；但如果她并未吸引所有人，那么此事就不是出于本性或美貌，而是出于不贞洁的眼睛。因为当你过于好奇地注视，赞赏并陷入恋慕时，你就中了箭。\Quote{那么，}有人说，\Quote{当一个人看到美丽的女子时，谁能不赞赏他所见到的呢？如果赞赏此类事物不是出于刻意的选择，那么爱慕就不取决于我们自己。}住口，人啊！你为何将一切混淆，东奔西跑，不愿看到邪恶的根源？因为我看见许多人赞赏并颂扬，却没有陷入恋慕。\Quote{人怎么可能赞赏却又不恋慕呢？}不要喧嚷（因为我接下来就要谈到这一点），但要等待，你将会听到\Person{梅瑟}赞赏\Person{雅各伯}的儿子，说：\BibleQuote{\Person{若瑟}是美男子，容貌俊秀。}\BibleRef{创 39:6}说这话的人是否恋慕呢？绝非如此。\Quote{因为，}他说，\Quote{他甚至没有见过他所赞赏的人。}然而，我们对于所描述的美丽也会有类似的感受，不仅是亲眼所见的。但为避免你对此质疑——\Person{达味}难道不是极其俊美，且眼睛有光彩吗？\BibleRef{撒上 16:12}而眼睛的美丽尤其是构成美貌最具威力的部分。那么是否有人恋慕他呢？绝非如此。所以，恋慕并非必然伴随赞赏。因为许多人的母亲容貌异常美丽。那么，他们的子女是否恋慕她们呢？不要有这种念头！但他们赞赏所见到的，却没有陷入可耻的爱慕。\Quote{不，这又是自然的美善安排。}怎么说自然的？告诉我。\Quote{因为她们是母亲，}他说。那么难道你没有听说波斯人毫无约束地与自己母亲行淫，而且不是一两个人，而是整个民族吗？但撇开这些，由此也显明这种疾病不是由于个人的青春美貌，而是由于懈怠游荡的灵魂。至少可以肯定，许多人常常略过成千上万的佳丽，而把自己交付给相貌平庸的人。由此可见，爱恋不取决于美貌：否则，那些先前的人必定会先被俘虏。那么原因何在？\Quote{因为，}他说，\Quote{如果不是美貌引起爱恋，那它有怎样的开端和根源呢？来自一个邪恶的魔鬼吗？}的确，它也来自那里，但这并不是我们要探究的；我们要探究的是我们自己是否也是原因。因为阴谋不仅是他们的，也首先有我们自己的参与。因为这种邪恶的疾病主要是由习惯、谗言、闲暇、懒散和无所事事产生的。因为习惯的暴政极大，甚至能成为自然的必要。如果习惯能产生它，那么显然习惯也能消除它。至少可以肯定，许多人通过这种方式不再恋慕，即不再看见他们所恋慕的人。初时这似乎是一件痛苦且极为不快的事，但久而久之却变为愉快，甚至即使他们想要，也不能再恢复那种激情。\par

8. 那么，当一个人没有习惯却一见就被俘虏时，又是怎样呢？这里同样，是身体的懈怠，或放纵自我，不留意自己的职责，也不从事必要的事务。因为这样的人，像个流浪者一样游荡，就会被任何邪恶刺透；像一个被释放的孩子，任何喜欢的人都能使这样的灵魂成为他的奴隶。因为灵魂习惯活动，当你阻止它在善事上运作时，既然它不能闲散，就被迫产生其他的东西。正如土地，当不被播种和种植时，只长出杂草；同样，灵魂当没有必要之事可做时，由于渴望无论如何都要活动，就会将自己交给邪恶的行为。正如眼睛从不停止观看，因此当好的事物不摆在眼前时，就会看见邪恶的事物；同样，思想当远离必要之事时，此后就会忙碌于无益之事。因为即使第一次攻击，人的工作和思想也能击退，这从许多事可以明显看出。那么，当你看见一个美丽的女子，并向她动心时，不要再注视，你就得救了。\Quote{但我怎能不再注视，}他说，\Quote{当被那欲望牵引时？}把自己交给其他能分散灵魂注意力的事，交给书籍、必要的关怀、保护他人、帮助受伤害的人、祈祷、以及思考未来之事的智慧：用这些事来束缚你的灵魂。通过这些方法，你不仅能治愈新近的创伤，也能轻易消磨已根深蒂固的旧伤。因为如果按照谚语，一次侮辱能使恋爱者放弃他的爱情，那么这些精神魅力岂不更能战胜邪恶，只要我们决心远离它。但是如果我们总是与那些向我们射箭的人交谈来往，与他们谈话、听他们所说的，我们就是在滋养这种疾病。那么，当你每天煽动火焰时，你怎能指望火被熄灭呢？\par

愿我们关于习惯所说的这番话成为对年轻人的告诫；因为对那些已经长大并受教于天上智慧的人来说，比一切更有力的是对天主的敬畏、对地狱的记念、对天国的渴望；因为这些都是能够熄灭那火的。并且随同这些也要存想：你所看见的不过是黏液、血和腐烂食物的汁液。有人却说：\Quote{然而，容貌的鲜艳是一件令人喜悦的事。}但没有什么比地上的花朵更令人喜悦，而它们也会腐烂凋谢。因此，在这件事上，也不要留意那鲜艳，而要用你的思想向内深入，并在思想中剥去那美丽的皮，仔细查看它下面是什么。因为即使是水肿病患者的身体也闪闪发光，表面并无令人反感之处；但一想到体内积聚的体液，我们仍会感到震惊，无法喜爱这样的人。\Quote{但是眼睛柔媚而顾盼，眉毛优美地弯曲，睫毛浓黑，眼球柔软，神情宁静。}但请看，即便是这些本身也不过是神经、静脉、膜和动脉。也请想想，这美丽的眼睛，当它生病衰老时，因绝望而消瘦，因愤怒而肿胀，看起来多么可憎，它消逝得多么快，甚至比画中的眼睛更快磨灭。从这些事让你的心灵转向真正的美。他说：\Quote{但是，我看不见灵魂的美。}但如果你愿意，你就能看见它：正如不在场的美丽可以用心灵欣赏，尽管肉眼看不见，同样，不用眼睛也能看见灵魂的美。你不是常常勾勒一个美丽的形式，并被吸引去描绘它吗？现在也想象灵魂的美，陶醉在那可爱之中。他说：\Quote{但是，我看不见无形的事物。}然而我们用心灵看见这些，而非有形之物。因此，例如，尽管我们看不见天使和总领天使、品德的习性和灵魂的美德，我们仍然钦佩它们。如果你看见一个体贴而节制的人，你会更钦佩他胜过那张美丽的容貌。如果你看见一个人受侮辱却忍受，被亏待却让步，要钦佩并喜爱这样的人，即使他们年事已高。因为灵魂的美就是如此；即使在老年，也有许多人迷恋它，它永不褪色，而是永远绽放。为了我们也能获得这美，让我们去寻求那些拥有它的人，并迷恋他们。因为这样，当我们获得了这美，我们也将能够获得永恒的福乐，愿我们众人都有分于这些，因我们的主\Person{耶稣基督}对人的恩宠和慈爱，愿光荣与权能归于祂和圣父及圣神，至于无穷之世。阿们。\par

\SourceRecord{Translated by Talbot W. Chambers.；Nicene and Post-Nicene Fathers, First Series；Vol. 12.；Edited by Philip Schaff.；Buffalo, NY: Christian Literature Publishing Co.,；1889.；<http://www.newadvent.org/fathers/220207.htm>.}

% structure: p0001 source=h1 target=section reason=page-title

\section{关于\WorkTitle{格林多后书}的第八篇讲道}

% structure: p0002 source=h2 target=subsection reason=relative-heading-rank; base=h2

\subsection{\WorkTitle{格林多后书}第四章1–2节}

\BibleQuote{因此，我们既然蒙受了这职务，照我们所蒙受的慈悲，就不沮丧；反而弃绝了隐晦的耻辱之事。}\par

既然他说出了伟大的事，并在\Person{梅瑟}面前把自己和所有信徒都高举起来，意识到他所说的话的高超与伟大，请注意他如何再次缓和语气。因为为了对抗假宗徒，也有必要提升听众的士气，但又要平息那膨胀；然而不能完全消除，否则就是轻率之举。因此他用了另一种方式来处理，表明这不是他们自己的功劳，而全是天主的仁慈。所以他也说：\BibleQuote{因此，我们既然蒙受了这职务。}因为我们除了成为仆人，并使自己在天主所赐的事上服侍之外，没有贡献什么。所以他并没有说\Quote{恩赐}或\Quote{供给}，而是说\BibleQuote{职务}。即便如此他仍不满足，还加上了\BibleQuote{照我们所蒙受的慈悲。}他说，就连服务于这些事本身，也是出于慈悲和仁慈。然而，慈悲是拯救人脱离邪恶，而不是赐予如此多的善事；但天主的慈悲也包括这一点。\par

\BibleQuote{我们不沮丧。}这确实要归功于祂的仁慈。因为\BibleQuote{照我们所蒙受的慈悲}这句话，应被理解为同时指\BibleQuote{职务}和\BibleQuote{我们不沮丧}。注意他多么努力地贬低自己的事。他说：\Quote{一个人被选配得到如此伟大的事物，而且这仅仅是出于慈悲和仁慈，他应当表现出这样的劳苦，经历危险，忍受试探，这并非什么大事。因此，我们不仅不消沉，反而欢喜并放胆说话。}例如，说完\BibleQuote{我们不沮丧}后，他补充道：\par

\BibleRef{格后 4:2}\BibleQuote{反而弃绝了隐晦的耻辱之事，不行诡诈，也不混乱地处理天主的圣言。}\par

什么是\BibleQuote{隐晦的耻辱之事}？他说：我们不夸口说大话，行动却显出另一套，像他们那样；所以他也说：\BibleQuote{你们只看外表；}但我们就像我们所表现出来的那样，没有任何表里不一，也不说或做那些应当用羞耻和脸红来隐藏和遮掩的事。为解释这一点，他加上了\BibleQuote{不行诡诈。}因为他们视为光荣的，他证明是耻辱和应当蔑视的。什么是\BibleQuote{诡诈}？他们享有不索取任何东西的名声，但实际上他们拿了并保密；他们享有圣人和公认宗徒的品格，却充满无数邪恶的事。但是，他说，\BibleQuote{我们弃绝了}这些事（因为这些正是他所谓的\BibleQuote{隐晦的耻辱之事}）；我们正如我们所表现的那样，没有任何隐藏。不仅在这（现世）生活中，而且在宣讲本身也是如此。因为这就是\BibleQuote{不混乱地处理天主的圣言}。\par

\BibleQuote{而是借着真理的彰显。}\par

不是借着面容和外表，而是借着我们行为的实际证明。\par

\BibleQuote{在每个人的良心面前推荐我们自己。}\par

因为我们不仅向信徒，也向不信者显明；我们公开地让所有人检验我们的行为，随他们判断；借此我们推荐自己，而不是扮演角色、带着一副漂亮面具。我们说我们不拿任何东西，并请你们作见证；我们说我们良心无罪，又从你们那里得到证明，不像他们（即假宗徒），他们遮掩自己的事，欺骗许多人。但我们既在众人面前展示我们的生活，也公开宣讲，使所有人都能领会。\par

2. 然后，因为不信者不知道它的能力，他补充说，这不是我们的错，而是他们自己的麻木。所以他又说：\par

% structure: p0013 source=h2 target=subsection reason=relative-heading-rank; base=h2

\subsection{\WorkTitle{格林多后书}第四章3–4节}

\BibleQuote{如果说我们的福音是蒙住的，那是对丧亡的人蒙住的；他们不信，被这世上的神蒙蔽了心眼。}\par

正如他先前也说过：\BibleQuote{对某些人是出于死亡以致死亡的气味，对另一些人却是出于生命以致生命的气味。}\BibleRef{格后 2:16}此处他同样如此说。但\Quote{这世界的天主}指的是谁？那些受\OriginalTerm{玛尔基翁}{Marcion}观念感染的人断言，这话是指造物主，那位只公义而不良善者；因为他们说有一位\OriginalTerm{天主}{God}，是公义而不良善的。但\OriginalTerm{摩尼教徒}{Manichees}认为这里指的是魔鬼，他们想借此经文引入另一位创造者，取代真神，这极其愚蠢。因为圣经惯常使用\Quote{\OriginalTerm{天主}{God}}一词，并不关乎被如此称呼者的尊贵，而关乎那些臣服于它之人的软弱；正如它称玛门为主，称肚腹为神。但肚腹并非神，玛门也非主，除非对于那些向它们俯首称臣的人。然而我们主张，这段经文既非论及魔鬼，也非论及另一位创造者，而是论及宇宙的天主，并且应当这样读：\Quote{天主蒙蔽了这世界不信者的心思。}因为未来的世界没有不信者，只有现今的世界才有。但即使有人以另一种方式读它，例如\Quote{这世界的天主}，这也不提供任何把柄，因为这并不表明祂仅是这世界的天主。因为祂被称为\BibleQuote{天上的天主}\BibleRef{咏 135:26}，但祂并非只是天上的天主；我们也会说\Quote{现今时代的天主}，但这并非将祂的能力局限于这一时代。此外，祂也被称为\BibleQuote{亚巴郎的天主，依撒格的天主，雅各伯的天主}\BibleRef{出 3:6}，但祂并非只是他们的天主。在圣经中可以找到许多类似的见证。那么，祂怎能使他们\Quote{蒙蔽}呢？并非祂为此做工——愿无此念！而是祂容忍并允许。因为圣经惯于如此表达，如它说：\BibleQuote{天主就让他们陷于可鄙的私欲。}\BibleRef{罗 1:28}因为他们自己先不信，使自己不配看见奥秘，随后祂也允许了。但祂应当做什么呢？难道要强拉他们，向那些不愿看的人启示吗？那样他们只会更加轻视，也仍然看不见。因此他又加上：\par

\Quote{不让他们看见基督光荣福音的光。}\par

并非使他们不信天主，而是使不信者看不见内在的事物，正如祂吩咐我们，命令不要将\BibleQuote{珍珠投在猪前}\BibleRef{玛 7:6}。因为假若祂甚至向那些不信者启示，他们的病症反而会加重。如果强迫一个患眼疾的人注视阳光，那反而加剧他的虚弱。因此医生甚至将他们关在黑暗中，以免加重病情。所以在此我们也必须认为，这些人自己成了不信者，而成了不信者之后，他们便不再看见福音的奥秘，天主从此阻止光芒照射他们。正如祂也对门徒说的：\BibleQuote{所以我用比喻对他们讲，因为他们听是听不见。}\BibleRef{玛 13:13}我所说的也可以用一个例子来阐明：假设一个希腊人，视我们的宗教为神话。这个人怎样更能获益？是进去看见奥秘，还是留在外面？因此他说：\Quote{不让他们看见光}，仍以梅瑟的史事为喻。那被遮蔽、不给他们光照的是什么呢？听他说：\Quote{不让他们看见那光荣福音的光，基督本是天主的肖像。}也就是说，十字架是世界的救恩，是祂的光荣；这位被钉者自己要带着极大光辉降临；所有其他事物，现今的、将来的、可见的、不可见的、所盼望的不可言喻的光辉。因此他也说\Quote{光照}，使你不要期望在这里得着全部，因为这里所给予的只是圣神的一线曙光。所以，他在前面提到\Quote{气味}\BibleRef{格后 2:16}以表明这一点，又提到\Quote{抵押}\BibleRef{格后 1:25}，表明更大的部分仍在将来。然而所有这些事对他们都是隐藏的，但隐藏是因为他们先不信。然后，为表明他们不仅不知道基督的光荣，也不知道父的光荣，因为他们不认识祂的（光荣），他加上：\Quote{基督本是天主的肖像？}因为你不要只停留在基督那里。因为正如通过祂你看见父，同样，你若不知道祂的光荣，你也不会知道父的。\par

% structure: p0018 source=h2 target=subsection reason=relative-heading-rank; base=h2

\subsection{格后 4:5}

3. \Quote{因为我们不是宣讲我们自己，而是宣讲耶稣基督为主，并且我们因耶稣的缘故做了你们的仆人。}\par

这里的联系是什么？这与所说的有什么共同处？他或者暗示他们抬高自己，并说服门徒以他们命名，如他在前书所说：\Quote{我是属保禄的，我是属阿颇罗的}；或者指另一件最严重的事。那是什么呢？既然他们猛烈攻击他们，各方面图谋反对他们，他说：\Quote{你们与我们争战、开战吗？不，是与我们所宣讲的那位作战，因为我们不宣讲自己。我是仆人，我甚至是那领受福音之人的仆役，为另一个人办理一切，并为祂的光荣做我所做的一切。因此，你们反对我，就是推翻属于祂的。因为我如此远离将福音的任何部分归于我个人利益，甚至我愿为基督的缘故做你们的仆人；因为祂喜欢这样尊荣你们，祂如此爱你们，为你们做了万事。}因此他也说：\Quote{并且我们因耶稣的缘故做了你们的仆人。}你看见一个远离虚荣的灵魂吗？\Quote{因为实在，他说，我们不仅不从我们主那里取任何东西，甚至为祂的缘故也谦卑自己服侍你们。}\par

% structure: p0021 source=h2 target=subsection reason=relative-heading-rank; base=h2

\subsection{格林多后书 4:6}

\BibleQuote{因为那说\InnerQuote{光要从黑暗中照耀}的天主，曾照耀在你们心中。}\par

你看见了吗？他又向那些渴望看见那超越的光荣——我指的是梅瑟的光荣——的人，以加增的光彩显示它？他说：\Quote{正如在梅瑟的面容上，也照样照耀在你们心中。}首先，他提醒他们想起创世之初所发生的：可感觉的光和可感觉的黑暗，表明这受造物更伟大。祂在哪里命令光从黑暗中照耀？在起初和创世的序幕中；因为他说：\BibleQuote{黑暗笼罩深渊，天主说：\InnerQuote{有光！}就有了光。}然而那时祂说：\BibleQuote{要有，就有了}；但现在祂什么也没有说，而是祂自己成了我们的光。因为他没有说\Quote{也曾命令}，而是祂自己\Quote{照耀了}。因此，我们藉这光所见到的，不是可感觉的对象，而是透过基督看见天主自己。你看见了在天主圣三中的不变性吗？因为论到圣神，他说：\BibleQuote{我们众人以揭开的面容反映主的荣耀，逐渐变成同样的肖像，从光荣到光荣，乃是出于主——圣神。}\BibleRef{格后 3:18}论到圣子：\BibleQuote{免得基督——天主的肖像——荣耀福音的光照射到他们身上。}\BibleRef{格后 4:4}论到圣父：\BibleQuote{那说\InnerQuote{光要从黑暗中照耀}的天主，曾照耀在你们心中，为使我们认识天主的荣耀，那是显在基督的面容上的。}因为当他说过\Quote{基督荣耀的福音}之后，又加上了\Quote{祂是天主的肖像}，表明他们也失去了祂的荣耀；同样，在说了\Quote{认识天主}之后，他又加上了\Quote{在基督的面容上}，以表明我们藉着祂认识父，正如我们也藉着圣神被带到祂面前。\par

% structure: p0024 source=h2 target=subsection reason=relative-heading-rank; base=h2

\subsection{格林多后书 4:7}

\BibleQuote{但我们是在瓦器中存有这宝贝，为彰显那卓著的力量是属于天主，并非出于我们自己。}\par

因为他已经说过许多伟大之事，论及那不可言喻的荣耀，恐怕有人说：\Quote{既然我们享受如此伟大的荣耀，怎么仍居于必死的身体中？}他说，这正是最卓越的奇事和天主大能的极好例证：一只瓦器竟能承受如此大的光辉，并保存如此高的宝贝。因此，他赞叹地说：\BibleQuote{为彰显那卓著的力量是属于天主，并非出于我们自己}；再次暗指那些自夸的人。因为恩赐的伟大和领受者的软弱，都显示祂的大能；祂不仅赐下伟大之物，而且赐给微小的祂。因为他用了\Quote{瓦器}一词，暗指我们可朽本性的脆弱，并表明我们肉身的软弱。它不比陶器更好；它易受损，被死亡、疾病、温度变化和成千上万的其他事物轻易瓦解。他说这些话，既是为了抑制他们的自负，也是为了向所有人表明，我们所拥有的一切都不是出于人。因为当天主藉卑贱成就大事时，祂的大能就特别显明。因此，在另一处他说：\BibleQuote{因为我的能力在软弱中才显得完全。}\BibleRef{格后 12:9}确实，在旧约中，成群的蛮族被蚊蝇驱散，因此祂也称蝗虫为祂的大军；\BibleRef{岳 2:25}在起初，祂仅以混乱口音就阻止了巴比伦那座高塔。在他们的战争中，祂一时以三百人击溃了无数军队，一时以号角倾覆了城市，后来又藉一个渺小贫穷的青年达味，击溃了蛮族的全军。因此，在这里祂只派遣了十二个人，就征服了世界；这十二个人又是受迫害和争战的。\par

4. 让我们惊叹于天主的德能，赞美、朝拜它。我们问犹太人，问希腊人：是谁说服了全世界离弃祖传的习惯，转投另一种生活方式？是渔夫，还是制帐幕者？是税吏，还是不学无术和无知的人？这些事如何能合乎理性，除非是神圣的德能藉着他们成就了一切？他们又说了什么来说服他们？\Quote{奉被钉十字架者的名受洗。}是怎样的人呢？是他们未曾见过、也未曾看到的人。然而，他们宣讲这些事，却说服了他们，使他们相信那些赐给他们神谕、并从祖先那里传统领受的神明不是真神；而这位基督，被钉在木头上的那一位，却把众人全都吸引到自己跟前。而且，他的确被钉死并埋葬了，在某种程度上是众人都明了的；但他复活了，除了少数人，没有人看见。但尽管如此，他们仍然说服了那些没有看见的人；不仅说服他们他复活了，而且还升了天堂，末了要来审判生者死者。那么，告诉我，这些话的说服力从何而来？无非是天主的德能。因为，首先，创新本身就得罪了所有人；而当人在这样的事上创新时，事情就变得更加严重：人拔除了古老习俗的基础，把法律从它们的宝座上拔除。此外，这些宣称者也似乎不值得信任，他们不仅出自一个万民憎恨的民族，而且还胆小无知。那么，他们如何征服了世界？他们如何把你们，以及你们那些被视为哲学家的祖先，连同他们的神明一起赶走了呢？这难道不是很明显，是因为天主与他们同在吗？因为这些成功既不是人力的，而是某种神圣且不可言喻的德能的。\Quote{不，}有人说，\Quote{是藉着巫术。}那么，魔鬼的权势本该增长，偶像崇拜本该扩展。但它们如何被推翻而消失，而我们所事奉的却与之相反呢？因此，由此甚至表明所行之事是天主的旨意；不仅从宣讲可以知道，从生命本身的名称也可以知道。因为守贞何时在世上如此广泛地扎根？轻蔑财富、生命以及其余一切，又何时如此？因为那些邪恶者和行巫术的，根本不可能成就这样的事，反而会在各方面走向反面；而这些人却在我们中间引入了天使般的生活；不仅引入，而且在我们本国、在化外人之地、在地极都确立了。因此，显而易见，是基督的德能随处成就了一切，它随处照耀，比任何闪电更快地照亮人的内心。那么，思量这一切，并将所成就的事视为未来诸事应许的确凿证据，就请同我们一起朝拜被钉十字架者不可战胜的大能，这样你们既能逃脱难以忍受的刑罚，又能获得永生的国度；愿我们众人藉着我们的主耶稣基督的恩宠与对人的慈爱，都能分享这国度；愿光荣归于祂，永世无穷。阿们。\par

\SourceRecord{Translated by Talbot W. Chambers.；Nicene and Post-Nicene Fathers, First Series；Vol. 12.；Edited by Philip Schaff.；Buffalo, NY: Christian Literature Publishing Co.,；1889.；<http://www.newadvent.org/fathers/220208.htm>.}

% structure: p0001 source=h1 target=section reason=page-title

\section{\WorkTitle{格林多后书讲道集 第九篇}}

% structure: p0002 source=h2 target=subsection reason=relative-heading-rank; base=h2

\subsection{\BibleRef{格后 4:8-9}}

\Quote{我们四面受压，却不被困住；心里作难，却不至绝望；遭逼迫，却不被丢弃。}\par

他继续论述证明整个工作应归功于天主的德能，抑制那些自夸者的高傲。他说：\Quote{不仅我们把这宝贝放在瓦器里是奇妙的，而且即使我们忍受千般磨难、四面受击，我们仍能保存它而不失去。然而，即使是一个金刚石般的器皿，也不足以承受如此巨大的宝藏，也无法抵挡如此众多的阴谋；但如今，它既承载了宝藏，又未受损害，这都是藉着天主的恩宠。} 因为他说：\Quote{我们四面受压，却不被困住。} 什么是\Quote{四面}？\par

\Quote{在仇敌方面，在朋友方面，在必需品方面，在其他需要方面，由敌对者，由自己家中的人。} \Quote{却不被困住。} 看他如何用相反之事说话，为要借此显出天主的德能。因为他说：\Quote{我们四面受压，却不被困住；心里作难，却不至绝望。} 也就是说：\Quote{我们并未完全跌倒。因为我们常常在计算上出错，未达目标，但并未因此从摆在我们面前的事上坠落：因为这些事情是天主允许的，为训练我们，而非使我们失败。}\par

\BibleRef{格后 4:9} \Quote{遭逼迫，却不被丢弃；被打倒，却不至灭亡。} 因为这些试炼确实临到，但试炼的后果却未临到。这实在是藉着天主的德能和恩宠。他在别处说，这些事之所以被允许，是为了使他们谦卑，也为了他人的安全：因为他说：\Quote{恐怕我过于高举自己，就有一根刺加在我身上} \BibleRef{格后 12:7}；又说：\Quote{免得有人看我所见的或听我所说的，过于我本身}；在另一处又说：\Quote{使我们不依靠自己} \BibleRef{格后 1:9}；在此处，却是为彰显天主的德能。你看见他的试炼带来多大的益处吗？因为既显明了天主的德能，又更多地启示了他的恩宠。因为他说：\Quote{我的恩宠够你用的} \BibleRef{格后 12:9}。这也傅油使他们谦卑，预备他们压制其余，并使他们更加刚强。因为他说：\Quote{忍耐生老练，老练生盼望} \BibleRef{罗 5:4}。因为那些陷入万千危险、藉着对天主的希望而得救的人，更加学会在万事上持守这希望。\par

% structure: p0007 source=h2 target=subsection reason=relative-heading-rank; base=h2

\subsection{\BibleRef{格后 4:10}}

\Quote{身上常带着主耶稣的死，使耶稣的生也显明在我们身上。}\par

他们所带着的\Quote{主耶稣的死}是什么？就是他们每日的死亡，藉此也显出复活。他说：\Quote{如果有人不信耶稣死了又复活了，看见我们每日死去又复活，就让他从此相信复活。} 你看见他发现了试炼的另一个理由吗？那么这个理由是什么？\Quote{使他的生也显明在我们身上。} 他说：\Quote{藉着将我们从危险中拯救出来。因此，这看似软弱和匮乏的记号，却宣告了他的复活。因为他的德能若使我们不受苦，就不如现在这样明显，即我们虽受苦，却未被胜过。}\par

% structure: p0010 source=h2 target=subsection reason=relative-heading-rank; base=h2

\subsection{\BibleRef{格后 4:11}}

\Quote{因为我们这活着的人，为耶稣的缘故，常被交于死地，使耶稣的生，也在我们这必死的身上显明出来。}\par

因为每当他说了什么晦涩的话，他都会再次解释自己。此处他也如此行，对我所引用的经文给出了清楚的解释。他说：\Quote{因此，我们被交于死地}，换句话说，我们带着他的死，为使他的生命之德能显明出来，他既不容许必死的肉身在如此巨大的苦难中被这灾难的风暴所胜。也可以从另一角度理解。如何理解？正如他在另一处所说：\Quote{我们若与他同死，也必与他同活} \BibleRef{弟后 2:11}。\Quote{因为我们现在忍受他的死，并选择活着时为他而死；同样，他也会选择在我们死后，赐予我们生命。因为如果我们从生命进入死亡，他也必将从死亡引导我们进入生命。}\par

% structure: p0013 source=h2 target=subsection reason=relative-heading-rank; base=h2

\subsection{\BibleRef{格后 4:12}}

\Quote{这样看来，死是在我们身上发动，生却在你们身上发动。}\par

他不再讲严格意义上的死亡，而是讲试炼和安息。他说：\Quote{我们确实在危险和试炼中，而你们却在安息中；你们收获这试炼所产生的生命。我们忍受危险，你们却享受美物；因为你们没有经受如此大的试炼。}\par

% structure: p0016 source=h2 target=subsection reason=relative-heading-rank; base=h2

\subsection{\BibleRef{格后 4:13}}

2. \Quote{但我们既有同样信德的精神，正如经上所记：\InnerQuote{我信了，因此我说了话。} 我们也信，因此也说话，知道那叫主耶稣复活的，也必叫我们与耶稣一同复活。} \BibleRef{咏 115:10}\par

他提醒我们想起一篇富含天上智慧、特别适合在危险中鼓舞人心的圣咏。因为那位义人在极大危险中说了这句话，那时除了天主的援助别无得救的可能。既然相似的情境最能安慰人，因此他说：\Quote{拥有同一圣神}，即\InnerQuote{藉着那拯救他的同样援助，我们也得救；藉着他说话的圣神，我们也说话。}由此他表明，在新约和旧约之间存在着伟大的和谐，同一圣神在两者中运行；不仅我们处于危险中，古代所有人也都如此；我们必须通过信德和望德找到补救，而不是立刻寻求摆脱加于我们身上的重担。因为他已通过论证显示了复活和生命，表明危险并非无能或匮乏的标志；他随后又引入信德，并将一切都交托给它。但他仍为此提供了证据，即基督的复活，他说：\Quote{我们也相信，所以我们也说话。}我们相信什么？请告诉我。\par

% structure: p0019 source=h2 target=subsection reason=relative-heading-rank; base=h2

\subsection{\BibleRef{格后 4:14-15}}

\Quote{那使耶稣复活起来的，也要使我们复活起来，并与你们一同把我们呈献在祂面前。因为一切都是为你们，好使恩宠藉着众多人的增加，使感谢洋溢，归光荣于天主。}\par

他又以崇高的思想充满他们，使他们不觉得自己欠了人——我是指假宗徒——的债。因为一切都是出于天主，祂愿意把恩宠赐给许多人，好让恩宠显得更大。因此，复活和一切其他事物都是为你们而行的。因为祂做这些事不只是为一个人，而是为所有人。\par

% structure: p0022 source=h2 target=subsection reason=relative-heading-rank; base=h2

\subsection{\BibleRef{格后 4:16}}

\Quote{因此，我们不气馁；但我们外在的人虽然朽坏，内在的人却日复一日地更新。}\par

它如何朽坏？受鞭打、受迫害、遭受无数极端之苦。\Quote{然而内在的人却日复一日地更新。}它如何更新？藉着信德、望德、进取的意志，最终，藉着勇敢面对那些极端。因为身体遭受多少苦难，灵魂就有多么美好的望德，并变得更为明亮，如同金子被火炼净，越来越纯净。请看他是怎样使今生的痛苦变为虚无。\par

% structure: p0025 source=h2 target=subsection reason=relative-heading-rank; base=h2

\subsection{\BibleRef{格后 4:17-18}}

\Quote{因为这轻微的痛苦，}他说，\Quote{正在片刻之间，为我们成就极重无比的永远光荣；我们并不注视那看得见的，而是注视那看不见的。}\par

他以望德结束了这个问题，（正如他在\WorkTitle{罗马书}中所说的：\Quote{我们得救是在于望德，但看得见的望德就不是望德了；}\BibleRef{罗 8:24}，他在这里也确立了同样的论点，）他将现在与将来、短暂与永恒、轻微与沉重、痛苦与光荣并列。他还不满足于此，又加了一句，重复说：\Quote{极重无比。}接着他也展示了如此巨大的痛苦如何成为轻微。那么如何轻微呢？\Quote{我们并不注视那看得见的，而是注视那看不见的。}因此，如果我们从看得见的事物中抽身，那么这现世将是轻微的，而将来的将是伟大的。\Quote{因为看得见的是暂时的。}\BibleRef{格后 4:18}但看不见的才是永远的。所以冠冕也是如此。他没有说痛苦是这样的，而是说\Quote{看得见的事物}，所有这一切，无论是惩罚还是安息，这样我们就不会被其中之一所膨胀，也不被另一个所压倒。因此，当谈到将来的事物时，他没有说王国是永远的，而是说\Quote{看不见的事物是永远的}，无论是王国还是惩罚；这样既以一方警醒人，又以另一方鼓励人。\par

3. 既然 \BibleQuote{看得见的是暂时的，看不见的却是永远的}，就让我们注视那些看不见的事。因为如果我们选择暂时的而舍弃永远的，我们还有什么借口？因为即使现世是快乐的，它却不持久；而它所招致的灾祸却是持久且不可赦免的。因为那些曾被圣神视为相称并享有如此大恩赐的人，若变得心思卑下，跌落到地上，他们还有什么借口？因为我听见许多人说出这些完全值得蔑视的话：\Quote{给我今天，明天拿去。} 因为有人说：\Quote{因为如果确实有你所说的那些事，那么就是一对一；但若根本没有这回事，那就是二对零。} 还有什么比这些话更无法无天？或者还有更无聊的空谈？我们正在谈论天堂和那些不可言喻的美善；而你却向我们提出赛马场上的术语，还不羞愧，也不掩面，竟然说出适合狂人的话？你难道不脸红，竟如此执着于眼前的事？你还不停止神志不清、失去理智，在年轻时像个老糊涂？如果希腊人这样说，倒不足为奇；但信徒竟说出这样的糊涂话，这能获得什么宽恕？因为你真的怀疑那些不朽的希望吗？你认为这些事完全不可靠？这些事有什么值得原谅？有人说：\Quote{谁曾来过，并带回那里有什么的消息？} 的确，没有人，但天主——比所有人都更可信——已经宣告了这些事。但你却看不见那里有什么。你也看不见天主。那么你会因为看不见祂就否认有天主吗？他回答：\Quote{是的，我坚信有天主。} 那么如果一个外教人问你：\Quote{谁从天堂来带回这个消息？} 你将如何回答？你怎么知道有天主？他回答：\Quote{从可见的事物，从整个受造界的井然有序，从祂向所有人显明。} 那么你也同样接受审判的教义。他问：\Quote{如何？} 我要问你，你回答我。这位天主是公义的，祂会按各人的行为报应吗？或者相反，祂愿恶人快乐奢侈地生活，而善人则相反？他回答：\Quote{绝不，因为连人也不会这样想。} 那么，在这里行善的人将在哪里享受美善？而恶人将在哪里得到相反的报应，除非在来世有生命和报应？你看见目前是一对一，而不是二对一。但我将继续向你表明，甚至不是一对一，而是对义人是二对零；对罪人和在此放纵生活的人，则完全相反。因为在此放纵生活的人连一对一也没有得到；而那些度德性生活的人却是二对零。因为谁在安息？是滥用今生的人，还是追随天上智慧的人？也许你会说是前者，但我用后者来证明，并召那些享受过现世事物的人作证；他们不会无耻到否认我将要说的话。因为多次他们诅咒媒人和他们新房挂花环的那一天，并称那些没有结婚的人为有福。许多年轻人，即使可以结婚，也因为这件事的麻烦而拒绝。我说这些，并不是指责婚姻；因为婚姻是 \BibleQuote{可尊敬的} \BibleRef{希 13:4}；但指责那些滥用婚姻的人。如果那些过婚姻生活的人常常认为自己的生命不值得活，那么对那些堕入娼妓深坑、比任何俘虏更奴役和悲惨的人，我们该说什么？对那些在奢侈中腐烂、身体缠绕着千万疾病的人呢？\Quote{但受人尊敬是一种快乐。} 是的，更确切地说，没有什么比这种奴役更痛苦了。因为追求虚荣的人比任何奴隶都更奴性，愿意取悦任何人；但践踏虚荣的人却超越一切，不关心来自他人的荣耀。\Quote{但拥有财富是可取的。} 然而我们经常表明，那些摆脱财富、一无所有的人，享有更大的财富和安宁。\Quote{但醉酒是愉快的。} 但谁会这样说？那么，如果无财富比有财富更愉快，不结婚比结婚更愉快，不追求虚荣比追求虚荣更愉快，不奢侈生活比奢侈生活更愉快；即使在今世，那些不执着于现世事物的人也占有优势。我尚未提及，即使前者遭受千万种折磨，也有那美好的希望支撑他；而后者即使享受千百种快乐，也有对未来的恐惧扰乱和破坏他的快乐。因为这本身也是一种不轻的惩罚；相反，后者也不是享受和安宁。此外还有第三种。这是什么？就是世俗享乐的事物即使在眼前时也不显得那样，因为被自然和时间所驳斥；但其他事物不仅存在，而且坚定不动摇。你看见我们不仅能提出二对零，甚至三、五、十、二十、一万对零吗？但为了让你通过例子也了解这同一真理——富人和拉匝禄——前者享受了现世的事物，后者享受了来世的事物。\BibleRef{路 16:19} 那么，在你看来，这是否是一对一：永远受惩罚，和短暂饥饿？在可朽的身体上患病，而在不朽的身体上悲惨地焚烧？在那短暂疾病后戴冠冕并享受不朽的快乐，而在短暂享受他的财富后无尽地受折磨？谁会这样说？因为你想我们比较什么？数量？质量？等级？天主对每一方的决定？你要像永远在粪堆中打滚的甲虫那样说话多久？因为这些不是理性人的话语——为了虚无而抛弃如此宝贵的灵魂，只需付出少许劳动就能获得天堂。你愿意我也用另一种方式教导你那里有一个可怕的审判台吗？打开你良心的门，看看坐在你心里的审判官。既然你虽爱自己，却仍审判自己，且不能不行公义的判决，难道天主不会更完全地照顾公义，对所有人施行公正的审判吗？或者祂会让一切松散随意地进行？谁会这样说？没有人；希腊人和野蛮人，诗人和哲学家，甚至全人类都在这方面与我们一致，虽然细节不同，但都断言在阴间有某种审判台；这件事是如此明显和无可争议。\par

有人说：\Quote{为什么他不在这里惩罚呢？}那是为彰显祂的宽仁，并藉悔改为我们提供救恩，免得我们的族类被扫除，也不在救恩之前拔除那些能藉卓越的改变而得救的人。因为如果祂一犯罪就立刻惩罚并毁灭，\Person{保禄}又怎能得救？\Person{伯多禄}，世界的主要导师，又怎能得救？\Person{达味}又怎能因悔改而获得救恩？\Person{迦拉达人}又怎能？其他许多人呢？因此，祂既不在此世向所有人追讨（只从所有人中取一部分），也不在彼世向所有人追讨，而是对一个人在这里，对另一个人在那里；以便藉受惩罚的人唤醒那些极其麻木的人，并藉不受惩罚的人使他们期待将来的事。或者你们没有看到许多人在这里受惩罚，例如那些被塔楼倒塌压死的人\BibleRef{路 13:4}，那些比拉多将他们的血混入祭品的人，那些在格林多人中因不相称地领受圣事而夭亡的人\BibleRef{格前 11:30}，\Person{法郎}，以及在野蛮人手下被杀的犹太人，还有许多其他的人，无论当时、现在，还是不断如此？然而另有一些人，虽犯了许多罪，却未在此受罚便离世，如同\Person{拉匝禄}故事中的富人，以及许多其他人（\EditorialNote{路 16}）。祂这样做，既为唤醒那些完全不信将来之事的人，也使那些相信却疏忽的人更加殷勤。\BibleQuote{因为天主是正义的审判者，刚强而宽仁，每日不轻易发怒。}\BibleRef{咏 7:11}但若我们滥用祂的宽仁，终有一天祂将不再宽容，哪怕片刻，而是立刻施罚。\par

那么，我们不要为了片刻的享乐（因今生就是这样）而为自己招致永无穷尽的刑罚；反要劳苦片刻，为获得永远的冠冕。难道你们没有看见，即使在世物上，大多数人也这样做：选择短暂的劳苦以换取长久的安息，即使结局相反？因为在此生中，劳苦与报偿是相等的；甚至常常劳苦无尽而果实很少，或连很少也没有；但在天国的事上，正好相反：劳苦微小，而快乐巨大且无限。试想：农夫全年辛苦，到希望的尽头时，往往得不到多劳的果实。船主和士兵则直到极老，都忙于战争和劳苦；而且他们中常常一个失去了富有的货物，另一个连胜利与性命都失去。那么，请告诉我，如果我们今世尚且宁愿劳苦以求片刻安息（或者甚至不止片刻；因为希望是不确定的），在属灵的事上却反其道而行，为了一点懈怠便招致不可言喻的惩罚，我们将有何托辞？因此，我恳求你们众人，虽已迟延，但仍要从此狂热中醒悟。因为在那一日，没有人能拯救我们；没有兄弟，没有父亲，没有子女，没有朋友，没有邻人，没有其他任何人；如果我们行为不忠，一切都完了，我们必然灭亡。那富人何等哀哭，恳求族长，请求派遣\Person{拉匝禄}！但是请听\Person{亚巴郎}对他说的话：\BibleQuote{在我们与你们之间隔着深渊，以致愿意过去的人不能过去。}\BibleRef{路 16:26}那些童女为了一点油向她们的同伴作了多少请求！但请听她们所说的话：\BibleQuote{恐怕不够你们和我们用的；}\BibleRef{玛 25:9}没有人能领她们进入婚宴。\par

既思念这些事，我们也当谨慎对待我们的生命。因为无论你提出什么劳苦，加上什么惩罚，这一切加起来与将来的福乐相比，都算不了什么。因此，你若愿意，可以举出火、刀剑、野兽，以及比这些更残酷的任何东西，但与那些刑罚相比，连影子都算不上。因为这些东西在过度使用时反而变得特别轻松，使解脱迅速；因为身体不能同时承受剧烈的和持久的痛苦；但二者（延长与剧烈）却同时存在，无论是对善还是对恶。趁我们还有时间，\BibleQuote{让我们带着感恩到祂面前，}\BibleRef{咏 94:2}，为使我们在那一天能见到祂温柔和蔼，能完全逃脱那些携带威胁的权势。你们没有看见今世士兵如何执行当权者命令，拖曳人，捆绑他们，鞭打他们，刺穿他们肋旁，用火炬烧灼他们，肢解他们吗？然而这一切与那些刑罚相比，不过是儿戏和玩笑。因为这些刑罚是暂时的；但那里虫不死，火不灭；因为那将要复活的身体是不朽坏的。但求天主使我们永不会从经验中知道这些事；使这些可怕的事对我们的逼近不超过它们的提及；使我们不被交到那些刑吏手中，却能因此变得明智。那时我们将如何控告自己！将发出多少哀叹！多少呻吟！但此后便毫无用处。因为水手在船只破碎沉没之后，不能有任何作为；医生在病人去世之后也是如此；他们固然会常说应该这样做那样做，但一切都是徒劳无益。的确，只要还有改过的希望，就该说和做一切；但当我们不再有任何能力，一切都已败坏时，说和做就都没有意义了。因为连犹太人那时也会说：\BibleQuote{奉主之名而来的，当受赞美：}\BibleRef{玛 23:39}，但他们从这呼喊中得不到任何好处以逃脱刑罚，因为当该说的时候，他们没有说。为使我们在生活中不致如此，让我们现在并从此刻开始悔改，好使我们能全然坦然无惧地站在基督的审判台前。愿我们众人因我主耶稣基督的恩宠和爱世人，得以达到这目的。愿光荣与权能归于祂，及圣父及圣神，至于永远。阿们。\par

\SourceRecord{Translated by Talbot W. Chambers.；Nicene and Post-Nicene Fathers, First Series；Vol. 12.；Edited by Philip Schaff.；Buffalo, NY: Christian Literature Publishing Co.,；1889.；<http://www.newadvent.org/fathers/220209.htm>.}

% structure: p0001 source=h1 target=section reason=page-title

\section{论\WorkTitle{格林多后书}第十讲}

% structure: p0002 source=h2 target=subsection reason=relative-heading-rank; base=h2

\subsection{格林多后书 5:1}

\BibleQuote{因为我们知道：如果我们这地上的帐棚式的房屋被拆毁，我们必由天主获得一所房舍，一所非人手所造，永远在天上的房屋。}\par

他又激起他们的热忱，因为许多考验临近了。因为他们很可能因他不在而在这一点上（需要）更为软弱。那么他说什么？我们不应因遭受磨难而惊奇，也不应困惑，因为我们由此还获得许多益处。其中一些他前面已提过；例如，我们\BibleQuote{身上常带有耶稣的死状}，并显出祂德能的最大明证：因为他说，\BibleQuote{那超出寻常的德能是属于天主的}；我们且显明复活的确证，因为他说，\BibleQuote{为使耶稣的生命也彰显在我们有死的肉身上}。但既然他同时说我们的内在人因而变得更好；因为\BibleQuote{我们外在的人虽在败坏}，他说，\BibleQuote{我们内在的人却日日更新}；再次表明这受鞭打和迫害是相称有益的，他又说，当这事彻底完成时，无数的美事将为那些忍受这些事的人涌现出来。因为怕你们听到你们外在的人毁坏而悲伤；他说，当这事完全实现时，你们将最为喜乐，并将获得更好的产业。因此，人不仅不应因它部分地毁坏而悲伤，反而应热切寻求那毁坏的完成，因为这事最能使你走向不朽。因此他也加上：\BibleQuote{因为我们知道：如果我们这地上的帐棚式的房屋被拆毁，我们必由天主获得一所房舍，一所非人手所造，永远在天上的房屋。}因为他再次强调复活的道理——在这方面他们特别软弱——他也求助于听众的判断，以此确立它；然而并不像先前那样，而是仿佛从另一个主题进入：（因为他们已被纠正过）并说：\BibleQuote{我们知道，如果我们这地上的帐棚式的房屋被拆毁，我们必由天主获得一所房舍，一所非人手所造，永远在天上的房屋。}有些人确实说这里的\OriginalTerm{地上的房屋}{earthly house}是指这个世界；但我主张他更是指身体。但请观察，他如何藉着[他所用的]词汇显示未来事物对现在的优越性。因为他说了\Quote{地上的}，就把它与\Quote{天上的}相对；他说了\Quote{帐棚式的房屋}，从而表明它容易拆散且是暂时的，他就把它与\Quote{永远的}相对，因为\Quote{帐棚}一名常表示暂时性。因此他说：\BibleQuote{在我父的家里有许多住处。}\BibleRef{若 14:2}但如果他也在别处称圣人们的安息之所为帐棚，他并非简单地称它们为帐棚，而是加上一个形容词；因为他没有说\BibleQuote{他们接你们}到他们的帐棚里，而是说\BibleQuote{到永远的帐棚里。}\BibleRef{路 16:9}此外，他在说\Quote{非人手所造}时，也暗示了那些人手所造的。那么怎样？身体是人手所造的吗？绝非如此；但他或是指这里人手所造的房屋，或若非如此，那么他称那非人手所造的身体为\Quote{帐棚式的房屋}。因为他使用这个术语不是与此形成对立和对比，而是为了提升那些赞美和增加那些称誉。\par

% structure: p0005 source=h2 target=subsection reason=relative-heading-rank; base=h2

\subsection{格林多后书 5:2}

2. \BibleQuote{因为我们在这帐棚里叹息，切望穿上那属天的住处。}\par

什么住处？请告诉我。是不朽的身体。为什么我们现在叹息？因为那是更好的。他称它为\Quote{来自天上的}是因为它的不朽性。因为确实不是有一个身体会从上面降临到我们这里；而是藉着这个表达，他指明从那里所赐的恩宠。因此，我们远不应为这些部分的考验而悲伤，反而应寻求它们的满全，就好像他说：你叹息，因为你被迫害，因为你这个人正在败坏？你叹息这事没有做到极致，没有完全毁坏。你看见他如何将所说的话颠倒过来；他证明他们应该为那些事没有完全成就而叹息；而他们却因为那些事部分地成就而叹息。因此他从此不再称之为帐棚，而称之为房屋，并且很有理由。因为帐棚确实容易拆散；但房屋却永远长存。\par

% structure: p0008 source=h2 target=subsection reason=relative-heading-rank; base=h2

\subsection{格林多后书 5:3}

\BibleQuote{如果我们脱去这帐棚，也不会是裸体的。}\par

这就是说，即使我们脱去了身体，我们不会在那里被呈现为没有身体，而是以同一个身体成为不朽的。但有些人读到（另一种读法），并且非常值得采纳：\BibleQuote{如果我们穿上衣服，就不会赤身露体。}因为怕所有人因复活而自信，他说：\BibleQuote{如果我们穿上衣服}，即获得了不朽和不朽的身体，\BibleQuote{我们就不会在光荣和安全上成为赤身露体。}如他在前书也说过：\BibleQuote{我们都要复活；但各人按自己的次序。}又说：\BibleQuote{有属天的身体，也有属地的身体。}\BibleRef{格前 15:22}因为复活确实是众人共同的，但光荣不是共同的；有些人在光荣中复活，有些人在羞辱中复活，有些人复活进入王国，有些人进入惩罚。他在这里无疑也暗示了这点，当他说：\BibleQuote{如果我们穿上衣服，就不会赤身露体。}\par

% structure: p0011 source=h2 target=subsection reason=relative-heading-rank; base=h2

\subsection{格林多后书 5:4}

3. \Quote{我们在这一帐棚里的确呻吟，不是因为我们愿意脱去，而是因为我们愿意穿上。} 此处他又彻底而显明地堵住了异端者的口，表明他并不是在绝对地谈论身体本质的差别，而是在谈论腐朽与不朽：\Quote{因为我们并不因此呻吟，}他说，\Quote{乃是为要脱离那身体：因为对于这一点，我们并不愿脱去衣服；而是急于脱离其内的腐朽。} 因此他说：\Quote{我们不愿脱去身体，而是愿它被穿上不朽。} 随后他又这样解释：\Quote{为使那有死的被生命吞灭。} 因为脱去身体在许多人看来似乎是痛苦的事；而当他所说\Quote{我们呻吟}时，他与所有人的判断相矛盾，说并不愿从中解脱；（\Quote{因为如果，}有人说，\Quote{灵魂在分离时如此受苦哀叹，你怎说我们呻吟是因尚未分离呢？}）为了避免这一点被用来攻击他，他说：\Quote{我也不是说我们因此呻吟，以便脱去它；（因为没有人能不痛苦地脱去它，看基督甚至对伯多禄说：\InnerQuote{他们要 \InnerQuote{带你}，并领你 \InnerQuote{往你不愿去的地方}；}\BibleRef{若 21:18}）而是为叫它被穿上不朽。} 因为我们受身体的重担，是在于这一点；不是因为它是身体，而是因为我们被一个可朽坏的、易受苦的身体所包围，因为这正是使我们痛苦的原因。但当生命来到时，它毁灭并消耗了腐朽；我是说，腐朽，而非身体。\Quote{这事如何成就？}有人问。不必追问；是天主成就的；不要过于好奇。因此他又补充道：\par

% structure: p0013 source=h2 target=subsection reason=relative-heading-rank; base=h2

\subsection{格后 5:5}

\Quote{那为这事工成就我们的，是天主。}\par

由此他表明这些事是从起初预表的。因为这不是现在才决定的：而是在起初祂用尘土造我们、创造亚当时；祂造他并不是为叫他死亡，而是为使他甚至成为不朽的。然后为表明这事可信并提供证明，他补充道：\par

\Quote{祂也赐下了圣神的凭据。} 因为甚至那时祂就为这事工而陶造了我们；而现在祂藉着圣洗成就了这事，并赐给了我们并非微小的保障，即圣神。他不断称它为凭据，意在证明天主是整个事工的债务人，并借此使他对较为粗俗之人所说的话更为可信。\par

% structure: p0017 source=h2 target=subsection reason=relative-heading-rank; base=h2

\subsection{格后 5:6}

4. \Quote{所以我们时常勇敢，并且知道。}\par

\Quote{勇敢}一词是就迫害、阴谋和不断的死亡而言的：仿佛他说：\Quote{有人烦扰、迫害、杀害你吗？不要沮丧，因为这一切是为你的益处。不要害怕：但要勇敢。因为你所呻吟和悲伤的事，即你受制于腐朽，祂从此将其挪开，并早早地使你脱离这束缚。}因此他也说：\Quote{所以我们时常勇敢}，不仅是在安息之时，也是在患难之时；\Quote{并且知道}。\par

% structure: p0020 source=h2 target=subsection reason=relative-heading-rank; base=h2

\subsection{格后 5:7-8}

\Quote{因为我们住在身内时，是远离主的（因为我们凭信德而行，并非凭目睹）；我说，我们勇敢，且情愿离身与主同住。}\par

那比一切都更伟大的，他放在了最后，因为与基督同住比领受不朽的身体更好。但他的意思是：\Quote{祂并不熄灭那与我们争战并杀害我们的生命；不要害怕；即使被肢解也要勇敢。因为祂不仅使你脱离腐朽和重担，还迅速送你到主那里。}因此他也没有说\Quote{当我们\Emph{在}身内时}：如同那些在异国他乡的人一样。\Quote{所以我们知道，当我们住在身内时，是远离主的；我说，我们勇敢，且情愿离身与主同住。} 你看见吗？他将那些令人痛苦的名词——死亡与终结——隐去，代之以激发极大渴望的名词，称之为与天主同在；而越过那些通常被当作甜美之物——生命中的事物——他用痛苦的名词来表达，称此世的生活为远离主。他这样做，既是叫无人迷恋于现世之物，反而对此厌倦；也是叫无人将死时感到不安，反而欣喜，如同前往更大的美善。然后，免得有人听见我们说远离主，就说：\Quote{你为何这样说？难道我们在此世时就与祂疏远吗？}他预先纠正了这种想法，说：\Quote{因为我们凭信德而行，并非凭目睹。} 即使在此，我们确实认识祂，但并不那么清楚。正如他也在别处所说，\BibleRef{格前 13:12} \Quote{藉着镜子}，和\Quote{模糊不清}。\par

\Quote{我说，我们勇敢，且情愿。} 奇妙！他把话题引向了何处？引向了极度渴望死亡，表明了痛苦成为可喜，可喜成为痛苦。因为藉着\Quote{我们情愿}一词，他意思是\Quote{我们渴望}。我们渴望什么？渴望\Quote{离身与主同住}。他常常如此行，（如前所述）把对手的反对论点完全扭转过来。\par

% structure: p0024 source=h2 target=subsection reason=relative-heading-rank; base=h2

\subsection{格后 5:9}

\Quote{因此我们立定志向，无论是在家还是离世，都要讨祂喜悦。}\par

\Quote{因为我们所追求的正是这个，}他说，\Quote{不论我们在那里或在这里，都要按照祂的旨意生活；因为这是首要的事。因此你们藉此已经拥有天国，无需考验。}因为恐怕当他们已有了如此渴望在那里的强烈愿望时，又因还要等那么久而不安，他在这里已经将那美好事物的首要部分给了他们。这是什么？就是中悦。因为离开本身并非绝对的好，而是在祂的恩宠中离开，这才使离开也成为好事；同样，留在这里并非绝对的痛苦，而是留在这里冒犯祂。所以不要以为脱离身体就够了；因为德行总是必需的。因为正如当他谈到复活时，他不容许他们单靠复活就放心，说：\BibleQuote{如果穿上，就不会赤身露体。}同样，在显示了离开（身体）之后，免得你们以为这就足够救你们，他补充说必须中悦。\par

5. 既然他用许多好事说服了他们，从此他也用更阴郁的事来惊吓他们。因为我们的利益既在于获得好事，也在于避免恶事，换句话说，就是地狱和天国。但是既然这一点——避免惩罚——是更有力的动机；因为当惩罚只限于得不到好事时，大多数人会甘心接受；但如果延伸到遭受恶事，就不再如此了。（因为他们本当认为前者无法忍受，但由于多数人的软弱和卑下天性，后者在他们看来更难承受。）既然（我说）给予好事并不像惩罚的威胁那样能激起一般的听众，他不得不以此作结，说道：\par

% structure: p0028 source=h2 target=subsection reason=relative-heading-rank; base=h2

\subsection{\BibleRef{格后 5:10}}

\BibleQuote{因为我们众人都必须出现在审判台前。}\par

然后，他通过提及那审判台使听者震惊和战栗，即使在这里，他也没有只列出可怕的事而不提好事，而是掺入了一些喜悦，说道：\par

\BibleQuote{使每人按他在身体内所行的，}凡\BibleQuote{他所行的，无论是} \BibleQuote{善或恶。}\par

说这些话，他既以这些希望鼓舞了那些行善且受迫害的人，也以那恐惧使那些后退的人更加热切。他这样证实了他关于身体复活的言论。\Quote{因为确然，}他说，\Quote{那服务于善和恶的同一个身体，必不排除在报应之外：而是与灵魂一起，在一种情况下受罚，在另一种情况下得冠冕。}但有些异端者说，是另一个身体复活。怎么会呢？请告诉我。一个人犯罪，另一个人受罚？一个人行善，另一个人得冠冕？而你将如何回答保禄的话：\BibleQuote{我们不愿意脱去，而愿意穿上。}以及\BibleQuote{有死的被生命吞灭。}？因为他没有说，有死的或可朽坏的身体要被不可朽坏的身体吞灭；而是说，朽坏【应被】\BibleQuote{生命}吞灭。因为当同一个身体复活时，这才发生；但如果祂放弃那个身体而预备另一个，朽坏就不再被吞灭，而是继续主导。因此不是这样；而是\BibleQuote{这可朽坏的}，就是说身体，\BibleQuote{必须穿上不可朽坏的。}因为身体处于中间状态，目前在此状态，日后将在彼状态；因此先在此状态，因为不可朽坏不能被解散。\BibleQuote{因为朽坏不能承受不可朽坏，}他说，（因为，它怎能是不可朽坏的呢？）而是相反，\BibleQuote{朽坏被生命吞灭：}因为生命确实胜过朽坏，而不是朽坏胜过生命。正如蜡被火熔化，而蜡本身不能熔化火；同样，朽坏在不可朽坏之下熔化消失，却永不能胜过不可朽坏。\par

6. 让我们听从保禄的话，他说，\BibleQuote{我们必须出现在基督的审判台前；}让我们想象那法庭，并想象它现在就在眼前，要我们交账。因为我要更详细地谈论它。因为保禄看到他在谈论苦难，他不愿意再次折磨他们，所以没有详细讨论；而是简要地表达了它的严峻，\BibleQuote{各人将按他所行的领受报应，}就很快过去了。让我们现在想象它就在眼前，让我们各人凭自己的良心核算，并认为审判者已临在，一切都要显露出来。因为我们不只要站，还要被显现。你们不羞愧吗？你们不惊慌吗？但如果现在，当现实尚未临到，只是假设和想象时，我们就被良心谴责而消灭；那么当它真正来临时，当整个世界都出现，当天使和总领天使，当一排排的队伍，当所有人都同时匆忙，有人被提到云彩上，当充满战栗的行列；当号角接连吹响，那些不间断的声音时，我们该怎么办？\par

假设没有地狱，但在如此灿烂的光辉中被拒绝、羞耻地离去——那惩罚何等之大！因为即使现在，当皇帝乘车前行、随从环绕时，我们每个人默想自己的贫乏，从这景象中所得的快乐，远不及因无法分享皇帝周遭之事、不能亲近君王而感到的沮丧；那时又将如何呢？或者你认为，不被列入那队伍、不配得那不可言喻的荣耀、从那集会和无数的美善中被远远抛到某处，是轻微的惩罚吗？但那里还有黑暗、咬牙切齿、不可解开的锁链、不死之虫、不灭之火、痛苦、狭窄、像富翁那样灼烧的舌头；我们哀号，却无人听见；我们因痛苦呻吟切齿，却无人理会；我们四处张望，无处有安慰；那处于此境的人当被置于何处？有什么比那些灵魂更悲惨？更可怜？因为如果我们进入监狱，看见囚犯，有的肮脏，有的锁链加身并饥饿，有的又囚于黑暗，我们便动怜悯之心，战栗，尽一切努力以免落入那地方；当我们被牵引拖入地狱本身的刑狱时，又将如何呢？因为那些锁链不是铁的，而是永不熄灭的火；管辖我们的不是我们的同伴（常常可能软化），而是甚至连看一眼都不许的天使，他们因我们对主人的侮辱而极其愤怒。那里不像这里，能看见有人带来金钱、食物、安慰的话语，并得到慰藉；那里一切都不赦免：即使有诺厄、约伯或达尼尔，看见自己的亲属受罚，也不敢救助。因为连自然的同情到时也被废除了。因为既然有义父生恶子，义子生恶父，为使他们的快乐纯全无杂，享受美善者不会因同情所迫而悲伤，这同情，我断言，也熄灭了，他们自己也与主人一同愤慨自己的骨肉。因为如果常人看见自己的子女邪恶，便断绝关系、不认他们；那时义人更会如此。所以，不要有人若未行善事，即或有千万义人先祖，指望美善。\Quote{各人要按自己所行的，领受身体内所做的事。} 这里他似乎也暗指那些行奸淫的人，以那世界的恐惧为他们树立围墙，但不仅为他们，也为所有以任何方式犯罪的人。\par

7. 那我们也当听。你若有着欲火，以那另一种火相对，这火必即刻熄灭消失。你若想说些刺耳的言语，当想起咬牙切齿，恐惧必成为你的缰绳。你若想掠夺，当听那审判者命令说：\Quote{捆起他的手和脚，丢在外面黑暗里，} \BibleRef{玛 22:13} 你也必驱逐这贪欲。你若酗酒、常暴食，当听那富翁说：\InnerQuote{打发拉匝禄来，用指尖蘸水凉却我这灼舌，} \BibleRef{路 16:24} 却不得如愿；你必远离那病患。你若喜爱奢侈，当想那里的痛苦狭窄，必不再念此。你若是严酷残忍，当想那些灯熄的童女如何错失洞房，你必速成仁厚。你若是怠惰松懈，当想那埋藏塔冷通的人，你必猛烈如火。你若贪恋邻人的财物，当想那不死之虫，你必轻易除去此病，在一切事上行善。因为祂的命令没有一样是烦难或沉重的。那么为何祂的命令在我们看来烦难呢？由于我们的懒惰。因为若我们殷勤劳力，看似不可忍受的也轻省易行；若我们怠惰，连可容忍的也觉困难。\par

既思这一切，让我们不想那奢侈者的结局：这里确是污秽肥胖，那里是虫与火；不想那掠夺者的结局：这里忧虑恐惧，那里是不可解开的锁链；不想那好光荣者的结局：这里是奴役和虚伪，那里是不可忍受的损失和永火。因为若我们如此与自己谈论，若我们常以这些和类似的事遏制恶欲，必速速驱逐对现世之物的爱，点燃对来世之物的爱。让我们点燃它，使它炽烈。因为即使微弱的想象已带来如此快乐，那亲身经历本身将带给我们何等的喜乐！有福，三倍有福，是的，多次三倍有福，是那些享受美善的人；正如可怜的、三倍可怜的，是那些承受相反之物的人。为使我们不属后者而属前者，让我们选择德行。如此我们将获得将来的美善：愿我们众人因我们的主耶稣基督的恩宠和慈爱而获得；藉着祂，并与祂偕同，及与圣神，归于圣父，荣耀、权能、尊威，现在及永远，世世无穷。阿们。\par

\SourceRecord{Translated by Talbot W. Chambers.；Nicene and Post-Nicene Fathers, First Series；Vol. 12.；Edited by Philip Schaff.；Buffalo, NY: Christian Literature Publishing Co.,；1889.；<http://www.newadvent.org/fathers/220210.htm>.}

% structure: p0001 source=h1 target=section reason=page-title

\section{论《格林多后书》第十一篇讲道}

% structure: p0002 source=h2 target=subsection reason=relative-heading-rank; base=h2

\subsection{格后 5:11}

\BibleQuote{我们既然知道敬畏上主，就劝告人；但在天主面前是显明的，我也希望在我们的良心中是显明的。}\par

\Quote{既然知道这些事，}他说，\Quote{那可怕的审判座，我们便凡事谨慎，免得给你们把柄或绊脚石，或任何对我们恶行的虚假嫌疑。你们看这生活的严谨，和警醒灵魂的热忱！他说：\InnerQuote{我们不仅在有恶行时受谴责；即使没有行恶，却被怀疑，且有能力消除嫌疑却硬撑，也要受罚。}}\par

% structure: p0005 source=h2 target=subsection reason=relative-heading-rank; base=h2

\subsection{格后 5:12}

\BibleQuote{我们不是再向你们举荐自己，而是给你们机会因我们而夸耀。}\par

请看他是如何不断消除自夸的嫌疑。因为没有什么比一个人说自己伟大奇妙的事更令听众反感。既然他在所说的话中不得不涉及这个话题，便采用修正的说法：\BibleQuote{我们这样做是为了你们，而不是为我们自己，使你们有可夸耀的，而不是我们夸耀。}即便如此，也不是绝对的，而是因为假宗徒的缘故。因此他又加上：\BibleQuote{好叫你们在外貌夸耀、而不在内心夸耀的人面前有可夸耀的。}你们看到他是如何使听众远离那些人，并吸引他们归向自己；他表明连格林多人自己也渴望找到机会，以便能替他们说话，并向控告他们的人辩护。因为他说：\BibleQuote{我们说这些事不是为了自夸，而是为了使你们能够自由地为我们说话。}这是见证他们大爱的话：\Quote{并且不仅是为了夸耀，而是为了不使你们被拉走。}但他没有明说，而是用另一种更温和的方式措辞，不打击他们，说道：\par

\BibleQuote{好叫你们在外貌夸耀的人面前有可夸耀的。}但他也没有绝对地吩咐这样做，除非有原因存在，而是当他们自夸时才那样；因为他在一切事上都看准时机。他这样做并非为了显耀自己，而是为了阻止那些滥用此事并损害这些人的人。然而，\Quote{外貌}是什么？是可见的，是展示的。因为那些人就是如此，他们做一切事都出于爱荣誉，内心空虚，外表确有敬虔和可敬的样子，但缺乏善工。\par

% structure: p0009 source=h2 target=subsection reason=relative-heading-rank; base=h2

\subsection{格后 5:13}

2. \BibleQuote{如果我们疯狂，那是为了天主；如果我们清醒，那是为了你们。}\par

他说：\Quote{如果我们说了什么了不起的事}（他在这里称之为疯狂，正如在其他地方也称之为狂妄 \BibleRef{格后 11:1}），\BibleQuote{我们这样做是为了天主，免得你们以为我们无用而轻视我们，以致丧亡；或者如果我们说了谦卑的话，那是为了你们，使你们学习谦卑。}或者他又另有所指：如果有人认为我们疯狂，我们是从天主那里寻求赏报，因为我们为此受嫌疑；但如果有人认为我们清醒，让他从我们的清醒中获益。另外，另一种方式：有人说我们疯狂吗？我们是为了天主才如此疯狂。因此他又接着说：\par

% structure: p0012 source=h2 target=subsection reason=relative-heading-rank; base=h2

\subsection{格后 5:14}

\BibleQuote{因为天主的爱催迫我们，因为我们这样断定。}\par

他说：\Quote{不仅是将来事物的恐惧，而且已经发生的事也不允许我们懈怠或打盹，反而激励我们驱使我们为了你们而劳苦。}那么，已经发生的事是什么？\par

\BibleQuote{如果一人替众人死了，那么众人就都死了。}\Quote{既然如此，因为众人都丧亡了，他才为众人死。}因为除非众人都是死的，他不会为众人死。因为救恩的机会在此；但那里再也找不到了。因此，他说：\BibleQuote{天主的爱催迫我们，}不允许我们安歇。因为极其悲惨且比地狱还糟的是，当他行了如此大能的事，竟然有人在如此伟大的眷顾之后毫无获益。因为那爱的浩瀚极大，既为一个如此广阔的世界而死，又为处于那种状态的世界而死。\par

% structure: p0016 source=h2 target=subsection reason=relative-heading-rank; base=h2

\subsection{格后 5:15}

\BibleQuote{他们活着不再为自己活，而是为替他们死而复活的那位活。}\par

因此，如果我们不该为自己活，他说：\Quote{不要烦恼，也不要因危险和死亡临到而惊慌。}此外，他还提出一个无可辩驳的论据，说明这是当尽的义务。因为如果我们这些已死的人借着他而活，我们就该为我们所借以活的那位而活。这话看似只是一件事，但如果有人仔细考察，它其实是两件事：一是我们因他而活，二是他为我们死。这两者中任何一样本身都足以使我们归于他；但当两者联合时，想想债务有多大。不，这里还有第三件事。因为为了你们的缘故，初果也被提升，被带往天上；因此他又加上：\BibleQuote{他为我们死而复活。}\par

% structure: p0019 source=h2 target=subsection reason=relative-heading-rank; base=h2

\subsection{格后 5:16}

3. \BibleQuote{所以我们从今以后，不再按血肉认识人了。}\par

因为如果众人都死了，众人都复活了；并且他们以那种方式死了，是罪恶的暴政定他们死；但借着\BibleQuote{重生的洗和圣神的更新}\BibleRef{铎 3:5}复活了，他有理由说：\BibleQuote{我们不再按血肉认识}任何一个信友。因为即使他们在血肉之内，那血肉的生命已被摧毁，我们由圣神重生，并学会了另一种举止、规则、生活和状态，即天上的。而对此，他自己又指出基督是根源。因此他又加上：\par

\BibleQuote{即使我们曾按血肉认识基督，如今却不再这样认识他了。}\par

那么，告诉我。祂是否把肉体置之一旁，如今不再带着那身体？切勿有此想法，因为祂如今仍穿着肉身；因为这位从你们中被接升天的耶稣，将来还要怎样来。怎样呢？带着肉身，带着祂的身体。那么，祂怎能说：\BibleQuote{纵然我们曾按肉体认识基督，但如今不再这样认祂了？}\BibleRef{宗 1:11} 因为在我们身上，\Quote{按肉体}是指处于罪恶中，\Quote{不按肉体}是指不处于罪恶中；但在基督身上，\Quote{按肉体}是指祂顺从天性的情感，如渴、饿、疲乏、睡眠。因为\BibleQuote{祂没有犯过罪，口中也找不出诡诈来。}\BibleRef{伯前 2:22} 因此祂也说：\BibleQuote{你们中间谁能指证我有罪？}\BibleRef{若 8:46} 又说：\BibleQuote{这世界的领袖来了，他在我身上一无所能。}\BibleRef{若 14:30} 而\Quote{不按肉体}是指从此以后甚至从这些事中得释放，并非指没有肉身。因为祂也要带着这肉身来审判世界，祂是不受苦的、纯洁的。我们也将要到达那境地，当\BibleQuote{我们的身体}被\BibleQuote{改变得与祂光荣的身体相似。}\BibleRef{斐 3:21}\par

% structure: p0024 source=h2 target=subsection reason=relative-heading-rank; base=h2

\subsection{格后 5:17}

4. \BibleQuote{所以谁若在基督内，他就是一个新受造物。}\par

因为既然他从祂的爱中劝勉他们行善，现在他又从为他们实际所做的事上引导他们；因此他又加上：\BibleQuote{谁若在基督内，}他就是\BibleQuote{一个新受造物。}\Quote{谁若，}他说，\Quote{相信了祂，他就到了另一个创造中，因为他藉着圣神而重生。}因此，他说，我们也应当为祂而活；不仅因为我们不属于自己，不仅因为祂为我们而死，不仅因为祂使我们的初果复活，还因为我们进入了另一个生命。看，他提出了多少个行善的生活的正当理由。正因此，他也用更粗俗的名称来称呼这革新，以显示转变和改变之大。然后，他进一步阐述所说的话，并说明它如何是\Quote{一个新受造物，}就加上：\BibleQuote{旧事已过，看，一切都成了新的。}\par

什么是旧事？他是指罪恶和不虔敬，或是指一切犹太人的规条。更确切地说，他是指两者。\BibleQuote{看，一切都成了新的。}\par

% structure: p0028 source=h2 target=subsection reason=relative-heading-rank; base=h2

\subsection{格后 5:18}

\BibleQuote{一切都是出于天主。}\par

我们没有丝毫出于自己。因为罪过的赦免、义子的名分和说不尽的光荣都是由祂赐给我们的。因为他不再只从将来的事劝勉他们，也从现在的事。试想。他说，我们将要复活，进入不朽，并拥有永恒的住所；但因为现实的事比将来的事更能说服那些未能恰当相信将来事的人，他就指出他们已经领受了何等大的事，以及他们本身是何等。那么，他们领受时自己是何等？全都是死的；（因为他说，\Quote{众人都死了；}以及\Quote{祂为众人死了；}祂如此同等爱了众人；）全都是积习已深的，在恶习中衰老。但看，既有新的灵魂（因为它被洁净了），也有新的身体、新的敬拜、新的应许、盟约、生命、筵席、装束，一切都是全新的。因为代替下面的耶路撒冷，我们领受了那在上的母城\BibleRef{迦 4:26}；代替物质的圣殿，我们看见了属神的圣殿；代替石版，有肉心的版；代替割损礼，有圣洗圣事；代替玛纳，有主的身体；代替从磐石流出的水，有从祂肋旁流出的血；代替梅瑟或亚郎的棍杖，有十字架；代替应许之地，有天国；代替成千的司祭，有一位大司祭；代替无理性的羔羊，有一只属神的羔羊。他心中存着这些和类似的事，就说：\BibleQuote{一切都是新的。}但这些\BibleQuote{一切都是出于天主，}是藉着基督，是祂的恩赐。因此他又加上：\par

\BibleQuote{祂藉着基督使我们与祂自己和好，并将这和好的职务赐给了我们。}\par

因为一切美善都出于祂。因为那使我们成为朋友的那一位，祂自己也是天主赐给祂朋友的其他一切事物的原因。因为祂没有让我们继续作仇敌而赐下这些，而是使我们与祂自己和好。但当我说基督是我们和好的原因时，我是说圣父也是如此；当我说圣父赐予时，我说圣子也赐予。\Quote{因为万物是藉着祂造的；}\BibleRef{若 1:3} 这事的创始者也是祂。因为我们不是跑向祂，而是祂亲自召唤了我们。祂如何召唤了我们？藉着基督的祭献。\par

\BibleQuote{并将这和好的职务赐给了我们。}\par

在这里，他又一次表明宗徒们的尊位，说明何等重大的事交在他们手中，以及天主之爱的超乎寻常的伟大。因为即使他们不愿听从前来的使者，祂也没有恼怒，没有弃他们于不顾，而是继续亲自并藉着他人劝勉他们。谁能适当地惊叹这种关怀呢？前来和好的圣子，祂的真独生子，被杀害了，但即使如此，圣父也没有转离祂的凶手；也没有说：\Quote{我派遣我的儿子作使者，他们不但不听祂，反而杀了祂，钉死了祂，从此以后该让他们自生自灭：}而是恰恰相反，当圣子离去后，祂将这事委托给了我们；因为他说：\Quote{将这和好的职务赐给了我们。}\par

% structure: p0035 source=h2 target=subsection reason=relative-heading-rank; base=h2

\subsection{格后 5:19}

5. \BibleQuote{这就是说：天主在基督内使世界与自己和好，不将他们的过犯归到他们身上。}\par

你看见这爱超越一切表达、一切思维吗？谁是被冒犯的一方？祂自己。谁首先寻求和好？祂自己。然而，\Quote{祂派遣了圣子，祂自己并未前来。}祂确实派遣了圣子；但并非只有祂恳求，而是圣父与祂一起、并藉着祂恳求；故此他说：\Quote{天主在基督内使世界与自己和好。}因为他先前曾说过：\Quote{将这和好的职务赐给了我们；}他在这里加以纠正，说：\Quote{不要以为在这事上我们是凭自己的权柄行事：我们是仆人；而做一切的是天主，祂藉着\OriginalTerm{独生子}{Only-Begotten}使世界与自己和好。}祂如何使世界与自己和好？因为奇妙之处不在于世界成为了朋友，更在于成了朋友的方式。什么方式？就是赦免他们的罪；因为没有其他方式可行。因此他又加上：\Quote{不将他们的\OriginalTerm{过犯}{trespasses}归咎于他们。}因为倘若祂乐意追究我们所犯之事，我们都早已灭亡；因为\Quote{众人都死了。}然而尽管我们的罪如此巨大，祂不但没有要求补偿，反而和好如初；祂不但宽恕了，甚至不\Quote{归咎。}因此我们也该宽恕我们的仇敌，好使我们自己也能获得同样的宽恕。\par

\Quote{并将这和好的话交托了我们。}\par

因为我们现在不是带着令人憎恶的使命而来；而是要使众人与天主为友。因为祂说：\InnerQuote{既然他们没有被我说服，你们要继续恳求，直到说服他们为止。}因此他又加上：\par

% structure: p0040 source=h2 target=subsection reason=relative-heading-rank; base=h2

\subsection{格林多后书 5:20}

\Quote{所以我们是代基督作大使，好像天主藉着我们恳求；我们代基督请求你们：与天主和好吧。}\par

你看见他如何藉着将基督呈现为恳求者的形象而颂扬此事吗？不，不仅是基督，甚至是圣父？因为他所说的是这样的：\InnerQuote{圣父派遣圣子来恳求，并作祂派往人类的大使。当祂被杀害并离去后，我们接替了这大使的职务；我们代祂和圣父恳求你们。祂如此珍视人类，甚至舍弃了圣子，且明知祂会被杀害，并为了你们使我们成为宗徒；}因此他有理由说：\Quote{一切都是为了你们。}\BibleRef{格后 4:15}\Quote{所以我们是代基督作大使，}也就是说，代替基督；因为我们接替了祂的职务。但如果这对你来说显得重大，也请听接下来他在其中证明他们这样做不仅是代替祂，也代替圣父。因此他又加上：\Quote{好像天主藉着我们在恳求。}因为祂不仅藉着圣子自己恳求，也藉着我们这些接替了圣子职务的人。所以，\InnerQuote{不要以为你们是被我们恳求，}基督自己、基督的圣父，藉着我们恳求你们。有什么能达到这个恩宠的\OriginalTerm{满溢}{excess}呢？祂曾施予无数恩惠，却遭受侮辱；受辱后，祂不但没有追讨公义，反而赐下祂的儿子，使我们能与祂和好。那些接待祂的人没有和好，反而杀了祂。祂又派遣其他大使来恳求，虽然这些大使受差遣，但恳求的仍是祂自己。祂恳求什么呢？\Quote{与天主和好吧。}祂并没有说：\InnerQuote{使天主与你们和好；}因为不是祂怀有敌意，而是你们；因为天主从不怀有敌意。此外，他像奉命出使的大使一样力陈己见，说道：\par

% structure: p0043 source=h2 target=subsection reason=relative-heading-rank; base=h2

\subsection{格林多后书 5:21}

\Quote{祂使那\OriginalTerm{不知罪}{knew no sin}的，替我们\OriginalTerm{成了罪}{made to be sin}。}\par

\Quote{我从前不说甚么，只说你们曾凌辱了祂，祂没有亏负你们，反而恩待了你们，祂没有追究公道，祂虽首先受辱，却首先恳求；目前暂且不提这些事。难道你们不该只因祂如今为你们所做的这一件事，就应当与祂和好吗？}祂做了甚么？\BibleQuote{祂使那不识罪的，替你们成了罪。}因为即使祂没有成就任何别的事，只做了这一件，试想，为那些凌辱了祂的人赐下自己的儿子，是何等伟大的事！但如今，祂既成就了伟大的事，又使那无辜者代替犯罪者受罚。但这还不是祂所说的；祂提到了比这大得多的事。那么这是什么呢？祂说：\BibleQuote{那不识罪的}，那本身就是义的，\BibleQuote{成了罪}，即像一个罪人受定罪，像一个受诅咒的去死。\BibleQuote{因为被挂在木头上的是受诅咒的。}\BibleRef{迦 3:13}因为这样的死远大于普通的死；这也是他在别处所暗示的，说：\BibleQuote{祂服从至死，且死在十字架上。}\BibleRef{斐 2:8}因为这事不仅带来惩罚，也带来羞辱。因此，请思量祂赐给了你们何等的恩惠。因为即使是罪人为任何人死，也是一件大事；但当那经历这事的人既是义人，又为罪人而死；且不仅死，而且像一个受诅咒的；且不仅作为受诅咒的而死，还借此白白地赐给我们那些我们从未指望的宏恩；（因为他说：\BibleQuote{我们能在祂内成为天主的义；}）有什么言语，什么思想能充分领悟这些事呢？\Quote{因为义人，}他说，\Quote{祂使其成了罪人；为使罪人成为义人。}不仅如此，他反而说了远为重大的话；因为他用的词不是习性，而是本质本身。因为他不是说\Quote{使}祂成了罪人，而是说\BibleQuote{成了罪}；不仅说\BibleQuote{那没有犯罪的}，而是说\BibleQuote{连罪也不曾认识的}；还说\BibleQuote{我们}也\BibleQuote{能成为}，他没有说\BibleQuote{义人}，而是说\BibleQuote{义}，又说\BibleQuote{天主的义}。因为这义是天主的义，当我们成义不是由于行为（若是如此，就必须连一个污点也没有），而是由于恩宠，在此情况下，一切罪都被除去。这点同时使我们不致高抬自己（因为一切都是天主的白白恩赐），也教导我们那赐予的伟大。因为先前的是法律的义和行为的义，而这是\BibleQuote{天主的义}。\par

6. 默想这些事，让我们比畏惧地狱更畏惧这些话；让我们比尊重天国更尊重这些事所表达的；不要认为受苦是痛苦的，而犯罪才是痛苦的。因为如果祂不惩罚我们，我们应当自己报复自己，因为我们对我们的恩主如此忘恩负义。凡有所爱之人者，常因恋爱不遂而自杀；即使成功，若对所爱之人犯了过错，也认为不配活下去；我们冒犯了如此慈爱和温和的那一位，难道不应该把自己投入地狱之火吗？我要说一些奇怪、奇妙、也许对许多人来说难以置信的话吗？对于明白事理并按照应当爱主的方式爱主的人，在激怒了如此慈爱的主之后受到惩罚，比不受惩罚更能得到安慰。这一点可以从常例中看出。因为凡亏待了他最亲密的朋友的人，当他向自己报复并受苦时，才感到最大的解脱。为此\Person{达味}说：\BibleQuote{我牧者犯了罪，我牧者做了错事；这些羊群做了什么？愿你的手加在我和我父家身上。}\BibleRef{撒下 24:17}当他失去\Person{阿贝沙隆}时，他向自己施加了最严厉的报复，尽管他不是加害者而是受害者；尽管如此，因为他极其爱那已故者，他用痛苦折磨自己，以此安慰自己。因此，当我们得罪了那不应得罪的主时，我们也要报复自己。你没有看见那些失去亲生子女的人，他们因此捶胸顿足、撕裂头发吗？因为为了所爱的人惩罚自己带来安慰。但如果当我们没有给我们最亲爱的人造成伤害时，因他们遭遇的不幸而受苦带来安慰；那么当我们自己就是挑衅和犯错的人时，受惩罚岂不更加是一种解脱？不受惩罚岂不更惩罚我们？人人都能看出这一点。如果有人像应当爱基督那样爱祂，他就知道我说的是什么：即使祂宽恕了，他也不愿不受惩罚；因为你最大的惩罚就是激怒了祂。我知道我说的话不会被许多人相信；但事实确如我所说。如果我们像应当爱基督那样爱祂，当我们犯罪时，我们就会惩罚自己。因为对于任何有所爱之人者，因激怒了所爱之人而受苦并非不愉快；而最不愉快的是激怒了所爱的人。如果后者在发怒时不惩罚，反而更折磨了爱他的人；但如果他要求赔偿，反而安慰了他。因此，让我们不要害怕地狱，而要害怕得罪天主；因为得罪天主比地狱更痛苦：这比一切都更沉重，这比一切都更沉重。为了让你明白这是什么事，请考虑我所说的。如果有一位国王，看见一个强盗和罪犯受惩罚，就将他心爱的儿子、独生子和真儿子赐予被杀；并将死亡和罪责从那人转移到他的儿子身上（而他的儿子本身并无这样的品性），以便既能拯救那被定罪的人，又能洗刷他的恶名；然后，如果再使他升到高位，而在这样拯救并提升他到那说不出的荣耀之后，竟被那承受如此恩惠的人所冒犯：那人不若还有一点理智，岂不宁愿选择千万次死亡，也不愿显得犯下如此忘恩负义之罪吗？那么，让我们也常常这样思想，为我们对恩主所行的冒犯而痛苦呻吟；不要因为祂虽被冒犯却耐心忍受而自恃；反而正因此要充满懊悔。因为即使在人间，被人打右颊而转左颊给他，比起打他千万拳更报复了他；被人辱骂而不还口，反而祝福，比向他倾泻千万辱骂更重创了他。如果对人，当我们侮辱人而遇到忍耐时，我们感到羞愧；那么对天主，那些继续不断犯罪却不遭灾祸的人，更应当惧怕。因为，甚至恶事临到他们头上，那不可言喻的惩罚已为他们积存。那么，记住这些事，让我们首先要害怕犯罪；因为这就是惩罚，这就是地狱，这就是千万种灾祸。让我们不仅害怕，还要逃避它，努力不断讨天主喜悦；因为这就是天国，这就是生命，这就是千万种福乐。这样，我们在此世也能预尝天国和未来的福乐；愿我们都能达到，借着我们的主耶稣基督的恩宠和慈爱；愿光荣、权能、尊崇归于祂及圣父及圣神，从现在直到永远，无穷世之世。阿们。\par

\SourceRecord{Translated by Talbot W. Chambers.；Nicene and Post-Nicene Fathers, First Series；Vol. 12.；Edited by Philip Schaff.；Buffalo, NY: Christian Literature Publishing Co.,；1889.；<http://www.newadvent.org/fathers/220211.htm>.}

% structure: p0001 source=h1 target=section reason=page-title

\section{《格林多后书》第十二篇讲道}

% structure: p0002 source=h2 target=subsection reason=relative-heading-rank; base=h2

\subsection{《格林多后书》6:1-2}

我们与天主一同工作，也劝你们不要白白接受天主的恩宠。因为祂说：\Quote{在悦纳的时候，我俯允了你；在救恩之日，我援助了你。}\par

既然他说了：\Quote{天主劝勉，我们是使者，恳求你们与天主和好}，免得他们变得懈怠，他便在此再次警醒并激励他们，说：\Quote{我们劝你们不要白白接受天主的恩宠。}他说：\Quote{因此，我们不要因为祂恳求并派了使者就松懈；反之，正因如此，我们更要努力取悦天主，收集属灵的财货。}正如他上面所说的：\Quote{基督的爱催迫我们}\BibleRef{格后5:14}，那就是催促、驱策、推动我们，\Quote{免得你们在如此深情的关怀之后，因懈怠而不显高尚，错失如此大的福分。因此，不要因为祂派了人来劝勉，就认为这将是常态。直到祂第二次来临之前，祂一直在恳求，只要我们还在世上；但此后便是审判与惩罚。}因此，他说：\Quote{我们被催迫。}\par

他不仅从福分的伟大和祂的慈爱，也从时间的短暂来不断督促他们。因此，他在别处也说：\BibleQuote{因为现在我们的救恩比当初更近了。}\BibleRef{罗13:11}又说：\BibleQuote{主已经近了。}\BibleRef{斐4:5}但在此，他做得更多。因为他不只从剩余时间短少来激励他们，更从这是唯一可得救恩的时间来激励。\par

因为他说：\BibleQuote{看，现在正是悦纳的时候；看，现在正是救恩之日。}因此，我们不要错过这有利时机，而要展示与恩宠相称的热忱。正因如此，我们也向前迈进，既知时间的短暂与合宜。故此他也说：\Quote{我们一同工作，也劝你们。}\Quote{一同工作}是与你们一同工作；\Quote{因为我们与你们同工，而非与天主同工，我们为祂作使者。祂一无所缺，而救恩全部归于你们。}即使是指与天主同工，他也不拒绝这种理解；因为在别处他说：\BibleQuote{我们是天主的同工}\BibleRef{格前3:9}，意思是，以此方式拯救人。又说：\Quote{我们也劝勉。}因为祂在恳求时，并非仅仅恳求，而是陈述祂的公义要求：即祂赐下了祂的圣子，那义者、不知罪者，为使我们罪人成为义，使祂成为罪。拥有这些要求，又是天主，祂显示了如此良善。但我们所求的，是你们接受恩惠，不拒绝恩赐。因此，请听从我们，\Quote{不要白白接受恩宠。}免得他们以为单是相信那呼召者就是\Quote{和好}，他添加了这些话，要求那关乎生活的热忱。因为一个从罪中得释放、成为朋友的人，若再沉溺于从前之事，就是重归敌意，并且\Quote{白白接受恩宠}，就生活而论。因为从\Quote{恩宠}中，如果我们生活不洁，我们对救恩毫无益处；甚至有害，因为我们有更大的罪过加重，在如此知识和如此恩赐之后，又退回从前的恶习。然而，他暂时没有提及这一点，以免使作品变得严厉；只说我们毫无益处。然后他又提及一个预言，敦促并强迫他们振作起来，以抓住自己的救恩。\par

他说：\Quote{因为}，祂说：\par

在悦纳的时候，我俯允了你，\newline{}在救恩之日，我援助了你；\newline{}\Quote{看，现在正是悦纳的时候；看，现在正是救恩之日。}\par

\Quote{悦纳的时候。}这是什么？就是恩赐的时候、恩宠的时候，那时不要求我们结算罪债，也不追讨惩罚；反而，除了得释放，我们还享有万般美物：公义、圣化等等。为了获得这个\Quote{时候}，我们本应付出多少劳苦！可是看，我们没有丝毫劳苦，它已来临，带来了对从前一切的赦免。故此，祂也称它为\Quote{悦纳的}，因为祂不仅接纳了在万般事上犯罪的人，而且不仅接纳，还提拔他们至最高尊荣；正如一位君王驾临时，不是审判之时，而是恩宠与宽赦之时。因此，祂也称它为悦纳的。所以，当我们还在赛场上，还在葡萄园里劳作，还有第十一个小时，让我们亲近祂，显示生命；因为这也容易。因为在如此大恩赐倾泻、如此大恩宠之时，那努力争先的人，必将早得奖赏。因为在此世君王的情形中，在他们的节日里，当他们身着执政官礼服时，那奉献微小礼物者必得厚赏；但在他们坐堂审判之日，则需要许多严谨、许多筛选。所以，我们也要在这恩赐之时努力争先。这是恩宠之日，天主的恩宠；因此，我们甚至容易获得冠冕。因为若我们在背负如此大恶时，祂仍接纳并释放了我们；那么，当我们脱离一切并尽自己本分时，祂岂不更要接纳我们吗？\par

2. 然后，正如他常做的，即将自己置于他们面前，并让他们从此处取榜样，他在这里也这样做了。故此，他又加上：\par

% structure: p0011 source=h2 target=subsection reason=relative-heading-rank; base=h2

\subsection{《格林多后书》6:3}

\BibleQuote{不给人任何绊脚石，免得我们的职务受人诋毁，}劝诫他们，不仅是从考虑\Quote{时间}上，也从那些与他们一同成功劳苦的人身上。且看他是多么没有骄傲。因为他并没有说\InnerQuote{看看我们是如何如此这般}，而是，为当前的目的，他只是为除去毁谤而叙述自己的行为。他提到了无可指责生活的两个要点：\Quote{无}在\Quote{任何}事上。他没有说\InnerQuote{毁谤}，而是远为轻的\BibleQuote{绊脚石}；即不给人任何把柄来谴责、定罪，\BibleQuote{使我们的职务不受诋毁}；即无人能抓住它。他又说，不是\InnerQuote{使它不受控告}，而是使它没有丝毫过失，也没有任何人有能力在任何方面指责它。\par

% structure: p0013 source=h2 target=subsection reason=relative-heading-rank; base=h2

\subsection{\BibleRef{格后 6:4}}

\BibleQuote{但在一切事上，我们举荐自己为天主的仆役。}\par

这是远为伟大的。因为摆脱控告，与在一切事上显现为\BibleQuote{天主的仆役}这样的品格，并不是同一回事。因为摆脱控告与备受赞美也不是一回事。他没有说\Quote{显现}，而是\BibleQuote{举荐}，即\InnerQuote{证明}。然后他也提到他们是如何成为这样的。那么是如何的呢？他说：\BibleQuote{在极大的忍耐中}，奠定了那些善事的基础。所以他不是仅仅说\Quote{忍耐}，而是\BibleQuote{极大的}，并也显示出它有多么大。因为忍受一两件事并不算什么。但他甚至在话语中加上了试炼的暴风雪：\BibleQuote{在磨难中，在困境中。}这是对磨难的加强，当邪恶不可避免，并且人身上仿佛有一种几乎无法摆脱的不幸的必要性。\BibleQuote{在忧患中。}他要么是指饥饿和其他必需品的艰辛，要么单纯是指他们试炼中的艰辛。\par

% structure: p0016 source=h2 target=subsection reason=relative-heading-rank; base=h2

\subsection{\BibleRef{格后 6:5-6}}

\BibleQuote{在鞭打中，在囚禁中，在颠沛流离中。}\par

然而，这些中的每一件本身都是难以忍受的：仅仅被鞭打，仅仅被捆绑，以及因迫害而无法在任何地方安顿（因为这就是在\Quote{颠沛流离}中），但是当所有这些同时袭来时，想想他们需要怎样的灵魂。然后，除了外在的事物外，他还提到了他自己所加上的。\par

\BibleQuote{在劳苦中，在守夜中，在禁食中；在纯正中。}\par

但此处他藉\BibleQuote{纯正}所意指的，要么是再次指贞洁，要么是一般的纯洁，要么是不朽，要么甚至是自由地宣讲福音。\par

\BibleQuote{在知识中。}什么是\Quote{在知识中？}在来自天主的智慧中；那是真正的知识；不像那些看似有智慧、夸耀自己熟悉外邦学问，却在这方面有缺失的人。\par

\BibleQuote{在忍耐中，在慈惠中。}因为这也是高尚灵魂的伟大标记，尽管在各方面被激怒和刺激，仍能以忍耐承受一切。然后，为显示他是如何成为这样的，他补充说：\par

\BibleQuote{在圣神内。} \Quote{因为在祂内，} 他说，\Quote{我们做这一切善工。} 但请注意，他是在何时提到了圣神的助佑。在陈述了自身方面之后。此外，我觉得他在这里还说了另一件事。那是什么呢？即 \Quote{我们既充满了丰盛的圣神，也由此证明我们的宗徒职份，因为我们被算为堪当领受神恩的人。} 因为即使这也是恩宠，但仍是因他本人藉其善工和劳苦吸引了那恩宠。若有人断言，除上述外，他还表明自己在运用圣神恩赐时也未给人任何绊脚石，那也不会误解他的意思。因为他们中那些领受[说语言]恩赐、因而自高自大的人，曾受过责备。因为一个人即使在领受圣神的恩赐时，也可能不正当地运用它。他说：\Quote{但我们并非如此，} 他又说，\Quote{反而在圣神内，即在恩赐上，我们也无可指责。}\par

\BibleQuote{在无伪的爱德中。} 这是这一切善事的原因；这使他成为他所是的；这也使圣神常与他同在，藉着圣神的助佑，他所有的事都做得正确。\par

% structure: p0025 source=h2 target=subsection reason=relative-heading-rank; base=h2

\subsection{\BibleRef{格后 6:7}}

\BibleQuote{在真理之言中。}\par

这是他在许多地方所说的事，即 \Quote{我们并没有诡诈地处理天主的圣言，也没有使它变质。}\par

那正是他常做的，不归功于自己，而将一切归于天主，并将他所做的一切正确之事都归于祂，他在这里也是如此。因为他既说出了伟大的事，并宣称自己在一切事上显示了一个无可指责的生命和崇高的智慧，他便将这些归于圣神和天主。因为他所说的那些事并非平庸。因为若连一个在安静中生活的人要做到正确、无可指责尚且是难事，那么想想那个被如此巨大诱惑折磨、却仍照耀所有方面的人，他具有何等的灵性！而且，他不仅经历了这些，还经历了比这些更多的事，正如他接下来所提到的。真正奇妙的是，不是他在如此巨浪中航行却无可指责，也不是他高贵地承受一切，而是他甚至喜悦地承受一切。这一切，他都藉着接下来的话来向我们显明，说：\par

\BibleQuote{藉着左右正义的武器。}\par

3. 你看到他那灵魂的自持和坚韧的精神了吗？因为他表明，苦难不仅是不会击倒人的武器，反而能加固并使人更强健。他将那些似乎痛苦的事称为 \Quote{左}，因为这些事会带来赏报。那么他为何如此称呼呢？或是顺应大众的观念，或是因天主命令我们祈祷不陷于诱惑。\par

% structure: p0031 source=h2 target=subsection reason=relative-heading-rank; base=h2

\subsection{\BibleRef{格后 6:8}}

\BibleQuote{藉着光荣与羞辱，藉着恶名与美名。}\par

你说什么？你说你享受尊荣，并以此为大吗？他说：\Quote{是的。}何以如此？因为忍受羞辱实为大，但分沾尊荣却不需要刚强的灵魂。不，需要一个刚强且极大的灵魂，使享受尊荣者不致跌倒折断颈项。因此他既在这事上夸口，也在那事上夸口，因为他在两者中都同样闪耀。但这如何是义德的武器？因为教师们受人尊敬会引导许多人归向虔敬。此外，这也是善工的证明，且光荣天主。这又是天主奇妙安排的一个实例，即祂用相反的事引入福音。试想：保禄被捆绑？这也是为福音的缘故。因为他说：\BibleQuote{我所遭遇的事，反而使福音更见进展，以致多数弟兄因我的锁链而信赖主，越发勇敢地、毫无畏惧地宣讲天主的圣言。}\BibleRef{斐 1:12}再者，他享受尊荣？这又使他们更加积极。\BibleQuote{藉着恶名和美名。}因为他不仅高尚地忍受那些临到身体的事，即\Quote{苦难}和他所列举的一切，也忍受那些触及灵魂的事；因为后者也不轻率地搅扰人。至少耶肋米亚在忍受了许多试探后，在这事上屈服了；当他受辱骂时，他说：\BibleQuote{我不再以祂的名讲论，也不再以祂的名发言。}\BibleRef{耶 20:9}达味也在多处抱怨辱骂。依撒意亚在诸多事情后，也对此劝诫说：\BibleQuote{你们不要怕人的辱骂，也不要因他们的毁谤而惊慌。}\BibleRef{依 51:7}基督也对祂的门徒说：\BibleQuote{几时人为了我而辱骂迫害你们，捏造一切坏话毁谤你们，你们是有福的。你们欢喜踊跃罢！}\BibleRef{玛 5:11}祂说：\BibleQuote{因为你们在天上的赏报是丰厚的。}（参\BibleRef{路 6:23}）祂又说：\BibleQuote{并且踊跃欢喜。}但祂不会使赏报如此之大，除非这争战是大的。因为在折磨中，身体也与灵魂一同感受痛苦；因为痛苦既是身体的，也是灵魂的；但这里只是灵魂的痛苦。至少许多人仅因这些而跌倒，丧失了自己的灵魂。甚至对约伯而言，朋友的责难比虫子和疮痍更难以忍受。因为没有什么，没有什么比一个能刺伤灵魂的言语对受苦者更难忍受了。因此他既提到危险和辛劳，也提到这些，说：\BibleQuote{藉着光荣和羞辱。}至少许多犹太人也因从众人得来的光荣而不愿相信。他们害怕，不是怕受惩罚，而是怕被逐出会堂。因此祂说：\BibleQuote{你们既然彼此寻求光荣，而不寻求出于唯一天主的光荣，怎能相信我呢？}\BibleRef{若 5:44}我们也可以看见许多人确实藐视了一切危险，却败给了光荣。\par

\BibleQuote{似乎是迷惑人的，却是真诚的。}\par

这就是\BibleQuote{藉着恶名和美名}。\par

% structure: p0036 source=h2 target=subsection reason=relative-heading-rank; base=h2

\subsection{格后 6:9}

\BibleQuote{似乎不为人知，却是人所共知的。}这就是\BibleQuote{藉着光荣和羞辱}。因为对于一些人，他们是人所共知、大受追捧；而另一些人则故意不认识他们。\par

\BibleQuote{似乎要死，看，我们却活着。}\par

如同被判处死刑、定了罪；这本身也是羞辱的事。但他说这话，是为了显示天主不可言喻的大能和他们的忍耐。因为就那些谋害我们的人而言，我们死了；这也就是众人都以为的；但靠着天主的帮助，我们逃脱了危险。然后，为了表明天主为何允许这些事，他加上：\BibleQuote{似乎受管教，却不被处死。}\par

表明他们从试探中所获得的益处，甚至在赏报之前已是大得，且敌人被迫对他们有益。\par

% structure: p0041 source=h2 target=subsection reason=relative-heading-rank; base=h2

\subsection{格后 6:10}

\BibleQuote{似乎忧愁，却常常喜乐。}\par

因为外人确实怀疑我们绝望了；但我们不理睬他们；是的，我们充满喜乐。他不是只说\BibleQuote{喜乐}，而是加上其恒久性，因为他说\BibleQuote{常常喜乐}。有什么能比得上这种生命呢？在其中，尽管如此众多的危险袭来，喜乐却变得更大。\par

\BibleQuote{似乎贫穷，却使许多人富足。}\par

有些人确实认为这里指的是属灵的财富；但我认为也包括物质的；因为他们在这方面也是富足的，以一种新的方式，众人的家门都为他们敞开。\par

\BibleQuote{似乎一无所有，却拥有万有。}\par

这怎么可能？其实不如说，相反的情况怎么可能？因为拥有许多的人一无所有；而一无所有的人却拥有众人的财富。不仅在这里，在其他方面，相反的事也由它们的反面产生出来。但若你惊奇一个人一无所有如何能拥有万有，让我们把这人本身摆在中间，他掌管世界，不仅是他们的财物之主，甚至是他们的眼目之主。他说：\BibleQuote{假使可能的话，你们也会把自己的眼睛挖出来给我。}\BibleRef{迦 4:15}\par

他说这些话，是为教训我们不要因多数人的意见而困扰，虽然他们称我们为欺骗者，虽然他们不认识我们，虽然他们认为我们是被定罪的、被判死刑的，要我们忧愁、贫穷、一无所有，要我们（这喜乐的人）沮丧；因为太阳对于瞎眼的人并不明亮，快乐对于疯狂的人也不可理解。只有信友才是这些事的正确判断者，他们不为同样的事而喜乐或痛苦，如同其他人一样。因为若有一个对游戏一无所知的人，看见一个拳击手身上有伤却戴着冠冕，他会因伤口而认为他痛苦，而不明白冠冕给他带来的喜乐。所以这些人，因为他们知道我们所受的痛苦，却不知道我们为何受这些苦，自然怀疑除了这些之外别无他物；他们确实看见了角力与危险，却看不见奖赏、冠冕和竞赛的主题。那么，当保禄说\Quote{似乎一无所有，却拥有万有}时，他所拥有的\Quote{万有}是什么？是暂时的事，属神的事。因为那被各城如天使般接待的人，他们甚至愿意剜出眼睛给他，\BibleRef{迦 4:14}那为他们甘愿舍命的人，\BibleRef{罗 16:4}他如何不拥有他们的一切？但若你也要看见属神的事，你会发现他在这些事上尤其富有。因为那与万有之王如此亲近，甚至与天使之主分享不可言说之事的人，\BibleRef{格后 12:4}他如何不比所有人更富足，并拥有万有？否则魔鬼不会如此服从他，痛苦与疾病也不会如此逃遁。\par

5. 因此，当我们为基督的缘故受苦时，我们不仅应当高尚地忍受，也应当喜乐。若我们禁食，要如同享受奢侈般雀跃；若我们受辱，要如同受赞美般欢舞；若我们花费，要如同获利般感受；若我们施舍穷人，要如同自己领受：因为不这样想的人不会愿意施舍。所以，当你想要施舍时，不仅在施舍中看这一点，也在每一种德行中，不要只计算辛劳的严酷，也要计算奖赏的甜蜜；更要紧的，是这场角力之主耶稣基督；这样你将欣然进入竞赛，并终生生活在喜乐中。因为没有什么比良心的平安更能带来喜乐。\par

因此，保禄虽然每日受伤，却喜乐欢跃；但今天的人们，虽连他所做的一点影子也未忍受，却忧愁哀叹，这无非是因为他们没有一个充满天上智慧的心。因为，请告诉我，为何哀叹？因为你贫穷，缺乏必需品？当然，你更应该为此哀叹，[不是]因为你哭泣，不是因为你贫穷，而是因为你心胸狭窄；不是因为你没有金钱，而是因为你如此重视金钱。保禄每日死去，却不哭泣，反而喜乐；他不断与饥饿搏斗，却不忧伤，反而因此夸耀。而你，只因为没有储存全年的粮食，就忧伤自责？\Quote{是的，}他回答，\Quote{因为他只需照顾自己的需要，而我还要照顾仆人、儿女和妻子。}相反，他恰恰不是只照顾自己的需要，而是照顾全世界的需要。你确实[需要照顾]一个家庭，但他需要照顾耶路撒冷那么多穷人，马其顿的，各处贫穷的，以及那些施舍的人，和那些领受的人。他对世界的关怀是双重的：既使他们不缺乏必需品，又使他们在属神的事上富足。你的饥饿儿童不会像所有信徒的事那样使他忧伤。我为何说\Quote{信徒}？因为他甚至对非信徒也并非无忧，而是如此被吞噬，甚至愿意为他们成为诅咒；而你，即使饥荒猖獗万倍，也绝不会为任何人选择死亡。你确实照顾一个女人，但他照顾普世的教会。因为他说：\Quote{我对众教会的挂虑。}\BibleRef{格后 11:28}那么，人啊，你要戏弄到何时，将自己与保禄相比，而不停止你这极大的卑怯？因为应当哭泣的，不是我们贫穷时，而是我们犯罪时；这值得哀叹，而其他一切甚至可笑。\Quote{但是，}他说，\Quote{这还不是使我忧伤的全部；还有，某某大权在握，而我不受尊敬、被排斥。}这又是什么？因为有福的保禄在许多人眼中也似乎不受尊敬、被排斥。\Quote{但是，}他说，\Quote{他是保禄。}显然，不是事情本身的性质，而是你意志的软弱导致你的沮丧。所以，不要哀叹你的贫穷，而要哀叹你这种心志；不，不要哀叹你自己，而要悔改；不要寻求金钱，而要追求那比千万金钱更令人面容开朗的东西，即智慧与德行。因为哪里有这些，贫穷便无害；哪里没有这些，金钱也无益。请告诉我，一个人富有却灵魂卑贱，有什么好处？你对自己的哀叹，还不如那富人对自己的哀叹，因为他没有万有的财富。即使他不像你那样哭泣，但请揭开他的良心，你会看见他的哀泣与悲叹。\par

你们愿意我给你们展示你们自己的财富，好叫你们不再认为那些有钱的人是有福的吗？你们看见这天空、太阳、这颗明亮且照耀遥远的星，那悦人眼目的，难道不也是为众人所共有的吗？穷人和富人不是同样享有它吗？星斗的冠冕和月轮，不也是平等地留给众人吗？是的，更确切地说，若我不得不说些奇妙的事，我们穷人比他们更享有这些。因为他们大多沉溺于酗酒，在宴饮和沉睡中度日，甚至没有觉察到这些事物，他们总是生活在荫蔽之下；但穷人比任何人更能享受这些元素的美妙。此外，若你考察那到处弥漫的空气，你会看见穷人享有它更清新、更丰沛。因为行路人和农夫比城市居民更能享受这些奢侈；而在城市居民中，工匠又比那些终日酗酒的人更能享受。大地又如何？不也是留给众人共有的吗？他说：\Quote{不。}你怎么这样说？告诉我。\Quote{因为富人即使在城里，也为自己得到几块\OriginalTerm{普莱斯拉}{plethra}之地，并在四周筑起长篱；在乡下，他又为自己割取许多份。}那么，当他割取时，他独享这些吗？绝不然，尽管他极力争辩。因为农产品他不得不分配给众人，他为你耕种谷物、葡萄酒和油，处处为你服务。而那些长篱和建筑物，在他花费无数、劳苦奔波之后，他为你预备，供你使用，只从你那里收取一点银钱作为如此大服务的报酬。在澡堂和任何地方，都可以看到同样的情况：富人轻松地享受这一切。他对大地的享受并不比你多，因为他确实不能吃十张肚腹，而你只吃一张。\Quote{他岂非享用更昂贵的肉食？}确如所言，这并非什么了不起的优势；然而，即使在这里，我们也会发现你占了上风。因为这种昂贵之所以被你看作嫉妒之事，是因为其中的乐趣更大。然而在穷人的情况下，乐趣更大；不仅乐趣，健康也是如此；而富人唯一的优势是使自己的体质更虚弱，并收集更丰富的疾病源泉。因为穷人的饮食完全合乎自然，而富人的饮食因过度而导致败坏与疾病。\par

6. 但若你愿意，让我们也用一个例子来看同样的事。假如必须生起一个火炉，一个人往里扔许多数不清的丝绸衣服和细麻布来点火，另一个人则扔橡木和松木的木头，那这人比那人有什么优势呢？没有，反而更糟。但怎样的呢？（因为没有什么能阻止我们以另一种方式转用同一个比喻）如果一人扔木头，另一人用尸体来生火，你愿意站在哪个炉边？是木头的，还是尸体的？显然，是木头的。因为那火燃烧自然，对观者也是赏心悦目的景象；而尸体火炉因蒸汽、汁液、烟雾和骨头的恶臭，会赶走每个人。你耳闻时岂不战慄，岂不厌恶那火炉？富人的肚腹也类似。因为在其中，你会发现比那火炉更多的腐烂、臭气、污秽的体液，因为到处都充满因暴饮暴食而导致的消化不良。由于天然热力不足以消化全部，反被食物所淹没，它们在上面冒着烟，产生的 unpleasantness 很大。那么，把这种肚腹比作什么呢？但请不要因我所说的话而恼怒，若我说的不真实，请反驳我。那么，把它们比作什么呢？因为已经说过的还不足以表明他们的悲惨处境。我又找到了另一个类比。是什么呢？就像在下水道里，有垃圾、粪肥、干草、残茬、石头、泥土的堆积，经常发生堵塞；然后污秽的洪流从上面溢出。同样，这些人的肚腹也如此。因为下面被堵住，大部分恶劣的液体向上喷涌。但穷人却不是这样，他们就像那些涌出清流的泉水，浇灌花园和游乐场；同样，他们的肚腹也免于这类多余之物。但富人的肚腹，或者更确切地说，奢侈者的肚腹却不是这样；他们充满体液、痰、胆汁、败血、腐烂的\OriginalTerm{风湿（体液）}{rheums}，以及其他此类东西。因此，没有人若常过奢侈生活，能忍受片刻；他的生命将在不断的疾病中度过。因此，我很乐意问他们，肉食是为了什么目的？是为叫我们毁灭，还是为滋养？是为叫我们生病，还是为强壮？是为叫我们健康，还是为患病？非常清楚，是为了滋养，而不是为了毁灭。然而，奢侈者却由于丰盛的食物而招致疾病。但穷人不是这样；相反，他通过简朴的饮食为自己获得健康、活力和力量。因此，不要因贫穷（健康之母）而哭泣，反而要欢跃；若你想富有，就要轻视财富。因为富有，不是拥有金钱，而是不想要金钱，这才是真正的富足。若我们能成就此事，我们在此世将比所有富人都更富足，并且在来世将获得那未来的美善，愿我们众人能达到，借着我们的主耶稣基督的恩宠和对人的爱，祂与圣父及圣神，永受光荣、权能、尊崇，自今至永世，阿们。\par

\SourceRecord{Translated by Talbot W. Chambers.；Nicene and Post-Nicene Fathers, First Series；Vol. 12.；Edited by Philip Schaff.；Buffalo, NY: Christian Literature Publishing Co.,；1889.；<http://www.newadvent.org/fathers/220212.htm>.}

% structure: p0001 source=h1 target=section reason=page-title

\section{第十三篇讲道：论格林多后书}

% structure: p0002 source=h2 target=subsection reason=relative-heading-rank; base=h2

\subsection{格后 6:11-12}

\BibleQuote{我们的口向你们张开，你们格林多人啊，我们的心是宽大的，你们在我们里面并不狭窄，而是在你们自己的情感里狭窄了。}\par

他详述了自己的考验和苦难，因为\BibleQuote{在忍耐中，}他说，\BibleQuote{在患难中，在急需中，在困苦中，（第4-5节）在鞭打中，在监禁中，在骚乱中，在劳苦中，在不眠中；}并表明这是一件大好事，因为\BibleQuote{像是忧苦的，}他说，\BibleQuote{却常常喜乐；像是贫穷的，却使许多人富足；像是一无所有，却拥有万有；}\BibleRef{格后 6:10}并称这些事为\BibleQuote{兵器，}因为\BibleQuote{像是受责罚的，}他说，\BibleQuote{却不被处死；}并由此彰显了天主的丰盛眷顾和大能，因为他说：\BibleQuote{那卓绝的能力是属于天主，而非出于我们；}\BibleRef{格后 4:7}并叙述了他的劳苦，因为他说：\BibleQuote{我们常带着祂的死；}并表明这是复活的确证，因为他说：\BibleQuote{那耶稣的生命也彰显在我们有死的肉身中；}\BibleRef{格后 4:10}并论及他成了何等事的分施者，以及受托付何等事，因为他说：\BibleQuote{我们是基督的使者，}\BibleRef{格后 5:20}\BibleQuote{好像天主藉着我们劝勉一样；}并论及他是什么事的执事，即：\BibleQuote{不是属文字的，而是属圣神的；}\BibleRef{格后 3:6}并表明他不仅因此应受尊敬，也因他的考验，因为\BibleQuote{感谢天主，}他说，\BibleQuote{祂常带领我们凯旋。}他现在也打算责备他们，因为他们对他心意不正。但虽有此意，他并不立即触及此事，而是先讨论这些事。因为责备者即使因自己的善行而应受尊敬，但当他显示自己对被责备者的爱时，就会使言词不那么冒犯。因此，宗徒也暂离自己考验、劳苦和斗争的题目，转而谈论他的爱，并以此击中他们的要害。那么，他爱的标志是什么？\BibleQuote{我们的口向你们张开，你们格林多人。}这算什么样的爱之标志？这些话究竟有何含义？\Quote{我们不能忍受向你们缄默，}他说，\Quote{反而总是渴望并切愿与你们说话和交谈；}这是爱人的常情。因为正如握手之于身体，语言的交流之于灵魂。除此之外，他还暗示了另一件事。那是怎样的呢？即\Quote{我们讲论毫无保留。}因为此后他打算责备，便请求原谅，以大胆责备他们作为深爱他们的证明。此外，加上他们的名字也是大爱、温暖和亲情的标志；因为我们习惯不断重复所爱之人的名字。\par

\BibleQuote{我们的心扩大了。} 因为正如温暖之物惯于膨胀，同样，扩大也是爱的工作。这既开了\Person{保禄}的口，又扩大了他的心。因为，\Quote{我不仅以口爱，} 他说，\Quote{但我也有合一的内心；因此我坦率说话，用我全部的口，用我全部的理智。} 因为没有什么比\Person{保禄}的心更宽广，他以一个人可能对自己所爱对象所怀有的全部炽热爱着所有信友；他的爱没有被分开因而削弱，而是完整地停留在每个人身上。既然甚至对不信者，\Person{保禄}的心也拥抱了全世界，那么对信友如此，又有什么可惊奇的？因此他没有说\Quote{我爱你}，而是更强调地说，\BibleQuote{我们的口是敞开的，我们的心是扩大的，}我们将你们所有人容纳在里面，而且不仅如此，还有极大的空间。因为被爱的人在爱他的人心中无拘无束地行走。因此他说，\BibleQuote{你们在我们心中并不狭窄，而是你们在自己心中狭窄。} 看，这指责是以宽容施行的，正如深爱者所常做的。他没有说\Quote{你们不爱我们}，而是说\Quote{不在同等程度上}，因为他不想过于刺痛他们。事实上，处处可见他对信友的炽热，从每一封书信中选词。他对罗马人说，\BibleQuote{我渴望见你们；} 并且，\BibleQuote{我屡次打算到你们那里去；} 并且，\BibleQuote{无论如何，如今终于或许得蒙顺利到你们那里去。} \BibleRef{罗 1:11} 他对迦拉达人说，\BibleQuote{我的孩子们，我为你们再受产痛。} \BibleRef{迦 4:19} 对厄弗所人又说，\BibleQuote{为此我跪伏膝盖}为你们。\BibleRef{弗 3:14} 对斐理伯人说，\BibleQuote{我的希望、喜乐、欢跃的冠冕是什么？难道不也是你们吗？} 并且他说他常在心中、在锁链中带着他们。\BibleRef{斐 1:7} 对哥罗森人说，\BibleQuote{我愿意你们知道我为你们以及一切没有在肉身见过我面的人多么奋斗，为使你们的心受安慰。} \BibleRef{哥 2:1} 对得撒洛尼人说，\BibleQuote{就如乳母抚育自己的孩子，同样，我们既深情恋慕你们，就乐意将福音不仅给你们，连自己的灵魂也给你们。} \BibleRef{得前 2:7} 对\Person{弟茂德}说，\BibleQuote{纪念你的眼泪，使我充满喜乐。} \BibleRef{弟后 1:4} 对\Person{弟铎}说，\BibleQuote{致我亲爱的儿子；} \BibleRef{铎 1:4} 对\Person{费肋孟}也是如此。\BibleRef{费 1:1} 对希伯来人也写了许多类似的事，并不停地恳求他们说，\BibleQuote{还有一点点时候，那要来的就来，并不迟延。} \BibleRef{希 10:37} 如同母亲对闹别扭的孩子一样。对格林多人自己他说，\BibleQuote{你们在我们心中并不狭窄。} 但是他不只说他爱他们，也说他被他们所爱，为借以更赢得他们。而且为他们见证这点，他说，\Person{弟铎}来了，并\BibleQuote{告诉我们你们的渴望、你们的哀恸、你们的热情。} \BibleRef{格后 7:7} 对迦拉达人说，\BibleQuote{若可能，你们会挖出眼睛给我。} \BibleRef{迦 4:15} 对得撒洛尼人说，\BibleQuote{我们如何进到你们当中。} \BibleRef{得前 1:9} 对\Person{弟茂德}也说，\BibleQuote{纪念你的眼泪，使我充满喜乐。} \BibleRef{弟后 1:4} 并且在他的书信中，处处可以找到他向门徒们作此见证，既证明他爱他们，又证明他被爱，但并非平等。而在这里他说，\BibleQuote{虽然我更丰盛地爱你们，却较少被爱。} \BibleRef{格后 12:15} 这，不过是在近结尾处；但此刻更激烈地说，\BibleQuote{你们在我们心中并不狭窄，而是你们在自己心中狭窄，} \Quote{你们接受一个人，} 他说，\Quote{但我接受一整个城市，如此众多的人口。} 而且他没有说\Quote{你们不接受我们}，而是说\Quote{你们是狭窄的}，暗示同样的事，但带着宽容，不深深刺痛他们。\par

% structure: p0006 source=h2 target=subsection reason=relative-heading-rank; base=h2

\subsection{格林多后书 6:13}

\BibleQuote{现在作为同样性质的回报（我如同对孩子说话），你们也要扩大。}\par

然而，这并非平等的回报：先被爱，而后去爱。因为即使一个人贡献同等数量，他在次序上仍然居后。\Quote{但是尽管如此，} 他说，\Quote{我仍不打算严格计较；如果你们在从我这里接受了最初的恩惠之后，只是表现出同等程度，我也就满足了。} 然后，为表明这样做甚至是一种债务，而且他的话毫无奉承之意，他说，\BibleQuote{我如同对孩子说话。} \BibleQuote{如同对孩子说话} 是什么意思？\Quote{如果身为你们的父亲，我希望被你们爱，这并非什么重大要求。} 再看他的智慧和谦虚。他没有在此提到他为你们所冒的危险、他的辛劳和牺牲，尽管他有许多可以讲述（他是多么谦逊！）；而是提到他的爱，并以此要求被爱；\Quote{因为，} 他说，\Quote{我是你们的父亲，我极度为你们燃烧。} (因为)常常，当一个人列举自己对所爱之人的恩惠时，这本身会让被爱者感到冒犯；因为这似乎是一种责备。因此\Person{保禄}没有这样做；而是说，\Quote{如同孩子，爱你们的父亲}，这更出于天性；也是每一位父亲应得的。然后，为免他似乎是为自己说话，他表明他邀请这种爱也是为了他们的益处。因此他补充道：\par

% structure: p0009 source=h2 target=subsection reason=relative-heading-rank; base=h2

\subsection{格林多后书 6:14-16}

\BibleQuote{你们不要与不信者不等地同负一轭。}\par

祂没有说：\Quote{不可与不信者混杂}，反而严厉地对待他们，如同违反正道者，说：\Quote{不可偏离}，他说，\BibleQuote{正义与不法有什么相通？}在随后的内容中，他进行了比较，不是在他自己的爱与那些败坏他们的人的爱之间，而是在他们的高贵与另一方的羞辱之间。因为这样，他的言论变得更加庄严，更符合他自己，也会更容易赢得他们。正如一个人对一个轻视父母、沉溺于邪恶之人的儿子说：\Quote{你在做什么，孩子？你轻视你的父亲，而宁愿要那些充满无数罪恶的污秽之人吗？你不知道你比他们好多少、可敬多少吗？}因为这样，他更容易[更易于]使他们脱离他们的社交，而不是通过表达对他父亲的钦佩。因为如果他真的说：\Quote{你不知道你的父亲比他们好多少吗？}他不会产生那么大的效果；但如果他省略了对父亲的提及，而是把他自己呈现在他们面前，说：\Quote{你不知道你是谁，他们是谁吗？你不记得你自己的高贵出身和血统，以及他们的耻辱吗？你与他们，那些盗贼、那些奸夫、那些骗子，有什么相通？}通过用这些对他自己的赞美来提升他，他会很快使他们准备好与他们断绝关系。因为前一种说法，他实际上不会欣然接受，因为对他父亲的赞誉是对他自身的责备，当他被显示为不仅令父亲伤心，而且是这样的父亲；但在后一种情况下，他不会有那样的感受。因为没有人会不愿意受到赞美，因此，伴随着对听众的赞美，责备就变得易于消化。因为听众变得柔和，充满高尚的思想，并鄙弃与那些人为伍。\par

但不仅这一点值得赞美他，他这样进行了比较，而且他还想象了另一件更伟大、更令人惊叹的事：首先，他以询问的形式展开话语，这适用于清晰和公认的事物，然后通过术语的快速连续和众多来展开它。因为他不仅使用了一个、两个或三个术语，而是几个。此外，代替人物，他使用事物的名称，并在此描绘了高度的美德和极端的恶行；并表明它们之间的差异是巨大且无限的，以至于甚至不需要证明。他说：\BibleQuote{正义与不法有什么相通？}\par

\BibleQuote{光明与黑暗有什么相通？}（15—16节）\BibleQuote{基督与\OriginalTerm{贝里雅耳}{Beliar}有什么和谐？或者，信者与不信者有什么份额？或者，天主的殿与偶像有什么一致？}\par

你看见他如何只使用名称，以及如何恰当地达到劝阻的目的。因为他没有说\Quote{疏忽正义}，而是更强硬的\Quote{不法}；他也没有说那些属于光明的和那些属于黑暗的；而是使用了那些不能接受其对立面的对立面本身，\Quote{光明与黑暗}。他也没有说那些属于基督的与那些属于魔鬼的；而是，那相距更远的，基督与\OriginalTerm{贝里雅耳}{Beliar}，他用希伯来语这样称呼那个背逆者。\BibleQuote{或者，信者与不信者有什么份额？}这里，终于，为了不显得仅仅是在进行对恶行的谴责和对美德的赞扬，他也提及了人物而不具体指名。他没有说\Quote{相通}，而是谈论了赏报，使用了\Quote{份额}这个词。\BibleQuote{天主的殿与偶像有什么一致？}\par

现在他说的是这个：你们的君王与祂没有任何共同之处，因为\BibleQuote{基督与贝里雅耳有什么和谐？}事物之间也没有共同之处，因为\BibleQuote{光明与黑暗有什么相通？}因此，你们也不应该如此。他首先提到他们的君王，然后是他们自己；通过这一点最有效地将他们分开。然后，说了\BibleQuote{天主的殿与偶像}，并宣告\BibleQuote{因为你们是天主生活的殿}，他不得不也附上对此的见证，以表明这事并非奉承。因为赞美的人若不也出示证据，甚至显得在奉承。那么他的见证是什么？因为，\par

他说：\BibleQuote{我要住在他们中间，并在他们中行走。}他重复了\BibleQuote{我要住在}，如同在殿中，\BibleQuote{并在他们中行走}，以此表示对他们的更丰富的依恋。\par

\BibleQuote{他们将成为我的子民，我将成为他们的天主。}\Quote{什么？}他说，\Quote{你怀有天主在你内，却奔向那些与祂毫无共同之处的人？这怎能得赦免？要记住谁在你内行走、居住。}\par

% structure: p0018 source=h2 target=subsection reason=relative-heading-rank; base=h2

\subsection{\BibleRef{格后6:17}}

\BibleQuote{因此，你们要从他们中间出来，与他们分离，不要触摸不洁之物；我必收纳你们，主说。}\par

祂没有说\Quote{不要做不洁之事}；而是，要求更大的严格，他说\Quote{甚至不要触摸}，\Quote{也不要靠近它们}。但是什么是肉体的污秽？奸淫、淫乱、各种放荡。什么是灵魂的污秽？不洁的思想，如邪淫的眼神、恶意、欺骗，以及诸如此类的事。因此，他希望你们在这两方面都洁净。你看见赏报有多大吗？脱离邪恶，与天主合一。也请听下文。\par

% structure: p0021 source=h2 target=subsection reason=relative-heading-rank; base=h2

\subsection{\BibleRef{格后6:18}}

\BibleQuote{我要作你们的父亲，你们要作我的儿女，主说。}\par

你看见从起初先知就预告了我们现在的高贵出生，即因恩宠而得的再生吗？\par

% structure: p0024 source=h2 target=subsection reason=relative-heading-rank; base=h2

\subsection{\BibleRef{格后7:1}}

\BibleQuote{所以，亲爱的，我们既有这些应许。}\par

什么应许？就是我们要成为天主的殿、儿女；祂在我们内居住、行走；我们是祂的子民；祂是我们的天主和父亲。\par

\BibleQuote{让我们洁净自己，除去肉身和灵魂的一切污秽。}\par

我们不要触摸不洁之物，因为这是肉体的洁净；也不要触摸玷污灵魂之物，因为这是灵魂的洁净。然而他并不满足于此，还加上：\par

\Quote{在敬畏天主中成全圣德。}\newline{} 因为不触及不洁之物并不能使人洁净，还需要别的东西才能使我们成为圣洁：认真、谨慎、虔诚。他说得好：\Quote{在敬畏天主中。}因为有可能不是在敬畏天主中成全贞洁，而是出于虚荣。同时他通过说\Quote{在敬畏天主中；}暗示了另一件事，即成全圣德的方式。因为即使情欲是一股强横的力量，如果你用敬畏天主占据它的地盘，你便制止了它的狂乱。\par

4. 现在他在这里所说的圣德不仅指贞洁，而且指免除各种罪恶，因为圣洁的人就是纯洁的人。一个人将成为纯洁，不仅如果他免于淫乱，而且如果也免于贪财、嫉妒、骄傲和虚荣，是的，尤其要避免虚荣，因为在一切事情上都应当避免它，但在施舍上更应如此；因为如果施舍有了这种弊病，它就不再是施舍，而是炫耀和残忍。因为当你不是出于慈悲，而是出于炫耀而施与时，这种作为不仅不是任何施舍，甚至是一种侮辱；因为你使你的弟兄公开蒙羞。因此，施舍不在于给钱，而在于出于慈悲而给钱。因为人们也在戏院里给钱，既给娼妓男孩，也给其他在舞台上的人；但这样的行为不是施舍。那些侮辱娼妓身体的人也给钱；但这并非慈爱，而是无礼对待。虚荣的人也是如此。因为正如侮辱娼妓的人付钱换取那种侮辱一样，你也向接受你施舍的人索取代价，即你对他的侮辱以及使你和他都背上恶名。此外，损失是不可言喻的。因为如同野兽和疯狗扑向我们一样，这种恶疾和这种残忍掠夺了我们的善行。因为这样的做法是不人道和残忍的；是的，甚至比这更严重。因为残忍的人确实不会给乞求者；而你做得更多：你阻止那些想要施舍的人。因为当你炫耀你的施舍时，你既降低了接受者的声誉，又阻止了将要施舍的人（如果他是一个粗心的人）。因为此后他不会再给他，理由是他已经收到，所以没有需要；是的，他常常甚至会控告他，如果他在收到后又前来乞讨，会认为他厚颜无耻。那么，当你在天主面前行施舍，却又寻求同仆的眼目，并且违反天主所立的禁止这些事的法律时，这算什么施舍呢？\par

我本希望也把这事延伸到其他主题上，就是斋戒和祈祷，并显示虚荣在哪里也是有害的；但我记得在上一篇讲道中，我留下了一个必要的要点没有完成。那个要点是什么？我在论及健康和快乐时曾说，穷人在今世的事物上胜过富人；但这一点显示得不够清楚。那么，今天让我们来证明：不但在今世的事物上，就是在更高的事物上，他们也占优势。因为什么引导人进入天国呢？是财富还是贫穷？让我们亲耳聆听主在天上论及那些人时所说的话：\BibleQuote{它比骆驼穿过针孔还容易，为富人进入天国。} \BibleRef{玛 19:24}但对穷人则相反：\BibleQuote{你若愿意成为成全的，去，变卖你所有的，分施给穷人，然后来跟随我；你将在天堂拥有宝藏。} \BibleRef{玛 19:21}但若你愿意，让我们看看双方的说法。\BibleQuote{路是狭窄而狭小的}祂说，\BibleQuote{那通往生命的}\BibleRef{玛 7:14}那么，谁走这条狭窄的路呢？是享乐的人，还是贫穷的人？是无拘无束的人，还是背负万钧重担的人？是散漫放荡的人，还是深思焦虑的人？但这些论证有什么必要呢？最好还是亲身看看具体人物。\Person{拉匝禄}是贫穷的，是的，非常贫穷；而当他躺在门口时，从他身边经过的那个人是富有的。那么，哪一个进入了天国，在\Person{亚巴郎}的怀抱中享受喜乐？哪一个被烈火焚烧，连一滴水也得不到？但有人说：\Quote{许多穷人将要丧亡，而许多富人将享受那些不可言喻的美好事物。}不，恰恰相反，可以看到，得救的富人很少，而穷人却多得多了。因为，试想一想，准确衡量财富的障碍和贫穷的缺陷（或者更确切地说，这些障碍和缺陷既不是财富的，也不是贫穷的，而是属于拥有财富或贫穷的人；尽管如此），让我们至少看看哪一个是更有效的武器。那么，贫穷似乎有什么缺陷呢？撒谎。财富呢？骄傲，万恶之母；骄傲也使魔鬼成为魔鬼，它原本不是那样的。再者，\BibleQuote{爱钱是万恶的根源。}\BibleRef{弟前 6:10}那么，谁靠近这个根源呢？是富人还是穷人？不是非常明显是富人吗？因为一个人积聚得越多，就越贪求。虚荣再次损坏了成千上万的好事，而富人又居住在这附近。\Quote{但是}有人说，\Quote{你没有提到穷人的坏事，他的痛苦，他的困境。}不，但这些既为富人所共有，而且富人承受得比穷人更多；所以那些看似贫穷的坏事，其实是两者共有的；而财富的坏事，则只有财富才有。有人说：\Quote{但穷人因缺乏必需品而行了许多可怕的事，那又怎么说呢？}然而，没有一个穷人，没有，一个也没有，因缺乏而行了许多可怕的事，像富人为积聚更多和不失去已有的积蓄所做的那样。因为穷人不会像富人渴求多余之物那样急切地渴求必需品；他也没有对方那样大的力量去实施邪恶。如果富人既有意愿又有能力，那么显然他会犯下这样的罪，而且犯得更多。穷人不会像富人那样因担心失去已有的财产、且因尚未拥有所有人的财产而颤抖焦虑。既然他接近虚荣、傲慢和爱钱——万恶之根，那么他除非表现出极大的智慧，还有什么得救的希望呢？他怎能走狭窄的路呢？因此，我们不要到处传播众人的观念，而要查考事实。因为在金钱上我们不信赖别人，而诉诸数字和计算；但在计算事实时，我们却轻易被别人的观念所左右，这岂不是很荒谬吗？更何况我们拥有精确的天平、直角尺和规矩——神圣的法律的宣示？因此，我劝勉并恳求你们所有人，不要理会这个人或那个人对这事怎么想，而是从圣经中查考这一切；在得知什么是真正的财富之后，让我们追求它们，好让我们也能获得永恒的美善。愿我们所有人都能获得，借着我们的主\Person{耶稣基督}对人的恩宠和慈爱，与祂一起，向圣父和圣神，荣耀、权能和尊威，从今时直到永远，永世无尽。阿们。\par

\SourceRecord{Translated by Talbot W. Chambers.；Nicene and Post-Nicene Fathers, First Series；Vol. 12.；Edited by Philip Schaff.；Buffalo, NY: Christian Literature Publishing Co.,；1889.；<http://www.newadvent.org/fathers/220213.htm>.}

% structure: p0001 source=h1 target=section reason=page-title

\section{论\BibleBook{格林多}{后书}第十四篇讲道}

% structure: p0002 source=h2 target=subsection reason=relative-heading-rank; base=h2

\subsection{\BibleBook{格林多}{后书} 7:2–4}

\Emph{你们要容纳我们：我们没有亏负过任何人，没有败坏过任何人，没有占过任何人的便宜。我说这话，并不是要定你们的罪；因为我以前说过[正如我在上面也说过]，你们是在我们心里，一同死去，一同活着。}\par

他再次将话题转向爱，以缓和责备的严厉。因为他此前已指责他们，说他们固然是可爱的，却没有同样地爱他，反而脱离了他的爱，与那些有害的败类混在一起。于是他又缓和了责备的激烈语气，说道：\Quote{你们要容纳我们}，也就是说，\Quote{要爱我们}；他恳求他们赐予一个并不沉重的恩惠，这恩惠使施予者比领受者更受益。他没有用\Quote{爱}这个字，而是用更令人怜恤的\Quote{容纳}，意思是：\Quote{谁驱逐了我们？}他说，\Quote{谁将我们从你们心中赶出去了？我们怎么会在你们心中被局限？}因为他在上面说过：\Quote{你们在情感上是局限的}。这里他更清楚地表明，说：\Quote{你们要容纳我们}；这样他又赢得了他们。因为没有什么能像被爱者知道爱他者极其渴望他的爱那样，能产生爱。\par

\Quote{我们没有亏负过任何人。}请注意，他如何又不提自己所施的恩惠，而是以另一种方式措辞，既少冒犯，又更犀利。同时，他也暗示了假宗徒，说道：\Quote{我们没有亏负过任何人，没有败坏过任何人，没有占过任何人的便宜。}\par

什么是\Quote{败坏}？即迷惑人，正如他在别处所说：\BibleQuote{免得你们的心意，如同蛇用诡诈诱惑了厄娃一样，被败坏。}\BibleRef{格后 11:3}\par

\Quote{我们没有占过任何人的便宜；}我们没有掠夺或谋害任何人。他暂时不说\Quote{我们曾如此这般地恩待了你们}，而是用更令他们羞愧的措辞说：\Quote{我们没有亏负过任何人，}意思是：\Quote{即使我们丝毫没有恩待你们，你们也不该离弃我们；因为你们不能指控我们任何大小过错。}然后，因为他深感责备的分量，又再次缓和了语气。他既没有完全沉默（否则就不能唤醒他们），也没有让言语的严厉不加缓和（否则就会过分伤害他们）。他怎么说呢？\par

\BibleRef{格后 7:3}：\BibleQuote{我说这话，并不是要定你们的罪。}\par

这如何显明呢？他接着说：\BibleQuote{因为我以前说过，你们是在我们心里，一同死去，一同活着。}\BibleRef{格后 7:3}这是最大的亲情：即使被轻视，他仍选择与他们同死同生。\Quote{因为你们不仅在我们心里，}他说，\Quote{而且正如我所说的。因为人既能爱，也能规避危险，但我们并非如此。}请看这里也有说不尽的智慧。他没有提他为他们所做的事，以免显得再次责备他们，而是为将来许诺。\Quote{因为倘若危险来临，}他说，\Quote{为了你们的缘故，我已准备好承受一切；无论是死是生，对我来说本身都不算什么，但你们在哪里，那里对我就更可取：若你们在其中，死比生更好，生比死更好。}不过，死显然证明爱；而生，即使是那些不是朋友的人，谁不选择呢？那么为什么宗徒将这当作一件大事来提呢？因为这是极大的事。许多人确实同情不幸中的朋友，但当他们尊荣时，却不与他们一同喜乐，反而嫉妒他们。\Quote{但我们并非如此：无论你们在患难中，我们不怕与你们同受厄运；还是你们在顺境中，我们不被嫉妒所伤。}\par

2. 在他不断重复这些话之后，说：\Quote{你们并没有在我们身上受到限制；}和\Quote{你们是在自己的情感上受到限制；}和\Quote{请给我们腾出空间；}和\Quote{你们也要被扩大；}和\Quote{我们没有亏负过任何人；}所有这些似乎都在谴责他们：请看他是如何用另一种方式缓和这种严厉，说道：\Quote{我对你们大有胆量。}\Quote{因此我敢于说这些事，}他说，\Quote{并非要通过我的话谴责你们，而是出于我极大的胆量，}他又进一步表明这一点，说道：\Quote{我以你们大为夸耀。}\Quote{你们不要以为，}他说，\Quote{因为我这样说，就好像我完全谴责了你们；（因为我对你们极其自豪，并夸耀你们；）但既是出于温柔的关怀，又是渴望你们在德行上取得更大的进步。}因此，在严厉责备之后也对希伯来人说：\BibleQuote{可是，亲爱的诸位！我们虽这样说，但对你们，却确信你们将有更好的表现，且更接近救恩。我们切愿你们每一位都表现同样的热心，以达成望德的圆满，一直到底。}\BibleRef{希 6:9} 同样在这里也是如此，他说：\Quote{我以你们大为夸耀。}\Quote{我们在别人面前夸耀你们，}他说。你看到他给予了何等真实的安慰吗？\Quote{而且，}他说，\Quote{我不只是夸耀，而且是大大地夸耀。}因此他补充了这些话：\Quote{我充满了安慰。}什么安慰？\Quote{来自你们的安慰，因为你们经过改造，以你们的品行安慰了我。}这是爱者的考验：既埋怨不被爱，又怕因过度埋怨而造成痛苦。因此他说：\Quote{我充满了安慰，我洋溢着喜乐。}\Quote{但是，}有人说，\Quote{这些话似乎与之前的话矛盾。}然而并非如此，甚至完全和谐。因为这些使前者获得良好的接受；其中所传达的赞美使那些责备的益处更加真诚，通过悄悄去除其中令人痛苦的部分。因此他使用这些表达，但极其真诚和恳切。因为他没有说：\Quote{我充满了喜乐；}而是说：\Quote{我充盈；}或者更确切地说，不是\Quote{充盈}，而是\Quote{盈溢}；这样又再次显示了他的渴望，即使他如此被爱而欢喜雀跃，他仍然不认为自己得到了应有的爱，也没有得到完全的回报；他如此贪得无厌，出于对她们超乎寻常的爱。因为被自己所热切爱的人以任何程度所爱带来的喜悦，因我们极其爱他们而变得巨大。所以这再次证明了他的爱。关于安慰，他说：\Quote{我充满了；}\Quote{我已得到了所亏欠我的；}但关于喜乐，他说：\Quote{我盈溢；}也就是说，\Quote{我曾为你们沮丧；但你们已充分辩解并提供了安慰：因为你们不仅消除了我悲伤的根源，甚至增加了喜乐。}然后展示其伟大，他不仅通过说\Quote{我在喜乐中盈溢}来宣告，还补充说\Quote{在我们一切的苦难中。}\Quote{因为从你们而来的喜悦如此巨大，甚至没有被如此大的患难所遮蔽，而是通过其自身伟大的超卓，克服了困扰我们的悲伤，使我们感觉不到它们。}\par

% structure: p0011 source=h2 target=subsection reason=relative-heading-rank; base=h2

\subsection{格林多后书 7:5}

\BibleQuote{因为甚至当我们到了马其顿时，我们的肉身没有得到安息。}\par

因为既然他说了\Quote{我们的患难}，他既解释其种类，又用言语将其放大，以显示从他们所得的安慰和喜悦何等大，因为它驱散了如此大的悲伤。\BibleQuote{但我们四面受敌。}\par

如何四面受敌？因为\BibleQuote{外有争斗}，来自不信者；\BibleQuote{内有恐惧}；因为信徒中的软弱者，恐怕他们被诱离。因为不仅是在格林多人中发生这些事，在其他地方也是如此。\par

% structure: p0015 source=h2 target=subsection reason=relative-heading-rank; base=h2

\subsection{格林多后书 7:6-7}

\BibleQuote{然而，那安慰谦卑人的天主，藉着弟铎的到来安慰了我们。}\par

因为既然他在所说的话中为他们作了大的见证，为显得不是在奉承他们，他引证弟兄弟铎为证人，他是在第一封书信之后从他们那里来到保禄这里，向他报告他们悔改的详情。但是，我请你们注意，他如何处处强调弟铎的到来。因为他在前面也说：\BibleQuote{此外，当我为了福音来到特洛阿时，我的心灵没有安息，因为我没有找到我的弟兄弟铎；}\EditorialNote{参阅 格后 2:12-13} 而在这里又说：\BibleQuote{我们得了安慰，}他说，\BibleQuote{因着弟铎的到来。}因为他也渴望使此人在他们中得到信任，并使他极其亲切可爱。请看他是如何兼顾这两件事的。因为一方面通过说\Quote{我的心灵没有安息，}他展示了他的美德之伟大；另一方面通过说，在我们的苦难中，他的到来足以安慰；然而\BibleQuote{不仅因着他的到来，也因着他在你们身上所得到的安慰，}他使这个人受到格林多人的爱戴。因为没有什么比说某人好话更能产生和巩固友谊。而他所见证的弟铎正是如此；当他说因着他的到来，他欣然给了我们翅膀，关于你们他报告了这样的事。因此他的到来使我们欢喜。因为我们并非只\Quote{因他的到来，也因着他在你们身上所得到的安慰}而欢喜。而他如何得了安慰？因你们的德行，因你们的善行。因此他也补充说，\par

\BibleQuote{他告诉我们你们的渴望，你们的哀恸，你们对我的热诚。}这些使他欢喜，他说，这些安慰了他。你看他是如何表明他也是他们热切的恋人，因为他把他们的好消息视为对自己的安慰；当他来到保禄那里时，他夸耀，如同为了自己的好事。\par

请注意他以何等热烈的言辞叙述这些事：\Quote{你们的渴望，你们的哀恸，你们的热忱。}因为很可能的，他们会哀恸悲伤，为何真福保禄如此不悦，为何他长久疏远他们。因此他没有简单地说泪水，而是说\Quote{哀恸}；也没有说愿望，而是说\Quote{渴望}；也没有说愤怒，而是说\Quote{热忱}；然后又提到\Quote{对他（保禄）的热忱}，这热忱是他们在关于那犯奸淫的人以及关于那些指控他的人身上所表现的。\Quote{因为，}他说，\Quote{你们在收到我的信后，就激愤起来，火焰般燃烧。}由于这些事，他满溢喜乐，由于这些事，他充满安慰，因为他使他们有了感觉。然而，在我看来，这些话不仅是为了缓和先前所说的，也是为了鼓励那些在这些事上行事有德的人。因为虽然我想某些人对先前的指控是有过错的，不配得到这些赞扬；然而，他没有将他们区分开来，而是使赞扬和指责都成为共同的，让听众的良心去选择属于他们自己的部分。这样，一方不会得罪人，另一方则引导他们趋向更大的心灵热忱。\par

4. 如今，受责备的人也应有这样的感受；他们应当这样哀伤悲痛；这样渴慕他们的导师；这样，比父母更甚地寻求他们。因为藉着父母，生活得以延续，但藉着导师，美好的生活得以实现。他们应当这样承受父亲的责备，这样因那些犯罪的人而与他们的长上同感。因为责任并不全在他们身上，也在你们身上。\newline{}因为如果犯罪者感到他确实被父亲责备了，但被弟兄们奉承了，他就变得更加心安理得。但当父亲责备时，你也当一同发怒，无论是出于对兄弟的关心，还是为了与父亲的义怒联合；只是你表现出的热忱要巨大；并且你要哀恸，不是因为受责备，而是因为他犯了罪。\newline{}但如果我建造，你拆毁，我们除了劳苦还有什么益处？\BibleRef{德 34:23}不仅如此，你的损失还不止于此，你还给自己招来惩罚。因为阻止伤口愈合的人，所受的惩罚不小于造成伤口的人，甚至更大。因为伤害和阻止受伤者愈合，并不是同等的过犯；因为后者必然导致死亡，而前者则不必然。\newline{}我现在对你们说这些，是要你们在你们的统治者义怒时，与他们的愤怒联合；当你们看见任何人受责备时，你们都要比教师本人更远离他。让犯罪的人害怕你们，甚于害怕他的统治者。因为他如果只害怕教师，他就会轻易犯罪；但如果他畏惧这么多的眼睛，这么多的舌头，他就会更安全。因为如果我们不这样做，我们将遭受最严厉的惩罚；如果我们行这些事，我们将分享从他悔改而来的益处。\newline{}因此，让我们这样行；如果有人会说：\Quote{对你的兄弟仁慈吧，这是基督徒的本分。}让他受教：那发怒（对他）的人才是仁慈的，而不是那过早让他松懈、甚至不允许他意识到自己过犯的人。因为，请告诉我，哪一个更怜悯一个发烧且神志不清的人：是那个让他躺在床上、捆绑他、不让他吃不合适的食物和饮料的人，还是那个允许他狂饮烈酒，吩咐他自由行动，一切如同健康人一样的人？难道后者不是加重了疾病，而那个看似仁慈的人却矫正了它吗？\newline{}确实，我们在这件事上也应当这样判断。因为仁慈的一部分，不是凡事迁就病人，也不是奉承他们不合时宜的欲望。在格林多人中，没有人像保禄那样爱那个犯奸淫的人，他命令将他交给撒殚；没有人像那些鼓掌讨好他的人那样恨他；而结果也证明了这一点。因为那些人确实使他骄傲并加重了他的炎症；但（宗徒）却降低了它，并一直不放弃，直到他把他带到完全的健康。而那些人则增添了已有的祸患，他则根除了最初存在的祸患。\newline{}因此，让我们也学习这些人道的法则。因为如果你看见一匹马正冲向悬崖，你会给马戴上嚼子，用力勒住它，并频繁鞭打它；虽然这是惩罚，但惩罚本身却是安全的母亲。同样，对待犯罪的人也要这样行。捆绑那犯过的人，直到他平息了天主的忿怒；不要让他松绑，免得他被天主的怒气捆绑得更紧。如果我捆绑，天主就不会锁住；如果我不捆绑，那解不开的锁链在等着他。\BibleRef{格前 11:31}因为如果我们判断自己，我们就不会被判断。\newline{}不要以为这样做是出于残忍和不人道；相反，这是出于最高的温柔、最巧妙的医术和极大的关怀。但是，有人说，他们已经受了很长时间的惩罚。多长时间？告诉我。一年、两年、三年？然而，我不要求这个，不要求时间长短，而是要求灵魂的改善。那么就请展示这一点：他们是否心痛，是否悔改，那么一切都完成了；因为如果没有这个，时间就没有益处。因为我们不询问伤口是否多次包扎，而是询问包扎是否有用。因此，如果有用，即使在短时间内，就不应再保留包扎；但如果无用，即使过了十年，仍应保留；并且让释放的期限由被捆绑者的益处来决定。\newline{}如果我们这样谨慎对待自己和他人，不顾虑人的荣誉和羞辱；而是铭记那未来的惩罚和耻辱，以及更重要的，对天主的激怒，努力施行悔改的良药：我们就会立刻获得完全的健康，并获得将来的美善；愿我们众人因我们的主耶稣基督的恩宠和慈爱，得以获得这一切；愿光荣、权能、尊崇归于圣父，偕同圣神，从今直到永远，万世无疆。阿们。\par

\SourceRecord{Translated by Talbot W. Chambers.；Nicene and Post-Nicene Fathers, First Series；Vol. 12.；Edited by Philip Schaff.；Buffalo, NY: Christian Literature Publishing Co.,；1889.；<http://www.newadvent.org/fathers/220214.htm>.}

% structure: p0001 source=h1 target=section reason=page-title

\section{《格林多后书》讲道集 第十五篇}

% structure: p0002 source=h2 target=subsection reason=relative-heading-rank; base=h2

\subsection{\BibleBook{格林多}{后书} 7:8-9}

\BibleQuote{\Emph{因此}，虽然我以那封信使你们忧愁，我并不后悔，虽然我曾后悔过。}\par

\Person{保祿}接着为他的书信道歉，因为（罪过既已纠正）温和地对待他们已无危险；他并显示此事的好处。因为他从前也这样做过，当他说：\BibleQuote{因为我怀着许多忧愁和心中痛苦，给你们写了信，不是为使你们忧愁，而是为使你们知道我向你们所怀有的爱。}\BibleRef{格后 2:4} 他现在也这样做，用更多的话确立同一要点。他并没说：\Quote{我从前确实后悔，但现在我不后悔}，而是怎样？\BibleQuote{我现在不后悔，虽然我曾后悔。}\Quote{即使我所写的，}他说，\Quote{超出了责备的应有分寸，以致让我后悔；然而从中产生的巨大益处，不容我后悔。}他说这话，不是好像他责备他们过了分，而是为更加赞扬他们。\Quote{因为你们所显示的改善是那么巨大，}他说，\Quote{即使我碰巧打击你们过重，以致我甚至谴责自己，我现在也因结果而赞扬自己。}如同对待小孩子，当他们经历了痛苦的医治，例如切开、烧灼或苦药，过后我们不怕安抚他们；\Person{保祿}也是如此。\par

\BibleQuote{因为我看见那封信使你们忧愁，不过是一时的。现在我不因你们忧愁而欢喜，而是因你们忧愁以致悔改。}\par

他说了\BibleQuote{我不后悔}之后，也说明理由；引述他的信所带来的益处；并巧妙地借说\BibleQuote{不过是一时的}来自我辩护。因为那痛苦的事是短暂的，但有益的事是永久的。而自然接着要说的是\Quote{即使它使你们一时忧愁，却使你们永远喜乐和受益。}但他没有这样说：而是在提及益处之前，他转而再次赞扬他们，并证明他多么关心他们，说：\BibleQuote{现今我不因你们忧愁而欢喜，（因为我从你们忧愁中得到什么益处呢？）而是因你们忧愁以致悔改，}这忧愁带来了益处。因为父亲看见儿子在刀下，并不因他受苦而欢喜，而是因他被治愈而欢喜；\Person{保祿}也是如此。但请观察他如何将此事的一切成就归功于他们自己；而将任何痛苦的事都归于那封信，说：\BibleQuote{它曾使你们一时忧愁；}而由此产生的益处，他却说成是他们自己的善功。因为他没有说\Quote{那封信纠正了你们}，虽然事实如此；而是说\BibleQuote{你们忧愁以致悔改。}\par

\BibleQuote{因为你们按照天主的心意而忧愁，好使你们因我们而一无所损。}\par

你看到无法言喻的智慧吗？\Quote{因为我们若不这样做，}他说，\Quote{我们就伤害了你们。}他肯定成就属于他们，但若他保持沉默，损害则属于他自己。因为他们如果可能因严厉的责备而被纠正，那么，若我们不严厉责备，我们就伤害了你们；这伤害不仅关乎你们，也关乎我们。因为正如那不给商人航行必需品的人，就是他造成损失；同样，我们若不给你们那悔改的机会，就给你们造成了损害。你看见不责备犯罪的人，对主人和门徒都是损害吗？\par

% structure: p0009 source=h2 target=subsection reason=relative-heading-rank; base=h2

\subsection{\BibleBook{格林多}{后书} 7:10}

2. \BibleQuote{因为按照天主的心意而忧愁，能产生悔改以致救恩，这是一种不令人后悔的悔改。}\par

\BibleQuote{因此，}他说，\BibleQuote{虽然我以前因未见果实与收获而感到后悔，但它们如此之大，我现在并不后悔。}因为天主所悦纳的忧愁就是如此。然后他对此加以哲理阐发，表明忧愁并非在所有情况下都是痛苦之事，只有当它属于世俗时才是。什么是世俗的？如果你为金钱、为名誉、为逝者而忧愁，这一切都是世俗的。因此它们也带来死亡。因为为名誉而忧愁的人会嫉妒，并常常导致毁灭：该隐的忧愁就是如此，厄撒乌的忧愁也是如此。这里他所说的世俗忧愁，是指对忧愁者有害的那种。因为只有在罪过方面的忧愁才是有益的；这一点可以这样清楚地看出：为损失财富而忧愁的人无法弥补那损失；为逝者而忧愁的人不能使死者复生；为疾病而忧愁的人不仅不会康复，反而加重病情；只有为罪过而忧愁的人，才能从忧愁中获得一些益处，因为他使自己的罪过减少并消失。既然药物是为这种情况准备的，那么只有在这种情况下它才有效力并显示其益处；在其他情况下甚至是有害的。\Quote{然而，}有人说，\Quote{该隐忧愁是因为他没有蒙天主的悦纳。}并非为此，而是因为他看见自己的兄弟在荣耀中受尊崇；如果他为此悲伤，他本应效仿并与兄弟一同喜乐；但事实却是，他忧愁表明他的忧愁是世俗的。但达味并非如此，伯多禄并非如此，任何义人也都并非如此。因此，当他们为自己的罪过或他人的罪过而忧愁时，他们蒙受了悦纳。然而，还有什么比忧愁更令人压抑呢？但是当它符合天主的心意时，它比世上的喜乐更好。因为世上的喜乐最终归于虚无；而那种忧愁\BibleQuote{产生导致救恩的悔改，这救恩是永不后悔的。}其中令人赞叹的是，一个如此忧愁的人永远不会后悔，而这正是世俗忧愁的一个特别特征。因为还有什么比亲生儿子更令人后悔的呢？还有什么比这种死亡更沉重的悲伤呢？然而，那些在极度悲伤中无法忍受任何人并疯狂捶打自己的父亲们，过了一段时间后会后悔，因为他们过度悲伤；因为他们这样做不仅对自己毫无益处，反而加重了痛苦。但天主所悦纳的忧愁并非如此；它拥有两个优点：一是人不会因所忧愁的事而受谴责，二是这种忧愁最终带来救恩；这两点世俗忧愁都欠缺。因为他们既忧愁以致受害，又在剧烈忧愁之后谴责自己，从而显示出最大的证据，表明他们这样做是出于害己。但天主所悦纳的忧愁则相反：因此他也说\BibleQuote{产生导致救恩的悔改，这悔改是永不后悔的。}因为一个人若为罪过而忧愁，若哀恸并克苦己身，他绝不会谴责自己。当蒙福的保禄说了这话之后，他不需要从别处引证他所说的话的证明，也不需要引述旧约历史中忧愁的人，而是举出你们格林多人自己；他用你们所做的事来提供证明；以便在赞美中既教训他们，又更加赢得他们归向自己。\par

% structure: p0012 source=h2 target=subsection reason=relative-heading-rank; base=h2

\subsection{格林多后书 7:11}

\BibleQuote{因为看，}他说，\BibleQuote{你们依天主的心意忧愁这事，在你们身上产生了何等热切的自省。}他说，\Quote{不但你们的忧愁没有使你们陷入自我谴责，仿佛你们这样做是徒然的；反而使你们更加谨慎。}然后他列举了那谨慎的某些标志：\par

\BibleQuote{是的，}何等\BibleQuote{的自白，}向着我。\BibleQuote{是的，何等的愤慨，}向着那犯罪的人。\BibleQuote{是的，何等的畏惧。}\BibleRef{格后 7:11}因为如此大的谨慎和非常迅速的改正，是属于那些极其畏惧的人。为了使自己不显得在抬高自己，请看他是如何迅速地缓和语气，说道：\par

\BibleQuote{是的，何等的渴望，}向着我。\BibleQuote{是的，何等的热忱，}向着天主。\BibleQuote{是的，何等的惩罚：}因为你们也惩罚了被冒犯的天主的法律。\par

\BibleQuote{在一切事上，你们证明了自己在这事上是纯洁的。}不仅因为没有犯下这事，因为这之前是明显的，而且因为没有赞同它。因为他在前书中说过\BibleQuote{你们还自高自大，}\BibleRef{格前 5:2}他在这里也说：\Quote{你们也清除了这个嫌疑；不仅没有称赞，反而责备并愤慨。}\par

% structure: p0017 source=h2 target=subsection reason=relative-heading-rank; base=h2

\subsection{格林多后书 7:12}

3. \BibleQuote{所以我虽曾写信给你们，} 我写信 \BibleQuote{不是为那行不义者，也不是为那受不义者。} 为免他们或许说：我们若是 \Quote{在这事上清白？} 你为何还要责斥我们？他更提前预备好应对此事，并预先处理了它，便说了他所说的，即 \BibleQuote{我不后悔，虽然我曾后悔。} \Quote{因为我如今远不后悔我当时所写的，那时我后悔的程度更甚于现在你们已证明了自己之后。} 你是否又见到他的强烈与恳切的争论，他如何将所说的话转向了完全相反的一面？因为他们原以为因他们的悔改，这会使他羞愧地收回成命，好像他草率地指责了他们；他却以此为证据，证明他直言不讳是正确的。因为当他发现自己能这样做时，他并不拒绝后来无畏地迁就他们。因为那在前文说了这样的话的人：\BibleQuote{与娼妓联合的，便是与她成为一体，} \BibleRef{格前 6:16} 以及 \BibleQuote{将这样的人交给撒殚，败坏他的肉体，} \BibleRef{格前 5:5} 以及 \BibleQuote{凡人所犯的罪，都在身体以外，} \BibleRef{格前 6:18} 以及类似的话；他怎能在这里说 \BibleQuote{不是为那行不义者，也不是为那受不义者？} 他并非自相矛盾，反而是极其一致。如何一致呢？因为向他表明他对他们的爱是他非常重要的一点。因此他并没有放弃对他的关心，但同时，如我所说，显示了他对他们的爱，以及一种更大的恐惧激动着他——即对整座教会的恐惧。因为他曾担心这邪恶会进一步蔓延，继续前行而吞噬整座教会。因此他也说了：\BibleQuote{一点面酵能使全团面发起来。} \BibleRef{格前 5:6} 然而那是他当时说的；但如今他们已做得好，他不再这样说，而是换了方式：确实暗示了同样的事，但更和悦地措辞，说道：\par

\BibleQuote{使我们为你们的关心，在你们面前显明出来。}\par

这就是说：\Quote{使你们知道我是多么爱你们。} 这其实与前文是一回事，但因表达方式不同，似乎传达了另一层意思。因为（要说服自己）这是同一回事，你只要展开他的思路，就会发现毫无差别。他说：\Quote{因为我非常爱你们，所以我担心你们会在那方面受到任何伤害，并且你们自己也会承受那忧愁。} 因此，正如他说 \BibleQuote{天主岂关心牛吗？} \BibleRef{格前 9:9} 时，并非指祂不关心（因为任何存在之物若被天主眷顾所抛弃，便不可能存留），而是指祂并非主要为牛立法；同样，他在这里的意思是：\Quote{我首先确实是为你们写的，但其次也是为他写的。} 而且我心中确有那爱，甚至不依赖我的书信；但我渴望通过那写作向你们，总而言之向所有人，显明这爱。\par

% structure: p0021 source=h2 target=subsection reason=relative-heading-rank; base=h2

\subsection{格后 7:13}

\BibleQuote{因此我们得到了安慰。}\par

因为我们既显明了我们对你们的关心，又完全成功了。正如他在另一处也说过：\BibleQuote{现今我们活着，若你们在主内站稳。} \BibleRef{得前 3:8} 以及：\BibleQuote{我们的希望、喜乐或所夸的冠冕是什么呢？不就是你们吗？} \BibleRef{得前 2:19} 因为对于一个有悟性的教师来说，这就是生命，这就是安慰，这就是慰藉；即他的弟子的成长。\par

4. 因为再没有什么比统治者对被统治者的慈父之情更能表明他是个统治者。因为单单生育并不构成父亲；而是在生育之后，还要有爱。如果在本性的事情上有如此大的爱的需要，那么在恩宠的事情上就更需要了。古时的所有杰出人物都是这样显明的。例如，凡在希伯来人中获得好名声的，都是因此彰显的。\Person{撒慕尔}就是这样显为大，说：\BibleQuote{但我决不因停止为你们祈祷而得罪天主。}\BibleRef{撒上 12:23} \Person{达味}如此，\Person{亚巴郎}如此，\Person{厄里亚}如此，每一位义人，新约和旧约中的，都是如此。因为\Person{梅瑟}为他所统治的人民的缘故，放弃了那样巨大的财富和说不尽的宝藏，\BibleQuote{宁愿选择与天主的百姓一同受苦。}\BibleRef{希 11:25} 并且在他被任命之前，他已通过行动成为人民的领袖。因此，那个希伯来人对他说的\BibleQuote{谁立了你作我们的首领和判官？}\BibleRef{出 2:14} 是极其愚蠢的。你说什么？你看见了行动，却怀疑称号？就好像一个人看见一位医生极其出色地使用手术刀，救治了身体上有病的肢体，却说：\Quote{谁立了你作医生，命令你使用手术刀？} \Quote{技艺，我的好先生，以及你自己的病患。} 同样，他的知识使他（即\Person{梅瑟}）成为他所声称的人。因为统治是一门艺术，不仅仅是一种尊严，而且是超越一切艺术的艺术。因为如果世俗的统治是一门超越其他一切的艺术和科学，那么这一统治就更是如此。因为这一统治比那一统治优越多少，那一统治就比其余的艺术优越多少；是的，甚至更多。如果你愿意，让我们更准确地考察这个论点。有农业、纺织、建筑等艺术；它们都非常必要，并极大地有助于维持我们的生命。因为其他的艺术无疑只是它们的附属：铜匠、木匠、牧羊人的技艺。但进一步，在所有艺术中，最必要的是农业，这甚至是天主在创造人之后首先引入的。因为没有鞋子和衣服，人还可以生存；但没有农业，就不可能生存。据说像\OriginalTerm{哈马克索比人}{Hamaxobii}、\OriginalTerm{斯基泰人}{Scythians}中的\OriginalTerm{游牧民族}{Nomads}，以及\OriginalTerm{印度裸智派}{Indian Gymnosophists}就是这样。因为他们不为房屋建造、纺织、制鞋这些艺术而烦恼，只需要农业。你们这些需要那些多余艺术的人——厨师、糕点师、刺绣工以及成千上万这样的人——应当羞愧，为了你们能生活；你们这些将虚浮的精致引入生活的人，应当羞愧；你们这些不信者，在面对那些不需要艺术的野蛮人时，应当羞愧。因为天主使本性极其自足，只需要很少的东西。然而，我并不强迫你们，也不立规矩叫你们这样生活；而是像\Person{雅各伯}所求的那样。他所求的是什么？\BibleQuote{若主赐给我食物吃，衣服穿。}\BibleRef{创 28:20} \Person{保禄}也同样吩咐说：\BibleQuote{只要有食物和遮盖，我们就当知足。}\BibleRef{弟前 6:8} 因此，首位是农业；第二位是纺织；第三位是建筑；制鞋是最后一位；因为在我们中间，毕竟有许多仆人和劳工，没有鞋也能生活。因此，这些是有用且必要的艺术。来吧，让我们把它们与统治的艺术相比。我之所以提出这些在所有艺术中最重要的，是为了当它被证明优于它们时，它胜过其余艺术就无可置疑了。那么，我们如何表明它比所有艺术更必要呢？因为没有它，这些艺术就毫无益处。如果你同意，我们就不提其他的，只把那个比任何艺术都更高、更重要的农业搬上舞台。如果你的劳动者互相争斗、抢夺彼此的财物，那么他们的众多人手又有何益处呢？因为实际上，对统治者的恐惧约束了他们，保护了他们所生产的东西；但如果你拿走了这个，他们的劳动就徒劳了。但如果有人仔细考察，他还会发现另一种统治，它是这一统治的父母和纽带。这可能是哪一种呢？那就是每个人应当控制和统治自己，惩罚自己不配的激情，同时小心地滋养和促进一切美德的萌芽。\par

因为统治有两种类型：一种是人统治人民和国家，规范政治生活；\Person{保禄}提到这一点时说：\BibleQuote{人人都应服从上级权力，因为没有权力不是出于天主的。}\BibleRef{罗 13:1} 之后，为表明这一点的益处，他接着说，统治者\BibleQuote{是天主的仆役，是为你们行善的}；又说，\BibleQuote{他是天主的仆役，是执行愤怒、惩罚作恶者的报复者。}\par

第二种统治是，每个有理智的人都统治自己；宗徒也进一步指出了这一点，说：\BibleQuote{你愿意不惧怕权力吗？你行善吧；}\BibleRef{罗 12:3} 这是指那统治自己的人。\par

5. 然而，这里还有另一种治理，比政治治理更高。这是什么？那就是在教会内的治理。保禄自己也提到这一点，说：\BibleQuote{应服从你们的领袖，且要顺服他们，因为他们为了你们的灵魂时刻警醒，好像要交账的人。}\BibleRef{希 13:17} 因为这种治理比政治治理，如同天高于地一样；是的，甚至更高更多。因为，首先，它主要考虑的并非如何惩罚已犯的罪，而是如何使罪恶根本不被犯下；其次，当罪已犯下时，不是如何去除死去的成员，而是如何将它们涂抹。它确实不太在意今世的事物，而是专注于天上的事物。\BibleQuote{因为我们的家乡是在天上。}\BibleRef{斐 3:20} 而我们的生命在这里。\BibleQuote{因为我们的生命，}他说，\BibleQuote{与基督一同藏在天主内。}\BibleRef{哥 3:3} 我们的奖赏在那里，我们的赛跑也是为了那里的冠冕。因为今世的生命在终结后并非消解，而是那时更加闪耀。因此，事实上，担负此种治理的人，手中被托付了更大的尊荣，不仅比总督更大，甚至比那些戴王冠的人更大，因为他们塑造人，向着更伟大的事物，为更伟大的事物。但是，无论追求政治治理的人，还是追求属神治理的人，都不能妥善管理它，除非他们首先妥善治理了自己，并严格遵守各自治理的法律。因为对众人的治理在某种程度上是双重的，同样，每个人对自己所施行的治理也是双重的。再者，在这一点上，属神的治理也超越政治的治理，正如我们所说的话所证明的。但也可以观察到某些技艺模仿治理，特别是农艺。因为正如耕作者在某种程度上是植物的治理者，修剪和抑制一些，使另一些成长并培育它们：同样，最好的治理者惩罚和剪除那些邪恶并损害众人的人，而提拔良善和守秩序的人。为此，圣经也将治理者比作葡萄园丁。尽管植物不会像国家中的受害者那样发出呼喊，但它们仍然通过外表显示所受的伤害，枯萎，被无价值的杂草挤压。正如邪恶被法律惩罚，同样，在这里，通过这种技艺，土壤的恶劣、植物的退化与野性也被纠正。因为我们会在这里找到人类性情的一切变种：粗鲁、软弱、胆怯、冒进、稳重：其中有些通过财富不合时宜地繁茂，损害邻人；另一些则贫困受损；例如，当树篱以牺牲邻近植物为代价而茂盛时；当其他不结果实和野生的树木长得极高，阻碍下面植物的生长时。正如统治者和君王有那些以暴行和战争扰乱其统治的人，同样，耕作者也遭受野兽的攻击、天气异常、冰雹、霉病、大雨、干旱以及所有这类事情。但这些事情发生，是为了让你不断仰望天主助佑的希望。因为其他技艺确实也依靠人的勤奋而前进；但这技艺的成败却由天主决定平衡，而且几乎完全依赖于此；它需要从上而来的雨水、天气的调和，以及超越一切的天主眷顾。\BibleQuote{因为栽种的不算什么，浇灌的也不算什么，只在乎那使之生长的天主。}\BibleRef{格前 3:7}\par

在此处也有死亡和生命，有阵痛和生育，正如在人中一样。因为在这里既有被剪除、结果实、死去、死而复生的事例，大地以此种种方式向我们清楚讲述复活。因为当根结果实，当种子发芽，岂不是一个复活吗？若逐一细查，便可在这一治理中领悟到天主眷顾和智慧的广大程度。但我所要说的是：此治理关乎土地和植物，而我们的治理却关乎灵魂的照料。植物与灵魂的差别有多大，此治理就优于彼治理多少。今世生活的统治者又多么逊于那治理，正如以自愿者比非自愿者更好管理。因为这也是一种自然的治理；因为在那情形下，一切都是出于害怕和强制而行；但在这里，正确行事是出于选择和意愿。不仅在此点，此治理胜过彼治理，而且在于它不仅是治理，而且可以说是父道；因为它有父亲的温柔；在命令更大事情的同时，仍加以说服。因为世俗统治者确实说：\Quote{你若犯奸淫，就丧了命。}但此治理，若你以不贞的眼光观看，就威胁最重的惩罚。因为这可畏的审判法庭，是为灵魂的纠正，不止是为身体。因此，灵魂与身体的差别有多大，此治理与彼治理之间的差别也多大。一个治理者确实审判公开的事；是的，甚至也不是所有公开的事，而是那些能被充分证明的事；而且常常，甚至在这些事上也背信弃义，但这法庭教导进入它的人说，那审判我们案件的主，将把\BibleQuote{一切赤露敞开的事物，}带到世界的公共剧场前，隐藏是不可能的。\BibleRef{希 4:13} 所以基督信仰维系我们的生活远胜于世俗法律。因为若因隐秘之罪而战栗，比因公开之罪而恐惧更使人安全；若因那些较小的过犯而追究责任，比只惩罚较重的罪更激励人走向美德；那么便容易看出，此治理比所有其他治理更能将我们的生命焊接在一起。\par

6. 但若你愿意，让我们也思考一下选立统治者的方式；因为在此你也会看见巨大的差异。因为获得这权柄并非藉着献金，而是要展现出极高的道德品行；并且不是作为世人的荣耀和自身的安逸，而是作为劳苦、操劳和众人的福祉，这样（我说）被任命者才被引入这统治。因此他所获得的来自圣神的帮助也是丰沛的。在那种情况下，统治权不过是宣告该做什么；但在此，它还添加了来自祈祷和圣神的帮助。更进一步；在那种情况下，根本没有关于哲学的谈论，也没有人坐下教导什么是灵魂、什么是世界、我们将来会成为什么、我们要从今世离开往何处去、以及我们如何成就德行。然而，关于契约、债券和金钱，却谈得很多，而对于那些事却毫不在意；而在教会中，却可以看到每一篇讲论的主题都是这些。因此，也完全有理由用所有这些名称来称呼它：审判庭、医院、哲学学校、灵魂的苗圃、以及通往天堂的那场赛跑的训练场。此外，这统治也是最温和的，尽管要求更大的严格性，从这一点就可以明显看出。因为世俗的统治者若抓住通奸者，立刻惩罚他。但这样有什么益处呢？因为这并不是要消除情欲，而是让灵魂带着它的创伤离去。但这统治者，当他察觉时，考虑的不是如何报复，而是如何根除情欲。因为你的所作所为，就好像当头部患病时，你不停止疾病，却砍掉头。但我不是这样：我砍掉疾病。我确实将他排除在奥迹和神圣的场所之外；但当我使他复原后，我重新接纳他，他立刻从邪恶中解脱并通过忏悔得以矫正。有人说：\InnerQuote{怎么可能根除通奸呢？}是可能的，是的，非常可能，如果一个人服从这些法律的话。因为教会是一个属灵洗浴，它不仅洗去身体的污秽，更洗去灵魂的污点，通过它的多种忏悔方法。因为你的做法，无论是放任一个不受惩罚的人使他变得更坏，还是惩罚他使他带着未治愈的创伤离去；而我既不让他不受惩罚，也不像你那样惩罚他，而是既要求符合我身份的补赎，又矫正了所行的恶事。你还要从另一个方式学习，你是如何虽然拔剑并显火焰给犯罪者，却没有带来什么明显的治愈；而我却不用这些，就引导他们达致完全的健康？但我无需论证或言词，而是用大地和海洋以及人类本性本身（作为见证）。请调查，在这审判庭设立之前，人类的事务状况如何；如何，甚至连现在所做的善功的名字都从未听说过。因为谁曾勇敢面对死亡？谁曾轻视钱财？谁曾淡漠荣誉？谁曾逃离生活的喧嚣，欢迎群山和独处，这天上智慧的母？贞洁的名号在哪里？因为这一切，以及比这些更多的，都是这审判庭的善功，这统治的作为。那么，知道这些，并清楚了解我们生活的一切益处和世界的改造都由此而来，就请常来聆听天主圣言，参与我们的集会和祈祷。因为如果你们这样调整自己，你们将能够，展现出相称于天国的态度，获得所应许的美善；愿我们所有人都能获得，藉着我们主耶稣基督的恩宠和祂对人的爱，愿光荣归于祂，至于无穷之世。阿们。\par

\SourceRecord{Translated by Talbot W. Chambers.；Nicene and Post-Nicene Fathers, First Series；Vol. 12.；Edited by Philip Schaff.；Buffalo, NY: Christian Literature Publishing Co.,；1889.；<http://www.newadvent.org/fathers/220215.htm>.}

% structure: p0001 source=h1 target=section reason=page-title

\section{《格林多后书》第十六篇讲道}

% structure: p0002 source=h2 target=subsection reason=relative-heading-rank; base=h2

\subsection{\BibleRef{格后 7:13}}

我们因你们的安慰，更因\Person{提多}的喜乐而格外欢喜，因为他的心从你们众人得了安息。\par

看，他又一次抬高对他们的赞美，并显示他们的爱。因为说过\Quote{我欢喜，不是因你们忧愁，而是因你们忧愁以致悔改}，并显示他自己的爱，因为他说\Quote{虽然我写信给你们，不是为那亏负人的，也不是为那受人亏负的，而是为要你们看出我对你们的关切}，他又一次提到他们善意的另一个标记，这为他们带来极大的赞美，并显示他们感情的真挚。因为他说\Quote{我们因你们的安慰，更因\Person{提多}的喜乐而格外欢喜}。然而，这似乎并不表明他极其爱他们，因为他为\Person{提多}欢喜胜过为他们欢喜。他回答说：\Quote{是，因为我的欢喜不是为他，而是为你们}。因此他附上理由说\Quote{因为他的\OriginalTerm{心肠}{\GreekText{σπλάγχνα}}从你们众人得了安息}。他说的不是\Quote{他}，而是\Quote{他的心肠}，即他对你们的爱。它们是如何得了安息的？\Quote{从众人}。因为这也是极大的赞美。\par

% structure: p0005 source=h2 target=subsection reason=relative-heading-rank; base=h2

\subsection{\BibleRef{格后 7:14}}

\Quote{因为若我在某事上曾为你们向他夸耀。}\par

这是极高的赞美，当师长夸耀时说\Quote{我并没有蒙羞}。所以我欢喜，因为你们显明自己已改正，并用行为证明了我所说的话。因此，我所获得的荣誉是双重的：首先，你们取得了进步；其次，我没有显出不符事实。\par

\BibleRef{格后 7:14}\Quote{但我们始终在真理中向你们说话，照样，我在\Person{提多}面前所夸耀的也显明是真理。}\par

这里他暗示了更进一步的事。如同我们在你们中间所讲的一切都是真理（很可能他也曾向他们大大赞美这个人），同样，我们对\Person{提多}所说关于你们的话也被证实为真。\par

% structure: p0010 source=h2 target=subsection reason=relative-heading-rank; base=h2

\subsection{\BibleRef{格后 7:15}}

\Quote{并且他内在的情感越发涌向你们。}\par

以下是称赞他（\Person{提多}）如何极度被爱充满并依恋他们。他没有说\Quote{他的爱}。然后，为了不显得在奉承，他处处提到这种情感的原因；正如我说，既逃避奉承的指责，又通过将赞誉归于他们来鼓励他们，并表明是他们向他注入了这巨大爱情的开端与基础。因为说了\Quote{他内在的情感越发涌向你们}后，他补充道：\par

\Quote{当他记起你们众人的服从。}这既表明\Person{提多}感激他的恩人，因为他回来后心中装着所有人，不停地记念他们，将他们留在唇边和心中；也表明\Person{格林多人}更受尊敬，因他们的表现使他心服。然后他也提到他们的服从，夸赞他们的热忱；因此他接着说：\par

\Quote{你们怎样以敬畏战栗的心接待了他。}不仅以爱，也以极大的尊敬。你们看见他如何为他们的双重美德作证吗？他们既爱他如父亲，又敬畏他如掌权者，不因敬畏而遮蔽爱，也不因爱而放松敬畏。他在上面也表达了这点：\Quote{你们依照天主的旨意忧愁，在你们身上产生了何等热忱，甚至恐惧，甚至渴慕。}\par

% structure: p0015 source=h2 target=subsection reason=relative-heading-rank; base=h2

\subsection{\BibleRef{格后 7:16}}

\Quote{所以我欢喜，因我在一切事上对你们放心。}你们看见他更是为他们欢喜吗？他说：\Quote{因为你们没有在任何事上使你们的师长蒙羞，也没有显出配不上我的见证。}因此，他的欢喜不单为\Person{提多}获得如此荣耀，更为他们自己展现了如此美好的情感。为了不让人以为他更替\Person{提多}欢喜，请注意他如何在此处也陈述了原因。正如他上面所说\Quote{若我在某事上曾为你们向他夸耀，我并没有蒙羞}，同样，此处\Quote{我在一切事上对你们放心}。\Quote{若有需要责备，我不怕你们会疏远；或再夸耀，我不怕被揭穿说谎；或称赞你们服从缰绳、有爱心、充满热忱，我对你们有信心。我命你们剪除，你们就剪除；我命你们接纳，你们就接纳。我在\Person{提多}面前说你们是伟大而可敬佩的人，知道如何尊敬师长：你们以行为证明这些话属实。而他得知这些，不完全是我告诉的，更多是从你们自己。无论如何，当他返回时，已成为你们的热烈爱人：你们的举止超出他所听闻的。}\par

% structure: p0017 source=h2 target=subsection reason=relative-heading-rank; base=h2

\subsection{\BibleRef{格后 8:1}}

2. \Quote{再者，弟兄们，我们告诉你们天主在\Place{马其顿}各教会中所赐的恩宠。}\par

用这些颂词鼓励他们之后，他又尝试劝诫。因为正是为此，他将这些赞美与责备混合在一起，免得从责备直接进入劝诫，使他要说的话不受欢迎；而是在安抚了他们的耳朵之后，借由此法为他的劝诫铺路。因为他打算谈论施舍；因此他事先说：\BibleQuote{我喜乐，因为在一切事上我对你们放心。}藉着他们过去的好行为，使他们也对此义务更加乐意。他没有立即说\Quote{因此，你们要施舍}，但请看他的智慧，如何从远处和高处为他的讲论预备。因他说：\BibleQuote{我将天主在马其顿各教会所施予的恩宠告诉你们。}为使他们不致自高自大，他将他们所行的称为\Quote{恩宠}；并在述说他人所为时，藉着对别人的赞美在他们心中激发更大的热忱。他同时提到马其顿人的两个赞美，更确切地说，是三个：即他们高贵地忍受试炼；他们知道如何怜悯；以及虽然贫穷，却在施舍上表现出慷慨，因为他们的财产也曾被掠夺。当他写信给他们时，正是为此意义说：\BibleQuote{因为你们成了在犹大天主的各教会的效法者，你们也从本族人受了同样的苦，如同他们从犹太人受的苦一样。}\BibleRef{得前 2:14} 请听他在后来写给希伯来人的信中所说的话：\BibleQuote{你们甘心忍受了家产的剥夺。}\BibleRef{希 10:34} 但他称他们所行的为\Quote{恩宠}，不仅是为了使他们谦卑；更是为了激发他们的竞争心，并防止他的话引起嫉妒。因此他又加上\Quote{弟兄}的称呼，以消解一切妒嫉情绪；因为他即将以华丽的言辞赞美他们。请至少听听他的赞美。因他说了：\Quote{我将天主的恩宠告诉你们}，他并未说\Quote{在这城或那城所施予的}，而是赞美整个民族，说：\Quote{在马其顿各教会。}随后他也详细说明这同一恩宠。\par

% structure: p0020 source=h2 target=subsection reason=relative-heading-rank; base=h2

\subsection{格林多后书 8:2}

\BibleQuote{他们在患难的大试炼中，仍充满喜乐。}\par

你看见他的智慧吗？因为他先不说他所愿望的，而是先说另一件事，免得显得是故意如此，而是由不同的关联引出。\BibleQuote{在患难的大试炼中}。这正是他在写给马其顿人的信中所说的：\BibleQuote{你们成了主的效法者，在患难中接受了圣言，怀着圣神的喜乐}；又说：\BibleQuote{主的圣言从你们那里传扬出去，不只传到马其顿和阿哈雅，更是传到了各处，你们对天主的信德已传遍了。}\BibleRef{得前 1:6} 但\BibleQuote{在患难的大试炼中，仍充满喜乐}是什么意思？他说，两者都过度临到他们，患难和喜乐都过度。因此，如此大的喜乐从患难中为他们产生出来，这是非常奇妙的。事实上，患难不仅没有成为忧愁的根源，反而成为他们喜乐的因由，而且这患难还是\Quote{大的}。他这样说，是为了准备他们在试炼中变得高贵和坚定。因为他们不仅受了苦，而且藉着耐心得以证明为精炼；更确切地说，他说的不是耐心，而是比耐心更甚的\Quote{喜乐}。他也不是简单地说\Quote{喜乐}，而是说\Quote{充盈的喜乐}，因为伟大的、说不出的喜乐在他们内涌出。\par

3. \BibleQuote{他们极度的贫穷，竟涌出慷慨的富足。}\par

再者，这两者都带有过度性。因为正如他们的大患难产生了大喜乐，甚至\Quote{充盈的喜乐}，同样，他们的大贫穷产生了施舍的大富足。为此他表明说：\BibleQuote{涌出慷慨的富足}。因为慷慨不是由所给之物的数量决定，而是由施与者的心意决定。\par

因此他绝不说\Quote{礼物的丰厚}，而是说\BibleQuote{慷慨的富足}。现在他所说的大意是：\Quote{他们的贫穷不仅没有阻碍他们慷慨，反而成为他们丰裕的机会，正如患难成为喜乐的机会一样。因为越贫穷，他们就越慷慨，越乐意奉献。}因此他也极其佩服他们，因为在如此极度的贫穷中，他们显示出如此大的慷慨。因为\BibleQuote{他们的深度}，即\Quote{他们伟大而无法言说的}，\BibleQuote{贫穷}，显示了他们的\BibleQuote{慷慨}。但他不是说\Quote{显示了}，而是说\BibleQuote{涌出}；他也不是说\Quote{慷慨}，而是说\BibleQuote{慷慨的富足}；也就是说，与贫穷的伟大相平衡，或更确切地说，大大超越它的，是他们所显示的慷慨。然后他甚至更清楚地解释说：\par

% structure: p0026 source=h2 target=subsection reason=relative-heading-rank; base=h2

\subsection{格林多后书 8:3}

\BibleQuote{因为他们按照自己的能力，我作证，}这见证是可信的。\BibleQuote{而且超过了自己的能力。}也就是说，它\BibleQuote{涌出慷慨的富足}。或者更确切地说，他不仅通过这一表达，也通过随后的一切来表明这一点；因为他说道：\BibleQuote{出于自愿}。看！这又是另一种过度。\par

% structure: p0028 source=h2 target=subsection reason=relative-heading-rank; base=h2

\subsection{格林多后书 8:4}

\Quote{恳切地请求} 看！还有第三次和第四次。\Quote{恳求我们} 看！甚至第五次。当他们在患难和贫穷中时，这里就有第六次和第七次。他们更是超额给予。既然这正是他最想在格林多人身上所供应的，即出于自愿的给予，所以他特别强调这一点，说：\Quote{恳切地请求}，和\Quote{恳求我们}。'我们不是恳求他们，而是他们恳求我们。'恳求我们什么？\Quote{恩宠和共融，即服务于圣徒的职务}。你看见他如何再次高举此行为，以可敬的名称称呼它吗？因为他们渴慕神恩，他便以恩宠之名称呼它，好使他们热心追求；又以\OriginalTerm{共融}{fellowship}之名，好使他们明白他们不仅给予，也领受。'因此他们恳求我们，'他说，'要我们接受这样的职务。'\par

% structure: p0030 source=h2 target=subsection reason=relative-heading-rank; base=h2

\subsection{格林多后书 8:5}

\Quote{And} 这，\Quote{不是我们所希望的}。这是就数量和他们的患难而说的。'因为我们绝想不到，'他说，'他们在这么大的患难和贫穷中，竟会恳求我们，并如此大力请求我们。'他也指出他们在其他方面的生活谨慎，说：\par

\Quote{但首先他们将自己献给主，并藉天主的旨意献给我们。}\par

'因为在一切事上，他们的顺服超出了我们的期望；并非因为他们行了怜悯，就忽略了其他德行，' \Quote{但首先他们将自己献给主}。什么是\Quote{将自己献给主}？'他们奉献了自己；他们表明自己在信德上蒙悦纳；他们在试炼中显示了极大的刚毅、秩序、良善、爱德，在一切事上既有准备又有热忱。'\Quote{并献给我们}是什么意思？'他们顺从管束，爱我们，服从我们；既遵守天主的法律，又因爱与我们相连。'请注意他在这里又怎样显示他们的热诚，说：\Quote{将自己献给主}。他们没有在一些事上服从天主，在另一些事上服从世俗；而是在一切事上服从祂；并将自己完全交给了天主。因为他们没有因行了怜悯就充满愚昧的骄傲，而是显示极大的谦卑、极大的顺服、极大的敬畏、极大的属天智慧，从而也实现了他们的施舍。但\Quote{藉天主的旨意}是什么意思？既然他说他们\Quote{献给我们}，但这并非按人的方式\Quote{给我们}，而是他们也按天主的旨意行了这事。\par

% structure: p0034 source=h2 target=subsection reason=relative-heading-rank; base=h2

\subsection{格林多后书 8:6}

4. \Quote{因此我们劝勉了弟铎，他既已开始，也当在你们中完成这恩宠。}\par

这有什么联系？关系很大，且紧密联系前文。'因为我们看见他们热烈，' 他说，'在一切事上热心：在试炼中，在施舍中，在他们对我们的爱中，在其他方面的纯洁生活中：为使你们也能与他们相等，我们派遣了弟铎。' 但他没有这么说，尽管他暗示了。看那爱的满溢！'因为虽然他们恳求并渴望我们，'他说，'我们却为你们的情况焦虑，免得你们不如他们。因此我们也派了弟铎，好使你们也因此被激发和提醒，效法马其顿人。' 因为弟铎恰好在写这封信时在那里。然而他表明他在保禄劝告之前已开始此事；\Quote{他既已开始，}他说。因此他也大大赞扬他；例如，在书信开头：\Quote{因为我找不到我的兄弟弟铎，我的心灵不得安宁：} \BibleRef{格后 2:13} 以及这里他所说的一切，也包括这一点。因为这也是不小的赞美，即甚至先开始：这显示出热切而炽热的灵魂。因此他也派遣他，借此在他们中间注入一个极大的激励，即弟铎的在场。为此他也用赞美来推崇他，希望使他对格林多人更加亲切。因为这在劝服上也大有分量，当劝告者处于亲密关系时。他一次、两次、三次提到施舍，并称之为'恩宠'；现在确实说：\Quote{此外，弟兄们，我将天主赐给马其顿众教会的恩宠告诉你们；} 又说：\Quote{他们出于自愿，以恳切请求恳求我们关于这恩宠和共融：} 再说：\Quote{要怎样开始，也要怎样在你们中完成这恩宠。}\par

5. 这实在是一件伟大的善事，是天主的恩赐；若善行之，便尽可能地使我们相似天主；因为这样的人，才是至高意义上的一个人。至少有一位，给出了人的模范，提到了这一点，因为他说：\Quote{人，}他说，\Quote{是一件伟大的事物；而仁慈的人是一件可敬的事物。}\BibleRef{箴 20:6}这项恩赐比复活死人更大。因为喂养饥饿的基督，远大于因耶稣之名复活死人：在前者，你向基督行善；在后者，祂向你行善。赏报当然来自于行善，而非受善。因为在此处（我指的是奇迹）你是天主的债务人；而在施舍中，你使天主成为你的债务人。施舍，乃是出于甘心乐意，出于慷慨，当你认为自己不是给出而是领受，当你好像受益了，好像得到了而非损失了；因为若非如此，这就不是恩宠了。因为向他人显示仁慈的人，应当感到喜乐，而非恼怒。如果你在消除他人沮丧的同时，自己却沮丧，这岂不荒谬？因为这样，你就不再使它成为施舍了。如果你因将他人从沮丧中解救出来而沮丧，你就提供了极端残酷和非人道的榜样；因为不救他，比这样救他更好。而你为何还沮丧呢，人啊？是怕你的金子减少吗？若你这样想，就根本不要给；若你不确定它会在天堂里为你增多，就不要施予。你却在这里寻求报酬。为什么？让你的施舍是施舍，而非交易。如今许多人确实在此生也得到了报酬；但他们得到的方式，并非好像比那些在此生未得报酬的人更有优势；而是他们中的一些人，因软弱过度，因为他们没有被那永恒之事深深吸引。正如那些贪婪、无礼、被肚腹所奴役的人，被邀请参加皇家宴会，却无法等待适当时机，像小孩子一样，预先进食，用劣质菜肴撑饱自己，破坏了乐趣：同样，这些在此生寻求并接受报酬的人，确实减少了他们在天上的赏报。此外，当你借贷时，你希望经过更长的时间后收回本金，甚至也许根本不收回，以便通过延迟使利息更大；但在这种情况下，你却立即要求归还；而且你还即将不在这里，而在那里永远；你即将不在这里受审判，而是呈上你的账目？如果真有人为你建造宅邸，而你却不打算留下，你会认为这是损失；但现在，你渴望在此地富有，而这里你甚至可能在晚上之前就要离开？你不知道我们住在异乡，如同陌生人和旅居者吗？你不知道旅居者的命运就是在意想不到、毫无预料的时候被驱逐吗？这也是我们的命运。因此，我们预备的一切，都留在了这里。因为主不允许我们带着它们离开，如果我们建造了房屋，买了田地，买了奴隶，买了器具，或任何其他类似的东西。主不仅不允许我们带着它们离开此地，甚至不将它们的价值计入你的账上。因为祂警告过你，不应建造，也不应花费在属于别人而非你自己的事物上。因此，为什么要离开你自己的，在别人的事物上劳苦花费，以致丧失你的劳苦和报酬，并遭受最严厉的惩罚呢？不要这样，我恳求你；但既然我们本性是旅居者，也让我们选择做旅居者；这样我们在那里就不会成为被驱逐、蒙羞辱的旅居者。因为如果我们执意要在此地做公民，那么我们在此地和彼地都将不是公民；但如果我们继续做旅居者，并以旅居者当有的方式生活，我们就会在此地和彼地都享有公民的自由。因为义人虽然一无所有，却会在此地住在所有人的财产中，如同自己的；而且当他离世升天堂时，将看到那些永恒的居所。他在此地也不会遭受任何不适，（因为没有任何人能使他成为陌生人，他把每片土地都当作自己的城市）；当他被接回自己的家乡时，将得到真正的财富。为了使我们获得此世和彼世的事物，让我们善用我们所拥有的。因为这样，我们就会成为天堂的公民，享有极大的胆量；愿我们都能藉着我们的主耶稣基督的恩宠和爱人之心，与圣父及圣神，永享光荣与权能。阿们。\par

\SourceRecord{Translated by Talbot W. Chambers.；Nicene and Post-Nicene Fathers, First Series；Vol. 12.；Edited by Philip Schaff.；Buffalo, NY: Christian Literature Publishing Co.,；1889.；<http://www.newadvent.org/fathers/220216.htm>.}

% structure: p0001 source=h1 target=section reason=page-title

\section{《格林多后书》讲道集第十七篇}

% structure: p0002 source=h2 target=subsection reason=relative-heading-rank; base=h2

\subsection{\BibleRef{格后 8:7}}

\BibleQuote{好使你们在一切事上洋溢；在信德、宣讲、知识和一切热忱上。}\par

看，他再次将劝勉与称赞结合，更大的称赞。他没有说\Quote{你们给予}，而是说\BibleQuote{你们洋溢}：即，在恩赐的\Quote{信德}上，在智慧言语的\Quote{宣讲}上，在道理的\Quote{知识}上，在达成其他一切德行的\Quote{所有热忱}上。\par

\BibleQuote{并在你们的爱上}，就是我先前所说、也曾证明过的爱。\par

\BibleQuote{好使你们也在这一恩宠上洋溢。}你们看见吗？正因如此，他才先以这些赞美开始，好更进一步引导他们在这些事上也同样勤勉。\par

% structure: p0007 source=h2 target=subsection reason=relative-heading-rank; base=h2

\subsection{\BibleRef{格后 8:8}}

\BibleQuote{我说这话，不是出于命令。}\par

看，他何等随时顾及他们的感受，何等避免冒犯，既不严厉也不勉强；或者不如说，他的话同时兼具这两者，且带着自愿之事的无冒犯性。因为他反复劝勉他们，又大大称赞了马其顿人，为免这看似构成一种必要，他说：\par

\BibleQuote{我说这话，不是出于命令，而是藉着别人的热忱考验你们之爱的真诚。}\par

\BibleQuote{我说这话，不是出于命令，而是藉着别人的热忱考验你们之爱的真诚。} \Quote{并非怀疑它}（因为他不是这个意思），\Quote{而是为了使其得到认可，展示它并使之更加坚强。因为我之所以说这些，是要激发你们同样奋勇。我提及他们的热诚，是为了照亮、鼓舞、激励你们的心志。}接着，他从这里进到另一个更重大的论点。因为他不放过任何劝服的方式，在处理他的论述时几乎翻天覆地。他既藉着对别人的赞美劝勉他们，说你们知道\BibleQuote{天主在马其顿众教会中所赐的恩宠}，也藉着他们自己的赞美：\BibleQuote{因此你们在一切事上洋溢，在宣讲和知识上。}因为这更令人刺痛：一个人发现自己不如自己，比发现自己不如别人更甚。然后他继续来到他劝勉的头等和冠冕。\par

% structure: p0012 source=h2 target=subsection reason=relative-heading-rank; base=h2

\subsection{\BibleRef{格后 8:9}}

\BibleQuote{因为你们知道我们主耶稣基督的恩宠：祂本是富有的，为了你们却成了贫穷的，好使你们因着祂的贫穷而成为富有的。}\par

\Quote{因为你们要记念}，他说，\Quote{思考并默想天主的恩宠，不要轻易放过，而要以认识其深厚与本质为目标，那么你们就不会吝惜自己的任何东西。祂倒空了自己的光荣，为叫你们不是因祂的富有，而是因祂的贫穷而成为富有的。如果你们不信贫穷能产生富有，那就记念你们的主，你们便不再怀疑。因为如果祂没有成为贫穷，你们就不会成为富有的。这正是惊人之处：贫穷使富有变得富有。}他这里所说的富有，是指虔诚的知识、罪恶的涤除、成义、圣化，以及祂已赐给我们和决意要赐给我们的无数美善。所有这一切都通过祂的贫穷而临到我们。什么贫穷？就是祂取了肉身，成了人，并承受了祂所受的苦难。然而祂并不欠这债，而是你欠祂。\par

% structure: p0015 source=h2 target=subsection reason=relative-heading-rank; base=h2

\subsection{\BibleRef{格后 8:10}}

\BibleQuote{我在这事上给你们建议，是为了你们的益处。}\par

看，他如何再次谨慎避免冒犯，并通过两件事缓和自己的话：一是说\BibleQuote{我建议}，二是说\BibleQuote{为了你们的益处}。他说：\Quote{因为我既不催促也不强迫你们，也不要求不情愿的臣民；我说这些话，与其说是着眼受惠者的利益，不如说是着眼你们的利益。}接着，他所举的例子也来自他们自己，而不是别人。\par

\BibleQuote{你们是首先开始的人，一年前不仅已开始实行，而且已开始愿意。}\par

看，他如何表明他们自己曾甘心情愿，并且毫无劝服就已下定决心。因为他曾向得撒洛尼人作证说，他们是\BibleQuote{自愿且再三恳求}推行这项施舍；他也渴望表明这些人的善工也是出于他们自己。因此他说：\BibleQuote{不仅已开始实行，而且已开始愿意}，又说不是\BibleQuote{开始}，而是\BibleQuote{一年前就已开始}。所以我劝勉你们在这些事上，就是你们先前已热切自发去做的事。\par

% structure: p0020 source=h2 target=subsection reason=relative-heading-rank; base=h2

\subsection{\BibleRef{格后 8:11}}

\BibleQuote{如今你们也当完成这事。}\par

他没有说\Quote{你们做了}，而是说\Quote{你们完成了这事}；\par

\BibleQuote{好使既然有愿作的热心，也按你们所有的得以完成。}\par

好使这项善工不止于热心，也领受伴随行为而来的赏报。\par

% structure: p0025 source=h2 target=subsection reason=relative-heading-rank; base=h2

\subsection{\BibleRef{格后 8:12}}

2. \BibleQuote{因为人只要有愿作的热心，是按他所有的蒙悦纳，不是按他没有的。}\par

看，何等的不可言喻的智慧。他指出了那些超过自己能力行事的人——我说的是得撒洛尼人——并为此称赞他们，说\BibleQuote{我给他们作证，他们甚至超过了自己的力量}，然后劝勉格林多人只\Quote{按}自己的力量去做，让榜样自行发挥作用；因为他知道，激发效法同样的行为，与其说是劝勉，不如说是竞争。因此他说：\BibleQuote{因为人只要有愿作的热心，是按他所有的蒙悦纳，不是按他没有的。}\par

\Quote{不要害怕}，他的意思是，\Quote{因为我说的这些话，不过是对他们慷慨的赞扬，但天主是按照人的力量要求}，即\BibleQuote{按他所有的，不是按他没有的}。因为\BibleQuote{蒙悦纳}这个词在此的意思是\Quote{被要求}。他大大缓和了语气，自信地依靠这个榜样，并通过给他们留自由而更稳妥地赢取他们。因此他又加上了：\par

% structure: p0029 source=h2 target=subsection reason=relative-heading-rank; base=h2

\subsection{\BibleRef{格后 8:13}}

\BibleQuote{我说这话，并不是要别人轻松，你们受累。}\par

然而基督却赞扬了那寡妇相反的行为，她将自己全部的生活费都倒空，并出于她的匮乏而施予。\BibleRef{谷 12:43} 但由于他是在向格林多人讲论，他宁愿在其中忍饥挨饿；\Quote{因为我宁可死，也不愿有人使我这夸耀落空。}\BibleRef{格前 9:15} 因此他采用温和的劝诫，确实赞扬那些做了超出能力之事的人，但并不强迫这些人也这样做；不是因为他不想这样，而是因为他们有点软弱。他为什么赞扬那些人呢？因为\Quote{在患难大考验中，他们的喜乐和极度的贫穷，竟涌出他们慷慨的富足：}并且他们施予\Quote{超过了自己的能力？}这岂不是很明显，是为了引导这些人也这样做吗？所以，即使他看起来允许较低的标准，他这样做是为了借此将他们提升到这个标准。例如，请看他在接下来的话中是如何暗中为此铺路的。说了这些之后，他补充道：\par

% structure: p0032 source=h2 target=subsection reason=relative-heading-rank; base=h2

\subsection{格林多后书 8:14}

\Quote{你们的富足补足他们的缺乏。}\par

因为他不仅通过先前的话，也通过这些话，希望使这诫命变得轻省。而且不仅从这一考虑，也从报酬的考虑，他又使之更容易；并说出比他们应得的更高的事，说：\Quote{这样，在此刻就有平等，他们的富足}补足\Quote{你们的缺乏。}那么，他说的是什么？\Quote{你们在金钱上富足；他们在生命和在天主前的胆量上富足。所以，你们要给他们，你们富足而他们缺乏的金钱，好使你们得到他们富足而你们缺乏的胆量。}看他是如何暗中准备让他们超出能力和出于匮乏地施予。\Quote{因为，}他说，\Quote{如果你们渴望得到他们的富足，就要拿出你们的富足；但如果你们想为自己赢得全部，你们就会出于匮乏并超出能力地施予。}然而，他没有这样说，而是留给听众自己去思考；同时他本人则实现他的目标，并给予恰当的劝诫，接着他按照所显示的内容，加上这些话：\Quote{在此刻就有平等。}怎样平等？你们和他们互相给出你们的富余，并填补你们的缺乏。这种平等是什么，用属灵的事交换属肉身的事？因为那一边有很大的优势；那么他为什么称之为\Quote{平等？}或者他是说，就各自富足和缺乏而言，这种[平等]发生；或者只是就现世生命而言。因此，在说了\Quote{平等}之后，他补充说：\Quote{在此刻。}他说这话，一方面是为了压制富人的高傲，另一方面是表明在我们离开之后，属灵的人拥有更大的优势。因为在此世，我们大家都享受许多平等的尊荣；但那时将有巨大的差别和极大的优越，当义人比太阳更明亮时。然后，既然他表明他们不仅要施予，也要接受，而且更多，作为回报；他尝试通过进一步的考虑使他们积极，表明如果他们不将自己的财物施予他人，他们也不会因将一切聚集在内而得利。他引述了一个古老的故事，这样说：\par

% structure: p0035 source=h2 target=subsection reason=relative-heading-rank; base=h2

\subsection{格林多后书 8:15}

\BibleQuote{正如经上所载：多收的没有剩余，少收的也没有缺乏。}\par

这事发生在玛纳的案例上。因为无论是收得多的还是收得少的，都被发现具有同样的数量，天主以此方式惩罚贪婪。他说这话，一方面是以当时发生的事警告他们，另一方面是劝说他们永远不要渴望更多，也不为少有而悲伤。现在，人们可以看到这在现世事物中发生，不仅在玛纳上。因为如果我们都只填满一个肚腹，活同样的年岁，穿同样的衣服；那么，富人不会因他的富余而获得什么，穷人也不会因他的贫穷而损失什么。\par

3. 那么，你为什么对贫穷战栗呢？为什么追求财富呢？有人说：\Quote{我害怕，}有人说，\Quote{我怕被迫到别人门口去向邻人乞讨。}我也常听到许多人这样祈祷说：\Quote{求祢任何时候都不要让我需要他人。}我听到这些祈祷时甚感可笑，因为这恐惧甚至如同孩童一般。因为每天、在每件事上，我们都可以说是彼此需要。所以说，这是出自一个鲁莽和傲慢的心灵，未能清楚辨识事物的本质。难道你看不见我们所有人都彼此需要吗？兵士需要工匠，工匠需要商人，商人需要农夫，奴隶需要自由人，主人需要奴隶，穷人需要富人，富人也需要穷人，不工作的人需要施舍者，施舍者也需要接受者。因为接受施舍的人满足了一个非常大的需要，比任何需要都大。如果没有穷人，我们救恩的大部分就会被推翻，因为我们将无处施舍我们的财富。所以，即使看起来比任何人都无用的穷人，却是最有用的。但如果需要他人是羞耻的，那么就只有死路一条；因为惧怕这个的人无法生存。有人说：\Quote{我受不了别人傲慢地扬起的拳头。}你指责别人傲慢，却用这个指控使自己蒙羞吗？因为不能忍受骄傲灵魂的膨胀，这本身就是傲慢的。你为什么恐惧这些事，为这些不值一提的事战栗，也因此惧怕贫穷呢？如果你富有，你将需要更多，而且更多更卑微的东西。因为你的财富越大，你就越受这个诅咒的束缚。你竟如此无知，当你祈求财富以便不需要任何人时，就像一个人来到大海，需要水手、船只和无数装备，却祈祷自己什么都不需要。如果你渴望极度独立于所有人，就祈求贫穷吧；然后如果你依赖任何人，也仅仅是依赖面包和衣服；但在另一种情况下，你将需要他人，涉及土地、房屋、赋税、工资、地位、安全、荣誉、长官及其下属，城里的和乡下的，商人和店主。你看见这些话是极度粗心的话吗？简言之，如果彼此需要在你看来是可怕的事，要知道根本无法完全逃脱；但如果你要避开纷扰（因为你可以逃入贫穷的无波之港），就斩断你事务的众多纷扰，并且不要认为需要他人是可耻的；因为这是天主不可言喻的智慧之作。如果我们彼此需要，即使这需要的强迫也未能使我们彼此联合于爱中；如果我们独立自主，岂不是成为未驯服的野兽吗？天主不得已强制地使我们彼此隶属，每天彼此碰撞。如果祂除去这缰绳，谁会轻易渴望邻人的爱呢？那么，我们既不要认为这是可耻，也不要祈祷反对它说：\Quote{求祢不让我们需要任何人；}而要祈祷说：\Quote{求祢不容我们在需要时拒绝那些能帮助我们的人。}需要他人并非可悲的，攫取他人之物才是可悲的。但现在我们从未为此祈祷，也没有说：\Quote{求祢不让我贪图别人的东西；}反而认为需要是一件值得祈求免去的事。然而\Person{保禄}曾多次处于需要之中，并不羞愧；反而以此为荣，并称赞那些服侍过他的人，说：\BibleQuote{你们曾一两次为我的需要送去财物;}\BibleRef{斐 4:16}又说：\BibleQuote{我剥削了别的教会，拿了工资来服侍你们。}\BibleRef{格后 11:8}因此，以此为耻不是宽宏大量的表现，而是软弱和卑下愚钝的心灵的表现。因为彼此需要正是天主的旨意。所以不要让你的哲学超越中道。有人说：\Quote{我受不了被人多次恳求却仍不答应。}那么天主多次恳求你，你却不听从，而且恳求你的还是对你有益的事，祂怎能忍受你呢？他说：\BibleQuote{我们作\Person{基督}的大使，}他说，\BibleQuote{好像天主借着我们劝勉你们；你们与天主和好吧。}\BibleRef{格后 5:20}\Quote{然而我是祂的仆人，}他说。那又怎样呢？当你这个仆人醉酒，而祂主人饥饿、连必需的食物都没有时，你\Quote{仆人}的名号能为你做什么呢？不，这反而会更压垮你，当你住在三层楼房，而祂连像样的住所都没有；当你躺在柔软的床上，而祂连枕头都没有。有人说：\Quote{我给了。}但你不应该停止这样做。因为只有当你没有东西可给、一无所有时，你才有借口；但只要你还有（即使你给了一万人），而别人还在挨饿，你就没有借口。当你囤积粮食、抬高价格、想出其他不寻常的交易诡计时，你从今以后还有什么得救的希望呢？你曾被命令免费给饥饿的人食物，但你连合理价格都不给。祂为了你倒空了如此大的荣耀，你却不认为祂配得一块面包；你的狗吃得饱饱的，而\Person{基督}却在挨饿；你的仆人因过饱而腹胀，而你的主和祂的主却缺少必需的食物。这难道是朋友的行为吗？\BibleQuote{你们与天主和好吧，}\BibleRef{格后 5:20}因为这是敌人和敌对者的行为。\par

4. 让我们怀着羞愧思想我们已经领受的巨大恩惠，以及我们还将领受的巨大恩惠。如果一个穷人来到我们面前乞求，让我们怀着极好的善意接待他，用言语安慰他、扶起他，好使我们自己也能从天主和人那里得到同样的待遇。\BibleQuote{凡你们愿意人给你们做的，你们也要照样给人做。}\BibleRef{玛 7:12} 这条法律不含任何沉重或冒犯之处。它说：\Quote{你愿意接受的，你就去做。}回报是相等的。而且它没有说\Quote{你不愿接受的，就不要做}，而是说了更进一步的。因为前一种确实是戒恶，后一种则是行善，后者包含了前者。而且祂没有说\Quote{你们也愿意那样做，而是去做}，而是说\Quote{去做。}那么好处是什么？\BibleQuote{这就是法律和先知。} 你愿意人怜悯你吗？那么你就怜悯人。你愿意得到宽恕吗？那么你就宽恕人。你愿意人不毁谤你吗？那么你就不要毁谤人。你渴望得到称赞吗？那么你就给予称赞。你愿意不受冤屈吗？那么你就不要掠夺人。你看见祂如何表明德性是自然的，我们不需要外在的法律和教师吗？因为在我们愿意从邻人那里接受或不接受的事情上，我们为自己立法。所以，如果你不愿意接受一件事，却去做它，或者如果你愿意接受它，却不去做，你就成了自我定罪的，从此没有任何借口以无知和不知道该做什么为理由。因此，我恳求你们，在我们自己内为自己设立这条法律，并阅读这写得如此清晰简洁的文字，让我们对待邻人，如同我们希望他们对待我们一样；好使我们既享受眼下的平安，又获得将来的福乐，借着我们的主耶稣基督的恩宠和祂对人的爱，愿光荣、权能、尊崇归于祂和圣父及圣神，从现在直到永远，永世无尽。阿们。\par

\SourceRecord{Translated by Talbot W. Chambers.；Nicene and Post-Nicene Fathers, First Series；Vol. 12.；Edited by Philip Schaff.；Buffalo, NY: Christian Literature Publishing Co.,；1889.；<http://www.newadvent.org/fathers/220217.htm>.}

% structure: p0001 source=h1 target=section reason=page-title

\section{格林多后书第十八篇讲道}

% structure: p0002 source=h2 target=subsection reason=relative-heading-rank; base=h2

\subsection{格林多后书 8:16}

\BibleQuote{\Emph{感谢归于天主，} \Emph{因为祂将同样的热切关怀放入了弟铎的心中，使他对你们关心。}}\par

他再次赞扬弟铎。因为既然他已谈论了施舍，随后也谈论那些将要接受他们钱财并带走的人。因为这有助于这项收集，并增加捐献者的热忱。因为那对管理钱财者怀有信心且不怀疑接收的人，会以更丰满的慷慨给予。为使那时的情况也是如此，请听他如何称赞为此目的而来的人，其首者是弟铎。因此他说：\Quote{感谢归于天主，祂将（字面意为\InnerQuote{赐予}）同样的热切关怀放入了弟铎的心中。} 什么是\Quote{同样的}？就是他对得撒洛尼人所怀有，或是\Quote{同样的}如同我所怀有。请注意此处的智慧。他表明这是天主的作为，也感谢那赐予者，借此激发他们。\Quote{因为如果天主激动了他并派遣他到你们那里，祂便借着他要求。所以不要以为所发生的是出于人。}从何可见是天主激动了他？\par

% structure: p0005 source=h2 target=subsection reason=relative-heading-rank; base=h2

\subsection{格林多后书 8:17}

\BibleQuote{因为他固然接受了我们的劝告，但因自己非常热切，便自愿前往。}\par

请注意他如何也描绘他尽了自己的本分，无需他人催促。并且提及天主的恩宠后，他并未将一切完全归于天主；同时，为借此赢得他们更大的爱，他又说他也是出于自己而被激动。因为，\BibleQuote{他非常热切，便自愿前往，}\Quote{他抓住了机会，他冲向宝藏，他将你们的服事视为自己的益处；因为他极其爱你们，不需要我给予的劝告；虽然他也接受我的劝告，但并非因之而被激动，而是出于自己及天主的恩宠。}\par

% structure: p0008 source=h2 target=subsection reason=relative-heading-rank; base=h2

\subsection{格林多后书 8:18}

\BibleQuote{我们还同他一起派去了那位兄弟，他在福音上的名声传遍了众教会。}\par

这位兄弟是谁？有些人说是路加，因为他所写的史记；但有些人说是巴尔纳伯，因为他也将未成文的宣讲称为福音。为什么他没有提他们的名字，却既指名提到弟铎\BibleRef{格后 8:23}，并称赞他在福音上的合作（因为他如此有用，以致由于他的缺席，甚至保禄也不能做什么伟大的事；因为\Quote{我没有找到我的兄弟弟铎，我的心灵不得安宁}\BibleRef{格后 2:13}），并且他被教会所选，经过了全体人民的投票？什么是\Quote{归荣耀于同一的主，并显明你们的乐意}？\Quote{为叫天主得荣耀，并使你们更乐意，这些接受钱财的人是有验证的，无人能对他们产生任何虚假的怀疑。因此我们寻找这样的人，没有将全部交给一个人，以免他也遭遇这种怀疑；但我们派了弟铎和另一个人同去。}然后，为解释这同一句话，\Quote{归荣耀于主，并显明你们乐意的心}，他补充说：\par

% structure: p0011 source=h2 target=subsection reason=relative-heading-rank; base=h2

\subsection{格林多后书 8:20}

\BibleQuote{我们避免这样，免得有人在我们经手的这丰富施舍中，责备我们。}\par

这是什么意思？一件相称于保禄德行的事，并显示了他极大关怀和谦抑自下。\Quote{因为}他说，\Quote{为避免有人怀疑我们，或对我们有丝毫挑剔，以为我们私吞了经手的金钱，因此我们派出这样的人，不仅一个，甚至两三个。}你看见他是怎样消除一切怀疑的吗？不是因福音的缘故，也不是单因他们被选，而是因他们是有验证的人，（正因此）被选出来，以免被怀疑。他并没有说\Quote{免得你们责备}，而是说\Quote{免得别人责备}。然而，他是为他们才这样做的；他在说\Quote{归荣耀于同一的主，并显明你们的乐意}时也暗示了这点；然而，他不愿伤害他们，所以换了一种方式表达：\par

\Quote{避免这样。} 他还不满足于此，又加上话安慰说：\par

\Quote{在我们经手的这丰富施舍中，} 将严厉与赞美相融合。为免他们感到受伤而说：\Quote{难道他必须暗中监视我们吗？我们竟这样可怜，以致一直被怀疑这些事吗？}他也对此提供了纠正，说：\Quote{你们所送出的钱数额巨大，这种丰盛，即钱的巨大数额，足以使恶毒的人产生怀疑，如果我们没有提供那种安全保障的话。}\par

% structure: p0016 source=h2 target=subsection reason=relative-heading-rank; base=h2

\subsection{格林多后书 8:21}

因为\BibleQuote{我们关心的是在主的眼中也在人的眼中做体面的事。}\par

有什么能与保禄相比？因为他没有说：\Quote{自取灭亡，那愿意怀疑此类事的人有祸了：只要我的良心不谴责我，我就不在乎那些怀疑的人。}相反，他们越软弱，他就越谦抑自下。因为应当帮助病人，而不是生气。然而，我们远离任何罪正如他远离这种怀疑一样？因为甚至魔鬼也不能怀疑那位有福的圣人会有这种不忠。但是，虽然远离这种恶毒怀疑，他仍做一切事，用尽一切方法，甚至不给那些可能想以任何方式怀疑错误的人留下阴影；而且，他不仅避免指控，还避免责备和最轻微的指责，甚至是单纯的怀疑。\par

% structure: p0019 source=h2 target=subsection reason=relative-heading-rank; base=h2

\subsection{格林多后书 8:22}

2. \BibleQuote{我们又同他们一起派去了我们的兄弟。}\par

看，他又加上了另一位，并对他也有称赞：既有他自己的判断，又有许多其他证人[对他]。\par

\BibleQuote{他，}他说，\BibleQuote{我们在许多事上屡次证明他是热心的，但现在更加热心。}而且因他自己的善行称赞了他，又因他对他们的爱而表扬他；他论\Person{弟铎}所说的，就是\BibleQuote{他出于非常的热心自愿前往；}这一点他也论及这人，说：\BibleQuote{但现在更加热心；}预先为他们存留了[证明]他们对格林多人之爱的种子。\par

然后，在显出他们的德行之后，他也劝勉他们代表他们，说：\par

% structure: p0024 source=h2 target=subsection reason=relative-heading-rank; base=h2

\subsection{\BibleRef{格后 8:23}}

\BibleQuote{若有人询问关于\Person{弟铎}；他是我的同伴，在你们中也是我的同工。}\par

\Quote{若有人问关于\Person{弟铎}的事？}\Quote{如果，}他说，\InnerQuote{有必要说什么，我要说，}\BibleQuote{他是我的同伴和我在你们中的同工。}因为他或是这意思；或，\Quote{如果你们愿意为\Person{弟铎}做什么，你们将不是对一个普通人做，因为他是\BibleQuote{我的同伴。}}而且当看似称赞他时，他赞扬他们，表明他们对他如此友善，以致于任何人显出是他的同伴，就足以成为他们中间的荣耀。但是，尽管如此，他并不满足于此，又加了另一件事，说：\Quote{在你们中的同工。}不仅是\Quote{同工，}\Quote{而且是在有关你们的事中，在你们的进步中，在你们的成长中，在我们的友谊中，在我们对你们的热情中；}最后一点最容易使他赢得他们的喜欢。\par

\BibleQuote{或论我们的弟兄：}\Quote{或是否你们愿意，}他说，\Quote{听到关于其他人的事：他们也有很大的理由被推荐给你们。}\Quote{因为他们也是，}他说，\Quote{是我们的弟兄，而且，}\par

\BibleQuote{众教会的使者，}也就是说，由众教会所派遣。然后，比一切更大的，\par

\BibleQuote{基督的光荣；}因为凡为他们所做的一切，都归于他。\Quote{既然如此，你们愿意接待他们如同弟兄，或如同众教会的宗徒，或如同为基督的光荣而工作；你们有许多向善的动机对他们怀有善意。因为就\Person{弟铎}而言，我要说，他既是\BibleQuote{我的同伴}，又是爱你们的人；就这些人而言，他们是\BibleQuote{弟兄}，他们是\BibleQuote{众教会的使者}，他们是\BibleQuote{基督的光荣}。}你们看，由此也很明显，他们是不为他们所知的那种人？否则他就会用那些他也用来称赞\Person{弟铎}的事来称赞他们，即他对他们的爱。但是既然他们还不为他们所知，\Quote{接待他们，}他说，\Quote{如同弟兄，如同众教会的使者，如同为基督的光荣而工作。}因此他接着说：\par

% structure: p0030 source=h2 target=subsection reason=relative-heading-rank; base=h2

\subsection{\BibleRef{格后 8:24}}

\BibleQuote{所以，请你们向他们，在众教会的面前，显出你们爱的证明，以及我们为你们所夸耀的证明。}\par

\BibleQuote{现在显明，}他说，\BibleQuote{你们如何爱我们；以及我们如何不是轻率地、也不是虚荣地以你们夸口：而你们要显明这一点，如果你们向他们显出爱。}然后他也使自己的话更加庄重，说：\BibleQuote{到众教会的面前。}他意指，到众教会的光荣和尊荣。\Quote{因为如果你们尊敬他们，你们就尊敬了派遣他们的教会。因为尊敬不只传到他们身上，也传到那些派遣他们、任命他们的人身上，而且更传到天主的光荣上。}当我们尊敬那些服侍他的人时，这接纳的善意就传到他身上，传到教会全体身上。现在这也不是小事，因为那个聚会的势力很大。\par

3. 至少可以肯定的是，教会的祈祷解开了\Person{伯多禄}的锁链，开启了\Person{保禄}的口；他们的声音在相当程度上装备了那些达到属神治理的人。\newline{}因此，确实如此：将要行按手礼的人那时也呼求他们的祈祷，他们以受启示者所知的欢呼来投票表示赞同：因为在未受启示者面前揭露一切是不允许的。\newline{}然而有些场合，司铎与属下之间毫无区别；例如，当我们参与神圣奥迹时；因为我们都同样配得同样的事：不像在\WorkTitle{旧约}之下，司铎吃某些东西，属下吃别的，民众不得分享司铎所吃的东西。但现在不是这样，而是摆在众人面前的是同一个身体和同一个杯。\newline{}在祈祷中，也可以看到民众贡献良多。因为为了附魔者，为了补赎者，司铎与他们共同祈祷；所有人同念一篇祈祷，那充满怜悯的祈祷。\newline{}再者，当我们把那些不能分享圣桌的人从神圣场所排除时，应当献上另一篇祈祷，我们全体一同俯伏在地，全体一同起身。\newline{}再者，在最神圣的奥迹本身中，司铎为民众祈祷，民众也为司铎祈祷；因为\Quote{与你的心灵同在}这句话，无非就是这个意思。\newline{}感恩的奉献也是共同的：因为他不仅独自感恩，全民众也一同感恩。因为他们首先发声，然后当他们赞同\Quote{适合而正确}时，他便开始感恩。\newline{}你为何惊奇民众在何处与司铎一同发声呢？因为实际上他们甚至与\Person{革鲁宾}及高天神圣一同献上那些神圣颂歌。\newline{}我说这一切，是为了让每位平信徒也警醒，使我们明白我们都是一个身体，彼此之间有如肢体与肢体之间的差异；并且不要把一切都推给司铎，而是我们自己也同样关心整个教会，如同关心我们共同的身体。\newline{}因为这条路将为我们提供更大的安全，并为你们提供更大的德性成长。\newline{}至少在此，就宗徒们而言，他们多么频繁地让平信徒参与他们的决定。因为他们拣选那七人时\BibleRef{宗 6:2}，先与民众商量；\Person{伯多禄}拣选\Person{玛弟亚}时，也与当时在场的所有人，无论男女，一同商议\BibleRef{宗 1:15}。\newline{}因为这里没有统治者的骄傲，也没有被统治者的奴性；而是属神的管理，在这方面最是独揽，承担更多为你们的辛劳和关怀，而不是追求更大的尊荣。\newline{}因为教会应当如此居住，如同一家；如同一个身体，如此安排；正如只有一个圣洗、一个圣桌、一个泉源、一个受造物和一个圣父。\newline{}那么，既然如此伟大的事物联合我们，我们为何分裂？为何被撕裂？因为我们不得不再次哀叹同样的事，我常为之悲叹。\newline{}我们现在的处境令人哀叹；我们彼此隔绝如此之深，而我们本应体现一个身体的合一。\newline{}因为这样，那更大的也能从较小的受益。因为如果\Person{梅瑟}从他岳父那里学到了他自己未曾觉察的合宜之事\BibleRef{出 18:14}，那么在教会中更可能发生。\newline{}那么，为何一个不信者所觉察的，属神的人却没有觉察？这是为了让当时所有人都明白他是人；尽管他分开海，劈开磐石，他仍需要天主的感动，那些行为不是出于人的本性，而是出于天主的权能。\newline{}所以，让另一个人起来发言；同样现在，如果某人不说合宜的话，就让另一个人起来发言；即使他是较卑微的人，但如果他说出切中要害的话，就要认可他的意见；即使他是最卑贱的人，也不要对他不敬。\newline{}因为这些人中没有谁与邻居相距如此遥远，如同\Person{梅瑟}的岳父与他相距之远，但他并没有不屑于听他，反而接受了他的意见，被说服并记录下来；并且不羞于将情况载入史册；从而打败了众人的骄傲。\newline{}因此，他将这个故事留给世界，如同刻在柱子上，因为他知道它将有益于许多人。\newline{}那么，我们不要忽视那些给我们有益劝告的人，即使他们是较低微的人；也不要坚持我们自己所提出的意见必须占上风；而是哪个主意看起来最好，就让大家认可它。\newline{}因为许多视力迟钝的人，藉着勤勉和专注，比那些视力敏锐的人更早察觉事情。\newline{}不要说：\Quote{你为何召我来议事，却不听从我的话？}这些指责不是参谋的，而是暴君的。\newline{}因为参谋只有权说出自己的意见；但如果另外的主意显得更有利，他却仍要执行自己的意见，他就不再是参谋，而是暴君，如我所说。\newline{}那么，我们不要这样做；而是使我们的灵魂摆脱一切傲慢和骄傲，让我们考虑，不是我们的主意如何成立，而是最好的意见如何胜出，即使它不是由我们提出的。\newline{}因为即使我们没有发现合宜之事，如果我们接受了他人所指出的，我们也不会没有大收获；我们必将从天主那里得到丰厚的赏报，同样也最会获得光荣。\newline{}因为那说出合宜之言的，是智慧的；同样，我们这些接受它的人，也将收获明智与坦诚的称赞。\newline{}如此，如果家庭和城邦如此，同样教会如果如此治理，她必获得更大的增长；我们自己也一样，既如此妥善安排了今生，必将领受将来的福乐：愿我们众人到达那里，因我们的主耶稣基督的恩宠与对人的爱，愿光荣归于祂，直到永远。亚孟。\par

\SourceRecord{Translated by Talbot W. Chambers.；Nicene and Post-Nicene Fathers, First Series；Vol. 12.；Edited by Philip Schaff.；Buffalo, NY: Christian Literature Publishing Co.,；1889.；<http://www.newadvent.org/fathers/220218.htm>.}

% structure: p0001 source=h1 target=section reason=page-title

\section{\WorkTitle{格林多后书第九篇讲道}}

% structure: p0002 source=h2 target=subsection reason=relative-heading-rank; base=h2

\subsection{\BibleRef{格后 9:1-2}}

关于为圣人服务的捐助，我实在不必写信给你们。\par

虽然他对此已经说了很多，但这里却说：\Quote{我实在不必写信给你们。}他的智慧不仅表现在：尽管他说了这么多，仍说\InnerQuote{我实在不必写信给你们，}而且还在于他再次提起这事。因为他稍前所说的话，是关乎那些接受捐款的人，要确保他们享有极大的尊荣；而他之前所说的（关于马其顿人的事，说\Quote{他们极度的贫穷，竟涌出慷慨的富足，}以及其他一切），是关于仁爱和施舍的。但尽管如此，尽管他之前说了那么多，并且还要再说，但他仍说：\Quote{我实在不必写信给你们。}他这样做，更是为了赢得他们的心。因为一个名誉如此之高，以至于不需要劝告的人，若显得低于或不及人们对他的评价，便会感到羞愧。他也常在控告中运用这种修辞手法——省略法，因为这非常有效。因为审判者看到控告者的宽宏大量，甚至不会产生怀疑。他会这样推理：\Quote{他既然可以说很多却不说，又怎会捏造不实之事呢？}他还让人推测出比他说出的更多的东西，并使自已享有良好品质的推定。他在劝告和赞美时也这样做。他说了\Quote{我实在不必写信给你们}之后，请注意他是如何劝告他们的。\par

\Quote{因为我知道你们的热情，我常以此在马其顿人面前夸耀你们。}他自己知道这件事已经非同小可，但更伟大的是，他也向别人宣扬了这事：因为这样带来的效力更大：他们不愿意如此广泛地蒙羞。你看见他智慧的目的了吗？他以他人的榜样（马其顿人）来劝勉他们，因为他说：\Quote{我告诉你们，天主在马其顿各教会所赐的恩宠。}他以他们自身的榜样来劝勉，因为他说：\Quote{你们不但从一年前就开始实行，而且也有愿意的心。}他以主的榜样来劝勉，因为他说：\Quote{你们知道}\BibleRef{格后 8:9}\BibleQuote{我们主耶稣基督的恩宠，祂本是富有的，却为了我们而成为贫穷的。}他又回到那有力的论点，即他人的行为。因为人有好胜之心。诚然，主的榜样本应最能吸引他们；其次是报酬的考虑；但因他们稍微软弱，故此事最能吸引他们。因为没有什么比好胜心更有效。但请注意他是如何以一种新颖的方式引入的。因为他没有说：\Quote{效法他们；}而是说什么？\par

\Quote{并且你们的热情激励了许多人。}你说什么？方才你说（他们）\Quote{出于自愿，再三恳求我们，}那么现在怎么又说\Quote{你们的热情}呢？\Quote{是的，}他说，\Quote{我们没有劝告，没有鼓励，我们只赞美了你们，只夸耀了你们，这就足以激励他们。}你看见他是如何用彼此激发他们的吗？用这些激发那些，用那些激发这些，并且在好胜心中也掺入了极高的赞美。然后，为了不使他们骄傲，他以温和的语气接着说：\Quote{你们的热情激励了许多人。}试想，那些已经给他人带来这种慷慨的人，若在这场捐助中自己却落后，那是何等的事。因此，他没有说：\Quote{效法他们，}因为这样不会激起如此大的好胜心，而是怎样？\Quote{他们效法了你们；那么你们作为教师，不可显得不如自己的学生。}\par

你看，在进一步激发和煽动他们时，他假装站在他们一边，仿佛在某种竞争和争论中支持他们。因为正如他上面所说：\Quote{他们出于自愿，再三恳求我们，以致我们劝勉弟铎，既然他早已开始，也当完成这一恩惠的事工；}这里他也这样说：\par

% structure: p0008 source=h2 target=subsection reason=relative-heading-rank; base=h2

\subsection{\BibleRef{格后 9:3}}

\Quote{为此，我打发众弟兄去，免得我们在这事上对你们的夸耀落了空。}\par

你看见他是何等焦虑和恐惧，恐怕他所说的话只是为劝勉的缘故吗？\Quote{但因事实如此，}他说，\Quote{我打发了众弟兄去；}我如此热切为你们，\Quote{免得我们的夸耀落了空。}他始终显得与格林多人同属一方，虽然他对所有人都同等关心。他说的话意即：\Quote{我极为你们自豪，我在众人面前夸耀，我甚至向他们夸耀了你们，所以如果你们有所欠缺，我就分担这羞愧。}但他这话是有条件的，因为他加了：\par

\Quote{在这一方面，}并非在一切方面；\par

\Quote{好使你们照我所说的，预备妥当。}\Quote{因为我并未说\InnerQuote{他们计划}，而是说\InnerQuote{一切齐备；如今他们方面毫无欠缺。}所以，}他说，\Quote{我愿你们用实际行动证明这事。}然后他甚至加剧了焦虑，说道：\par

% structure: p0013 source=h2 target=subsection reason=relative-heading-rank; base=h2

\subsection{\BibleRef{格后 9:4}}

\Quote{免得如果有什么从马其顿来的人与我同去，我们（更不必说你们）会在这自信上蒙羞。}身旁的观望者众多，甚至那些亲耳听过他夸口的人，这羞辱就更大了。他没有说\Quote{我正带着马其顿人来}\Quote{有马其顿人同我来}，免得显得刻意；而是怎么说呢？\Quote{免得如果有什么从马其顿来的人与我同去？}他说：\Quote{这事可能会发生，只是可能的。}这样使他的话不受怀疑；若他用另一种方式表达，反而会使他们更加争辩。看他如何引导他们，不仅出于属神的动机，也出于属人的动机。他说：\Quote{因为即使你们不大看重我，并且自信我会原谅你们，但也要考虑到马其顿人，免得他们来时发现你们。}他没有说\Quote{不情愿地}，而是说\Quote{没有预备好}，即尚未完全齐备。如果迟迟不奉献已是羞辱，那么根本不奉献或奉献不及应尽之分，将是何等大的羞辱！接着他以温和而尖锐的措辞说出后果，如此说：\Quote{我们（更不必说你们）会蒙羞。}他又缓和语气，说\Quote{在这自信上}，并非使他们更懈怠，而是表明他们在其他各方面都受认可，在这件事上也当有极大的坦然无惧。\newline{}\newline{}\newline{}\newline{}\par

% structure: p0015 source=h2 target=subsection reason=relative-heading-rank; base=h2

\subsection{\BibleBook{格林多}{后书} 9:5}

2. \Quote{所以我认为有必要恳请弟兄们，让他们预先预备好你们这祝福，使同一件事能以祝福的方式准备妥当，而非出于勒索。}\par

他又以不同的方式重新开始这个话题：为了不显得说这些话毫无目的，他断言这次旅行的唯一理由就是他们不致蒙羞。你看见他的话\Quote{我再写信给你们也是多余的}如何成了劝告的开端？至少你看见他为这项服务讲论了多少事。此外，还可以注意到，为了不显得自相矛盾（他曾说\Quote{多余的}，却详加论述），他转而谈论奉献的快速、丰盛和乐意，以此确保这一点。因为他要求这三件事。事实上，他在起初就提出了这三个要点：当他说\Quote{在患难大试炼中，他们洋溢的喜乐和极度的贫穷，涌现出他们慷慨的富足}时，他无非是说他们奉献得多，既乐意又快速；而且，不仅奉献不使他们痛苦，连患难也不使他们痛苦——而患难比奉献更难受。还有\Quote{他们将自己献给了我们}，这句话也显示了他们的乐意和信德的伟大。在这里他又论述了这些要点。因为由于慷慨和乐意是彼此对立的，一个人可能奉献得多却悲伤，另一个人为了不悲伤而奉献得少，请注意他如何以智慧关心二者。他没有说\Quote{少量而情愿的奉献胜于大量而被迫的奉献}，因为他希望他们既大量又情愿地奉献；但他怎么说呢？就是叫他们预先预备好你们这祝福，使同一件事能以祝福的方式准备妥当，而非勒索。他首先从最愉快、最轻松的一点开始，即\Quote{非出于必要}，因为他称之为\Quote{祝福}。请注意，在劝告的形式中，他同时描绘了果效的涌现和施予者充满祝福。他用这个措辞赢得了他们，因为没有人会痛苦地给予祝福。但他还不满足于此，还加上\Quote{不是出于勒索}。他说：\Quote{不要以为我们是勒索者，而是为给你们带来祝福。}因为勒索属于不情愿的人，所以不愿施舍的人就是出于勒索。然后他从这一点又转到慷慨奉献。\par

% structure: p0018 source=h2 target=subsection reason=relative-heading-rank; base=h2

\subsection{\BibleBook{格林多}{后书} 9:6}

\Quote{但我要说：}即，此外我还要说。说什么？\par

\Quote{那少种少收的，也要少收多种多收的。}他没有说\Quote{吝啬}，而是用了较温和的表达，采用\Quote{少种}这个词。并且他称这事为\Quote{播种}，为要你盼望报酬，想到收获，感到你所得的超过你给的。因此他没有说\Quote{那给予的}，而是\Quote{那播种的}；他没有说\Quote{你们若播种}，而是将这话说得一般化。他也没有说\Quote{大量地}，而是说\Quote{多收的}，这比\Quote{大量}更甚。他又回到原先喜乐的那一点，说道：\par

% structure: p0021 source=h2 target=subsection reason=relative-heading-rank; base=h2

\subsection{\BibleBook{格林多}{后书} 9:7}

\Quote{各人要照心里所决定的。}因为人若自主行事，比受强迫更为甘愿。因此他也强调这一点：因为说了\Quote{照他所决定的}，他加上：\par

\Quote{不是出于勉强，也不是出于不得已。}但他还不满足于此，还从圣经引证，说道：\par

\BibleQuote{天主喜爱乐意捐献的人。}你看他多么频繁地强调这一点？\BibleQuote{我这样说不是出于命令：}以及\BibleQuote{我在这事上给出我的意见：}以及\BibleQuote{是出于慷慨，而不是出于勉强，}又再次说\BibleQuote{不是出于勉强，也不是出于不得已；因为天主喜爱乐意捐献的人。}在这段话中，我认为指的是大量捐献的人；但宗徒却将其理解为乐意奉献。因为马其顿人的榜样和所有其他事物都足以产生慷慨，他并没有在那方面多说，而是强调不勉强奉献。因为如果这是德行的行为，而一切出于不得已的都被剥夺了赏报，那么他在此点上努力也是有理由的。他不仅劝告，还如他常做的那样，加上一个祈祷，说：\par

% structure: p0025 source=h2 target=subsection reason=relative-heading-rank; base=h2

\subsection{\BibleRef{格后 9:8}}

\BibleQuote{愿天主，那有能力者，使一切恩宠充满你们。}\par

通过这个祈祷，他消除了潜伏在慷慨面前的顾虑，这顾虑如今也阻碍了许多人。因为许多人害怕施舍，说：\BibleQuote{恐怕我变得贫穷，} \BibleQuote{恐怕我需要别人的帮助。}因此，为了消除这种恐惧，他加上这个祈祷，说：愿\BibleQuote{祂使一切恩宠丰富地赐予你们。}不仅仅是满足，而是\BibleQuote{使之丰富。} \BibleQuote{使之丰富}是什么意思？\BibleQuote{充满你们，}他意思是，\BibleQuote{以如此大的事物，使你们能在这慷慨上丰富。}\par

\BibleQuote{使你们在一切事上常常十分充足，能多行各种善事。}\par

看，即使在他的祈祷中，他也显示出伟大的智慧。他并非祈求财富或丰裕，而是祈求\BibleQuote{充足}。不仅如此，他并未紧逼或强迫他们从匮乏中给予，而是迁就他们的软弱；他只求\BibleQuote{充足}，同时表明他们不应滥用从天主领受的恩赐。\BibleQuote{使你们能多行，}他说，\BibleQuote{各种善事。} \BibleQuote{我祈求这个，}他说，\BibleQuote{使你们也能施予他人。}然而他并没有说\BibleQuote{施予}，而是说\BibleQuote{多行}。因为在物质事物上，他为他们祈求充足；但在属灵事物上，他甚至祈求丰裕；不仅在施舍上，也在所有其他事上，\BibleQuote{多行各种善事。}于是，他为他们引荐先知作为谋士，找出一段鼓励慷慨的证言，说：\par

% structure: p0030 source=h2 target=subsection reason=relative-heading-rank; base=h2

\subsection{\BibleRef{格后 9:9}}

\BibleQuote{正如经上所载：\newline{}他慷慨施舍，施舍给穷人；\newline{}他的正义永远存留。}\par

这就是\BibleQuote{多行}的含义；因为\BibleQuote{他慷慨施舍}一词无非表示大量给予。即使事物本身不存留，它们的结果却存留。这正是令人惊叹之处：当它们被保留时，就失去了；而当它们被分散时，却存留，甚至永远存留。这里\BibleQuote{正义}一词，他指的是对人的爱。因为当这爱大量倾注时，它使人成义，像火一样焚烧罪恶。\par

3. 所以，我们不要精打细算，而要慷慨撒种。难道你没有看见别人给戏子和妓女多少钱吗？至少把她们给舞女的一半给\Person{基督}。他们因虚荣而给台上的人的，至少你也要同样给饥饿的人。因为他们的确用数不清的金子给娼妓穿衣；但你连一件破旧的衣服也不给\Person{基督}的肉身，尽管你看见祂赤裸。这该得到什么宽恕？是的，这难道不该受到多大的惩罚？当他真的把那么多给那毁坏他、羞辱他的女人时，你却连一点点也不给那拯救你、使你更光耀的祂？但只要你把钱花在肚腹上、醉酒和放荡上，你就从不想贫穷；但当需要救济穷人时，你却变得比谁都穷。当你供养食客和谄媚者时，你高兴得好像有流不尽的财源；但如果你碰巧看到一个穷人，那么对贫穷的恐惧就缠绕着你。所以在那一天，我们必定会被定罪，既被我们自己定罪，也被别人定罪；既被那些行善的人定罪，也被那些作恶的人定罪。因为祂会对你说：\Quote{为什么你在那些合宜的事上不这样慷慨呢？但这里有一个人，当他给妓女时，什么也不考虑；而你，把财物献给你的主，祂曾吩咐你\BibleQuote{不要忧虑}\BibleRef{玛 6:25}，你却充满恐惧战栗。}那时你该得到什么宽恕呢？因为如果一个接受了恩惠的人不会忽视，反而会回报恩情，何况\Person{基督}呢？因为那连没有收到也给予的，收到后岂不更要给予吗？有人会说：\Quote{那么，如果有些花了很多钱的人需要别人的帮助呢？}你说的是那些倾尽所有的人；而你自己连一分钱也不给。答应我，你放弃一切，然后再问关于那些人的问题；但只要你是一个吝啬鬼，只给一点点你的财物，为什么对我抛出一堆借口和托词？因为我并没有引导你走上完全贫穷的崇高顶峰，而是目前我只要求你砍掉多余之物，只渴求知足。如今，知足的界限就是使用那些没有就不可能生活的东西。没有人阻止你使用这些；也不禁止你每日的食物。我说的是食物，不是宴饮；衣服，不是装饰。是的，如果要准确询问，这从最好的意义上才是宴饮。因为，请想一想。我们应该说谁更真正地享受了宴饮呢？是那个饮食为蔬菜，身体健康、毫无不适的人呢？还是那个拥有\Person{昔巴利人}餐桌、充满无数疾病的人呢？显然前者。所以，如果我们想同时过得奢华又健康，就不要追求更多；让我们设定这些知足的界限。让那能满足于豆类并能保持健康的人，不要再追求更多；但让那身体更弱、需要以蔬菜调养的人，不要被阻止这样做。但如果有人比这更弱，需要适量肉食的支持，我们也不会禁止他。因为我们劝告这些事，不是要杀死或伤害人，而是要砍掉多余之物；多余的就是超过我们需要的。因为当我们即使没有某物也能健康体面地生活时，那东西的增加肯定就是多余。\par

4. 同样，让我们也以这种态度来看待衣着、饮食、住所以及我们其他一切需要；凡事只求必需。因为多余的也是无用的。当你学会靠足够的东西生活之后，如果你有心效法那位寡妇，我们将引导你走向比这些更伟大的事。因为你还没有达到那位妇人的\OriginalTerm{哲学}{philosophy}，你还在为足够的东西忧虑。因为她甚至超越了这个；她把本应维持生计的全部都投了进去。那么，你还要为这些必需品烦恼吗？难道你不因被一个妇人战胜而羞愧吗？你不以效法她为耻，反而远远落后于她吗？她没有说我们所说的那些话——\Quote{如果我花光了所有，岂不是要被迫向他人乞讨？}——而是慷慨地舍弃了她所有的一切。关于旧约中厄里亚先知时代的那位寡妇，我们又能说什么呢？她所冒的风险不是贫穷，而是死亡和灭绝，不仅是自己的，还有她孩子的。因为她也没有期待从别人那里得到什么，而是即将死去。但有人说：\Quote{她看见了先知，这才使她慷慨。}但你们不是也看见了无数的圣人吗？我为什么要说圣人呢？你们看见了先知们的主在乞求施舍，却仍然没有变得仁慈；尽管你们的钱柜一个接一个地满溢，却连多余的都不肯分施。你说什么？来到她面前的那位是先知，这使她如此大度吗？这本身就值得极大的钦佩，因为她相信他是一个伟大而奇妙的人物。她怎么没有像野蛮的异族妇人所可能推理的那样说：\Quote{如果他是先知，他就不会向我乞讨。如果他是天主的朋友，天主就不会无视他。就算犹太人因罪遭受这惩罚，但这个人为何受苦？}但她没有这些念头；她向他打开了家门，更在开门之前就向他敞开了心扉；把她所有的一切都摆在他面前；她把天性放在一边，不顾自己的孩子，把陌生人置于一切之上。那么，如果我们竟比不上一个妇人、一个寡妇、穷人、外邦人、蛮族、有孩子的母亲，而且她不知道我们所知道的这些事，我们该受多大的惩罚呢！因为我们有身体的力气，并不因此就是刚毅的人。因为只有那内心有力量的人才有这种美德，即使他卧床不起；没有这美德，即使有人靠体力移山，我也会说他并不比一个女孩或老妇人更强。因为前者与无形的灾祸搏斗，而后者甚至不敢正视它们。为使你明白这是刚毅的标准，请从这例子中体会。还有什么能比那位妇人更刚毅呢？她既抵抗了天性的暴政，又抵抗了饥饿的力量，还抵抗了死亡的威胁，坚强地站立，战胜了这一切。请听听基督如何称扬她。因为祂说：\Quote{在厄里亚的日子里，有许多寡妇，但先知没有被派到她们中的任何一个，只派到她那里。}\BibleRef{路 4:25}我要说一件伟大而惊人的事吗？这位妇人比我们的父亲亚巴郎待客更慷慨。因为她没有像他那样\Quote{跑到牛群里}\BibleRef{创 18:7}，而是用那\Quote{一把面}\BibleRef{列上 17:12}超越了所有以好客著称的人。因为亚巴郎的卓越在于他亲自去做那一服务；而她的卓越在于为了陌生人，她连自己的孩子都不顾惜，而且那时她并不指望来世。但我们，尽管有天国存在，有地狱在威胁，尽管（比一切都更重要的是）天主为我们行了如此伟大的事，并对这些事感到喜悦和欢欣，我们却慵懒地往后躺。我恳求你们，不要这样！让我们\Quote{广施}，让我们\Quote{济贫}，按照应当的施与。因为天主定义多与少，不是根据所施与的量，而是根据施与者的家资。确实，常常你们这些投入一百金币的人，反不如那投入一文小钱的人，因为你们是从多余中投入。但即使只做这一件，你们也会很快达到更大的慷慨。分散财富，好使你们积聚正义。因为财富拒绝与我们同来；然而通过它，虽不与之同来，却使我们拥有它。因为贪财和正义不可能同住；它们的帐篷是分开的。所以，不要固执地试图调和不能相容的事，而是驱逐篡位者贪婪，如果你们想得到天国。因为正义是合法的王后，她使奴隶成为自由的人，而贪婪则相反。因此，让我们极其热忱地逃避一个，欢迎另一个，以便在此生获得自由，并得到天国，藉着我们的主耶稣基督的恩宠和慈爱，祂与圣父及圣神，同享光荣、权能、尊崇，从今世直到永远，世世无穷。阿们。\par

\SourceRecord{Translated by Talbot W. Chambers.；Nicene and Post-Nicene Fathers, First Series；Vol. 12.；Edited by Philip Schaff.；Buffalo, NY: Christian Literature Publishing Co.,；1889.；<http://www.newadvent.org/fathers/220219.htm>.}

% structure: p0001 source=h1 target=section reason=page-title

\section{《格林多后书》第二十篇讲道}

% structure: p0002 source=h2 target=subsection reason=relative-heading-rank; base=h2

\subsection{\BibleRef{格后 9:10}}

\BibleQuote{那供给撒种者种子、又供给食粮作你们食物的主，必将加增你们的种子，增多你们仁义的果实。}\par

由此尤可钦佩保禄的智慧：他先以属神的理由，又以现世的理由劝勉之后，在论及报偿时，又同样做，使他所提及的报偿兼具双重性质。例如，\BibleQuote{他散施，他给了穷人，他的仁义永远常存}属于属神的报偿；而\BibleQuote{加增你们的种子}则属于现世的报偿。然而他并不停留于此，反而又转回属神之事，将两者不断并列；因为\BibleQuote{增多你们仁义的果实}是属神的。这样做是为使他的言论富于变化，根除他们那怯懦而气馁的推理，并用许多论据驱散他们对贫穷的恐惧，正如他现在所举的例子。因为如果连播种于田地的，天主都赐予；连供养肉体的，祂都赐予丰盛；那么对那些耕耘天国之田、关怀灵魂的人，祂岂不更加赐予？因为祂更愿意这些事享受祂的眷顾。然而，他并不以推理的方式陈述，也不像我那样说，而是以祈祷的形式；如此，既使所指明显，又更引导他们盼望，不仅从通常发生的事，也从他自己的祈祷：因为他说：\BibleQuote{愿祂供给并加增你们撒种的种子，增多你们仁义的果实。}这里他又以不惹人怀疑的方式暗示了慷慨施予，因为\Quote{加增并增多}这些词正是指明这一点；同时，他允许他们只求必需之物，说\Quote{食物作食粮。}因为这一点在他身上尤其值得钦佩（这也是他之前成功确立的一点），即：在必需之事上，他只允许他们求所需之用，不逾本分；但在属神之事上，他劝勉他们为自己获得大量丰裕。因此他上文也说\BibleQuote{使你们有余，可以多行各种善事}；这里又说\BibleQuote{那供给食粮作食物者，加增你们撒种的种子}；即属神的种子。因为他所要求的不仅是施舍，而是慷慨的施舍。所以他不断称其为\Quote{种子。}因为如同撒在地上的谷粒长出茂盛的庄稼，同样，施舍也产生许多把的仁义，并结出不可言喻的果实。然后，他为求得大丰收而祈祷后，又指出他们应用之于何处，说：\par

% structure: p0005 source=h2 target=subsection reason=relative-heading-rank; base=h2

\subsection{\BibleRef{格后 9:11}}

\BibleQuote{使你们凡事富足，以致十分慷慨，这慷慨藉着我们使天主得到感谢。}\par

不是叫你们把它消费在不合宜的事上，而是用在能带给天主许多感谢的事上。因为天主使我们能支配大事，却把那较小的留给自己，将更大的赐给我们。因为肉体的食物完全由祂掌管，但属神的食物祂却交给了我们；因为我们能自己支配，看我们所要呈现的庄稼是否茂盛。这里不需要雨水和气候的变化，只需要意志，它们便直达天上。他在这里所称的慷慨，就是大方施舍。\BibleQuote{这慷慨藉着我们使天主得到感谢。}因为所行的不只是施舍，也是许多感谢的缘由；不仅如此，还不只是感谢，还有其他许多益处。他继续提及这些，为指出这是许多善工的原因，从而激励他们更加热忱。\par

2. 那么这些许多善工是什么呢？听他说道：\par

% structure: p0009 source=h2 target=subsection reason=relative-heading-rank; base=h2

\subsection{\BibleRef{格后 9:12-14}}

\BibleQuote{因为这供职的事，不但补足了圣徒的缺乏，而且因着许多感谢也丰盛地归于天主；他们因着这供应所证明的你们的服从，便为你们对福音的宣认的服从，以及你们慷慨的捐助对他们和众人所显的，而荣耀天主；同时他们也以恳求为你们祈祷，因天主在你们身上所赐的极丰富的恩宠而记念你们。}\par

他所说的是：\Quote{首先，你们不仅补足了圣徒的缺乏，而且甚至有余；}即\Quote{你们供给他们超过所需；其次，藉着他们，你们向天主献上感谢，因为他们为你们对宣认的服从而荣耀天主。}为免他使人以为他们仅为此（即因为他们收到了什么）而感谢，看他使他们的心志多么高尚，正如他亲自对斐理伯人说：\BibleQuote{我并非寻求馈赠。}\BibleRef{斐 4:17}\Quote{我也为他们作证同样的事。因为他们确实因你们供给他们的所需、减轻他们的贫穷而喜乐；但更因你们如此顺服福音而喜乐；你们的慷慨捐助就是这事的证据。}因为福音正命令这个。\par

\BibleQuote{以及你们慷慨的捐助对他们和众人所显的。}他说：\Quote{为此他们荣耀天主，因你们如此慷慨，不仅对他们，也对众人。}这也成为对他们的赞美，因为他们甚至为施予别人的事而感谢。\Quote{因为，}他说，\Quote{他们不仅尊重自己的事，也尊重别人的事，尽管他们处于极端的贫困中；这证明他们的大德。}因为没有什么比穷人的嫉妒心更强烈的了。但他们却摆脱了这种激情；不仅不因你们赐予他人之物而感到痛苦，反而为此喜乐，不亚于为自身所领受的。\par

\Quote{他们自己也以恳求。}\Quote{因为这些事，}他说，\Quote{他们感谢天主，但为你们的爱和你们的聚会，他们恳求祂，使他们堪当见到你们。因为他们渴望这个，不是为金钱的缘故，而是为能亲眼看见那赐予你们的恩宠。}你看到保禄的智慧吗？他如何高举他们之后，又将一切归于天主，称那事为\Quote{恩宠}？因为他既说了他们的大事，称他们为服役者并将他们高举到高处（因为他们献上服务，而他自己只是管理），又称他们为\Quote{被验中的}，他表明天主是这一切的根源。他自己再次与他们一起献上感谢，说：\par

% structure: p0014 source=h2 target=subsection reason=relative-heading-rank; base=h2

\subsection{\BibleRef{格后 9:15}}

\Quote{感谢天主，为祂不可言喻的恩赐。}\par

这里他称\Quote{恩赐}，甚至那些由施舍所成就的许多善事，无论是对于领受者还是给予者；或者毋宁说，那些不可言喻的美善事物，祂借着降临以大慷慨赐予全世界，这可能是最可能的。因为他要同时使人清醒并使他们更慷慨，便提醒他们从天主所领受的恩惠。因为这在激励人修一切德行上非常有效；因此他以它结束他的讲论。但如果祂的恩赐是不可言喻的，那么那些对祂的本体提出好奇问题的人的狂热，还有什么能与之相比呢？不仅祂的恩赐不可言喻，而且那\Quote{平安}也\Quote{超越一切理解，}\BibleRef{斐 4:7}，祂借此使天上与地下的事物和好。\par

3. 既然我们享有如此大的恩宠，让我们努力表现出与之相称的美德生活，并重视施舍。我们若避免过度、醉酒和贪食，就会做到这一点。因为天主赐予食物和饮料，不是为了过度，而是为了滋养。因为不是酒导致醉酒，因为如果是这样，人人都必须醉酒。但有人说：\Quote{若能大量饮用而无害，就更好了。}这是醉汉的话。因为若大量饮用确有害处，而你仍不停止你的过度；若这是如此可耻和有害，而你仍不停止你那堕落的渴望；若有可能大量饮用而毫无损害，你在过度中会止步何处？你岂不会渴望连河流都变成酒吗？你岂不会毁灭和败坏一切吗？若饮食有节制，超过便受损害，而即使如此你仍不能忍受约束，反而折断之，夺取别人所有，以满足这贪食的邪恶暴政；若这自然的节制被废除，你岂不会做何事？你岂不将全部时间花费在这事上？那样巩固一种如此不合理的欲望，而不防止过度产生的害处，岂是好的？这又会生出多少其他害处？\par

但愚昧的人啊！他们沉溺于醉酒和一切其他放荡中如同淤泥，稍微清醒一点后，坐下来说的只是这样的话：\Quote{为什么这事以这种方式结束？}而此时他们本该谴责自己的过犯。因为，与其像你现在说的\Quote{为什么祂设定了界限？为什么一切不都毫无秩序地进行？}不如说：\Quote{为什么我们不停上醉酒？为什么我们总不满足？为什么我们比无理性的受造物更愚昧？}因为这些事他们该互相询问，并听从宗徒的声音，学习他见证格林多人从施舍中得来的许多好事，并抓住这宝藏。因为轻视金钱使人蒙悦纳，如他所说；并使天主受光荣；温暖爱德；在人心成就灵魂的高超；并使他们成为司铎，甚至是一种带来大赏报的司祭职。因为仁慈的人不穿着垂到脚的长袍，也不摇铃，不戴冠冕；但他是披着仁慈的外袍，比圣衣更神圣；他被傅油，不是由物质元素构成，而是由圣神所产生；他戴着怜悯的冠冕，因为经上说：\Quote{祂以怜悯和慈爱给你作冠冕；}\BibleRef{咏 102:4}，而且代替佩戴刻有天主名的金牌，他自己相似天主。因为如何？\Quote{你们，}祂说，\Quote{要像你们的天父一样。}\BibleRef{玛 5:45}\par

你愿也看见祂的祭台么？不是\Person{贝匝肋耳}建造了它，也不是任何其他的人，而是天主自己；不是用石头，而是用一种比天空更光亮的材料，即理性的灵魂。但是司铎进入了至圣所。当你奉献这祭献时，你可以进入更可畏的地方，那里除了\BibleQuote{你那在暗中看见的父}\BibleRef{玛 6:4}之外，无人同在，也没有别的观看者。有人说：\Quote{怎么可能无人观看，既然祭台是公开设立的？}因为这就是令人惊叹之处：在古代，双重的门和帷幔造成了隔绝；但现在，即使你在公开的视野中奉献你的祭献，你也可以如同在至圣所内、以更可畏的方式行事。因为当你不是为在人前炫耀而做时，即使全世界都看见了，也没有人看见，因为你是如此做的。因为祂不是简单地说：\BibleQuote{不要}在人\BibleQuote{前}做，而是加上：\BibleQuote{为叫人看见。}\BibleRef{玛 6:1}。这个祭台是由基督的肢体本身组成的，主的身体成了你的祭台。那么，你当敬畏那祭台；你在主的肉身上祭献牺牲。这个祭台甚至比我们现在所使用的这个更可畏，不仅比古代所使用的那个更可畏。不，不要喧嚷。因为那个祭台因放在其上的牺牲而可敬；但那个仁慈之人的祭台不仅为此缘故，而且因为它本身就由使另一个祭台可敬的牺牲所组成。再者，这个祭台在本性上只是石头，但因接受基督的身体而成为圣洁；但那祭台是圣洁的，因为它本身就是基督的身体。因此，你这位平信徒所站立的这个祭台，比那个更可畏。那么，与这个相比，\Person{亚郎}在你眼中还算什么呢？或是他的冠冕、他的铃铛、或\OriginalTerm{至圣所}{holy of holies}呢？因为从今以后，何需将我们的比较引到\Person{亚郎}的祭台呢？然而，你确实尊敬这个祭台，因为它接受基督的身体；但你却侮辱那本身就是基督身体的人，当他处于困境时，你忽视他。这个祭台你可以随处看见，躺在小巷里和市场上，并且可以每时每刻在它上面献祭；因为在这个祭台上也举行祭献。正如司铎站着呼求圣神，你也同样呼求圣神，不是用言语，而是用行为。因为没有什么比大量倾注这油更能维持和点燃圣神的火。但如果你也想看见祭台上所放之物变成了什么，请到这里来，我将指给你看。那么，这个祭台的烟是什么？这甘饴的香气是什么？是赞美和感恩。它上升多远？直到天上么？绝非如此，而是它越过了天本身、以及天上的天，甚至达到君王的宝座。因为经上说：\BibleQuote{你的祈祷和你的施舍已经上升到天主面前。}\BibleRef{宗 10:4}。感官所感知的甘饴香气不能穿透空气很远，但这香气却打开了穹苍本身。而你虽然沉默，但你的行为在说话：一种颂谢的祭献被奉献，不是宰杀牛犊或焚烧皮革，而是一个属神的灵魂呈上她自身的奉献。因为这样的祭献比任何一种仁爱更蒙悦纳。因此，当你看见一位贫穷的信徒时，你当思想你看见了一个祭台；当你看见这样一个人——一个乞丐时，不仅不要侮辱他，反而要尊敬他，如果你看见别人侮辱他，要阻止、斥退这事。因为这样你便能使天主对你和善，并获得所应许的福乐；愿我们众人藉着我们的主耶稣基督的恩宠和慈爱，得以达到那里，祂和父及圣神，是荣耀、权能、尊威，从今时直到永远，世世无穷。阿们。\par

\SourceRecord{Translated by Talbot W. Chambers.；Nicene and Post-Nicene Fathers, First Series；Vol. 12.；Edited by Philip Schaff.；Buffalo, NY: Christian Literature Publishing Co.,；1889.；<http://www.newadvent.org/fathers/220220.htm>.}

% structure: p0001 source=h1 target=section reason=page-title

\section{\WorkTitle{论 \BibleBook{格林多}{后书} 讲道词 21}}

% structure: p0002 source=h2 target=subsection reason=relative-heading-rank; base=h2

\subsection{\BibleRef{格后 10:1-2}}

\BibleQuote{我这保禄，就是当着你们面时是谦卑的，不在你们那里时，对你们却是勇敢的，现在因基督的良善与温和恳求你们。我恳求你们，不要让我在你们那里时，不得不勇敢地用我自以为该大胆对待某些人的信心，他们以为我们是按血肉行事的。}\par

他既已按恰当的方式完成了关于施舍的论述，又已表明自己爱他们多于被他们所爱，并讲述了自己忍耐和试炼的种种境遇，如今适时地进入更带责斥意味的要点，暗示假宗徒，并以更令人不悦的内容和自夸来结束他的论述。因为他在整封书信中也以此为务。他觉察到这一点，便时常纠正自己，直言不讳地说：\Quote{我们岂是又开始举荐自己呢？}\BibleRef{格后 3:1} 随后又说：\Quote{我们不是又举荐自己，而是给你们一个荣耀我们的机会}\BibleRef{格后 5:12}，之后又说：\Quote{我成了在夸耀上的愚妄人，是你们逼我的}\BibleRef{格后 12:11}。他用了许多这样的纠正。将这部书信称为保禄的\OriginalTerm{颂词}{eulogium}也不算错，因为他如此多次提及他的恩宠和忍耐。因为在他们当中有些自视甚高，将自己置于宗徒之上，并指控他为自夸者、一无是处、不教导纯正道理；（而这本身就是他们自己败坏的最确凿证据；）请看他是如何开始责备他们的：\Quote{我这保禄本人。}你看到这里有何等严厉、何等威严吗？因为他想说：\Quote{我恳求你们不要逼我，不要让我使用我的权能反对那些轻视我们、视我们属血肉的人。}这比在先前那封书信中对他们发出的威胁更为严厉：\Quote{你们愿意我带着棍棒到你们那里去呢，还是怀着爱和温柔的精神？}\BibleRef{格前 4:21} 接着又说：\Quote{有些人以为我不会到你们那里去，就自高自大起来；但我必要快去找寻，并且要知道的，不是那些自高自大之人的言语，而是他们的能力。}\BibleRef{格前 4:18-19}因为在此处他显示了两点，既显示了他的权能，又显示了他的智慧和宽容；他如此恳求他们，且如此热切，以免被迫展现他所有的报应之权，击打惩罚他们，并索要极端的刑罚。因为他要说这话时暗示了这一点：\Quote{但我恳求你们，使我不要当着你们的面时，不得不勇敢地用我自以为该大胆对待某些人的信心，他们以为我们是按血肉行事的。}现在，我们暂且谈谈开头。\Quote{我这保禄本人。}这里大有强调、大有分量。他在别处也说：\Quote{我保禄告诉你们}\BibleRef{迦 5:2}，又说：\Quote{保禄，老人}\BibleRef{费 1:9}，在另一处又说：\Quote{她曾帮助过许多人，也帮助过我}\BibleRef{罗 16:2}。同样在这里，\Quote{我这保禄本人。}这本身是一件大事，因为他亲自恳求；但他随后所加的更是伟大，他说：\Quote{因基督的良善与温和。}他为了大大羞辱他们，就提出了\Quote{良善与温和}，以此使他的恳求更有力；仿佛在说：\Quote{你们要尊敬基督的温和，因我藉此恳求你们。}他说这话，同时也表明，尽管他们对他施加了如此强的压力，他自己却更倾向于温和：他不攻击他们，是出于温和，而非缺乏权能；因为基督也照样行了。\par

\Quote{当着你们面时是谦卑的，不在你们那里时，对你们却是勇敢的。}请问这是什么？他必定是用讥讽的话，引用他们的言辞。因为他们说：\Quote{他当面时确实微不足道，既贫穷又卑贱；但不在时却膨胀自夸，高抬自己来反对我们，并且恐吓。}因此他至少在后文也暗示了，他们引用说：\Quote{他的书信}，他们说，\Quote{是严厉的，但他亲身在场却是软弱的，言语不足介意。}\BibleRef{格后 10:10}他要么用讥讽的语气，显示严厉之意，说道：\Quote{我，卑贱的我，卑微的我，当面时（如他们所说）和不在时，却是高傲的；}要么意味着，即使他说了大事，也不是出于骄傲，而是出于对他们的信任。\par

\BibleQuote{但是，我恳求你们，免得我临在时，凭着那我认为敢对那些以为我们按血肉行事的人所拥有的自信，而显勇敢。} 你们看到他是多么愤慨，以及如何彻底驳斥了他们的那些说法吗？因为他说：\Quote{我恳求你们，不要迫使我显示，即使我临在时，我也是刚强且有能力的。} 既然他们说：\Quote{当他不在这里时，他公然对我们大胆并抬高自己}，他就用了他们的话：\Quote{因此我恳求，不要迫使我使用我的能力。} 因为这就是\BibleQuote{自信}的意思。他没有说\Quote{我准备用}，而是\Quote{我认为用}。\Quote{因为我尚未对此下定决心；但他们给了我足够的理由，但即使如此，我也不希望这样。} 然而他这样做不是为了为自己辩护，而是为了福音。如果在他必须为信息辩护时，他不严厉，反而退缩和拖延，并恳求不要有这样的必要；那么他更不会为自己辩护而这样做。\Quote{那么，请赐给我这个恩惠，}他说，\Quote{就是不要迫使我显示，即使我临在时，我也能对任何必要的人刚强；也就是说，惩罚他们。} 你们看到他多么没有野心，他做事不为了炫耀，因为即使是在必要的情况下，他也毫不犹豫地称此行为为勇敢。因为他说：\BibleQuote{我恳求你们，}\BibleQuote{免得我临在时，凭着那我认为要敢用的自信而显勇敢} 反对一些人。因为以下尤其是老师的本分：不急于报复，而是改过自新，并且总是勉强而缓慢地施行惩罚。请问，他如何描述他所威胁的人？\BibleQuote{那些以为我们按血肉行事的人}：因为他们指责他是伪君子、邪恶、自夸。\par

% structure: p0007 source=h2 target=subsection reason=relative-heading-rank; base=h2

\subsection{\BibleRef{格后 10:3}}

2. \BibleQuote{因为我们虽然在血肉中行走，却不按照血肉作战。}\par

在此，他继续用他所用的比喻来警告他们，因为他说：\Quote{我们确实被血肉包围；我承认，但我们不靠血肉生活；}或者更确切地说，他连这个也没有说，而是暂时保留它，因为它属于对他生活的赞词：但他先谈论宣讲，并表明它不是出于人，也不需要来自下面的帮助。因此，他没有说\Quote{我们不按血肉生活}，而是\BibleQuote{我们不按血肉作战}，也就是说，\Quote{我们从事了一场战争和战斗；但我们不凭血气的武器作战，也不靠任何人的帮助。}\par

% structure: p0010 source=h2 target=subsection reason=relative-heading-rank; base=h2

\subsection{\BibleRef{格后 10:4}}

\BibleQuote{因为我们的武器不是属血肉的。}\par

因为什么样的武器是属血肉的？财富、荣耀、权力、流利、聪明、诡计、谄媚、虚伪，以及其他类似的一切。但我们的不是这一类的：那么它们是怎样的呢？\par

\BibleQuote{在天主前是有能力的。}\par

并且他没有说\Quote{我们不是属血肉的}，而是\BibleQuote{我们的武器}。因为正如我所说，目前他谈论宣讲，并将全部能力归于天主。他没有说\Quote{属神的}，尽管这本来是\BibleQuote{属血肉的}的恰当对立词，而是\BibleQuote{有能力的}，这样一来也暗示了前者，并表明他们的武器是软弱无力的。注意他毫无骄傲；因为他没有说\Quote{我们是有能力的}，而是\BibleQuote{我们的武器在天主前是有能力的}。\Quote{我们并没有使它们如此，而是天主自己。}因为他们被鞭打、被逼迫、遭受无数不可治愈的伤害，这些都是软弱的证明：为了显示天主的力量，他说\BibleQuote{但它们在天主前是有能力的}。因为这尤其显示祂的力量，即祂借着这些事获得胜利。所以即使我们被它们包围，然而却是祂在作战并借着它们工作。然后他长篇赞扬它们，说：\par

\BibleQuote{攻破堡垒}。并且免得你们听到堡垒时想到任何物质的东西，他说：\par

% structure: p0016 source=h2 target=subsection reason=relative-heading-rank; base=h2

\subsection{\BibleRef{格后 10:5}}

\BibleQuote{攻破计谋。}\par

首先用比喻来强调，然后用这额外的表达来宣示战争的神性特征。因为这些堡垒围攻的是灵魂，不是身体。因此它们比其他的更强，所以它们所需的武器也更强大。但堡垒指的是希腊人的骄傲，以及他们的诡辩和推理的力量。然而，\Quote{这些武器，}他说，\Quote{挫败了一切起来反对它们的东西；因为他们攻破计谋，}\par

\BibleQuote{并且一切高张起来反对天主之知识的事。} 他坚持这个比喻，以便使强调更强烈。\Quote{因为即使有堡垒，}他说，\Quote{即使有防御工事，即使有任何其他东西，它们都会在这些武器面前屈服让步。}\par

\Quote{并将一切思想掳去，使之服从基督。}然而，\Quote{\OriginalTerm{被掳}{captivity}}这个名称听起来不好，因为它是自由的毁灭。那么，他为何使用它？它是带着自身的含义，涉及另一点。因为\Quote{被掳}一词传达两个观念：失去自由，以及被如此猛烈地压倒以致不再起来。因此，他是在这第二个意义上使用它。正如当他说\Quote{我抢夺了别的教会}\BibleRef{格后 11:8}，他指的不是暗中偷取，而是剥取并拿走他们的一切；同样在这里，当他说\Quote{掳去}时也是如此。因为战斗并未对等维持，而是他轻易获胜。因此，他不仅说\Quote{我们战胜并占上风}，而且说\Quote{我们甚至将他们\InnerQuote{掳去}}。正如上面，他并没有说\Quote{我们向\InnerQuote{\OriginalTerm{堡垒}{strongholds}}推进器械}，而是说\Quote{我们将其摧毁，因为我们的武器远胜}。因为他说：\Quote{我们不是以言语争战，而是以行为对抗言语，不是凭血肉的智慧，而是凭良善和德能的精神。}那么，我怎么会追逐荣誉、口中夸耀、以书信威胁呢？（正如他们指控他，说\Quote{他的书信是严厉的}），当我们的能力并不在这些事上时？但是，既然说了\Quote{将一切思想掳去，使之服从基督}，因为\Quote{被掳}的名称令人不快，他随即终止了比喻，说\Quote{使之服从基督}：从奴役到自由，从死亡到生命，从毁灭到救恩。因为我们来不只是击倒，而是将反对我们的人带到真理面前。\par

% structure: p0021 source=h2 target=subsection reason=relative-heading-rank; base=h2

\subsection{格林多后书 10:6}

3. \Quote{并且我们已经预备好了，等你们完全服从的时候，要惩罚一切的不服从。}\par

这里他也警戒这些，不单是那些人：因为他说：\Quote{我们等候你们，等我们用劝诫和警告改造了你们，把你们从他们的团体中净化分别出来；那时，当只剩下那些不可救药的人时，我们才予以惩罚，因为我们看到你们确实已经与他们分离了。因为即便现在你们确实服从，但并不完全。}有人说：\Quote{如果你现在就这样做，会带来更大的益处。}\Quote{绝不可，因为如果我现在就做，我也会连累你们受罚。然而，惩罚他们固然应当，却要宽恕你们。但如果我宽恕了，似乎是因为偏爱这样做：而这是我不愿的，我要先纠正你们，然后再处理他们。}有什么比宗徒的心肠更温柔的呢？因为他看到自己人混杂在外人之中，确实想施加打击，却忍耐着，克制自己的怒气，直到这些人退去，好让他只打击这些人；甚至也不是这些人。为此，他如此威胁，并说他渴望只将他们分开来惩罚，好使他们因畏惧而改正，而他不向任何人发泄愤怒。因为他就像一位最优秀的医生、共同的父亲、保护者和监护人，他如此行事，如此关心所有人，移除障碍，遏制为害者，四处奔走。因为他不靠争战如此成就事业，而是如同迈向一场现成的、轻易的胜利，他树立了战利品，暗中破坏、推倒、推翻魔鬼的堡垒和恶魔的器械，并将他们的全部战利品运到基督的营中。他甚至没有稍作喘息，从这些人跳到那些人，从那些人再到别人，如同一位极为能干的将军，每天甚或每小时都树立战利品。因为当他身穿一件小外衣进入战场时，保禄的舌头攻取了敌人的城市，连同他们的士兵、弓、矛、标枪以及一切。\par

因为他只说话；他的话语比任何火焰更猛烈地扑向敌人，驱赶了魔鬼，并将它们附身的人带到他面前。因为当他驱逐那个魔鬼（那恶者）时，五万个巫师聚集起来焚烧了他们的魔法书，并归向了真理\BibleRef{宗 19:19}。如同在战争中，当一座塔楼倒塌或一个暴君被推翻，他所有的党羽都丢弃武器，奔向敌方的将军；那时也确实如此发生。因为当魔鬼被驱逐后，他们全都被围困，并丢弃，甚至毁灭了他们的书籍，跑到保禄的脚前。但他将自己置身于对抗世界如同对抗一支军队，任何地方都不停止他的行进，行事如同一个被赋予翅膀的人：时而治愈瘸子，时而复活死人，时而弄瞎另一个人（我指的是那个巫师）；即使被关在监狱里也不休息，甚至在狱中也使狱卒归向自己，实现了我们所谈论的良善的\OriginalTerm{被掳}{captivity}。\par

4. 我们也当按我们的能力效法他。我为什么说按我们的能力呢？因为凡愿意的，甚至可以接近他，看见他的英勇，并效法他的英雄气概。因为他现在仍在做这工作，\BibleQuote{推翻想象及一切反抗天主知识的自高之事}尽管许多异端者试图将他碎尸万段，但即使被肢解，他仍然显出巨大的力量。因为\Person{马尔西翁}和\Person{摩尼}确实使用他，但却是将他割裂之后；即使如此，他们仍被他的各个肢体所驳斥。因为这位勇士的一只手落在他们中间，就使他们彻底溃败；一只脚留在别人中间，就能追逐并打倒他们，为的是叫你学习他力量的丰盈，即使被砍去四肢，他仍能毁灭所有仇敌。\Quote{然而，有人会说，这是一个颠倒的例子，那些彼此争斗的人都使用他。}这诚然是颠倒，但不在\Person{保禄}身上（天主不许），而在使用他的人身上。因为他不是杂色的，而是纯粹清晰的，但他们将他的话曲解以适应自己的见解。\Quote{为什么，有人会说，它们这样说，以致给那些想要的人以把柄？}他并没有给出把柄，而是他们的狂热误用了他的话；因为整个世界也是奇妙伟大的，是天主智慧的可靠证明，\BibleQuote{诸天宣扬天主的荣耀，日与日发出言语，夜与夜传出知识；}\BibleRef{咏 19:1}然而，许多人却因此绊倒，并且彼此方向相反。有些人过于崇拜它，超过其价值，以致认为它是天主；另一些人则对它的美如此麻木，以致断言它不值得天主的创造之手，并将其中较大部分归因于某种邪恶的物质。然而天主为这两点作了预备：使它美丽伟大，以免被认为与他的智慧无关；又使它有缺陷，不能自足，以免被怀疑就是天主。尽管如此，那些被自己的推理弄瞎了眼的人，陷入矛盾的概念，彼此反驳，互相指控，甚至藉着使他们偏离的那些推理，证明天主的智慧。我为什么提及太阳和天空呢？犹太人亲眼看见这么多奇迹发生在眼前，却立刻崇拜牛犊。他们又看见基督驱逐魔鬼，却称祂附有魔鬼。但这并非对驱逐者的指控，而是对他们如此盲目之理智的谴责。那么，不要因为那些误用保禄之人的判断而责备保禄；而要好好理解他里面的宝藏，发掘他的富源，这样你就能因他的武装而坚固，高贵地对抗一切。这样你就能堵住希腊人和犹太人的口。\Quote{怎么堵住呢？有人会说，既然他们不信他？}藉着藉他所行的事，藉着世界上发生的革新。因为如此大事不是人力所能行的，而是被钉者的力量吹嘘在他身上，使他成为他所是，并显出他比演说家、哲学家、暴君、君王和所有人更强大。他不但能武装自己，击倒仇敌，还能使别人也像他一样。因此，为使我们能对自己和他人都有益，让我们不断将他握在手中，把他的著作当作喜乐的草地和花园。因为这样我们才能脱离恶习，选择德行，获得应许的美善。愿我们众人因我们的主耶稣基督的恩宠和爱人之心，都能达到，祂偕同圣父及圣神，享受光荣、权能、尊威，自今而始，以至永远，万世无疆。阿们。\par

\SourceRecord{Translated by Talbot W. Chambers.；Nicene and Post-Nicene Fathers, First Series；Vol. 12.；Edited by Philip Schaff.；Buffalo, NY: Christian Literature Publishing Co.,；1889.；<http://www.newadvent.org/fathers/220221.htm>.}

% structure: p0001 source=h1 target=section reason=page-title

\section{\WorkTitle{格林多后书讲道集 第二十二篇}}

% structure: p0002 source=h2 target=subsection reason=relative-heading-rank; base=h2

\subsection{格后 10:7}

\BibleQuote{你们看表面的事。若有人自信是属于基督的，他自己就该再想想：他怎样属于基督，我们也怎样。}\par

在保禄身上，除其他事外，尤其令人钦佩的是：当他遇到迫切需要夸赞自己的时候，他既能完成这一点，又不会因这种自我标榜而冒犯多数人；这一点我们尤其可以在他写给迦拉达人的书信中看到。因为他在那里遇到了这样的论证，便兼顾这两点；这是一件极其困难且需要极大明智的事；他既谦逊，又说了些关于自己的大事。请看他在此处也如何使这一点变得重要：\Quote{你们看表面的事。}看，这里也有明智。因为他责备了那些欺骗他们的人之后，并没有将话只限于他们，而是从他们跳到了这些人身上；他经常这样做。事实上，他不仅鞭策那些引人误入歧途的人，也鞭策那些被误导的人。因为如果他甚至让他们不受责备地过去，他们不会那么容易因对别人所说的话而悔改，反而会大为得意，以为自己不受指控。因此他也鞭策他们。不仅如此，更令人钦佩的是，他以适合各人的方式责备双方。请听他对这些人说的话：\BibleQuote{你们看表面的事。}这指责并不轻；而是表明人极其容易被欺骗。他所说的意思是：\Quote{你们以表面、以属肉体、以属身体的事来判断。} 所谓\Quote{表面}是什么意思？如果一人有钱，如果一人傲慢，如果一人被许多谄媚者围绕，如果一人夸耀自己，如果一人虚荣，如果一人没有德行却假装有德行，因为这正是\BibleQuote{你们看表面的事}的含义。\par

\BibleQuote{若有人自信是属于基督的，他自己就该再想想：他怎样属于基督，我们也怎样。}因为他不想一开始就严厉，而是逐渐加强并达到高潮。但请看这里有多少严厉和隐含的意思。他通过使用\BibleQuote{他自己}这个词来表明这一点。因为他说：\Quote{让他不要等着从我们这里得知这事，即通过我们对他的责备，而让他自己想想：他怎样属于基督，我们也怎样；} 并不是说他像别人那样属于基督，而是\BibleQuote{他怎样属于基督，我也怎样属于基督}。至此，共同点成立：因为的确，他属于基督，而我不是属于别人的。然后，在建立了这种平等之后，他又进一步超越，说：\par

% structure: p0006 source=h2 target=subsection reason=relative-heading-rank; base=h2

\subsection{格后 10:8}

\BibleQuote{因为即使我对主赐给我们的权柄——这权柄是为建设你们，而不是为破坏你们——稍稍夸耀，我也不会羞愧。}\par

因为他将要说出一些伟大的事，请看如何缓和语气。因为没有什么比自我赞美更让大多数听众反感。为了从根本上消除这种反感，他说：\BibleQuote{因为即使我稍稍夸耀。}他没有说：\Quote{若有人自信属于基督，就让他想他远不及我们，因为我从祂那里得了极大的权柄，可以随意惩罚和处死任何人。} 而是说什么？\BibleQuote{因为即使我甚至稍稍夸耀。}然而，他拥有无以言表的能力，但他用言语降低了它。他没有说\Quote{我夸耀}，而是\BibleQuote{即使我夸耀}，即使我选择这样做：既表现出谦逊，又宣告了他的优越。因此他说：\BibleQuote{即使我夸耀主赐给我的权柄。}他再次将一切归于祂，并使恩赐成为共同的。\BibleQuote{为建设，而不是为破坏。}你看见他如何再次平息因赞美可能引起的嫉妒，并通过提及他接受权柄的用途来争取听众吗？那么他为什么说：\BibleQuote{破坏人的计谋？}因为这本身就是建设的一种特殊形式：移除障碍，识别不健全的部分，并将真正的部分砌在一起。因此我们领受权柄正是为此，为了建设。但如果有人与我们争斗，并且无可救药，我们也会使用另一种能力，摧毁并推翻他。因此他也说：\BibleQuote{我不会羞愧，}即我不会被证明是说谎者或夸口者。\par

% structure: p0009 source=h2 target=subsection reason=relative-heading-rank; base=h2

\subsection{格后 10:9-11}

2. \BibleQuote{免得我好像恐吓你们：因为有人说，他的书信是沉重而有力的，但他的身体来临是软弱的，他的言语是可鄙的。这样的人该想想：我们不在时，在书信上所说的，我们来到时，在行事上也是这样的。}\par

他说的话是这样：\Quote{我固然可以夸耀，但为使人们不再说同样的话，即我在信中夸耀，而在场时却可鄙，所以我不说什么大事。}然而后来他确实说了大事，但并非关于这使他可畏的能力，而是关于启示以及更长的关于试炼的事。\Quote{所以，为免我似乎恐吓你们，\InnerQuote{让这样的人算算这事：我们不在时信中所呈现的，我们在场时行事也是如此。}}因为既然他们说：\Quote{他写信时把自己说得很伟大，但他在场时却毫不足道}，所以他这样说，而且再次以温和的方式。因为他并没有说：\Quote{正如我们写大事，我们在场时也做大事}，而是用更平缓的措辞。因为当他向其余的人说话时，他确实以激烈的方式陈述，说：\Quote{我恳求你们，使我当在场时，不必用那我认为要勇敢反对某些人的自信来表现勇气}；但向这些人说话时，他更为克制。因此他说：\Quote{我们在场时如何，不在场时也如何，那就是谦卑、温和，毫不夸耀。这从下文可以明显看出。}\par

% structure: p0012 source=h2 target=subsection reason=relative-heading-rank; base=h2

\subsection{格后 10:12}

\BibleQuote{因为我们不敢将自己与那些自夸的人同列或相比。}\par

这里他显示出那些假宗徒是自夸者，说自己的大事；并嘲笑他们自吹自擂。\Quote{但我们不做这样的事：即使我们做了大事，也全归于天主，并彼此比较。}因此他还说：\par

\BibleQuote{但他们自己以自己衡量自己，拿自己与自己比较，是无知的。}现在他所说的是：\Quote{我们不与他们相比，而是彼此相比。}因为往下文他说：\BibleQuote{我一点不在那些最大的宗徒以下。}\BibleRef{格后 12:11}并在前书说：\BibleQuote{我比他们众人更劳苦。}\BibleRef{格前 15:10}又说：\BibleQuote{借着一切忍耐，在你们中间确实行了宗徒的标记。}\BibleRef{格后 12:12}\Quote{所以我们与自己相比，不与那些一无所有的人相比：因为这样的狂妄出于愚昧。}然后他要么就自己说这话，要么就指着他们说：\Quote{我们不敢将自己与那些彼此相争、自夸大事而又不领悟的人相比；}也就是说，他们不察觉自己如此狂妄，又在彼此之间抬高自己，是何等可笑。\par

% structure: p0016 source=h2 target=subsection reason=relative-heading-rank; base=h2

\subsection{格后 10:13}

\BibleQuote{但我们不会过分地夸耀，}像他们那样。\par

因为很可能在他们的自夸中，他们说：\Quote{我们使世界归化，我们传到地极，}并吐露许多其他类似的大话。\Quote{但我们并非如此，}他说。\par

\BibleQuote{但按照天主给我们划分的范围的尺度，达到你们那里。}由此他的谦卑在两方面都明显可见：一方面他夸耀的没有超过他所成就的，另一方面他甚至把这本身也归于天主。因为他说：\BibleQuote{按照范围的尺度，}\BibleQuote{是天主给我们划分的，一个达到你们那里的尺度。}好比将葡萄园分给农夫一样，祂就这样分给了我们。因此，我们被算配得达到多远，就夸耀多远。\par

% structure: p0020 source=h2 target=subsection reason=relative-heading-rank; base=h2

\subsection{格后 10:14}

\BibleQuote{因为我们并没有越过自己的范围，好像达不到你们那里：因为我们甚至远至你们那里，借着传基督的福音，已经来到了。}\par

不单是\Quote{我们来了，}而且是\Quote{我们宣告，我们宣讲，我们劝服，我们成功了。}因为很可能他们只是来到宗徒的门徒那里，就因着他们只是在他们中间，便将一切归功于自己。\Quote{但我们并非如此：也没有人可以说我们未能来到你们那里，或者只在言语上把夸耀扬及你们；因为我们也将圣言传给了你们。}\par

% structure: p0023 source=h2 target=subsection reason=relative-heading-rank; base=h2

\subsection{格后 10:15-16}

3. \BibleQuote{不夸耀超过}我们的\BibleQuote{范围，}就是说，\BibleQuote{在别人的劳苦中，但怀着希望，当你们的信德增长时，我们将在你们中间按我们的范围被扩大，直至丰盛有余，以至将福音传到你们以外的地方，而不在别人范围内夸耀那已做成的事。}\par

他以此为依据对他们提出严厉的控诉，既因为他们夸耀范围之外的事和他人的劳苦，也因为在全部辛劳都是宗徒的时候，他人却以自己的劳苦自夸。\Quote{但我们，}他说，\Quote{在行为中显明了这些事。因此我们不会效仿那些人，而只会在行为为我们作证的地方说这样的话。为什么，}他说，\Quote{我说\Quote{你们}？}\BibleQuote{因为我希望，当你们的信德增长时，}因为他并没有绝对断言，保持自己的品格，而是说\Quote{我希望，}他说，\Quote{如果你们进步，我们的范围将延伸得更远，}\BibleQuote{在那些边界之外传福音。}\Quote{因为我们还会更进一步，}他说，\Quote{去宣讲和劳作，而不是夸耀别人劳苦的事。}他恰当地称之为\BibleQuote{范围和尺度，}好像他得了世界为产业，一份丰厚的遗产；并显示一切完全属于天主。\Quote{既然有这些工作，}他说，\Quote{并且期望更大的，我们不学那些一无所有的人夸耀，也不将任何部分归于自己，而将全部归于天主。因此他又说：}\par

% structure: p0026 source=h2 target=subsection reason=relative-heading-rank; base=h2

\subsection{格后 10:17}

\BibleQuote{夸耀的，当在主内夸耀。}\par

% structure: p0028 source=h2 target=subsection reason=relative-heading-rank; base=h2

\subsection{格后 10:18}

\BibleQuote{因为不是自夸的被认可，而是主所称赞的。}\par

他没有说我们是如此，而是\Quote{主所称赞的}。你看他说话多么谦逊？但是，如果他在继续时说出更高昂的话语，不要惊奇，因为这也是出于\Person{保禄}的明智。因为如果他在每一部分都只讲谦卑的话语，他就不会如此有效地打击这些人，也不会使门徒从他们的错误中解脱。因为不合时宜的谦逊可能造成伤害，而在适当的时候说一些关于自己的可称赞之事也能有益。因此他也这样做了。因为门徒被诱导对\Person{保禄}产生卑下的看法，这并非小危险。并非\Person{保禄}寻求从人而来的光荣。因为如果他寻求这个，他就不会那么久对那些\Quote{十四年前}的大事\BibleRef{格后 12:20}保持沉默；也不会在不得已时如此退缩和犹豫去谈论它们；显然，若非被迫，他甚至不会说出来。所以，他这样说，并非出于对人的光荣的渴望，而是出于对门徒的温柔关怀。因为他们指责他是夸夸其谈者，言语自夸却无法在行动中证明，于是他被迫随后提到那些启示。虽然他在说这些话时本可以通过行动说服他们，但他仍然坚持用言语警告。因为他特别远离虚荣；他的一生前后都证明了这一点。例如，正因如此他忽然改变；改变后，他使犹太人混乱，抛弃了从他们那里得来的一切光荣，尽管他自己曾是他们的首领和战士。但当他找到真理时，他不再考虑那些；反而接受了他们的侮辱和辱骂；因为他着眼于许多人的救恩，认为这就是一切。因为那个为了对基督的渴望而不顾地狱、天堂甚至万千世界的人，怎么会追逐从许多人来的光荣呢？绝不会；他甚至在可以谦卑的时候非常谦卑，并耻辱地标记自己以前的生活，称自己为\Quote{亵渎者、迫害者和凌辱者}\BibleRef{弟前 1:13}。而他的门徒\Person{路加}也说了许多关于他的事，显然是从他本人学来的，他自己充分展示了他以前的生活，不亚于他归化后的生活。\par

4. 我说这些话，不但是为叫我们听，也是为叫我们学习。因为如果他在圣洗之前还记得那些过犯，虽然它们都被涂抹了，那么我们在圣洗之后自己却忘记了那些过犯，还能得到什么宽恕呢？人哪，你说什么？你得罪了天主，而你却忘记了？这是第二次的过犯，第二次的仇视。那么你求宽恕的是什么罪呢？是你自己也不知道的那些吗？当然（难道不是吗？）你深深焦虑并思虑如何交代它们，你甚至不肯关心去记住它们，反而戏弄那不可戏弄的事。但总有一个时候，我们的戏弄不能再继续。因为我们必须要死（因为许多人的麻木不仁迫使我甚至说那些明显的事），必须复活，受审判，受惩罚；不，如果我们选择，这倒不是必须的。因为那些事不由我们支配：我们的结局、复活、审判，都不由我们，而是由我们的主支配；但我们是否受惩罚却由我们自己支配；因为这是可能发生也可能不发生的事。但如果我们选择，我们就使它成为不可能的事；正如保禄、伯多禄、以及所有圣人所作的那样；因为他们受惩罚甚至是不可能的。因此，如果我们有心，同样我们也不可能遭受任何事。因为即使我们得罪了千万件事，只要我们还在世，就有可能挽回自己。那么，让我们挽回自己吧：让老年人想想，不久之后他就要离世，因为他有生之年已经享乐够了（虽然这种享乐是什么，就是生活在邪恶中吗？但此时我按他的思维方式这样说）；此外，让他想想，他在短时间内就能洗去一切。青年人，也让他想想死亡的不确定性，常常许多年长者还留在此世时，青年人却先被带走。因为正是为此，使我们不致于为死亡做交易，死亡被留在不确定中。因此，一位智者劝告说：\BibleQuote{你不要迟延归向上主，不要一天又一天地拖延；因为你不知道明天会发生什么。}\BibleRef{德 5:7} 因为拖延带来危险和恐惧；而不拖延则带来确定和安全的救恩。所以，你要紧紧抓住德行。因为这样，即使你年轻离世，你也安全地离世；如果你到了老年，你将带着丰富的储备到达死亡，并且你一生将有双重的盛宴：既在于你戒除恶习，又在于你抓住德行。不要说：\Quote{会有时间可以回转的}，因为这种话极其激怒天主。为什么呢？因为祂应许了你无数个世代，而你甚至不愿意在此生劳苦，这短暂而只持续一段时间的生命；却如此懒惰和懦弱，甚至寻求比这更短的时间。难道每天不是同样的宴乐？难道不是同样的筵席、同样的娼妓、同样的戏院、同样的财富吗？你还要爱这些东西多久，好像它们是什么似的？你的恶欲还要多久才满足？想想看，你犯了多少次奸淫，就定了自己多少次罪。因为罪的性质就是这样：一旦犯了，审判者也已经宣判了。你醉酒、贪食、或抢劫了吗？现在就停住，立刻回转，向天主承认祂的仁慈，就是在你犯罪时没有把你夺走；不要再寻求另一个作恶的时间。许多人在贪婪中被夺走，离世而受明显的惩罚。要怕你也遭受这事，而毫无借口。\Quote{但天主给许多人直到极度老年悔改的时间。}那又怎样？祂也会给你吗？有人说：\Quote{或许祂会。}你为什么说\Quote{或许}、\Quote{有时}、\Quote{常常}？想想你是在为你的灵魂考虑，也要考虑相反的情况，并计算一下，说：\Quote{但如果祂不给呢？}他却说：\Quote{但如果祂给呢？}天主确实给了；但这个假设比那个更安全、更有益。因为如果你现在开始，无论你是否得到延长时间，你已赢得一切；但如果你总是拖延，正因如此，也许你得不到延长时间。当你上战场时，你不会说：\Quote{我不需要立遗嘱，也许我会平安回来。}当你考虑婚姻时，你不会说：\Quote{假设我娶了一个穷妻子，许多人这样出乎意料地致富了。}当你建房时，你不会说：\Quote{假设我打了腐烂的地基，许多房子也这样立住了。}然而在考虑灵魂时，你却依靠更腐烂的东西：坚持你的\Quote{或许}、\Quote{常常}、\Quote{有时}，并把自己托付给这些不确定的事。有人说：\Quote{不，不是不确定，而是依靠天主的仁慈，因为天主是仁慈的。}我也知道；但即便如此，这位仁慈的天主还是夺走了我所提到的人。如果你得到了时间，你仍然像原来一样呢？因为这种人即使在老年也会懒散。他说：\Quote{不，不是这样。}因为这种推理方式在八十年后渴望九十年，九十年后渴望一百年，一百年后更加不愿行动。这样，整个生命就白白消耗了，从前对犹太人所说的话也会临到你：\BibleQuote{他们的日子消耗于虚无。}\BibleRef{咏 77:33} 但愿只是虚无，而不是邪恶。因为当我们带着罪的沉重负担离开那里时，这也会成为邪恶。因为我们必带走火的燃料和虫子的盛宴。因此，我祈祷并恳求你们终于高尚地停步，停止作恶，使我们也能获得所应许的美善。愿我们众人得以达到，赖我们的主耶稣基督的恩宠和爱人，祂和圣父及圣神，永受光荣、权能、尊崇，自今而始，及世世无穷。阿们。\par

\SourceRecord{Translated by Talbot W. Chambers.；Nicene and Post-Nicene Fathers, First Series；Vol. 12.；Edited by Philip Schaff.；Buffalo, NY: Christian Literature Publishing Co.,；1889.；<http://www.newadvent.org/fathers/220222.htm>.}

% structure: p0001 source=h1 target=section reason=page-title

\section{第23篇讲道：论\BibleBook{格林多}{后书}}

% structure: p0002 source=h2 target=subsection reason=relative-heading-rank; base=h2

\subsection{\BibleBook{格林多}{后书} 11:1-2}

\BibleQuote{但愿你们能容忍我这一点愚妄；其实你们在容忍我。}\par

正要开始谈论自己的长处时，他事先作了许多纠正。他这样做不只一次两次，尽管话题的必要性以及他常说的话足以成为他的理由。因为他记念天主所不记念的罪，并因此声称自己不配宗徒之名，就连最无知的人也清楚看出，他接下来要说的话并非为了荣耀。因为若要说出令人震惊的话，那甚至是尤其损害他的荣耀的——谈论自己的一些事；并且对大多数人来说，这是令人反感的。然而，他毫不畏惧地看待这些事中的任何一件，他只关注一件事：听众的救恩。但为了不因此（我指的是说自己大事）给无知的人造成伤害，他出于极大的谨慎，用了许多预备性的纠正，并说：\Quote{但愿你们能容忍我，} 当我做一点愚妄之事时；是的，反倒是\Quote{你们确实容忍了我。} 你看到智慧了吗？因为当他说\Quote{但愿}时，像是把决定权交给他们；而当他甚至断言（他们容忍）时，则是极大地信任他们的情感，并表明他既爱他们也被爱。是的，不仅出于单纯的爱，而且出于一种热烈而近乎疯狂的情感，他说他们应当容忍他，即使他做愚昧事。因此他接着说：\Quote{因为我以天主的妒爱妒爱你们。} 他没有说\Quote{我爱你们}，而是用了比这强烈得多的词。因为那些灵魂是妒爱的，它们为所爱的人炽烈燃烧；而妒爱只能从强烈的感情中产生。然后，为了不让他们以为他是为了权力、荣誉、财富或任何其他类似的东西而渴望他们的感情，他加了\Quote{以天主的妒爱。} 因为天主也被说成是妒爱的，这不是要人认为祂有激情（因为天主性是无情的），而是让所有人都知道，祂做一切事并非出于别的关切，而是为了祂所妒爱之人的缘故；不是为使祂自己有所得，而是为拯救他们。在人中间，妒爱确实不是出于这个原因，而是为了自己的安宁；不是因为所爱者遭受了侮辱，而是为了不让爱他们的人受伤、在恩宠中被超越、在所爱者心中的情感中降低。但在这里并非如此。他说：\Quote{我不在乎这个，生怕我在你们眼中地位降低，而是生怕看见你们被败坏。因为天主的妒爱正是如此；我的妒爱也是如此，既强烈又纯洁。} 然后还有这个必要的理由：\par

\BibleQuote{因为我曾把你们许配给一个丈夫，如同把贞洁的童女献给基督。}\newline{} 所以，我妒爱，不是为我自己的缘故，而是为那一位，我曾把你们许配给祂。因为现在是订婚的时候，而婚礼的时候是另一个时候；那时他们歌唱：\Quote{新郎站了起来。} 哦，何等闻所未闻的事！在世界上，童女在婚前是处女，婚后就不再是了。但在这里却不是这样：即使他们在这次婚礼前不是处女，婚后他们却成为处女。所以整个教会是一个童女。因为他甚至对所有人——丈夫和妻子——说话，这样称呼他们。让我们看，他带来了什么，并用什么聘礼把我们许配了。不是金子，不是银子，而是天国。因此他也说：\BibleQuote{我们是基督的使者}，当他将要迎娶新娘时，他恳求他们。发生在亚巴郎身上的事是这一点的预像。\BibleRef{创 24:4} 因为他派他忠信的仆人去为儿子寻找外邦女子为妻；同样，天主派祂的仆人为自己的儿子寻找教会为妻，并派古代的先知说：\BibleQuote{女儿，请听，请看，要忘记你的民族和你的父家，君王就会恋慕你的美丽。}\BibleRef{咏 45:10} 你看到先知也在提亲吗？你看到宗徒自己也大胆地说同样的事，并说：\BibleQuote{我曾把你们许配给一个丈夫，要把你们如同贞洁的童女献给基督吗？} 你再次看到智慧了吗？因为说了\Quote{你们应当容忍我}之后，他没有说\Quote{因为我是你们的老师，我不是为我自己说话：} 而是用了这个使他们格外尊贵的说法，把自己放在促成婚姻的位置上，把他们放在新娘的地位上；他接着说：\par

% structure: p0006 source=h2 target=subsection reason=relative-heading-rank; base=h2

\subsection{\BibleBook{格林多}{后书} 11:3}

\BibleQuote{但我怕你们的心意，或许怎样被毒蛇以它的狡猾所迷惑，以致败坏了对基督所怀有的纯朴。}\par

因为虽然毁坏是你们（自己）的，但悲伤也是我的。请考虑他的智慧。因为他并没有断言他们已经被败坏；这样他在说\BibleQuote{当你们的服从完成时}\BibleRef{格后 10:6}和\BibleQuote{我要哀哭许多从前犯罪的人}\BibleRef{格后 12:21}时表明了这点；但他仍不让他们变得无耻。因此他说\BibleQuote{恐怕有时。} 因为这不谴责也不沉默；因为无论直说还是永远隐藏都不安全。所以他采用这种折中的方式，说\BibleQuote{恐怕有时。} 因为这话既不是完全不信赖，也不是完全信赖他们，而是介于两者之间。这样他缓和了语气，但通过提及那段历史，他使他们陷入无法形容的恐惧，并断绝了他们所有得赦免的希望。因为即使毒蛇是恶毒的，她也是愚昧的，但这些事中没有一件使那女人免于惩罚。他说：\Quote{所以你们要小心，免得你们遭遇同样的命运，而没有任何东西保护你们。因为他也是应许更大的事，如此欺骗了人。} 由此可知，这些人也是通过自夸和自吹自擂而受骗。这不仅可以从这里推测，也可以从他后面的话看出：\par

% structure: p0009 source=h2 target=subsection reason=relative-heading-rank; base=h2

\subsection{\BibleBook{格林多}{后书} 11:4}

\Quote{如果那来的人宣讲另一位耶稣，就是我们不曾宣讲过的，或者你们接受另一位圣神，就是你们不曾接受过的，或者另一个福音，就是你们不曾接纳过的，你们容忍他，倒也不错。}\par

他没有说：\Quote{免得恐怕如同\Person{亚当}被欺骗了一样。}而是指出那些人是妇女，他们这样被虐待，因为受欺骗是妇女的本分。他没有说：\Quote{你们也应当这样被欺骗。}而是保持比喻说：\Quote{你们的心意应当败坏，离开那向着基督的纯一。}\Quote{离开纯一，我说，不是出于邪恶；既不是出于邪恶，也不是出于你们不信，而是出于纯一。}但是，即便如此，受欺骗的人也不配得到宽恕，\Person{厄娃}显示了这一点。但如果这不配得到宽恕，那么当有人出于虚荣而如此行时，就更不配得到宽恕了。\par

2. \Quote{因为如果那来的人宣讲另一位耶稣，就是我们不曾宣讲过的：}由此表明，你们的欺骗者不是格林多人，而是从别处先前被败坏的人；因此他说\Quote{那来的人}。\par

\Quote{如果你们接受另一位圣神，如果你们接受另一个福音，就是你们不曾接纳过的，你们容忍他，倒也不错。}你说什么？你曾对迦拉达人说过：\Quote{若有人给你们宣讲另一个福音，与你们所接受的不同，他就该受\OriginalTerm{绝罚}{\GreekText{ἀνάθεμα}}。}现在你却说：\Quote{你们容忍他，倒也不错}？然而因这缘故，理应不容忍他们，而应回避他们；但如果他们说的是同样的事，就应容忍他们。那么你怎么说：\Quote{因为他们说的是同样的事，就不该容忍他们？}因为他说：\Quote{如果他们说了别的事，就应当容忍他们。}因此我们要留意，因为危险巨大，悬崖深邃，如果人粗心大意地跑过这里；这里所说的给所有异端打开了入口。那么这些话是什么意思呢？那些人如此夸耀，仿佛宗徒们教导得不完全，而他们引进了一些更多的东西。因为很可能，他们用许多空谈，引入无意义的垃圾来掩盖这些教义。因此他提到了蛇和\Person{厄娃}，她被期望获得更多的念头所欺骗。在前一封书信中也暗示这一点，他说：\Quote{你们已经富足了，已经不用我们而作王了；}又说：\Quote{我们为了基督成了愚妄的，你们在基督内却是聪明的。}\BibleRef{格前 4:8}既然很可能他们用外表的智慧说了许多空话，他所说的是这样：即\Quote{如果这些人说了更多的话，而且宣讲了理应宣讲的另一位基督，但我们却省略了它，\InnerQuote{你们容忍他，倒也不错}。}因为为此他加上了：\Quote{就是我们不曾宣讲过的。}\Quote{但如果信仰的主要点是一样的，你们比他们多得了什么呢？无论他们说什么，他们不会说超出我们所说的话。}且看他多么精确地陈述情况。因为他没有说：\Quote{如果那来的人说了更多的事；}因为他们确实说了更多的事，用更多的权威和优美的言辞高谈阔论；因此他没有说这话，而是说了什么？【如果】\Quote{那来的人宣讲另一位耶稣，}这是不需要那一套辞藻的事：\Quote{或者你们接受另一位圣神，}（因为这事也不需要言语；）也就是说，\Quote{使你们在恩宠上更丰富；}或者\Quote{另一个福音，就是你们不曾接纳过的，}（这也不需要言语，）\Quote{你们容忍他，倒也不错。}但是，我请求你们想想，他如何到处使用这样的定义，表明他们没有引进任何非常伟大的事，甚至没有更多的事。因为当他说道：\Quote{如果那来的人宣讲另一位耶稣，}他加上：\Quote{就是我们不曾宣讲过的；}以及\Quote{你们接受另一位圣神，}他接着说：\Quote{就是你们不曾接受过的；或者另一个福音，}他加上：\Quote{就是你们不曾接纳过的，}所有这些表明，应当留意他们，不是简单地说如果他们说了更多，而是如果他们说了任何理应被说而由我们省略的事。但如果那事不该被说，因此我们没说；或者他们只说和我们相同的事，为什么你们这样惊奇地张开嘴望着他们？\Quote{然而如果他们说着同样的事，}有人说，\Quote{你为什么阻止他们？}因为他们用伪善，引入陌生教义。但这暂时他没有说，而是在后面断言了，当他说：\Quote{他们伪装成基督的宗徒；}\BibleRef{格后 11:13}现在他通过较不冒犯的考虑，使门徒们从他们的权柄下退开；这不是出于嫉妒他们，而是为了保全这些门徒。否则，他为什么不阻止\Person{阿颇罗}，他可是\Quote{有学问的人，且精通圣经；}\BibleRef{宗 18:24}反而还恳请他，并许诺要送他来？因为随着他的学问，他也保存了教义的完整；但这些人相反。因此他同他们作战，并责备门徒们惊奇地张着嘴望着他们，说：\Quote{如果我们省略了任何应当说的事，而他们补充了，我们不阻止你们注意他们；但如果所有事都已由我们完全完成，没有留下任何不足，他们怎么抓住了你们？}因此他又说：\par

% structure: p0014 source=h2 target=subsection reason=relative-heading-rank; base=h2

\subsection{格林多后书 11:5}

\BibleQuote{因为我自以为一点也不在那些最大的宗徒以下}，他不再把自己与他们相比，而是与伯多禄和其他人相比。\Quote{所以，如果他们比我懂得多，那也超越他们。}请注意他在这里也表现出谦逊。因为他没有说：\Quote{宗徒们没有比我多说什麼}，而是说什么？\BibleQuote{我以为}，我这样认为，\BibleQuote{一点也不在那些最大的宗徒以下}。\par

% structure: p0016 source=h2 target=subsection reason=relative-heading-rank; base=h2

\subsection{\BibleRef{格后 11:6}}

3. \BibleQuote{但我虽然在言语上粗鲁，但在知识上却不然。}\par

因为那些败坏格林多人的在这点占优势，即他们不粗鲁；他也提到这点，表明他不但不以此为耻，反而以此为傲。他没有说：\BibleQuote{但我虽然在言语上粗鲁}，但他们也是如此，因为这会显得既指控他们也指控他自己，而抬高这些人；但他推翻这事本身，即外在的智慧。的确，在他前书里，他甚至激烈地争论这事，说它不仅对宣讲毫无贡献，反而遮蔽十字架的光荣；\BibleRef{格前 2:1} 因为他说：\BibleQuote{我来到你们那里，并没有用高超的言辞或智慧，以免基督的十字架落空；} \BibleRef{格前 1:17} 以及其他许多同样的话；因为\BibleQuote{在知识上}他们\BibleQuote{粗鲁}，这也是粗鲁的极端形式。因此，当需要在大事上作比较时，他把自己与宗徒相比；但当要显示那看似缺点的事情时，他不再这样做，而是直接处理这事本身，并表明它是一种优越性。的确，当没有必要时，他说他是\BibleQuote{宗徒中最小的}，甚至不配这个称号；但在这里，当情况需要时，他说他是\BibleQuote{一点也不在那些最大的宗徒以下}。因为他知道这最有益于门徒。因此他又加上：\par

\BibleQuote{不，我们在一切事上已在众人中向你们显明了。}因为在这里他再次指责假宗徒\BibleQuote{行事诡诈}。\BibleRef{格后 4:2}他先前也这样论自己，说他不是按外表生活，也不宣讲天主圣言而加以歪曲腐蚀。\BibleRef{格后 4:2}但那些人是表里不一。他却不是这样。因此他处处采取高调，仿佛不因人的意见而作任何事，也不隐瞒关于自己的任何事。正如他先前所说：\BibleQuote{借着真理的彰显，在众人良心前推荐自己}，\BibleRef{格后 4:2}所以现在他又说：\BibleQuote{在一切事上我们已向你们显明了。}但这意思是什麼？\Quote{我们粗鲁，}他说，\Quote{并且不隐瞒：我们接受一些人的东西，也不保密。我们从你们那里接受，我们不假装不接受，正如他们接受时那样，但我们使所行的一切向你们显明；}这是一种既对你们极其信任，又把一切如实告诉你们的行为。因此他也称你们为证人，现在说：\BibleQuote{在众人中向你们}，又先前：\BibleQuote{因为我们写给你们的，不外是你们所念或所承认的。} \BibleRef{格后 1:13}\par

4. 然后，在为自己行为辩护之后，他接着严厉地说：\par

% structure: p0021 source=h2 target=subsection reason=relative-heading-rank; base=h2

\subsection{\BibleRef{格后 11:7}}

\BibleQuote{我贬抑自己，为使你们受高举，难道我犯了罪吗？}\par

为了解释这点，他加上：\par

% structure: p0024 source=h2 target=subsection reason=relative-heading-rank; base=h2

\subsection{\BibleRef{格后 11:8}}

\BibleQuote{我掠夺了别的教会，取了他们的工资，为能服事你们。}\par

他所说的是：\Quote{我生活在窘迫中}；因为这是\Quote{贬抑自己}的力量。\Quote{那么你们能以此为罪责我吗？你们因此就高抬自己反对我吗，因为我借着乞讨、忍受窘迫、受苦、挨饿来贬抑自己，为使你们受高举？}他处于窘迫中，他们如何受高举呢？他们更加受到建立，没有跌倒；这甚至可能成为对他们的极大控告，指责他们的软弱；除非先贬抑自己，就没有其他方法引导他们。\Quote{那么你们是否以此为罪责我贬抑自己？但你们却因此受高举。}因为他在上面甚至说过他们指责他，因为他当面谦卑，背后大胆；在为自己辩护时，他在这里再次打击他们，说：\Quote{这也是为了你们。}\par

\Quote{我抢夺了别的教会。}这里他终于以责备的口吻说话，但他先前的话使这些话不显得刺耳；因为他曾说：\Quote{请你们容忍我一点愚妄吧。}并且在讲述其他成就之前，他首先以此事夸口。因为世俗之人尤其注重这一点，而他的对手们也对此极为自夸。因此，他并非首先谈及自己的危险，也不是自己的奇迹，而是谈及自己不贪图钱财的事，因为他们以此自夸；同时他也暗示他们富有。但值得钦佩的是，他本可以说自己甚至亲手做工维持生计，却没有这样说；而是说出那特别使他们羞愧、却又并非称赞自己的话，即\InnerQuote{我从别人那里拿了来。}他没有说\Quote{拿了}，而是说\Quote{抢夺了}，意即\InnerQuote{我剥夺了他们，使他们贫穷。}而更甚的是，这不是为奢侈，而是为必需；因为当他说\InnerQuote{工资}时，他指的是必要的生计。更为严重的是，\Quote{来服事你们。}我们向你们宣讲；而我本应由你们供养，却从别人那里得到了。这指控是双重的，或者更确切地说是三重的：即当他在他们中间生活并服事他们，寻求必要支持时，却有别人供应他的需要。一方的疏忽与另一方的热忱，差距何等之大！因为这些人即使远在他方也送给他，而那些人连他在他们中间时也不供养他。\par

5. 然后，因他严厉地责备了他们，便又平静地缓和了斥责的激烈语气，说道：\par

% structure: p0029 source=h2 target=subsection reason=relative-heading-rank; base=h2

\subsection{\BibleRef{格后 11:9}}

\Quote{当我在你们那里时，虽然有所缺乏，却没有拖累任何人。}\par

因为他没有说\InnerQuote{你们没有给我}，而是说\InnerQuote{我没有拿}；此时他仍宽容他们。然而即便在他的话中带着克制，他也巧妙地再次击打他们；因为\Quote{在你们那里}一词极其有力，\Quote{有所缺乏}也是如此。免得他们说：\InnerQuote{那么，如果你有〔足够的〕，又算什么呢？}他补充说：\Quote{并且有所缺乏。}\par

\Quote{我没有拖累}你们。这里他再次轻轻地敲打他们，暗示他们不情愿地做这样的供给，视之为负担。接着是理由，也充满指责并引起嫉妒。因此他引入这理由，并非作为主要论点，而是为了告诉他们他从何处、由何人得到供养，以便以一种不引人怀疑的方式再次激励他们关于施舍的事；\par

\Quote{因为我的缺乏，}他说，\Quote{从马其顿来的弟兄们供应了。}你看见他如何再次激怒他们，提出那些曾服事他的人吗？因为他先激发他们想知道这些人是谁，当他说\Quote{我抢夺了别的教会}时，然后他指名提及他们；这也会激励他们施舍。因为他说服那些在支持宗徒这件事上已被击败的人，不要在救济穷人上也被击败。他也在写给马其顿人自己的书信中说了这话：\Quote{因为在我的需要中，你们曾一再地送给我，甚至在福音开始的时候；}\BibleRef{斐 4:16}这一点也是对他们极大的称赞，因为他们从一开始就发光。但请注意，他处处提到自己的\Quote{需要}，而从未提到奢侈。因此，现在通过说\Quote{在你们那里}和\Quote{有所缺乏}，他表明自己本应由格林多人供养；而通过\Quote{他们供应了我的缺乏}，他表明他并没有要求。他给出的理由并不是真实的理由。这理由是什么呢？即他从别人那里领受了；\Quote{因为}他说，\Quote{那些来的人供应了我的缺乏。}\InnerQuote{为此原因}，他说，\InnerQuote{我没有拖累你们；并非因为我不信任你们。}然而他正是因为这个原因才如此行，并在后面表明这一点；但他没有明说，而是将其暗藏，留给听众的良心。他在后面隐约地证明了这一点，说：\par

\Quote{我在一切事上保守自己不成为负担，并且将来也必保守自己。}\Quote{因为不要以为，}他说，\Quote{我说这些是为了要得着什么。}如今\Quote{将来也必保守自己}这话更加严厉，因为即或到了现在，他仍不信任他们；而是一次决定不再从他们那里接受任何东西。他还表明，他们甚至把这视为一种负担；因此他说：\Quote{我保守自己不成为负担，并且将来也必保守自己。}他在前封书信中也说：\Quote{我写这些，并非要这样待我；因为我宁愿死，也不愿有人使我这夸耀落了空。}\BibleRef{格前 9:15}这里又再说：\Quote{我保守自己不成为你们的负担，并且将来也必保守自己。}\par

6. 然后，为免他显得说这些话是为了更好地赢取他们，他说：\par

% structure: p0036 source=h2 target=subsection reason=relative-heading-rank; base=h2

\subsection{\BibleRef{格后 11:10}}

\Quote{基督的真理在我里面。}\InnerQuote{你们不要以为我因此说了话，是为了得着什么，是为了吸引你们；因为}他说，正如真理在我里面，\par

\BibleQuote{在阿哈雅各地，没有人能阻止我以此夸耀。}因为怕有人再以为他为此伤心，或说这些话是出于愤怒，他甚至称这事为\Quote{夸耀}。在他前书里，他也用同样的措辞修饰。为使他在那里也不致伤害他们，他说：\BibleQuote{那么，我的报酬是什么？} \BibleQuote{就是当我传福音时，我能使人无偿地获得基督的福音。}\BibleRef{格前 9:18}正如他在那里称它为\Quote{报酬}，他在这里也称它为\Quote{夸耀}，免得他们因他的话过分羞愧，好像他是在索取而他们却没有给他。\Quote{因为，即使你们愿意给，}他说，\Quote{我也不接受。}而\Quote{没有人能阻止我}这个说法，是取自河流的比喻，或者像流言一样，到处传播他不收任何东西的事。\Quote{你们不能以你们的给予来阻止我这言论的自由。}但他没有说\Quote{你们不能阻止}，那样会太伤人，而是说\Quote{在阿哈雅各地，没有人能阻止我。}这又是给他们致命的一击，极易使他们沮丧和痛苦，因为他们是唯一被他拒绝接受的人。\Quote{因为如果他以此为荣，他应该在各地都这样做；但如果他只在我们中间这样做，也许这是由于我们的软弱。}因此，免得他们这样推理而沮丧，请看他是如何纠正这一点的。\par

% structure: p0039 source=h2 target=subsection reason=relative-heading-rank; base=h2

\subsection{格林多后书 11:11}

\BibleQuote{为什么呢？是因为我不爱你们吗？天主知道。}\par

［解决得很快］，方法也很简单。但即便如此，他也没有完全消除那些指控。因为他既没有说\Quote{你们不软弱}，也没有说\Quote{你们是坚强的}，而是说\Quote{我爱你们}，这反而大大加重了对他们的控诉。因为不接受他们的馈赠，是由于他们觉得这极其委屈，这恰恰是特别爱他们的证明。所以，他出于爱采取了两种相反的行为：既接受，也不接受；但这种矛盾是由于给予者的心态。他没有说\Quote{我因此不接受你们的东西，是因为我极其爱你们}，因为这会包含对他们软弱的控诉，使他们陷入痛苦；而是将他的话转向另一个理由。那么这个理由是什么呢？\par

% structure: p0042 source=h2 target=subsection reason=relative-heading-rank; base=h2

\subsection{格林多后书 11:12}

为要断绝那些找机会之人的机会，使他们在所夸的事上，被显明和我们一样。\par

\BibleQuote{为要断绝那些找机会之人的机会，使他们在所夸的事上，被显明和我们一样。}因为他们极力寻求找把柄攻击他，所以也必须消除这一点。因为这是他们引以为傲的惟一点。因此，为使他们在任何方面都不能占优势，必须纠正这一点；因为他们在其他方面是逊色的。因为，正如我所说，没有比不接受他们的馈赠更能教化世俗之人的了。因此，魔鬼在狡猾中特别投下了这个诱饵，想要在其他方面损害他们。但在我看来，这甚至是出于虚伪。因此他没有说\Quote{在他们做得好的地方}，而是说\Quote{在他们所夸的事上}；这也是在嘲笑他们的夸耀；因为他们也夸耀他们本来没有的东西。但高尚的人不仅不应该夸耀他没有的东西，甚至也不应该夸耀他拥有的东西；如同这位有福的圣人常做的，如同圣祖亚巴郎所做的，说：\BibleQuote{我只是灰土。}\BibleRef{创 18:27}因为既然他没有可说的罪，却闪耀着善工；他在各处奔跑，找不到什么大的把柄反对自己，他便转向自己的本性；既然\Quote{土}这个名称在某种意义上是尊严的，他便加上了\Quote{灰}。因此，另一个人也说：\BibleQuote{灰土为什么要骄傲？}\BibleRef{德 10:9}\par

7. 请你们不要对我讲面容的红润，昂起的颈项，披风、骏马和随从；但请思考这一切的结局，并将这一点放在他们身上。如果你告诉我眼睛所见的事物，我也将告诉你图画中的事物，比这些更加光彩夺目。但我们不因它们的外表而赞美它们，因为我们知道它们的本性全是泥土；所以，让我们也不赞美这些人，因为他们也不过是泥土。不仅如此，甚至在它们被分解成尘土之前，请给我看这昂起的颈项被热病折磨、喘着最后一口气的样子；那时我将与你交谈，并问你：那所有过分的装饰都到哪里去了？那群谄媚者、跟随的奴隶、丰裕的财富和产业都消失到哪里去了？什么风来把它们都吹走了？不，甚至躺在棺材上，他还带着那财富和骄傲的记号：一件华丽的外衣盖在他身上，穷人和富人护送着他，聚集的人群说着吉祥的话。这确实也是一种极大的嘲弄；然而，这很快也显为虚无，如同凋谢的花朵。因为当我们穿过城门，将尸体交给虫蛆后返回时，我会再问你：那大群人都到哪里去了？那喧闹和嘈杂变成了什么？火把在哪里？妇女的队伍在哪里？这些岂不都是梦吗？那些呼喊声又在哪里？那些许多嘴唇喊叫，命令他\Quote{放心，因为没有人是不死的}？这些话现在不应该对一个听不见的人说，而应该在他掠夺别人、欺诈别人的时候说，那时稍加修改应是对他说：\Quote{不要放心，没有人是不死的；抑制你的疯狂，熄灭你的贪欲；}而\Quote{放心}是对受害者说的。因为现在在这个人身上唱这样的歌，不过是人们在他面前夸耀并说反话；因为他现在不该为此放心，而该恐惧战栗。\par

即使这劝告如今对他已无用，因为他已经走完他的路程，但至少让那些同样染此病症的富人们，并跟随他到坟墓的人，听到它。因为虽然先前因财富的陶醉，他们心中毫无此念，但在那时，当看见那躺卧之人证实了所言，让他们清醒，让他们受教：反省不久之后，将有人来把他们带走，去面对那可怕的账目，并为他们的强夺和敲诈行为受罚。有人会说：\Quote{这对穷人有什么益处？}其实，对许多人来说，这也是一种满足，看到那冤屈他们的人受罚。\Quote{但对我们来说，这不是满足，而是自己免于受苦。}我极其称赞并认可你们，因为你们不以别人的灾祸为乐，只求自身的安全。那么来吧，我也要确保你们这一点。因为如果我们从人手中受害，我们通过高贵地承受所发生的事，就削减了我们债务的不小部分。因此我们并未受伤害；因为天主将所受的苦算作对我们的债务，不是按公义的原则，而是按祂的慈爱；并且祂没有救助那受害的人。有人会说：\Quote{这从何见得？}犹太人曾遭受巴比伦人的迫害，天主并没有阻止；他们被掳走，连孩童和妇女；但后来这被掳却成了他们赎罪的慰藉。因此祂对\Person{依撒意亚}说：\BibleQuote{你们要安慰，安慰我的百姓，你们司铎：对耶路撒冷说知心话，因为她从主的手中已为罪恶受到了双倍的惩罚。}\BibleRef{依 40:1}又说：\BibleQuote{求你赐给我们和平，因为你已经偿还了我们一切。}\BibleRef{依 26:12}\Person{达味}说：\BibleQuote{请看我的仇敌何其众多，求你赦免我所有的罪过。}\BibleRef{咏 25:19}当他忍受\Person{史米}咒骂他时，他说：\Quote{由他吧，愿主看见我的屈辱，并在这天用善来回报我。}\EditorialNote{389 11,12}因为当我们在受冤屈时，祂若不帮助我们，那我们最得益；因为祂把这算作我们罪债的偿还，只要我们感恩地承受。\par

8. 因此，当你看见富人掠夺穷人时，撇下那受苦的人，为掠夺者哭泣。因为前者除去污秽，后者却沾染更多污秽。\Person{厄里叟}的仆人在\Person{纳阿曼}的故事中就遭遇了这样的事\BibleRef{列下 5:20}。虽然他并非用暴力夺取，但他行了不义；因为用欺骗得钱就是不义。结果如何？他连同不义还得了癞病；那受欺的人蒙受了益处，而行不义的人却遭受了极大的损害。现在灵魂方面也是如此。这力量如此之大，常常单凭它本身就能使天主息怒；是的，尽管那受害的人不配得援助，但是当他如此过度受苦时，仅凭这一点，他就吸引天主宽恕他自己，并惩罚那行不义的人。因此天主昔日对外邦人说：\BibleQuote{我固然把他们交在少数人手中，但他们却一同趋向邪恶；}\BibleRef{匝 1:15}他们必遭受无法补救的灾难。因为没有什么，不，没有什么像强夺、暴力和敲诈那样激怒天主。为什么呢？因为戒除这种罪非常容易。因为在此扰乱心神的并非任何自然的欲望，而是出于故意的疏忽。那么，宗徒为何称它为\BibleQuote{万恶的根源}\BibleRef{弟前 6:10}？是的，我也这样说，但这根源出于我们，并非出于事物的本性。如果你们愿意，让我们作个比较，看看哪一个更难以驾驭，是贪财的欲望还是美色的欲望；因为那曾绊倒伟人的欲望更难制伏。那么，让我们看贪财曾缠住过哪位伟人？没有一个；只缠过极其可怜卑微的人，如\Person{革哈齐}、\Person{阿哈布}、\Person{犹达斯}、犹太人的司祭们；但美色的欲望却战胜了伟大的先知\Person{达味}。我说这话，并非宽恕那些被这类情欲征服的人，而是为了让他们警醒。因为当我显出这情欲的力量时，我就尤其显明他们被剥夺了一切求恕的权利。因为如果你们从前不认识这野兽，你们还可以有这作为托词；但现在，既然认识了，却仍然陷入其中，你们将毫无借口。在他之后，这情欲更完全地占据了其子。然而，从未有人比他更智慧，他获得了所有其他的德行；但他仍被这情欲猛烈地抓住，以致在要害处受了伤。那父亲确实再次起来，重新战斗，并再次获冠；但儿子却毫无此类表现。\par

因此保禄也说：\BibleQuote{与其欲火中烧，倒不如结婚为妙} \BibleRef{格前 7:9}；基督也说：\BibleQuote{谁能领受，就领受吧。} \BibleRef{玛 14:12} 但论及金钱，祂却没有这样说，而是说：\BibleQuote{凡舍弃} 自己财物的，\BibleQuote{必要获得百倍。} \BibleRef{玛 14:29} 有人会说：\Quote{那么，祂论及富人时，为什么说他们难以进入天国呢？} 这再次指出他们品格的软弱，并非金钱的专横，而是他们彻底的奴役。这一点从保禄所给的劝告中也显而易见。因为他使人完全远离那种欲望，说道：\BibleQuote{那些想致富的人，陷于诱惑；} \BibleRef{弟前 6:9} 但论及另一种欲望时，却不是这样；而是仅使他们\BibleQuote{暂时}分开，且要\BibleQuote{两厢情愿}，他劝告\BibleQuote{再行同居} \BibleRef{格前 7:5}。因为他害怕情欲的浪潮，唯恐它们导致惨重的沉船。这种情欲甚至比愤怒更为猛烈。因为当没有激起愤怒的事实时，人不可能发怒；但即使未见引起欲望的面容，人仍不禁有欲念。因此，主并没有完全断绝这种情欲，而是加上了\BibleQuote{无缘无故}这几个字 \BibleRef{玛 5:22}。祂也没有废除一切欲望，而只是非法的欲望，因为祂说：\BibleQuote{但为了避免欲望，男人当各有自己的妻子。} \BibleRef{格前 7:2} 然而，积蓄钱财，无论是出于原因还是无缘无故，祂都不允许。因为那些情欲植根于我们的本性，是为必要的目的：欲望是为了生育子女，愤怒是为了救助受屈者；但贪财之心并非如此。因此，这种情欲也不是我们本性所固有的。所以，如果你被它俘虏，你将遭受更可耻的惩罚。正因如此，保禄甚至允许第二次婚姻，但在金钱方面却要求极大的严格，说道：\BibleQuote{为什么不情愿受欺呢？为什么不情愿吃亏呢？} \BibleRef{格前 6:7} 而当论及童贞时，他说：\BibleQuote{我没有主的命令，} \BibleRef{格前 7:25} 又说：\BibleQuote{我说这话是为了你们的益处，并不是要牢笼你们；} \BibleRef{格前 7:35} 但当论及金钱时，他说：\BibleQuote{只要有衣有食，就当知足。} \BibleRef{弟前 6:8} 有人会说：\Quote{为什么更多人被这种欲望胜过，而不是被另一种呢？} 因为他们没有像防范淫乱和奸淫那样防范它；因为如果他们同样认为它危险，或许就不会被它俘虏了。同样，那些愚拙的童贞女也被赶出婚筵，因为她们虽然击倒了强敌，却被更弱小的、微不足道的敌人所伤 \BibleRef{玛 25:1}。此外，还可以说：如果一个人制服了情欲，却被金钱所胜，那么他实际上并没有制服情欲，只是从本性上领受了不受那种强烈不安的恩赐；因为并非所有人都同等倾向情欲。既然知道这些，并常常在心中思量童贞女的榜样，让我们避开这邪恶的野兽。因为，如果童贞于她们毫无益处，反而在无尽的劳苦之后因贪财而丧亡，那么如果我们陷入这种情欲，谁能拯救我们呢？因此，我恳求你们尽力而为，既不要被它俘虏，又若被俘虏，也不要继续受奴役，而要挣断那些坚硬的锁链。这样，我们才能稳立天上，获得那不可计数的福乐；愿我们众人赖我们的主耶稣基督的恩宠和慈爱，得以达到此境。愿光荣、权能、尊崇归于祂，偕同圣父及圣神，自今至永远，及世之世。阿们。\par

\SourceRecord{Translated by Talbot W. Chambers.；Nicene and Post-Nicene Fathers, First Series；Vol. 12.；Edited by Philip Schaff.；Buffalo, NY: Christian Literature Publishing Co.,；1889.；<http://www.newadvent.org/fathers/220223.htm>.}

% structure: p0001 source=h1 target=section reason=page-title

\section{论\BibleBook{格林多}{后书}讲道词第二十四篇}

% structure: p0002 source=h2 target=subsection reason=relative-heading-rank; base=h2

\subsection{\BibleRef{格后 11:13}}

\BibleQuote{这样的人是假宗徒，是欺诈的工人，冒充基督的宗徒。}\par

你说什么呢？那些宣讲基督、不收钱、不传别的福音的人，\Quote{假宗徒吗？}\Quote{是的，}他说，而正因如此，因为他们假装这一切都是为了欺骗。\Quote{欺诈的工人，}因为他们确实工作，却拔除所栽种的。因为他们深知，否则就不会被接纳，于是戴上真理的面具，上演错误的戏剧。\Quote{然而，}有人说，\Quote{他们不收钱。}那是为了获取更大的；为了毁灭灵魂。不如说，那也是谎言；他们收钱，只是暗中进行：他在下文指出了这一点。事实上，他早已暗示了这一点，他说：\Quote{叫他们所夸耀的，能与我们一样：}\BibleRef{格后 11:12}然而在后文中，他更清楚地暗示了这一点，说：\Quote{若有人吞没你们，若有人掳掠你们，若有人自高自大，你们就容忍他。}\BibleRef{格后 11:20}但此时他以另一理由指控他们，说：\Quote{他们冒充。}他们只有一种\Emph{外表}；羊皮只是外在的装扮。\par

% structure: p0005 source=h2 target=subsection reason=relative-heading-rank; base=h2

\subsection{\BibleRef{格后 11:14-15}}

\BibleQuote{这并不稀奇，因为即使撒殚也装作光明的天使，那么他的仆役装作正义的仆役，又算什么大事呢？}\par

因此，若有人应当惊奇，这正是他应当惊奇之处，而不是他们的转变。因为当他们的师傅敢做任何事时，门徒也跟随，这并不稀奇。但什么是\Quote{光明的天使？}它能够自由发言，站在天主近旁。因为也有黑暗的天使，那些属于魔鬼的，那些黑暗而残酷的。魔鬼曾这样欺骗了许多人，他装作\Quote{成为，}而不是成为\Quote{光明的天使}。同样，这些人也带着宗徒的外形，而不是能力本身，因为他们不能。但没有任何事比为了炫耀而行事更像魔鬼。什么是\Quote{正义的职分？}那就是我们这些向你们宣讲拥有正义的福音的人。因为他或是指这个，或是说他们自己披上正义之人的外表。那么，我们如何认出他们？正如基督所说的：\Quote{凭他们的果实。}因此，他不得不将自己的善行与他们的恶行并列，以便假冒者通过比较而显明。当他即将再次进入自我赞扬时，他先指控他们，以表明这样的辩论是迫不得已的，免得有人因他谈论自己而指责他，他说：\par

% structure: p0008 source=h2 target=subsection reason=relative-heading-rank; base=h2

\subsection{\BibleRef{格后 11:16}}

\Quote{我再说。}因为他已经使用了许多预备性的纠正：\Quote{然而我对我所说的仍不满意，但我再说一次，}\par

\Quote{不可有人以为我是愚昧的。}因为这是他们所做的——无故自夸。——但是，我请你们注意，他多少次在即将开始自我赞扬时克制自己。他说：\Quote{因为夸耀确实是愚蠢的行为，但我这样做，不是出于愚妄，而是由于不得已。如果你们不相信我，即使看到有必然性还要谴责我，我也不会推辞这事。}\par

% structure: p0011 source=h2 target=subsection reason=relative-heading-rank; base=h2

\subsection{\BibleRef{格后 11:17}}

\BibleQuote{我所说的，不是依照主说的，而是如同在愚昧中，在这自夸的自信中说的。}\par

你们看见夸耀不是\Quote{依照主吗？}因为主说：\BibleQuote{你们做完了一切，要说：我们是无用的仆人。}\BibleRef{路 17:10}但是，就其本身而言，确实不是\Quote{依照主}，但因意向而成为如此。因此他说：\Quote{我所说的，}不是指责动机，而是词语。因为他的目的是如此值得称赞，以至于也尊崇了这些词语。因为如同杀人犯，虽然他的行为是最严格禁止的，却常常因意向而被认可；又如割礼，虽然它不是\Quote{依照主}，却因意向而成为如此；夸耀也是如此。那么他为什么不使用如此严格的表达呢？因为他正急于转向另一点，他慷慨地满足那些想找把柄反对他的人，甚至过度满足，使他只说得益的事；因为这些话足以消除所有怀疑。\Quote{但是如同在愚昧中。}在此之前，他说：\Quote{但愿你们能容忍我一点愚昧，}\BibleRef{格后 11:4}但现在说\Quote{如同在愚昧中；}因为越往前走，他越澄清自己的言语。为使你们不以为他在所有点上都是愚妄，他加上了\Quote{在这自夸的自信中。}在此特定语境中，他的意思是：正如在另一处，他说：\Quote{使我们不致蒙羞，}并加上了\Quote{在这自夸的自信中。}\BibleRef{格后 9:4}又在另一处，他说：\Quote{我决定的事，难道是按肉体决定的，以致在我有是而又非吗？}\BibleRef{格后 1:17}并表明他不能在所有情况下都履行他所承诺的，因为他不按肉体决定，免得有人将这种怀疑延伸到教义上，于是加上：\BibleQuote{但天主是信实的，我们对你们所说的话，并不是是而又非的。}\BibleRef{格后 1:18}\par

2. 注意，在说了那么多之后，他又进一步提出其他理由，这样说：\par

% structure: p0015 source=h2 target=subsection reason=relative-heading-rank; base=h2

\subsection{\BibleRef{格后 11:18}}

\BibleQuote{既然有许多人按肉体夸耀，我也要夸耀。}\par

\Quote{按血肉}是什么？是指外在的事物：出身高贵、财富、智慧、受割损、希伯来血统、声望。看哪，还有智慧。他列出那些他显示为虚无的事物，然后还有愚妄。因为若在真正美好的事物上夸耀是愚妄，那么在虚无的事物上夸耀更是愚妄。这就是他所谓的\Quote{不按主的意思}。因为身为希伯来人，或任何此类事物并无益处。'所以，不要以为我把这些当作德行；不；而是因为那些人夸耀，我也不得不在这几点上进行比较。'他在另一处也这样做，说：\BibleQuote{若有人以为自己可以信赖血肉，我更可以}\BibleRef{斐 3:4}；在那里，也是因为那些信赖血肉之人的缘故。就像一个出身显贵却选择了哲学生活的人，看到别人因出身高贵而极其自负，为了消除他们的虚荣，被迫谈论自己的显赫；不是为了自夸，而是为了压低他们；同样，\Person{保禄}也确实这样做了。然后，他离开那些话题，将全部责备倾泻在格林多人身上，说：\par

% structure: p0018 source=h2 target=subsection reason=relative-heading-rank; base=h2

\subsection{\BibleBook{格林多}{后书} 11:19}

\Quote{因为你们欣然容忍愚妄人。}'所以，这要归咎于你们，而且超过他们。因为如果你们没有容忍他们，并且在他们那里受到损害，我就不会说一句话；但我这样做是出于对你们得救的深切关怀，并宽容迁就。看哪，他甚至把责备和赞美结合在一起。因为说了\Quote{你们欣然容忍愚妄人}之后，他加上：\par

\Quote{你们自己是明智的。}因为夸耀，尤其是在这种事上夸耀，是愚妄的标记。然而，本该责备他们说，\Quote{不要容忍愚妄人}；他却以更有利的方式这样做。因为那样他会显得是因为自己缺乏这些优势而责备他们；但如今他显示自己在这些方面也超越他们，并且视之为虚无，从而更有效地纠正他们。目前，在进入他自己的夸耀和比较之前，他也斥责格林多人大大的奴性，因为他们对他们过分顺从。请注意他如何讽刺他们。\par

% structure: p0021 source=h2 target=subsection reason=relative-heading-rank; base=h2

\subsection{\BibleBook{格林多}{后书} 11:20}

\Quote{因为你们容忍一个人，}他说，\Quote{如果他吞吃你们。}\par

那么，你们怎么说：\BibleQuote{使他们在所夸耀的事上，可以显明与我们一样？}\BibleRef{格后 11:12}你看见他显示他们确实拿了他们的东西，而且不单是拿，甚至是过分地拿：因为\Quote{吞吃}一词清楚表明这一点。\par

\Quote{若有人使你们作奴隶。}'你们已经交出了你们的钱财、你们的身体和你们的自由。'他说，\Quote{这比拿你们的东西更甚；他们不仅成为你们钱财的主人，而且成为你们自身的主人。}他甚至在之前就表明了这一点，他说：\BibleQuote{若别人在你们身上分享这权利，我们岂不更有分吗？}\BibleRef{格前 9:12}然后他加上更严厉的话，说：\par

\Quote{若有人自高自大。}'因为你们的奴役并非温和，你们的主人也不温和，而是令人厌烦和憎恶。'\par

\Quote{若有人打你们的脸。}你看，这里又是暴政的进一步延伸？他说这话，不是指他们真的被打脸，而是指他们被唾弃、受侮辱；因此他加上：\par

% structure: p0027 source=h2 target=subsection reason=relative-heading-rank; base=h2

\subsection{\BibleBook{格林多}{后书} 11:21}

\Quote{我带着羞愧说这话，}因为你们所受的丝毫不亚于被打脸的人。还有什么比这更强烈呢？什么压迫比这更苦呢？当从你们那里拿走了你们的钱财、自由和荣誉之后，他们甚至对你们也不温和，也不让你们停留在仆人的行列中，反而比任何买来的奴隶更侮辱地对待你们。\par

\Quote{好像我们原是软弱的。}这个表达很隐晦。因为话题令人不快，所以他这样表达，通过隐晦来消除不快。他想说的是：\Quote{难道我们不能也做这些事吗？是的，但我们不做。那么，为什么你们容忍这些人，好像我们不能做这些事呢？如果你们甚至容忍那些装愚妄的人，那已经是你们的不是了；但你们容忍他们，尽管他们如此藐视你们、掠夺你们、自高自大、打你们，这是无可原谅，也毫无理由的。因为这是一种新型的欺骗方式。因为欺骗的人既给予又奉承；但这些人既欺骗，又拿走东西并侮辱你们。因此，你们完全没有借口，因为你们唾弃那些为你们而谦卑自下、为使你们被高举的人，却钦佩那些自高自大、为使你们被贬低的人。难道我们不能也做这些事吗？是的，但我们不愿意，是为了你们的益处。因为他们牺牲你们的利益追求自己的利益，而我们牺牲自己的利益追求你们的利益。}你们看，他在每一次的明白说话中，也借着所说的话警告他们。\Quote{因为，}他说，\Quote{如果是为了这个缘故你们尊敬他们，因为他们打你们、侮辱你们，我们也能这样做：使你们作奴隶、打你们、自高自大反对你们。}\par

3. 你看到他是如何将整个责任，既包括他们愚蠢的骄傲，也包括他自己看似愚拙的事，都归在他们身上吗？\Quote{因为我被迫稍作夸耀，并非为了显扬我自己，而是为使你们摆脱这痛苦的奴役，我被迫稍作夸耀。}但不应只考察所说的话，还当考察其原因。因为\Person{撒慕尔}在傅油立\Person{撒乌耳}为王时，也对自己作了一番崇高的赞辞，说：\BibleQuote{谁的驴我取过，或牛犊，或鞋？或我曾压迫过你们中的任何人？}\BibleRef{撒上 12:3}然而没有人责备他。原因在于他这样说并非为了显耀自己，而是因为他要立一位君王，便藉着自辩的形式教导他要谦和温良。请看先知的智慧，不如说是天主的慈爱。因为他想使他们转离（他们的计划，）便汇集了许多严厉的事，用以指称他们未来的君王，例如他要使他们的妻子推磨，\BibleRef{撒上 8:11}使男人做牧人和赶骡的；他详尽地列举了王权的一切职事。但当他看见这些事都不能阻止他们，他们已病入膏肓时，他便这样既宽容了他们，又使他们的君王变得温和。\BibleRef{撒上 12:5}因此他也请他作证。因为当时并没有人控告或指责他需要自辩，他说这些话是为了使他变得更好。因此他也为了挫其骄气而加上说：\BibleQuote{如果你听从，你和你的君王，}\BibleRef{撒上 12:14}如此这般的美好事物必属于你们；\BibleQuote{但如果你不听从，则相反的一切。}\Person{亚毛斯}也说：\BibleQuote{我不是先知，也不是先知的儿子，只是牧人，采集野无花果的。天主召叫了我。}\BibleRef{亚 7:14}但他这样说并非为了抬高自己，而是为了堵住那些怀疑他不是先知的人的口，并表明他不是骗子，也不是凭自己的心意说这些话。又如另一位，为表明同样的事，说：\BibleQuote{但我确实藉上主的神和力量充满能力。}\BibleRef{米 3:8}\Person{达味}在述说有关狮子和熊的事时，\BibleRef{撒上 17:34}也不是为了荣耀自己，而是为达致一个伟大而可赞叹的目的。因为人们不信他赤手空拳能战胜那蛮人，他自己甚至不能拿起武器；他被迫提供自己勇力的证据。当他割下\Person{撒乌耳}的衣边时，他所说的话不是为了炫耀，而是为了驳斥一种恶意的猜疑，有人散布说他想要杀害\Person{撒乌耳}。\BibleRef{撒上 24:4}因此，我们处处都当寻求原因。因为凡顾及听众益处的人，即使他称赞自己，非但不该受责备，反而当受加冕；如果他不说话，则当受责备。因为如果\Person{达味}当时在\Person{哥肋雅}一事上保持沉默，他们就不会允许他出战，也不会竖起那辉煌的胜利标记。因此他才被迫说话；而且不是对他的兄弟们说，尽管他们也不信任他，正如君王不信任他一样；但嫉妒堵住了他们的耳朵。因此他撇下他们，向那个尚未嫉妒他的人讲述了他的事迹。\par

4. 嫉妒真是一件可怕、可怕的事，它诱使人轻视自己的救恩。加音就是这样毁灭了自己，而在他之前的魔鬼——他父亲的毁灭者——也是如此。撒乌耳也是这样邀请恶灵来攻击自己的灵魂，邀请之后，又嫉妒那位医治他的人。因为嫉妒的本性正是如此：他知道自己得了救恩，却宁愿自己灭亡，也不愿看见救他的人受到尊荣。还有什么比这情欲更令人痛心呢？称它为魔鬼的后代也不算错。其中也包含着虚荣的果实，或者更确切地说，它的根；因为这两种恶习往往相互产生。确实，当人们说：\Quote{达味杀了万万}\BibleRef{撒上 18:7}，撒乌耳就是这样嫉妒的，还有什么比这更愚蠢的呢？因为，告诉我，你嫉妒什么？\InnerQuote{因为某人称赞了他？}其实你应当欢喜；况且，你甚至不知道这称赞是否真实。难道你因一个本不值得钦佩的人被这样称赞而悲伤吗？其实你应当同情。因为如果他是好人，你在别人称赞他时不应当嫉妒，而应当与那些赞美他的人一同赞美他；但若他不是如此，你又何必恼火呢？为什么把剑刺向自己？\InnerQuote{因为他被人钦佩。}但人今天在，明天就不在了。\InnerQuote{因为他享受荣耀？}告诉我，什么样的荣耀？就是先知称之为\Quote{草上的花}\BibleRef{依 40:6}的那种。那么，你是不是因此嫉妒，因为你没有背负重担，也没有随身带着这么一大堆草？但如果他因此在你眼中值得嫉妒，那为什么樵夫们每天背着东西进城，你不嫉妒他们呢？因为那负担并不比这好多少，甚至更糟。他们的负担只是使身体疲乏，而这往往伤害灵魂，带来比快乐更多的忧虑。即使有人通过口才获得声誉，他所受的恐惧也大于他所获得的好名声；不仅如此，前者是短暂的，后者却是长久的。\InnerQuote{但他得到当权者的宠爱？}在这方面也有危险和嫉妒。因为你对他的感受，也正如许多其他人对他的感受一样。\InnerQuote{但他不断受赞扬？}这导致痛苦的奴役。因为他不敢无所畏惧地做任何他认为应该做的事，以免得罪那些赞扬他的人，因为那种尊荣对他是一种沉重的束缚。所以他越被人认识，就有越多的主人，他的奴役就越大，因为到处都有他的主人。的确，一个仆人，当从主人的视线中解脱出来时，既能喘口气，也能自由地生活；但这个人处处遇到主人，因为他是在广场上出现的所有人的奴隶。即使有某些必要的事情催促，他也不敢踏入广场，除非有仆从跟随、有马匹和所有其他排场，以免他的主人谴责他。如果他看到某个真正朋友，他也不敢平等地与他交谈；因为他害怕他的主人，怕他们剥夺他的荣耀。所以他越显赫，就越被奴役。如果他遇到什么不愉快的事，侮辱就更加恼人，因为目击者更多，而且似乎有损他的尊严。这不仅是一种侮辱，而且是一种灾难，因为他有许多人为此幸灾乐祸；同样，如果他享受到什么好事，就有更多的人嫉妒、诽谤，并竭力要毁灭他。那么，告诉我，这是好事吗？这是荣耀吗？绝不是，而是不光彩、奴役、枷锁，以及一切可说的令人负担的东西。但如果人间的荣耀在你眼中如此值得羡慕，并且当某人受到众人称赞时让你如此不安；那么当你看到他在享受那种掌声时，就要转念想到来世和那里的荣耀。正如你匆忙逃避野兽攻击时，会进入小屋并关上门；同样，现在也要逃到来世和那说不尽的荣耀中。这样，你既能践踏今世的荣耀，也能轻易把握那来世的荣耀，享受真正的自由和永恒的美好。愿我们众人藉着吾主耶稣基督的恩宠和慈爱，与圣父及圣神，获享这一切，光荣、权能、尊崇，自今世至永远，永无穷尽。阿们。\par

\SourceRecord{Translated by Talbot W. Chambers.；Nicene and Post-Nicene Fathers, First Series；Vol. 12.；Edited by Philip Schaff.；Buffalo, NY: Christian Literature Publishing Co.,；1889.；<http://www.newadvent.org/fathers/220224.htm>.}

% structure: p0001 source=h1 target=section reason=page-title

\section{第二十五篇讲道：论格林多后书}

% structure: p0002 source=h2 target=subsection reason=relative-heading-rank; base=h2

\subsection{\BibleRef{格后 11:21}}

\BibleQuote{然而，不论别人在何事上大胆——我是在愚妄中说——我也大胆。}\par

看他又在退缩，预先使用贬抑和纠正，尽管他已经说过许多此类的话：\BibleQuote{但愿你们能容忍我一点愚妄；}\BibleRef{格后 11:1} 又说：\BibleQuote{不要有人以为我是愚妄的；若是这样，也要以愚妄来接纳我。}\BibleRef{格后 11:16} \BibleQuote{我所说的，不是按主说的，乃是像在愚妄中说的。}\BibleRef{格后 11:17} \BibleQuote{既然有许多人按着血肉夸耀，我也要夸耀；}\BibleRef{格后 11:18} 这里又说：\BibleQuote{不论别人在何事上大胆——我是在愚妄中说——我也大胆。} 他把说自己任何大事称为大胆和愚妄，即使有必需，也教导我们要格外避免这类事。因为若我们做了一切后，仍该自称是无用的仆人；那么，在没有理由催促时，高举自己并自夸的人，能配得什么宽恕呢？因此，法利塞人也遭到了那样的命运，甚至在港湾中触礁沉船，因为他撞上了这块暗礁。同样，保罗即使看到充分的必要，也依然退缩，并不断指出这样的说话是愚妄的标记。然后最后，他才冒险，以必需为借口，说道：\par

% structure: p0005 source=h2 target=subsection reason=relative-heading-rank; base=h2

\subsection{\BibleRef{格后 11:22-23}}

\BibleQuote{他们是希伯来人吗？我也是。他们是以色列人吗？我也是。}\par

因为并非所有希伯来人都是以色列人，因为阿孟人和摩阿布人都是希伯来人。因此他加上一点，以澄清他血统的高贵，并说：\par

\BibleQuote{他们是亚伯拉罕的后裔吗？我也是。他们是基督的仆役吗？——我像发狂的人说话——我更是。}\par

他不满足于先前推卸的言词，而在这里又再次使用。\BibleQuote{我像发狂的人说话，我更是。} 我是他们的上级，比他们更优秀。事实上，他有明确的证据证明他的优越，但即使如此，他仍称这事为愚妄。然而，如果他们是假宗徒，他就不必以比较的方式提出自己的优越性，而只要完全否定他们被称为\Quote{仆役}的资格。确实，他否定了，说：\BibleQuote{假宗徒，欺诈的工人，装成基督的宗徒。}\BibleRef{格后 11:13} 但现在他并没有这样做，因为他的讨论即将进入严格的审查；而在审查时，没有人只是单纯地断言。但他先以比较的方式陈述情况，然后表明事实否定了它，这是一个非常强有力的否定。此外，当他说\BibleQuote{他们是基督的仆役吗？}时，他给出的是他们的看法，而非他自己的断言。说了\BibleQuote{我更是}之后，他继续进行比较，并表明他不是靠空言，而是通过提供事实的证据来维持宗徒身分的印记。并且，他撇开所有的奇迹，从自己的考验开始，于是说道：\par

\BibleQuote{劳碌更频繁，鞭打超过限度。} 后者比前者更甚；即既被棍打又被鞭打。\par

\BibleQuote{坐监更频繁。} 这里同样有增加。\BibleQuote{屡次冒死。}\BibleRef{格前 15:31} 因为他说：\BibleQuote{我是天天冒死。} 但这里，甚至是实际的死亡；因为我常被置于死亡的险境中。\par

% structure: p0012 source=h2 target=subsection reason=relative-heading-rank; base=h2

\subsection{\BibleRef{格后 11:24}}

\BibleQuote{从犹太人那里，我五次接受了四十减一下的鞭刑。}\par

为什么是\Quote{减一}？有一条古法，凡受鞭超过四十下者，在他们中间被视为可耻。因此，为免行刑者的猛烈和冲动在施刑时超过数目而使人蒙羞，他们规定应施行\Quote{减一}，这样即使行刑者超过，也不会超过四十，而留在规定的数目之内，不会使受鞭者蒙受耻辱。\par

% structure: p0015 source=h2 target=subsection reason=relative-heading-rank; base=h2

\subsection{\BibleRef{格后 11:25}}

\BibleQuote{三次被棍打，一次被石头砸，三次遭遇海难。}\par

这与福音有什么关系？因为他远行；而且走海上。\par

\BibleQuote{在深海中度过了一昼一夜。} 有人说这意味在开阔的海面上，有人说是在海水中游泳，这解释更正确。至少前者并不稀奇，他也不会将其置于比海难更大的位置。\par

% structure: p0019 source=h2 target=subsection reason=relative-heading-rank; base=h2

\subsection{\BibleRef{格后 11:26}}

\BibleQuote{河流的危险。}\par

因为他也被迫渡河。\BibleQuote{盗贼的危险，城中的危险，旷野的危险。} \Quote{到处都有争战摆在我面前，在地方、在地区、在城巿、在旷野。}\par

\BibleQuote{外邦人的危险，假弟兄中的危险。}\par

看，另一种交战。因为不仅敌人攻击他，连那些假意的人也是如此；他需要极大的坚定和极大的谨慎。\par

% structure: p0024 source=h2 target=subsection reason=relative-heading-rank; base=h2

\subsection{\BibleRef{格后 11:27-28}}

2. \BibleQuote{劳碌和辛苦。}\par

危险接着劳碌，劳碌接着危险，一个接一个，从不间断，甚至不允许他稍作喘息。\par

\BibleQuote{旅行频繁，饥渴赤身，此外还有那些外面的事。}\par

略去的比列出的更多。其实，即使是列出的，也无法数算其数目；因为他没有具体列出，而只提到那些数目少且容易理解的，说\BibleQuote{三次}和\BibleQuote{三次}\BibleRef{格后 11:25}，以及\BibleQuote{一次}；至于其他的，他没有提及次数，因为他经常忍受。他不叙述它们的成果，比如他归化了多少人，而只叙述他为宣讲所受的苦；这既是出于谦虚，也是表明即使除了劳苦一无所获，他获得酬报的权利也已经满足了。\par

\BibleQuote{那每日压在我身上的事。}骚动、扰乱、群众的攻击、城池的围攻。因为犹太人最攻击此人，因为他最使他们混乱，他突如其来的转变是对他们疯狂的最有力的驳斥。一场大战向他袭来，来自本族、外人、假弟兄；处处是波涛和悬崖，在有人居住之地，在无人之域，陆上、海上、外在内里。他甚至连必需的食物和薄衣都没有，但世界的冠军赤身搏斗、饥饿奋战；他远未使自己富足。然而他不抱怨，反而为这些事感谢战斗的裁判者。\par

\BibleQuote{对众教会的挂虑。}这是所有事中最主要的，他的灵魂也分心，思绪分散。因为即使没有外在的攻击，内在的战争、一波又一波的浪潮、忧虑的冰雹、思想之战也已足够。若一个只管理一座房屋、有仆人和管家的人，常常因挂虑而喘不过气，虽无人搅扰；那么他不仅管理一座房屋，而是城市、民族和整个世界，且面对如此重大的关切，被如此多人恶待，独自一人，忍受如此多事，又如此温柔地关心他人，甚至超过父亲对儿女——请想想他所忍受的。为使你不说：他固然挂虑，但挂虑轻微，他更进一步说明挂虑的强度，说道：\par

% structure: p0031 source=h2 target=subsection reason=relative-heading-rank; base=h2

\subsection{\BibleRef{格后 11:29}}

\BibleQuote{谁软弱，我不软弱呢？}他没有说\Quote{我不分享他的沮丧？}而是\Quote{我如此烦恼不安，仿佛我自己正受着那种情感、那种软弱之苦。}\par

\BibleQuote{谁跌倒，我不燃烧呢？}看，他如何再次以\Quote{燃烧}一词向我们展示他忧伤的过度。他说：\Quote{我着火}、\Quote{我燃起火焰}，这无疑比他所说的一切都更强烈。因为其他事虽猛烈，却迅速过去，并带来不褪色的快乐；但这却使他痛苦和窘迫，彻底刺透他的心灵：为每一个软弱的人遭受这些事，无论他是谁。因为他不仅为较大者感到痛苦而轻视较小者，而且将最卑微的人也视为密友。因此他也说\BibleQuote{谁软弱？}无论他是谁；他好像自己就是全世界的教会，为每一个肢体而忧伤。\par

% structure: p0034 source=h2 target=subsection reason=relative-heading-rank; base=h2

\subsection{\BibleRef{格后 11:30}}

\BibleQuote{若必须夸耀，我将夸耀那些关乎我软弱的事。}\par

你们看到，他从不夸耀奇迹，而是夸耀他的迫害和试炼？因为这就是\Quote{软弱}的意思。他表明他的战斗具有多样性。因为犹太人攻击他，外邦人反对他，假弟兄与他作战，弟兄们因软弱和跌倒而令他忧伤——他从朋友和外人各方面都遇到麻烦和困扰。这是宗徒的特殊标志，福音由此编织而成。\par

% structure: p0037 source=h2 target=subsection reason=relative-heading-rank; base=h2

\subsection{\BibleRef{格后 11:31-33}}

\BibleQuote{主耶稣的天主和圣父知道我不说谎。在阿勒达王手下的总督把守了大马士革城，想要捉拿我。}\par

为什么他在这里强烈确认并保证[他的真实]，而之前任何事上都没有这样做？也许因为这件事年代较久，不太为人所知；而其他事实，他对教会的挂念等等，他们自己都清楚。请看战事多么巨大，以致全城因他而被\Quote{把守}。当我说战事时，我说的是保禄的热忱；除非这热忱强烈燃烧，它不会把总督煽动到如此疯狂。这些是宗徒灵魂的作为：忍受如此大事，却丝毫不转向，而是高尚地承受一切；但也不去冒险，不冲向危险。例如这里，他满足于\Quote{被人用篮子从窗户缒下}以躲避围困。因为他虽也渴望离世，但他仍热切地关心人的救恩。因此他常采取这类办法，为宣讲而保存自己；当情况需要时，他也不拒绝使用人的计谋，如此清醒和警醒。因为在无法避免的灾祸中，他只需要恩宠；但在有限度的试炼中，他自己也筹划许多事，却又将一切归于天主。正如一颗不灭的火星，若落入海中，会被层层波涛淹没，却会再次闪耀升到水面；同样，有福的保禄也时而因危险被淹没，时而又从其中潜出，更加光彩夺目，以受苦战胜恶事。\par

3. 这就是辉煌的胜利，这就是教会的凯旋，魔鬼就是这样被推翻：当我们受苦时，他被战胜。因为当我们受苦时，他被俘虏；当他想要加害我们时，自己反受其害。这事也发生在保禄身上：他越是遭遇危险，就越是得胜。魔鬼不仅用同一种试炼来攻击他，而是用各种不同的试炼。有些涉及劳苦，有些涉及忧愁，有些涉及恐惧，有些涉及痛苦，有些涉及挂虑，有些涉及羞辱，有些则同时包含这一切；然而他在一切事上都得胜了。就好像一个单独的士兵，面对全世界与他为敌，却穿过敌人的行列而毫发无损：保禄正是这样，独自一人，在野蛮人中，在希腊人中，在每一块土地上，在每一片海洋上，都屹立不倒。又如一点火星落在芦苇和干草上，使被点燃的东西变成自己的性质；同样，这个人扑向一切，使万物转变成真理；如同冬日的洪流，席卷一切，冲垮每个障碍。又如一位角力、赛跑、拳击的冠军；或一名士兵轮流进行攻城、步战、海战；他如此轮流尝试各种战斗，口吐火焰，无人能接近；他以一己之身夺取世界，以一己之舌驱散众人。那些众多号角落在耶里哥城墙上将其摧毁的力量，也不及这人声音的响声，既击倒魔鬼的营垒，又将反对他的人争取过来。当他聚集了一大批俘虏，又武装了他们，使他们成为自己的军队，并借着他们征服敌人。达味以一块石头击倒哥肋雅，令人惊叹；但若你细查保禄的成就，那简直是小孩子的把戏；牧羊人与将军之间的巨大差异，正是你在这里要看到差距。因为这人不是投石击倒哥肋雅，而是只凭言语就驱散了魔鬼的全军；他如狮子般咆哮，从舌头上射出火焰，所有人都感到他不可抗拒；他不停地四处跳跃；他跑到这些人跟前，又来到那些人面前，转向这些人，又跳到那些人面前，比风更快；管理整个世界，如同管理一座房屋或一艘船；拯救下沉者，稳定头晕者，鼓舞水手，坐在舵旁，注视船头，收紧帆桁，划桨，拉桅，观察天空；他自己成为一切，既是水手，又是舵手，又是舵手助手，又是帆，又是船；他忍受一切，为要解除他人的苦难。试想：他经历船难，为要制止世界的船难；\BibleQuote{他曾在深渊里度过一日一夜}，为要把它从错误的深渊中拉上来；他\BibleQuote{处在劳苦中}，为要安慰劳苦的人；他忍受鞭打，为要医治那些被魔鬼击打的人；他身在监狱，为要带领那些坐在监狱和黑暗中的人进入光明；他\BibleQuote{多次受死亡的危险}，为要拯救他们脱离死亡的痛苦；\BibleQuote{五次受了四十下鞭刑，减去一下}，为要释放那些施行鞭刑的人脱离魔鬼的鞭笞；他\BibleQuote{被棍杖击打}，为要带领他们归入基督的\BibleQuote{棍杖和牧杖}\BibleRef{咏 22:4}；他\BibleQuote{被石头砸}，为要拯救他们脱离无知的石头；他\BibleQuote{在旷野里}，为要把他们领出旷野；\BibleQuote{行路}，为要止住他们的漂泊，并打开通往天堂的道路；\BibleQuote{在城中有危险}，为要显示那上面的城；\BibleQuote{饥渴}，为要拯救他们脱离更严重的饥饿；\BibleQuote{赤身露体}，为要以基督的长袍遮盖他们的不雅；被群众围攻，为要解救他们脱离魔鬼的围困；他焚烧，为要熄灭魔鬼燃烧的箭；\BibleQuote{他从窗口被用筐子从城墙上缒下}，为要把那些卧在地上的人从下面升上去。我们还能再说什么呢？既然我们连保禄受了什么苦都不知道？我们还能再提财产、妻子、城市或自由吗？当我们看见他千万次地连生命本身都视若无睹？殉道者只死一次；而这位蒙福的圣人在他一个身体、一个灵魂内，承受了如此多的危险，足以动摇铁石般的灵魂；众圣人一起在许多身体里所承受的一切，他独自一人在一个身体里承受了；他进入世界如同进入一个竞技场，脱去一切，然后如此英勇地站立。因为他知道与他角力的魔鬼。因此，他从一开始，从起跑线就大放光芒，一直坚持到底；不，他甚至越接近奖赏就越加奋力。最令人惊叹的是，尽管他经历和成就了如此伟大的事，他却懂得保持极其谦卑。当他被迫叙述自己的善行时，他匆匆带过一切；尽管如果他愿意详细展开他所提到的每件事，他可以写成无数的书卷；如果他说出他所挂虑的教会、他的监狱和在其中的成就，以及其他事一件件地——围攻、袭击。但他不愿意。既然知道这些，让我们也学习谦卑，任何时候不要以财富或世间其他事物夸耀，只以我们为基督所受的凌辱夸耀，而且只在必要时才这样做；如果没有必要，我们连这些也不要提（免得我们骄傲），只提我们的罪过。因为这样，我们既容易从罪过中得释放，也会使天主对我们仁慈，并得到未来的生命；愿我们所有人都能借着我们的主耶稣基督的恩宠和慈爱达到这生命，愿光荣、权能、尊崇归于圣父及圣神，从今世直到永远。阿们。\par

\SourceRecord{Translated by Talbot W. Chambers.；Nicene and Post-Nicene Fathers, First Series；Vol. 12.；Edited by Philip Schaff.；Buffalo, NY: Christian Literature Publishing Co.,；1889.；<http://www.newadvent.org/fathers/220225.htm>.}

% structure: p0001 source=h1 target=section reason=page-title

\section{格林多后书讲道 第26篇}

% structure: p0002 source=h2 target=subsection reason=relative-heading-rank; base=h2

\subsection{格林多后书 12:1}

\BibleQuote{我自夸固然无益，[因为]我必来到主的异象和启示。}\par

这是什么？说了如此伟大之事的人，难道会说\Quote{我自夸固然无益？} 好像他什么也没说？不，并非好像他什么也没说：而是因为他将转向另一种夸耀，这种夸耀确实并非因如此大的赏报而有意为之，但在众人（虽非谨慎的检视者）看来，却似乎更显精彩，他说：\Quote{我自夸固然无益。} 对于真正伟大的夸耀理由，是他所数算的那些，即他的考验；然而他还有其他事要讲述，即关于异象和不可言传的奥秘。那么，他为什么说\Quote{我自夸固然无益}呢？他的意思是，\Quote{免得它使我骄傲。} 你说什么？因为即使你不讲论它们，难道你不知道它们吗？我们知道自己知道它们，并不及向他人公开宣扬那样使我们骄傲。因为使人骄傲的通常不是善行的本质，而是被众人见证和知晓。为此，他说\Quote{我自夸固然无益}；并说，\Quote{免得我在听者心中植入对我过高的想法。} 因为那些假宗徒，甚至说自己不真实的事；而此人连真实的事也隐藏，尽管有如此大的必要临到他身上，却说\Quote{我自夸固然无益}；教导每一个人，甚至过度地避免这类事。因为这种事没有益处，反而有害，除非有某种必要而有益的理由促使我们为之。那么，在讲述了他的危险、考验、陷阱、沮丧、海难之后，他转向另一种夸耀，说：\par

% structure: p0005 source=h2 target=subsection reason=relative-heading-rank; base=h2

\subsection{格林多后书 12:2-5}

\BibleQuote{我认识一个人，十四年前（或在身内，我不知道；或在身外，我不知道；天主知道）被提到第三层天。我也知道他被提到乐园里，（或在身内，我不知道；或在身外，我不知道；）听见了不可言传的话，是人不可说出的。为这样的人，我要夸耀；但为我自己，我不夸耀。}\par

这个启示确实是伟大的。但这并非唯一的：还有其他的，但他只提其中之一。因为有许多，听他怎么说：\BibleQuote{免得我因启示的宏大而过于高抬自己。} 然而，有人会说：\Quote{如果他想要隐藏它们，他就不应该给予任何暗示或说任何这类话；但如果他想要谈论它们，就应该坦白地讲。} 那么，为什么他既不坦白地讲，也不保持沉默？是为了借此表明他是不情愿地做这事。因此他也说明了时间，\Quote{十四年。} 因为他不是无目的地提及，而是表明那长久以来克制自己的人，除非有极大的必要，否则不会现在说出来。但他仍会保持沉默，除非他看到弟兄们正在灭亡。如果保禄从最初就是这样一个人，被堪当如此大的启示，虽然他还未行出这样的善行；那么，试想十四年后他成长了多少。并且留意在这件事上他是如何表现谦逊的，他说了一些事，却承认对其他事无知。因为他确实被提了，他宣告了，但他说不知道是\Quote{在身内}还是\Quote{在身外}。然而，如果他只告诉自己被提而保持沉默，那也足够了；但实际上，他因谦逊又加上了这一点。那么，是心神和灵魂被提，而身体留在死亡中吗？还是身体被提？无法断定。因为保禄自己被提，遭遇了如此多不可言传的大事，尚且无知，更何况我们。确实，他当时在乐园里，他知道；他在第三层天，他并非不知；但方式他并不清楚。再从一个方面看他如何远离骄傲：在他关于\Quote{大马士革城}的叙述中\BibleRef{格后 11:32}，他证实了自己所说的，但在这里却没有；因为他的目的不是要强烈地确立这一事实，而只是提及和暗示它。所以他也接着说：\BibleQuote{为这样的人，我要夸耀；} 意思不是被提的是另一个人，而是他尽可能以最好的方式措辞，以便既提及事实，又避免公开谈论自己。当谈论自己时，引入另一个人还有什么逻辑呢？那么他为什么这样表达呢？说\Quote{我被提了}，和说\Quote{我认识一个被提的人}；以及\Quote{我要夸耀自己}，和\Quote{我要夸耀这样的人}，这两者是不一样的。如果有人问：\Quote{没有身体怎么可能被提呢？} 我将问他：\Quote{有身体怎么可能被提？} 因为，如果你用推理来考察而不给信心留余地，后者比前者更难以解释。\par

2. 但他为什么也被提呢？我想，是为了使他看起来不逊于其他宗徒。因为他们曾与基督交往，而保禄没有；所以基督也把他提到荣耀中来。\BibleQuote{进入乐园。} 因为这个地方的名声很大，到处被人传颂。因此，基督也说：\BibleQuote{今天你就要同我在乐园里。} \BibleRef{路 23:43}\par

\Quote{我愿为这样的人夸耀吗？} 为什么呢？因为如果另一个人被提起，你为什么夸耀？由此显然他是在说自己的事。而且如果他加上，\Quote{但我不夸耀我自己}，他不外乎是说，\InnerQuote{在没有必要时，我不会无益地随意说那样的话}；或者他又在尽量模糊他所说的话。因为后面的话也清楚表明，整个论述是关乎他自己的；因为他接着说：\par

% structure: p0010 source=h2 target=subsection reason=relative-heading-rank; base=h2

\subsection{格林多后书 12:6}

\BibleQuote{但我若愿意夸耀，也不算愚妄，因为我要说实话。}\par

那么，你先前怎么说，\BibleQuote{但愿你们容忍我一点愚妄} \BibleRef{格后 11:1} 而且，\BibleQuote{我所说的，不是按主说的，而是像在愚妄中} \BibleRef{格后 11:17} 但这里却说，\BibleQuote{即使我愿意夸耀，也不算愚妄？} 这不是关于夸耀，而是关于说谎；因为如果夸耀是愚妄，那么说谎更是如此？\par

因此他接着说，\BibleQuote{我不算愚妄。} 为此他又加上：\par

\BibleQuote{因为我要说实话；但我住口，免得有人把我看得超过他所看见的，或从我所听见的。} 这里你有了公认的理由；因为他们甚至把他们当作神，由于他们奇迹的伟大。正如在元素的情况下，天主做了两件事：同时创造了它们软弱而又荣耀；一是彰显祂自己的大能，二是防止人类的错误：同样，在这里他们也既奇妙又软弱，以便不信的人通过事实本身受到教导。因为如果他们在继续只显奇妙而没有显示软弱的情况下，试图用言语劝许多人不要对他们有过高的想法；他们不仅不会成功，反而会适得其反。因为言语上的劝阻似乎更出于谦卑，反而会使他们更受钦佩。因此，在行动和事实中，他们的软弱被揭示了。在旧律法下生活的人中也可以看到这一点。因为\Person{厄里亚}是奇妙的，但有一次他被证明是胆怯的；\Person{梅瑟}是伟大的，但他也在同样的情绪影响下逃跑了。这些事发生在他们身上，是因为天主退避，允许他们的人性显露出来。因为如果因为\Person{梅瑟}带领他们出来，他们说\Quote{梅瑟在哪里？} 如果他们也在里面，他们会说什么呢？因此，[保禄]自己说，\BibleQuote{我住口，免得有人把我看得}。他没有说\Quote{说我看}，而是\BibleQuote{免得有人甚至把我看得}超过我所应得的。由此也明显看出，整个论述是关乎他自己的。因此，甚至当他开始时，他说，\BibleQuote{我固然不必夸耀}，如果他是在说别人的事，他就不会这么说。因为为什么夸耀别人是\Quote{不必}的呢？但正是他自己被算为配得这些事的；因此他接着说：\par

% structure: p0015 source=h2 target=subsection reason=relative-heading-rank; base=h2

\subsection{格林多后书 12:7}

\BibleQuote{又恐怕我因那启示的超越性而过于自高，所以有一根刺加在我身上，就是撒殚的使者来攻击我。}\par

你说什么？他把天国算作虚无，甚至地狱也不算什麼，比起他对基督的渴望；难道他把从许多人得来的尊荣看作什麼，以致既会自高，又需要那样的缰绳不断地约束？因为他不是说\Quote{它}要\Quote{攻击我}，而是\Quote{它}要\Quote{攻击我}。但是谁能这样说呢？那么，所说的话是什么意思呢？当我们解释了所谓的\Quote{刺}和这个\Quote{撒殚的使者}到底是谁以后，我们也将说明这一点。有些人说，他指的是一种头痛，是由魔鬼所施加的；但天主禁止！因为\Person{保禄}的身体绝不可能被交在魔鬼手中，因为魔鬼本人也服从\Person{保禄}的单纯命令；当\Person{保禄}把犯奸淫者交出来毁灭肉体时，他给魔鬼设定了律法和界限，魔鬼也不敢越界。那么，所说的话是什么意思呢？在希伯来语中，敌人被称为撒殚；在\BibleBook{列王}{纪上}，圣经就这样称呼那些敌对者；论到\Person{撒罗满}，说：\BibleQuote{在他的日子没有撒殚}，也就是说，没有敌人、仇敌或对手。\BibleRef{列上 5:4} 那么，他所说的是：天主不允许宣讲进展，为了抑制我们的骄傲思想；却允许敌对者攻击我们。因为这一点足以打消他的骄傲思想，而不是头痛。所以，通过\Quote{撒殚的使者}，他指的是铜匠\Person{亚历山大}、\Person{希默纳约}和\Person{费肋托}一党，以及圣道的所有敌对者；那些与他争斗对抗的人，那些把他投入监狱、鞭打他、将他带往死亡的人；因为他们做撒殚的工作。正如他称那些效法魔鬼行为的犹太人为魔鬼的儿女，同样他也称每一个对抗者为\Quote{撒殚的使者}。所以他说：\BibleQuote{有一根刺加在我身上来攻击我；} 不是天主把武器交在这样的人手中，天主禁止！不是祂来惩罚或责打，而是祂在时间和场合上允许和准许。\par

% structure: p0018 source=h2 target=subsection reason=relative-heading-rank; base=h2

\subsection{格林多后书 12:8}

3. \BibleQuote{关于这事，我三次求过主。}\par

也就是说，多次。这也是极其谦卑的标志，他不隐瞒自己不能忍受那些隐藏的阴谋，他在这些事下晕倒，并被迫祈求解脱。\par

% structure: p0021 source=h2 target=subsection reason=relative-heading-rank; base=h2

\subsection{格林多后书 12:9}

\BibleQuote{他对我说：我的恩宠够你用的，因为我的能力在软弱中得以完全。}\par

就是说：\Quote{为你来说，使死人复活，使瞎子复明，使癞病人洁净，行那些其他的奇迹，已经足够了；不要还寻求免于危险和恐惧，以及在无困扰中宣讲。但你是否因着许多人对你图谋、殴打、骚扰、鞭打，而痛苦沮丧，怕这显得是由于我的软弱？这正是显示了我的能力。因为他说：\InnerQuote{我的能力在软弱中才得以完全。}当你们被迫害而胜过迫害者时；当你们受骚扰而胜过骚扰者时；当你们被捆绑而使捆绑你们的人归化时。所以不要寻求超出需要的事。}你看见他自己提出一个理由，而天主提出另一个吗？因为他自己说：\Quote{免得我过于高举自己，就给了我一根刺}；但他说天主允许这事是为了显示祂的能力。所以你所寻求的不仅是不需要的，而且会遮蔽我能力的荣耀。因为藉着\Quote{为你来说已经足够}这句话，祂要表明的就是：无需再加添别的，一切已经完备。由此可见，他所说的并非头痛；事实上，他们生病时并不宣讲，因为生病不能宣讲；但遭受骚扰和迫害时，他们却战胜了一切。他说：\Quote{在听了这话之后，}\par

\Quote{所以我极其欣然夸耀我的软弱。}为使他们不致消沉，当那些假宗徒因着相反的命运而夸耀，而这些人却受迫害时，他表明他因此更为光耀，天主的德能也因此更为彰显，所发生的事正是夸耀的缘由。因此他说：\Quote{所以我极其欣然夸耀。}\Quote{所以我刚才列举的事，或是我刚说的\InnerQuote{有一根刺加给了我}，并非出于忧伤，而是以此自豪，并为自己吸引更大的能力。}因此他又加上：\par

\Quote{为使基督的能力覆庇我。}在这里他也暗示了另一件事，即：考验越强烈，恩宠也按同样比例增加并持续。\par

% structure: p0026 source=h2 target=subsection reason=relative-heading-rank; base=h2

\subsection{格林多后书 12:10}

\Quote{因此我为许多软弱而喜悦。}什么样的软弱？请告诉我。\Quote{在凌辱中，在迫害中，在困乏中，在艰难中。}\par

你看见他现在怎样以最清楚的方式揭示了吗？因为列举软弱的种类时，他说的不是热病，也不是任何这类的复发，也不是任何其他身体的疾病，而是\Quote{凌辱、迫害、艰难}。你看见一个纯朴的灵魂吗？他渴望脱离那些危险；但当他听到天主的答复说这并不适宜时，他不但不因祈祷失望而难过，反而欢喜。因此他说：\Quote{我喜悦}，\Quote{我欢喜，我渴望为基督的缘故受凌辱、受迫害、受艰难。}他说这些话，既是为了抑制那些人，也是为了鼓舞这些人的精神，使他们不以保禄的受苦为耻。因为那理由足以使他们比所有人都更光辉。然后他又提出另一个理由。\par

\Quote{因为我软弱时，正是我刚强时。}\Quote{你为何惊奇天主的德能在那时显明呢？我也在\InnerQuote{那时}刚强}；因为那时恩宠最丰富地降临到他身上。\Quote{因为基督的苦难怎样在我们身上增多，我们的安慰也怎样藉着基督而增多。}\BibleRef{格后 1:5}\par

4. 哪里有苦难，哪里也有安慰；哪里有安慰，哪里也有恩宠。例如，当他被投入监狱时，就在那时行了那些奇异的事；当他遭遇船难并被抛到那野蛮之地时，就在那时他更为光荣。当他被捆绑着进入审判厅时，就在那时他甚至胜过了判官。在旧约中也是如此；义人因着考验而兴盛。三个青年如此，达尼尔如此，梅瑟和若瑟也是如此；他们因此都大放光彩，并堪配巨大的荣冠。因为当灵魂为天主而受苦难时，它便被净化；那时它享受更大的助佑，因为它更需要帮助，也配得更多的恩宠。事实上，在天主为它预备的赏报之前，它已藉着成为哲学家而获得了丰富的善果。因为苦难撕裂骄傲，剪除一切懈怠，并锻炼忍耐；它揭示人世间事物的卑微，并引领人进入高深的哲学。因为所有的情欲都在它面前退却：嫉妒、竞争、贪欲、统治、贪财、贪色、夸耀、骄傲、忿怒，以及其余所有这些紊乱的癖好。若你愿意在实际运作中看到这一点，我将能够向你展示一个个体和一个民族，在苦难中与在安逸中，从而教导你前者带来多大的益处，后者带来多大的懈怠。\par

希伯来人在受磨难和迫害时，曾呻吟呼求天主，从而由上获得了极大的助佑；但一旦养尊处优，便踢跳反抗。而尼尼微人在安享太平之时，如此激怒了天主，以致天主威胁要从根基上拔除整座城市；但经过那次宣讲的谦卑之后，他们彰显了全部德性。若你愿意看一个个人，且看撒罗满。当他怀着焦虑和烦恼为治理那民族而思虑时，蒙受了那神视；但当他纵情享乐时，便滑入了邪恶的深坑。至于他的父亲达味，何时堪受钦佩、令人难以置信？岂不是在患难之时吗？阿贝沙隆不是在他流亡时头脑清醒，而返回后却变成了暴君和弑父者吗？约伯呢？他在顺境中固然闪耀，但在磨难后更显光辉。为何一定要谈论古老的事呢？因为若有人考察我们现今的状况，就会看到磨难有多大益处。的确，如今我们安享和平，变得懒散松懈，以无数的邪恶充满了教会；然而当我们受迫害时，我们更清醒、更仁慈、更热切，对于这些聚会和聆听也更积极。因为火之于金子如何，磨难之于灵魂也如何：它拭去污秽，使人洁净，使之明亮发光。磨难引人走向天国，而享乐引人走向地狱。因此一条路是宽的，另一条是窄的。为此，主亲自说：\Quote{在世界上你们要受苦难，}\BibleRef{若 16:33} 仿佛为我们留下了某种大福。所以，若你是门徒，就应行走那又直又窄的路，不要厌恶，不要气馁。因为即使你不以这种方式受磨难，也定会因其他无法获益的原因而受苦。嫉妒的人、贪财的人、为妓女痴狂的人、爱慕虚荣的人，以及所有随从邪恶的人，都遭受许多灰心和磨难，其痛苦不亚于哀伤的人。若他们不哭泣不哀伤，那是出于羞耻和麻木：因为若你细察他们的灵魂，便会看见它充满无数波涛。既然无论我们选择这种生活方式或那种，都必然要受苦，为何不选择这条在磨难中带来无数荣冠的道路呢？因为天主就是这样藉磨难和困苦引导了所有圣人，既造福他们，也使其余的人免于对他们产生不应有的过高评价。因为偶像崇拜最初得以立足，正是由于人们受到了超过其应得的钦佩。例如，罗马元老院曾将亚历山大封为第十三位神祇，因为该院有权选举和登记神明。当时，关于基督的一切已有所传闻，该国统治者派人询问，是否乐意也将基督选为神祇。但他们拒绝同意，愤怒而愤慨，因为在他们投票和法令之前，被钉者的能力已经闪耀四方，使全世界都皈依了对祂的崇拜。然而，即使他们不愿意，也如此安排：基督的神性不是藉人的谕令宣告的，祂也不被视为由他们所选举的众多神明之一。因为他们甚至将拳击手和哈德良的宠儿封为神祇，安提诺城即以其命名。既然死亡见证了他们的道德本性，魔鬼便发明了另一种方式，即灵魂不朽之说；并以此混入过分的谄媚，诱使许多人陷入不敬虔。请看何等邪恶的诡计！当我们为善的目的宣讲那教义时，他推翻我们的言语；但当他本人想为恶构建论据时，却热心地建立它。若有人问：\Quote{亚历山大如何是神？他不是死了吗？而且死得很惨？}他回答说：\Quote{是的，但灵魂是不朽的。}你看，你为不朽辩护并哲学化，是为了使人们远离统御万有的天主；但当我们宣称这是天主最大的恩赐时，你却说服你的受骗者，说人是卑贱低下的，与牲畜无异。若我们说：\Quote{被钉者活着}，立刻引来嘲笑；尽管普世都宣扬这事，古时和现今都是：古时藉奇迹，现今藉皈依者；因为这些成功确实不是一位死者的作为；但若有人说\Quote{亚历山大活着}，你便相信，尽管你提不出任何奇迹。\par

5. \Quote{是的，}有人回答说，\Quote{我有；因为他活着时行了许多伟大的事业；他征服了国家与城市，在许多战争和战役中得胜，并竖立了凯旋柱。}\par

如果我要显明一些东西，是他在世时从未梦想过的，他自己和任何其他在世的人都没有梦想过的，那么你还要求什么复活的证明呢？因为一个人活着时，作为国王，拥有军队，赢得战役和胜利，这并不奇妙，不，也不惊人或新奇；但经过\OriginalTerm{十字架}{Cross}和\OriginalTerm{坟墓}{Tomb}之后，竟能如此行大事遍及各地和海上，这才最充满惊奇，并宣告祂神圣而不可言喻的\OriginalTerm{德能}{Power}。而且\Person{亚历山大}去世后，确实再也没有恢复他那被分裂和完全灭亡的王国：的确，死了的他，又怎能做到呢？但基督却是在死了以后，最有力地建立祂的（王国）。我又何必说起基督呢？因为祂也赐予祂的门徒们，在死后发光。告诉我，\Person{亚历山大}的坟墓在哪里？指给我看，并告诉我他去世的日子。但是基督的仆人们，他们的坟墓是光荣的，因为它们占据了最忠信的城市；他们的纪念日为人熟知，使世界欢庆。而他的坟墓连自己人也不知道，但这人的坟墓连化外人也都知道。而被钉十字架者的仆人们的坟墓，比王宫更为辉煌；不是因为建筑物的大小和美丽（虽然在这方面也胜过它们），而是——更重要的——在于到访者的热情。因为穿紫袍的人亲自去拥抱那些坟墓，放下骄傲，站着恳求圣人们作他在天主面前的\OriginalTerm{中保}{Advocates}，戴王冠的人恳求那做帐篷的与渔夫，虽然死了，作他的\OriginalTerm{主保}{Patrons}。那么，告诉我，你敢称这些死者的主为谁吗？他们的仆人即使在死后，还是世上君王的\OriginalTerm{主保}{Patrons}。这件事不仅发生在\Place{罗马}，也发生在\Place{君士坦丁堡}。因为在那里，\Person{君士坦丁大帝}的儿子认为，若把他埋葬在渔夫的门廊里，便是给予他极大的尊荣；而看门人对于王宫中的君王是怎样的，君王对于渔夫的坟墓也就是怎样的。这些人确实作为该地的主人占据里面，而其他人仿佛只是寄居者和邻居，乐于分得门廊的门；以此向世上、甚至向不信的人表明，在\OriginalTerm{复活}{Resurrection}中，渔夫将更加是他们上面的。因为若在此世，在（各人的）埋葬中已是如此，那么在复活中更将如此。他们的地位互换；君王取了臣仆的，臣民取了君王的尊荣，甚至更加光明。事实本身就证明这不是奉承；因为由于那些人，这些坟墓变得更加著名。因为人们对这些坟墓的尊敬远超过其他的王陵；因为在那里是深深的孤寂，而在这里是庞大的人群。但若你拿这些坟墓与王宫相比，这里依然是它们得胜。因为在那里有许多人避而远之，但这里有许多人邀请并吸引到他们那里，富有的、贫穷的、男人、女人、奴隶、自由人；在那里有很多恐惧；在这里有说不出的喜乐。\Quote{但是，}有人说，\Quote{看到一位遍体金饰、头戴王冠的君王，以及站在他旁边的将军、统帅、骑兵和步兵指挥官、副将，是一幅悦目的景象。}然而，我们的这一景象却远为宏伟和可畏，以至于相比之下，那景象必须被视为舞台布景和儿戏。因为一旦你跨过门槛，那地方立刻将你的思想提升到天堂，到那至高之王，到天使的军队，到高耸的宝座，到不可接近的荣耀。在此世，祂确实将统治者的权柄——释放一个臣民、捆绑另一个——放在他们手中；但圣人们的骨头并不拥有这样可怜而卑微的权柄，而是远胜于此的权柄。因为他们召唤\OriginalTerm{魔鬼}{Demons}，加以折磨，并从那些最痛苦的束缚中解放被捆绑的人。还有什么比这个法庭更可怕的呢？虽然看不见任何人，虽然没有人堆积魔鬼的两侧，却仍有呼喊、撕扯、鞭打、折磨、灼烧的舌头，因为魔鬼不能忍受那奇妙的德能。那些一度穿着身体的人，胜过无身体的权势；他们的尘土、骨头和骨灰折磨那些不可见的本性。因此，确实没有人会出门远行去看王宫，但许多君王却常常旅行去看这一景象。因为圣人们的\OriginalTerm{殉道者纪念堂}{Martyries}彰显了将来审判的轮廓和预表；在那里魔鬼被鞭打，人受惩罚和释放。你看到圣人们的能力了吗？即使死了。你看到罪人的软弱了吗？即使活着。那么，逃避邪恶吧，好使你能有这样的权能；并尽力追求德行。因为若此世的情况已是如此，试想将来之世将会如何。并且，既然常为这种爱所充满，就抓住永生吧。愿我们都能达到永生，借着我们的主耶稣基督的恩宠和爱人之心，愿光荣、权能、尊崇，归于祂、圣父和圣神，从今世直到永远，永无穷尽。阿们。\par

\SourceRecord{Translated by Talbot W. Chambers.；Nicene and Post-Nicene Fathers, First Series；Vol. 12.；Edited by Philip Schaff.；Buffalo, NY: Christian Literature Publishing Co.,；1889.；<http://www.newadvent.org/fathers/220226.htm>.}

% structure: p0001 source=h1 target=section reason=page-title

\section{格林多后书讲道 第二十七篇}

% structure: p0002 source=h2 target=subsection reason=relative-heading-rank; base=h2

\subsection{格林多后书 12:11}

\BibleQuote{我在夸耀中成了愚妄，是你们逼我的；因为我本该受你们称赞。}\par

在充分完成了关于自己赞美的论述之后，他并未就此止步；反而再次为自己辩护，请求宽恕他所说的话，声称他这样做是出于必要，而非出于自愿。尽管如此，虽然有必要，他还是称自己为\Quote{愚妄}。当他刚开始时，他说：\Quote{请你们当作愚妄人接纳我}，和\Quote{如同在愚妄中}；但现在，省略了\Quote{如同}，他直接称自己为\Quote{愚妄}。因为在他通过所说的话确立了他所希望的观点之后，他随后大胆而无情地抨击了所有这类过失，教导所有人：在没有必要的情况下，绝不应赞美自己，因为即使存在理由，保禄也称自己为\Quote{愚妄}。然后，他将自己这样说话的责备不是转向假宗徒，而是完全转向了门徒。因为他说：\Quote{是你们逼我的}。\Quote{因为如果他们夸耀，但并没有因此误导你们或导致你们毁灭，我就不会被迫陷入这种讨论；但因他们正败坏整个教会，为了你们的益处，我才被迫成为愚妄。}他没有说：\Quote{因为我害怕，如果他们获得你们最高的评价，就会播撒他们的教义}，然而他确实在上文如此表达，当他说：\Quote{我害怕，免得如同蛇诱惑了厄娃一样，你们的心思被败坏。}\BibleRef{格后 11:3}然而在这里他没有那样表达，而是以更命令式的方式和更大的权威，因他所说的话获得了胆量：\Quote{因为我本该受你们称赞。}接着他也给出了理由；他再次提及的不仅是他的启示和奇迹，还有他的试探。\par

\Quote{因为在一切事上，我没有落后于那些最大的宗徒。}请看，他在这里如何又更加权威地说话。事实上，之前他说：\Quote{我自认毫不逊色}，但在这些证明之后，他如今大胆地直陈事实，如我所说，如此绝对。然而即使如此，他也没有偏离中庸之道，也没有脱离其本有特质。因为好像他说了什么伟大而超过自己应得的话，将自己与宗徒等同，他再次谦卑地说话，并加上：\par

% structure: p0006 source=h2 target=subsection reason=relative-heading-rank; base=h2

\subsection{格林多后书 12:12}

\BibleQuote{纵使我算不了什么，宗徒的征兆却在你们中间行了出来。}\par

他说：\Quote{不要看这一点}，\Quote{我是卑微还是微小，而要看你们是否没有享受到那些理应从宗徒那里领受的好处。}然而他没有说\Quote{卑微}，而是说更低的\Quote{算不了什么}。因为伟大而对他人毫无用处，又有什么好处呢？正如一位技艺高超的医生，若他不医治任何病人，则毫无益处。他说：\Quote{所以，不要查究我算不了什么，而要考虑到，在你们应受益的事上，我没有丝毫欠缺，而是证明了我的宗徒职分。那么，本无需我为自己说什么。}他这样说，并非想要受称赞（因为他将天堂本身都视为与他对基督的渴望相比微不足道，又怎么会呢？），而是渴望他们的救恩。然后，免得他们说：\Quote{即使你与那些最大的宗徒毫不逊色，与我们何干？}他因此加上：\par

\Quote{宗徒的征兆在你们中间，以一切忍耐，并借着征兆和奇事，行了出来。}令人惊叹！他在寥寥数语中跨越了何等浩瀚的善工之海！注意他首先提到的是什么：\Quote{忍耐}。因为这就是宗徒的标志：高贵地承担一切。然后他用一个词简短地表达了这一点；但对于奇迹，那些并非凭他自己成就的事，他用了更多言辞。想想看，他在\Quote{忍耐}一词中暗含了多少监牢、多少鞭打、多少危险、多少阴谋、多少试探的冰雹、多少内战、多少对外战争、多少痛苦、多少攻击！而通过\Quote{征兆}，又包含了多少死人复活、多少瞎子看见、多少癞病人洁净、多少魔鬼被驱逐！听到这些，我们当学习：如果我们偶然需要这样的叙述，就要像他那样，简略地提及我们的善行。\par

2. 然后，免得有人说：\Quote{好吧！既然你既伟大，又行了那么多事，但你还没有像宗徒们在其他教会中所行的那样伟大的事，}他加上：\par

% structure: p0011 source=h2 target=subsection reason=relative-heading-rank; base=h2

\subsection{格林多后书 12:13}

\BibleQuote{因为你们在什么事上比其他教会逊色呢？}\par

他说：\Quote{你们领受了与别人一样多的恩宠。}但也许有人会说：\Quote{为什么他把话题转向宗徒，而放弃了与假宗徒的争论？}因为他渴望进一步振奋他们的精神，并表明他不仅超越假宗徒，而且甚至不逊于伟大的宗徒。因此，的确，当他论及假宗徒时，他说：\Quote{我更多；}但当与宗徒们相比时，他认为不\Quote{落后}就是一件大事，尽管他比他们劳苦更多。由此他表明，他们侮辱了宗徒，因为他们将与他相等的人置于这些人之下。\par

\Quote{除非是我本人没有成为你们的负担？}他又一次极其严厉地责备了他们。接下来的话更加令人不悦。\par

\Quote{请原谅我这桩错事。}尽管如此，这种严厉仍包含爱的话语和对他们自身的称赞；也就是说，如果他们视之为一桩错事，即宗徒从他们那里什么都不肯收，也不够信任他们而由他们供养。他说：\Quote{如果你们为此责备我；}他没有说\Quote{你们责备我是错的}，而是极其温和地说：\Quote{我请你们原谅，请宽恕我这个过失。}请注意他的智慧。因为反复提及此事容易给他们带来羞辱，所以他不断缓和语气；例如前面他说：\Quote{基督的真理在我内，这夸耀必不会在我内停止；}\BibleRef{格后 11:10}然后又引用：\Quote{是因我不爱你们吗？天主知道……但为要断绝那些寻找机会之人的机会，使他们所夸耀的显明与我们一样。}\EditorialNote{格后 11:11, 12}在前一封书信中：\Quote{那么我的赏报是什么呢？}诚然，\Quote{就是我在传福音时，能使人无偿地得到福音。}\BibleRef{格前 9:18}而在这里：\Quote{请原谅我这桩错事。}因为他在各处都避免显出是因他们的软弱才不收，以免伤害他们。因此他在这里这样表达：\Quote{如果你们认为这是冒犯，我请求原谅。}他说这些话，既是为了刺痛，也是为了医治。请不要这样说，我恳求你们：\Quote{如果你要刺痛，为何道歉？但如果你道歉，为何刺痛？}因为这是智慧的一部分：既要切开伤口，又要包扎好。随后，为了不让人觉得他——正如他前面也说过——不断反复提这件事是为了从他们那里收受，他在这疑虑上作了补救，甚至在前一封书信中说：\Quote{我写这些话，并不是要人家这样待我；因为我宁愿死，也不愿任何人使我这夸耀落了空。}\BibleRef{格前 9:15}但在这里则更加温良柔和。如何温良柔和呢？以什么方式？\par

% structure: p0016 source=h2 target=subsection reason=relative-heading-rank; base=h2

\subsection{格林多后书 12:14}

\Quote{看，这是我第三次预备到你们那里去，并且我不会成为你们的累赘；因为我寻求的不是你们的财物，而是你们自己；因为儿女不该为父母积蓄，而父母该为儿女积蓄。}\par

他所说的是：\Quote{因为我不是不收你们的，就不来你们这里；不，我已经去过两次；我也预备这第三次前来，\InnerQuote{并且不会成为你们的累赘。}}理由也是高贵的。因为他没有说\Quote{因为你们卑微}、\Quote{因为你们为此受伤}、\Quote{因为你们软弱}；而是什么？\Quote{因为我寻求的不是你们的财物，而是你们自己。}\Quote{我寻求更宝贵的东西：灵魂而非财物；替代金子的是救恩。}然而，此事仍存有某种猜疑，仿佛他对他们不悦；因此他甚至提出一个论据。既然他们很可能说：\Quote{你难道不能既得我们，也得我们的财物吗？}他便极其恩慈地加上这个为他们辩护的理由，说：\Quote{因为儿女不该为父母积蓄，而父母该为儿女积蓄；}以父母和儿女的称呼代替教师和门徒，表明他所做的是出于本分，而非出于义务。因为基督并未这样命令，但他这样说正是为了体恤他们；因此他还进一步补充。他不仅说\Quote{儿女不该为父母积蓄}，而且也说父母该为儿女积蓄。因此，既然给予是合宜的，\par

% structure: p0019 source=h2 target=subsection reason=relative-heading-rank; base=h2

\subsection{格林多后书 12:15}

\Quote{我极愿意为了你们的灵魂而花费自己，并完全耗尽。}\par

\Quote{因为自然律确实命令父母为儿女积蓄；但我不仅这样做，还把自己也一并给出。}他的这种慷慨，不仅不收，反而还给予，不是寻常的方式，而是伴随着极大的大方，并且出于他自身的匮乏；因为\Quote{耗尽}一词正暗示这一点。\Quote{即使需要耗尽我的肉身，为了你们的得救，我也在所不惜。}随后的话语同时包含责备和爱：\Quote{即使我越多地爱你们，却越少被爱。}他说：\Quote{我这样做，是为了那些我所爱的人，但他们却没有同等地爱我。}那么，现在请观察，在这件事上有多少层次。他有权收受，却没有收受；这是第一件善行：而这是在他匮乏的情况下，第二件善行；并且他是向他们宣讲，第三件；他还给予，第四件；不是仅仅给予，而是慷慨地给予，第五件；不仅是财物，还有他自己，第六件；是为那些并不怎么爱他的人，第七件；是为他所深爱的人，第八件。\par

3. 那么，让我们也效法这人罢！因为不爱他人已是严重的罪责；但若被爱而不爱，就更严重了。因为若那爱他之人者爱他，并不比税吏强多少；\BibleRef{玛 5:46} 而连这点都做不到的人，就与野兽为伍；甚至比它们还低劣。人啊，你说什么？你不爱那爱你的人吗？那你活着是为了什么？你将来在何事上有用？在公事上？在私事上？绝不可能；因为没有什么比一个不知道爱的人更无用的。这法律连强盗也常尊重，凶杀者、破门盗窃者亦然；他们仅与人同食盐，便结为朋友，让餐桌改变他们的性情；而你不仅与人同食盐，且言语、行为、出入与共，你却不爱吗？不，那些生活不洁者尚且为娼妇倾尽家产；而你拥有值得的爱，却如此冷漠、软弱、无丈夫气，竟不愿付出分文去爱？\Quote{有人问：\InnerQuote{谁会如此卑劣，如此野蛮，竟转身憎恨那爱他的人？}}你确当不信，因这事悖于常理；但我若指出如此之人甚多，我们又怎能不羞愧？因为当你诽谤你爱的人，当你听别人诽谤他而不辩护，当你嫉妒他得好名声，这是什么样的情感？而且，这不嫉妒、不敌对、不争执，只是支持并抬举那爱你的人，尚不足以证明爱；但当一个人尽一切言行去摧毁邻人，还有什么比这种心态更可怜？昨日及前日还是他的朋友，你与他交谈、共食；然后因忽然看见自己的肢体被看重，便撕下友谊的面具，戴上敌意甚至疯狂的面具。因为为邻人的善行而烦恼，实是十足的疯狂；这是狂犬和疯狗的行为。因为像它们一样，这些人也冲所有人的脸扑去，被嫉妒激怒。与其让嫉妒在我们内爬行，不如让一条蛇缠住内脏。因为蛇常可借药物吐出，或以食物平息；但嫉妒不缠内脏，却寄居在灵魂的胸怀里，是一种难以消除的情欲。若体内真有这样的蛇，只要它食物充足，就不会伤及人身；但嫉妒，即使你为它摆开无尽的宴席，它却吞噬灵魂本身，从四面啃咬、撕裂、拉扯，无法找到任何缓和剂使其放弃疯狂，只有一条：兴盛者的逆境；这样它才得平息——不，即使这样也不能。因为即使这人遭遇逆境，他仍见他人兴盛，便受同样痛苦折磨；处处是伤口，处处是打击。因为活在世上不可能不见到受好评的人。而这病症如此过度，即使将其受害者关在家中，他仍嫉妒那些已死的古人。\par

世俗之人有此感觉，固然可悲，却并非十分可怕；但那些从忙碌生活的纷扰中解脱出来的人，竟被此病所困——这才是最可怕的。我确愿保持沉默；若沉默也能除去这些行为的耻辱，那么缄口不言倒是有益；然而，即使我沉默，这些行为也将比我的舌头喊得更响，那么我的言语不会因其在我们面前炫耀这些恶行而带来损害，反而可能有所增益和益处。因为这疾病已感染了教会，翻转了一切，割裂了身体的联系；我们彼此对立，嫉妒为我们提供武器。因此分裂甚大。因为若我们都建造，门徒尚难站立；当大家都在拆毁，结局将如何？\par

4. 人啊，你在做什么？你以为在拆毁邻人的，但在拆他的之前，你已拆了你的。你没有看见那些园丁、农夫，如何都同心协力？一人挖土，另一人栽种，第三人细心培土，一人浇灌所种的，另一人围篱加固，再一人驱赶牲畜；大家都着眼于一目的，即植物的安全。但这里却不如此：我自己固然栽种，另一人却摇动并扰动那植物。至少让它好好扎根，使其能有力抵抗袭击。你毁的不是我的工作，而是放弃了你自己的。我栽种了，你本应浇灌。你若摇动它，你已连根拔起，便没有地方显示你的浇灌。但你见栽种者受重视。不要怕：我算不得什么，你也不算不得什么。\Quote{因为栽种的和浇灌的，都算不得什么。}\BibleRef{格前 3:7} 一切工作都是天主的。所以你是在与祂争斗、作战，当你拔除所栽种的时。\par

那么让我们终于清醒过来吧，让我们警醒。因为我不太害怕外在的战争，却怕内在的争斗；因为树根若牢牢扎入土中，就不会被风损害；但若它自己摇动，一条虫从内部啃穿它，树就会倒下，即使无人打扰它。我们像虫子一样啃咬教会的根基要到几时呢？因为这类想象也是从土里生出来的，或者更准确地说，不是从土里，而是从粪里，以腐败为母亲；它们且不停止来自女人的可憎谄媚。让我们终于成为慷慨的人，让我们成为哲学的勇士，让我们击退这些邪恶的猛烈进程。因为我看见教会的大众现在俯卧着，如同一具尸体。如同在一个刚死去的身体上，人可以看到眼睛、手、脚、颈、头，却没有一个肢体执行其本分的职务；同样，确实，在这里，所有在场的人都是信友，但他们的信德没有运作；因为我们熄灭了它的温热，使基督的身体成了一具尸体。这话说出来若已可怕，在行为上显现出来时就更加可怕。因为我们确有兄弟之名，却行仇敌之实；人人都被称为肢体，我们却像野兽一样彼此分裂。我说这话并非想炫耀我们的处境，而是为使你们羞愧，并使你们止息。某某人进了一户人家，尊敬给予了他；你本当感谢天主，因为你的肢体得了荣耀，天主受了光荣；但你们却反其道而行：你对那个尊敬他的人说他的坏话，以致你使两人都跌倒，此外，你还使自己蒙羞。为什么呢，悲惨可怜的人？你听见你的兄弟受人称赞，无论是在男人中还是在女人中？你当增添对他的赞美，因为这样你也赞美了你自己。但若你推翻这称赞，首先，你对自己说了坏话，从而获得了坏名声，而且你把他抬得更高了。当你听见一个人受称赞时，你要在所说的话中成为同伴；如果不能在生活和德行上，至少也要因他的卓越而喜乐。某某人称赞了他吗？你也要赞叹：这样他也会称赞你为善良正直的人。不要害怕，好像你因称赞别人而损害了自己的利益：因为这是\Emph{反而}指责他的结果。因为人类有争竞的精神；当它看见你说别人的坏话时，它就会加倍地称赞他，想借此来羞辱你；并且谴责那些控告者，既在心里也对他人。你们看见我们给自己招来的耻辱了吗？我们怎样毁灭和撕裂羊群？让我们终于成为（彼此相联的）肢体，让我们成为一个身体。让那受称赞的人推拒赞美，并把赞词转给兄弟；让那听见别人受称赞的人，自己感到喜乐。若我们这样彼此团聚，我们也会将头吸引到我们这里；但若我们彼此分开生活，我们也会推开从那里来的援助；当那援助被推开时，身体就会受到巨大的伤害，因为没有从上面被结合在一起。为使这事不发生，让我们驱除恶意和嫉妒，轻视众人可能对我们的看法，拥抱爱和和睦。因为这样我们就会获得现世和来世的美物；愿我们众人能达到这一点，藉着我们的主\Person{耶稣基督}对人的恩宠和爱，愿光荣、权能、尊崇与祂（基督）一起归于圣父及圣神，从今直到永远，永无穷尽。阿们。\par

\SourceRecord{Translated by Talbot W. Chambers.；Nicene and Post-Nicene Fathers, First Series；Vol. 12.；Edited by Philip Schaff.；Buffalo, NY: Christian Literature Publishing Co.,；1889.；<http://www.newadvent.org/fathers/220227.htm>.}

% structure: p0001 source=h1 target=section reason=page-title

\section{格林多后书第二十八讲}

% structure: p0002 source=h2 target=subsection reason=relative-heading-rank; base=h2

\subsection{格后 12:16-18}

\BibleRef{格后 12:16}\par

\BibleQuote{然而，就算如此，我自己并没有拖累你们；但我既是诡诈的，就用诡计牢笼了你们。难道我藉着我所差遣到你们那里去的任何人，占了你们的便宜吗？我劝动了弟铎，又差了那位弟兄同去；难道弟铎占了你们的便宜吗？我们不是因同一圣神行事吗？我们不是同一步履吗？}\par

\Person{保禄}这些话说得非常隐晦，但不是没有意义或目的。因为当他谈论钱财并为自己的辩护时，他所说的话包裹在隐晦之中是十分合理的。那么他所说的到底是什么含义？他曾经说：我没有收取，反而甚至愿意加给并花费。在之前的书信和这一封书里，对此主题有许多论述。现在他说了另一件事，以反驳的形式引入主题并预先应对。他的话大意是这样的：我确实没有从你们身上获利；但也许有人会说，我自己没有收取，但我诡诈地安排我所差遣的人替他们自己向你们索取，通过他们我自己收取，同时让自己看起来没有收取，因为是通过他人收取。但没有人能说这一点；你们是见证。因此他继续以问话的方式说：\Quote{我劝动了弟铎，并同他差了那位弟兄。难道弟铎占了你们的便宜吗？} 他不是像我所行的那样行事吗？也就是说，他也没有收取。你们看到怎样一种严格的律己，不仅他自己远离收取，而且如此约束他所差遣的人，以至于一点借口都不给那些想攻击他的人。因为这远远超过圣祖\Person{亚巴郎}所做的。\BibleRef{创 14:24} 因为当他从胜利归来，王要把战利品给他时，他拒绝接受任何东西，只除了那些人已经吃过的；但保禄不仅自己未享用（他们供给的）必要饮食，而且也不允许他的同工分享这些；这样就彻底堵住了厚颜之人的口。因此他不作断言，也不说他们也没有收取；但远比这更甚的是，他引格林多人自己作见证，证明他们什么都没有收取，这样他就不像亲自作证，而是通过他们的判断；我们惯常对完全承认和确信的事情采取这种做法。\Quote{请告诉我，他说，我们差遣的人中，有谁对你们做了\OriginalTerm{不法的得利}{unfair gain}？} 他没有说：有谁从你们那里收取了什么？而是称这些事为\OriginalTerm{不法的得利}{unfair gain}，严厉责备并羞辱他们，表明从不情愿的（给予者）那里收取就是\OriginalTerm{不法的得利}{unfair gain}。他也没有说\Quote{难道弟铎？}，而是说\Quote{有谁？} \Quote{因为你们也不能这么说，}他说，\Quote{某某人确实没有收取，但另一个人收取了。没有一个来的人收取过。} \Quote{我劝动了弟铎。}这话也说得严厉。因为他不是说：我差遣了弟铎，而是我劝动了他；表明即使他收取了，也是正义的；但尽管如此，他仍保持纯洁。因此他再问他们说：\Quote{难道弟铎占了你们的便宜吗？我们不是因同一圣神行事吗？} \Quote{因同一圣神}是什么意思？他将一切归于恩宠，表明这一切赞美不是我们劳苦的结果，而是出于圣神和恩宠的恩赐。因为，虽然他们生活匮乏、饥饿，却不收取任何财物，这是恩宠极大的事例。\Quote{我们不是同一步履行事吗？} 也就是说，他们丝毫未偏离这种严格，而是保持了同一规则的整体。\par

% structure: p0006 source=h2 target=subsection reason=relative-heading-rank; base=h2

\subsection{格后 12:19}

2. \Quote{再者，你们想我们是向你们辩护吗？}\par

你们看到他是如何常常担心，唯恐被人怀疑谄媚？你们看到宗徒的谨慎，他如何经常提起这一点？因为他先前说过：\Quote{我们不是再举荐自己，乃是给你们机会夸耀；}\BibleRef{格后 5:12} 以及在这书信的开头：\Quote{我们岂用得着荐书吗？}\BibleRef{格后 3:1}\par

\Quote{但一切事都是为造就你们。} 他再次抚慰他们。他这里也没有明说：我们因此不收取，是因为你们的软弱；而是说\Quote{为了能造就你们}；话虽比先前说得更清楚，揭示了他心头所重担的，但并无严厉。因为他不是说\Quote{因为你们的软弱}；而是\Quote{为使你们得到造就。}\par

% structure: p0010 source=h2 target=subsection reason=relative-heading-rank; base=h2

\subsection{格后 12:20}

\Quote{因为我怕，怕我来到的时候，见你们不合我所想的，你们见我也不合你们所想的。}\par

他要说一些伟大而令人反感的话。因此，他也插入这一辩解，既说\BibleQuote{一切都是为了建树你们}，又加上\BibleQuote{我害怕}，以缓和将要说出的话的严厉性。因为这不是出于傲慢，也不是出于教师的权威，而是出于父亲般的慈爱关怀，他比罪人自己更恐惧战栗地面对能改造他们的事。即便如此，他也没有斥责他们或作出绝对的断言，而是怀疑地说：\BibleQuote{恐怕我来的时候，找不到你们如同我所愿意的。}他没有说\InnerQuote{不道德}，而是说\Quote{不如我所愿意的}，处处使用亲情的用语。而\BibleQuote{我找到}一词，表达的是出乎自然期望的事，同样还有\BibleQuote{我被你们找到}。因为这不是出于自主选择，而是出于你们方面的必要性。因此他说：\BibleQuote{我恐怕被你们找到，是你们所不愿意的。}他没有在这里说\BibleQuote{是我不愿意的}，而是更严厉地说\BibleQuote{是你们所不愿意的}。因为那样就会成为他自己的意愿，虽然不是他首先愿意的，但仍然是他的意愿。因为他本可以再次说\BibleQuote{是我不愿意的}，以此显示他的爱；但他不愿松懈听众的心。不，那样他的话甚至会更加严厉；但现在他既给了他们更重的打击，又显得更加温柔。这正是他智慧的特点：更深入地切中要害，却更温柔地击打。然后，因为他说的隐晦，他揭示了他的意思，说：\par

\BibleQuote{免得有纷争、嫉妒、恼怒、毁谤、谗言、自高自大。}\par

而他本可先说的，却放在最后：因为他们对他非常骄傲。因此，为不显得主要是在追求自己的利益，他先提到公共的事。因为这一切都是从嫉妒产生的：他们的毁谤、指控、分裂。正如某种恶根，嫉妒产生了愤怒、指控、骄傲和所有其他恶事，并且又被它们进一步滋长。\par

% structure: p0015 source=h2 target=subsection reason=relative-heading-rank; base=h2

\subsection{\Emph{格后12:21}}

和\BibleQuote{免得我再来时，我的天主使我在你们中间谦卑。}\par

而\BibleQuote{再来}一词也是打击他们。因为他的意思是：\InnerQuote{先前的事已经够了；}正如他在书信开头也说过：\BibleQuote{为了宽容你们，我至今没有再到格林多去。}\EditorialNote{格后1:18, 23}你看他如何同时显示愤怒和慈爱？但\InnerQuote{使我谦卑}是什么意思？其实这是光荣的事：控告、报复、追究、坐在审判官的位置上；但他却称之为谦卑。他如此不以那谦卑的原因羞愧，因为\BibleQuote{他的肉身在场是软弱的，言语是微不足道的}，以至于他甚至愿意永远处于那种情况，并反对相反的情况。他继续更清楚地说明了这一点；他认为特别谦卑的事，就是陷入当前这种必须惩罚和报复的处境。为什么他不说\Quote{恐怕我来时，我会谦卑}，而说\BibleQuote{恐怕我来时，我的天主会谦卑我}？\InnerQuote{若不是为了祂，我本不会在意，也不会担忧。因为我要求满足，不是由于拥有权威或为了自己的喜悦，而是因为祂的命令。}上面他是这样表达的：\BibleQuote{我被找到}；在这里他却放松了，采用了更温和、更温柔的语言，说：\par

\BibleQuote{我要为许多犯罪的人哀哭。}不是单纯地说\InnerQuote{犯罪的人}，而是\par

\BibleQuote{他们没有悔改。}他没有说\InnerQuote{所有人}，而是说\InnerQuote{许多人}；也没有明确是谁，从而使他们悔改的归途变得容易；并且为了表明悔改能够纠正过犯，他哀哭那些不悔改的人，那些不治之症的人，那些继续留在创伤中的人。请观察宗徒的德行：他自己虽无愧于心，却为别人的恶事哀哭，并为别人的过犯而谦卑。因为这是教师特有的标志：如此同情门徒的灾难，并为属下人的创伤而哀伤。然后他也提到了具体的罪。\par

\BibleQuote{关于他们曾行的邪淫和不洁。}在这里他确实暗指淫乱；但如果仔细考察，每一种罪都可以这样称呼。因为虽然犯邪淫者和奸夫特别被称为不洁，但其他罪也在灵魂中产生不洁。因此，基督也称犹太人为不洁，不仅指控他们犯奸淫，也指控其他恶行。因此祂也说他们洗净外面，并且\BibleQuote{不是入口的污秽人，而是出口的污秽人。}\BibleRef{玛 15:11} 又在另一处说：\BibleQuote{凡是心中傲慢的人，在上主面前都是不洁的。}\BibleRef{箴 16:5}\par

3. 因为没有任何事物比德行更纯洁，也没有任何事物比恶行更污秽；前者比太阳更明亮，后者比污泥更臭恶。那些在污泥中打滚、在黑暗中生活的人，自身将为此作证；只要稍微使他们清醒一点，他们就能看清。因为他们独自一人时，被情欲醉倒，如同活在黑暗中一般，以不当的姿势躺卧，带来许多耻辱，甚至那时他们也意识到自己的所在，虽然不完全；但一旦他们看到任何生活在德行中的人谴责他们，甚至只是显现自身，那么他们就更清楚地理解自己的可怜；如同阳光照射在他们身上，他们遮掩自己的不当，在知道他们行为的人面前脸红，即使一个是奴隶，另一个是自由人，即使一个是国王，另一个是臣民。因此，当阿哈布看见厄里亚时，他感到羞愧，尽管厄里亚还没有说什么；仅仅看见他就被定罪；当控告者沉默时，他宣布了对自己的定罪判决；说出被捉住者的话，说：\Quote{我的仇敌，你找到了我！}\BibleRef{列上 21:20} 因此厄里亚本人那时怀着极大的勇气与那暴君交谈。同样，黑落德无法忍受那些责备的羞耻，（那先知舌头的响声以强大而透明的清晰更明显地揭露了这羞耻）将若翰投入监狱：如同一个裸体的人试图扑灭光，好能再次处于黑暗；或者不如说他本人不敢扑灭它，而是把它放在斗底下；而那邪恶可悲的女人强逼了这事。但即便如此，他们也无法掩盖责备，反而使它更加明亮。因为那些问\Quote{若翰为何住在监狱里？}的人知道了原因，并且所有此后生活在陆地或海上的人，无论是那时活着，还是现在活着，以及将来的人，都清楚知道并将会清楚知道这些邪恶的悲剧，包括他们淫乱和流血的悲剧，没有任何时间能抹去对它们的记忆。\par

德行是如此伟大：它的记忆是如此不朽，甚至仅凭言语就能完全击倒它的对手。因为他为何将他投入监狱？为何不轻视他？难道他要把他拖到审判台前吗？难道他要求因他的奸淫而报复他吗？他当时所说的岂不只是一句责备？那他为何恐惧战栗？岂不只有言语和谈话吗？但这话比行动更刺痛他。他没有把他带到任何审判台前，而是把他拖到良心的另一个法庭；他让所有在心中自由宣判的人作他的法官。因此暴君战栗，无法忍受德行的光辉。你看到哲学是多么伟大吗？它使一个囚犯比国王更光彩，而后者在他面前害怕战栗。他确实只把他捆绑了；但那个污秽的女人也冲去杀害他，虽然责备主要是针对他（而不是她）。因为他当时没有遇见\Quote{她}并对她说：\Quote{你为何与国王同居？}不是说她没有罪（她怎能无罪呢？），而是他希望通过另一种方法纠正一切。因此他责备国王，但对他并不粗暴。因为他没有说：\Quote{啊，污秽的、全污秽的、无法无天、亵渎的人，你践踏了天主的法律，你藐视了诫命，你使你的强权成为法律。}没有说这些；但即使在责备中，这人的温和也很大，他的温良很大。因为他说：\Quote{你不可占有你兄弟斐理伯的妻子。}这些话更像是一位教导者而非责备者，一位训诫者而非惩罚者，一位整理秩序者而非揭露者，一位修正者而非践踏者。但是，正如我所说，光对贼是可恨的，看见义人就令罪人厌恶；\Quote{他使我们感到烦恼，连看也不愿看他。}\BibleRef{智 2:15} 因为他们无法忍受他的光辉，如同病眼无法忍受太阳的光辉。但对许多恶人来说，他不仅看见是烦恼，甚至听说也是烦恼。因此那污秽的、全污秽的女人，她女儿的教唆者，甚至可以说是她女儿的谋杀者，尽管从未见过他或听过他的声音，也冲去杀害他；并准备她在那淫荡中养大的女儿也去行凶，她如此过分地惧怕他。\par

4. 她又说了什么？\Quote{请把\Person{洗者若翰}的头放在盘子里给我。}\BibleRef{玛 14:8} 可怜不幸的人啊，你为什么要冲向悬崖？什么？控告者在你面前吗？他在眼前困扰你吗？另一些人曾说：\Quote{他甚至连看都令我们痛苦；}但对她来说，正如我所说，听他说话都是痛苦的。所以她要求说：\Quote{请把\Person{洗者若翰}的头放在盘子里给我。}然而，因为你，他身陷囹圄，戴着重锁，而你却可以放纵你的情欲，甚至说：\Quote{我完全制伏了王，以至于他公开受责骂却不屈服，也不中止与我的奸情，反而将责骂他的人投入了监牢。}你为什么疯狂、狂暴，即使经过对他罪恶的责备，你仍然保留你的情夫？你为什么寻求复仇女神的餐桌，准备报复魔鬼的宴席？你难道看不出恶行是多么无价值、多么懦弱、多么不男子气吗？当它最得逞时，它就变得更无力了吗？因为这个女人在把\Person{洗者若翰}投入监狱之前并没有如此不安，而在他被捆绑之后，她却更为焦虑，逼迫着说：\Quote{请把\Person{洗者若翰}的头放在盘子里给我。}这是为什么？\Quote{我怕，}她说，\Quote{怕他的谋杀被掩盖，怕有人把他从危险中救出。}为什么你不要求整具尸体，而只要头？\Quote{那使我痛苦的舌头，}她说，\Quote{我希望看到它永远沉寂。}但相反的事情会发生，正如它也确实发生了，可怜不幸的人！它将更响亮地呼喊，当它被割掉之后。因为那时它只在\Place{犹太}呼喊，但现在它要传到地极；无论你进入哪一座教堂，无论是在\Person{摩尔人}中，在\Person{波斯人}中，甚至到\Place{不列颠群岛}本身，你都会听到\Person{洗者若翰}呼喊：\Quote{你占有你兄弟\Person{斐理伯}的妻子是不合法的。}但那个女人，不知道理智思考，催促、逼迫着那个无知的暴君去杀人，怕他改变主意。但由此你也再次学到美德的力量。即使当他被关押、捆绑、沉默时，她仍然忍受不了那个义人。你看出恶行是何等软弱吗？何等污秽吗？因为它在肉食的位置上带来了一个人头放在盘子里。有什么比那个少女更污秽、更可咒、更无耻的呢？她在\Person{魔鬼}的剧场里、在\Person{恶魔}的宴席上发出了怎样的声音啊！你看到这个舌头和那个舌头吗？一个带来健康的药物，另一个带着毒药，成了\Person{魔鬼}宴席的供应者。但她为什么不命令他在那里、在宴席上被谋杀，那时她的快乐会更大？她怕如果他到了那里并被看见，他的面容和勇气会改变他们所有人。因此她一定要求他的头，想为淫乱立一个辉煌的凯旋，并把它交给她的母亲。你看到跳舞的代价了吗？你看到那个\Person{魔鬼}阴谋的战利品了吗？我说的不是\Person{洗者若翰}的头，而是她的情夫本人。因为如果有人仔细考察，这个战利品是针对王的，而胜利者被击败了，被斩首者得了冠冕，并被宣告胜利，甚至死后更剧烈地震撼了犯法者的心。我所说的不是夸张，问\Person{黑落德}本人吧；他听见\Person{基督}的神迹时说：\Quote{这是\Person{洗者若翰}，他从死人中复活了，所以这些异能在他身上运行。}\BibleRef{玛 14:2} 恐惧是如此鲜明，痛苦如此持久地留在他身上；没有人能推翻他良心的恐惧，那个不朽的审判者继续掐住他的喉咙，日复一日向他追讨杀人的赔偿。那么，知道这些事，我们不要害怕受苦，而要害怕作恶；因为前者确实是胜利，后者是失败。\par

\Person{保禄}也说：\Quote{为什么不宁愿受欺？为什么不宁愿被亏负？然而你们反倒欺压、亏负，而且正是亏负你们的弟兄。}因为通过受苦而获得那些冠冕、那些奖赏、那[胜利的]宣告。这在所有圣人身上都可以看见。既然他们都已经如此被加冕、被宣告，让我们也走这条路，并确实祈求我们不陷入诱惑；但如果诱惑来临，让我们以极大的男子气概站稳，并表现出应有的心灵准备，这样我们就能获得将来的美好事物，透过我们的主\Person{耶稣基督}的恩宠和慈爱，愿光荣、权能、尊崇归于祂、及\Person{圣父}、偕同\Person{圣神}，自今至永远，及无穷世。亚孟。\par

\SourceRecord{Translated by Talbot W. Chambers.；Nicene and Post-Nicene Fathers, First Series；Vol. 12.；Edited by Philip Schaff.；Buffalo, NY: Christian Literature Publishing Co.,；1889.；<http://www.newadvent.org/fathers/220228.htm>.}

% structure: p0001 source=h1 target=section reason=page-title

\section{《格林多后书》讲道集第二十九篇}

% structure: p0002 source=h2 target=subsection reason=relative-heading-rank; base=h2

\subsection{格林多后书 13:1}

\BibleQuote{这是我第三次要到你们那里去。凭两三个人的口，句句都可定准。}\par

\Person{保禄}的智慧和他极为温柔的深情，在许多其他情况下都可以观察到，但尤其在此，他的劝诫如此丰富而强烈，但惩罚却如此迟缓而拖延。因为他没有在他们犯罪后立即惩罚他们，而是一再警告；即使如此，在他们不加理会时，他也不执行惩罚，而是再次警告，说：\BibleQuote{这是我第三次要到你们那里去；} 并说：\Quote{在我来之前，我再写信。} 然后，为了使他的拖延不会导致漠不关心，请看他也如何纠正这一结果：通过不断威胁，将打击悬在他们头上，说：\BibleQuote{若我来时，必不宽容；} 以及 \Quote{恐怕我来时，要哀悼许多人。} 那么，他这样做和这样说，在这事上也效法万有的主：因为天主确实也不断威胁，常常警告，但不常责罚和惩罚。同样，\Person{保禄}也确实如此，因此他先前也说：\Quote{我之所以还未到\Place{格林多}去，是为了宽免你们。} 什么是\Quote{宽免你们}？恐怕我发现你们犯了罪且仍未悔改，就会用责罚和惩罚来对待你们。而在此，\BibleQuote{这是我第三次要到你们那里去。凭两三个人的口，句句都可定准。} 他将未写的话与已写的话结合起来，如同他在另一处所做的一样，说：\BibleQuote{与娼妓结合的，便是成为一体；因为祂说：二人}，祂说，\BibleQuote{要成为一体。} \BibleRef{格前 6:16} 然而，这话是关于合法婚姻的；但他却恰当地将这话应用到这件事上，以便更震慑他们。他在这里也是如此，以他的到来和警告为证人。他所说的是：\Quote{我曾在你们那里说过一次两次；现在也以书信说话。如果你们听从我，我所愿的就成就了；但如果你们不加理会，从今以后就必须停止说话，而要施加惩罚。} 因此他说，\par

% structure: p0005 source=h2 target=subsection reason=relative-heading-rank; base=h2

\subsection{格林多后书 13:2}

\BibleQuote{我曾预先说过，当第二次在你们那里时也预先说过，现在不在你们那里，就写信给那些以前犯了罪的和其余所有的人：若我来时，必不宽容。}\par

\Quote{若凭两三个人的口，句句都可定准，而我已经来过两次并说过话，现在又通过这封信说话；那么，我此后必须遵守诺言。因为请你们不要以为我的写信比不上我的亲临；因为就如我在时说话，现在我不在时也写信。} 你看见他的兄弟般的关怀吗？你看见相称于师傅的先见之明吗？他既不保持沉默也不惩罚，而是常常预告，不断威胁，延迟惩罚，如果他们继续不改悔，他便威胁要付诸行动。\Quote{但你先前在的时候对他们说了什么？不在的时候又写了什么？} \BibleQuote{若我来时，必不宽容。} 他先前已表明，除非被迫，他不愿这样做，并把这事称为哀痛和自卑；（因为他曾说过：恐怕我的天主在你们面前降卑我，使我为那些从前犯了罪而仍未悔改的人哀痛。\BibleRef{格后 12:21}）他在他们面前辩解说，他曾预先告诉他们一次、两次、三次，并且他尽一切努力抑制惩罚，通过言语的恐惧使他们变好，然后才用了这句令人不悦和恐惧的话：\BibleQuote{若我来时，必不宽容。} 他没有说：\Quote{我要报复、惩罚并索取补偿}；而是再次用慈父般的言语表达惩罚本身；显示他的柔情，并与他们一同悲伤；因为他总是为了\Quote{宽免}他们而拖延。为了使他们不再认为还会有拖延，只是言语上的威胁，因此他先前说：\BibleQuote{凭两三个人的口，句句都可定准；} 现在[他说]：\BibleQuote{若我来时，必不宽容。} 他的意思是：\Quote{如果（但愿不会）我发现你们仍未改悔，我将不再拖延；而必定要追究，实现我所说的话。}\par

2. 然后，带着极大的愤怒和强烈的义愤，针对那些嘲笑他软弱、讥讽他的亲临、并说：\BibleQuote{他亲临是软弱，言语是可鄙的；} \BibleRef{格后 10:10} 的人，他针对这些人努力，说道：\par

% structure: p0009 source=h2 target=subsection reason=relative-heading-rank; base=h2

\subsection{格林多后书 13:3-4}

\BibleQuote{既然你们寻求基督在我内说话的凭据。}\par

他说这话，既打击了这些人，同时也鞭打了那些人。他的意思是：\Quote{既然你们想要证明基督是否住在我内，并向我追究，并因此嘲笑我卑贱可鄙，仿佛我缺乏那大能；你们要知道，我们并不缺乏，如果你们给我们机会（但愿不会）。} 那么怎样？告诉我。难道你们因此就惩罚，因为他们寻求凭据？\Quote{不，}他说；因为如果他寻求这个，他早就他们在犯罪时惩罚他们，不会拖延了。但他在继续时更清楚地表明他不寻求这个，说：\BibleQuote{现在我祈祷你们不作任何恶事，不是要我们显得合格，而是要你们合格，即使我们像是被淘汰的。} \BibleRef{格后 13:7}\par

他不用那些话来作为理由，而是出于愤慨，更是在攻击那些轻视他的人。他说：\Quote{我实在不愿意给你们这样的证据，但如果你们自己引起原因并选择挑战我，你们将藉着事实知道。}又注意他所说的话有多严厉。因为他没有说\Quote{你们既然寻求我的证据}，却说\Quote{基督在我内说话的证据，表明他们得罪的是祂}。他不仅说\Quote{住在我内}，而是\Quote{在我内说话}，表明他的话是属神的。但如果他不显示祂的权能也不惩罚，（因为从此宗徒把话从自己转到基督，使他的威胁更可怕），不是由于软弱，因为祂能这样做，而是由于含忍。所以不要让任何人以为祂的宽容是软弱。为什么你奇怪祂现在不处罚罪人，不在祂的宽容和含忍中要求赔偿，而祂甚至忍受了被钉十字架，且受了这样的苦却未惩罚？因此他也补充了：\par

\BibleQuote{祂对你们不是软弱的，而是在你们中间大有能力的。因为祂虽因软弱被钉十字架，却因天主的德能活着。}\par

这些话很有隐晦性，使较弱的人感到困扰。因此有必要更清楚地阐明它们，并解释存在隐晦的那个表达的含义，以免即使较单纯的人也不被冒犯。那么，这里所说的究竟是什么意思，\Quote{软弱}一词指的是什么，它被用在什么含义上，这是必须学习的。因为这个词确实是一个，但有很多意义。因为身体上的疾病被称为\Quote{软弱}：由此甚至在福音中说到：\BibleQuote{看，祢所爱的人病了。}\BibleRef{若 11:3}，关于拉匝禄；祂自己说：\BibleQuote{这病不至于死。}；保禄论及厄帕夫辣说：\BibleQuote{他确实因软弱几乎死了，但天主可怜了他。}\BibleRef{斐 2:57}；论及弟茂德：\BibleQuote{为了你的胃和你的屡次软弱，用一点酒。}\BibleRef{弟前 5:23}。所有这些都表示身体的疾病。再者，在信德中没有坚定站稳被称为\Quote{软弱}，即不完全、不圆满。为此保禄说：\BibleQuote{在信德上软弱的，你们要接纳，但不要辩驳疑惑的事。}\BibleRef{罗 14:1}；又说：\BibleQuote{一个人信自己可以吃一切，而那在信德上软弱的却只吃蔬菜。}，指在信德上软弱的人。这里\Quote{软弱}一词有两个含义；还有第三种东西叫做\Quote{软弱}。那么这是什么？迫害、阴谋、侮辱、试探、攻击。为此保禄说：\BibleQuote{为这事，我三次求主，祂对我说：我的恩宠够你用，因为我的德能在软弱中才完全。}\EditorialNote{格后 12:8,9}什么是\Quote{在软弱中？}在迫害中、危险中、试探中、阴谋中、死亡中。为此他说：\Quote{因此，我喜好软弱。}然后显示他指的是哪种软弱，他没有说发烧，也没有说对信德的怀疑；而是什么？\BibleQuote{在凌辱中，在急需中，在磨难中，在鞭打中，在监禁中，为的是基督的德能常在我身上。因为我几时软弱，正是我有能力的时候。}\BibleRef{格后 12:10}这就是说\Quote{当我受迫害时，当我被驱赶时，当我被谋害时，我便刚强，此时我更胜过并胜过那些谋害我的人，因为恩宠更多地停留在我身上。}正是在这第三种意义上，保禄使用了\Quote{软弱}；这就是他所指的；他再次，如我先前所说，瞄准那一点，即他在他们眼中显得卑贱和可轻视。因为他确实不想夸耀，也不想显露出他真正的样子，更不想显示他惩罚和报复的能力；因此他也被看作卑贱。当他们如此看待他，继续在极大的冷漠和麻木中前行，并不悔改自己的罪时，他抓住一个好时机，在这些问题上也非常有力地讲论，并表明他不做什么并非由于软弱，而是出于含忍。\par

3. 然后，正如我所说，通过把论证从自己转到基督，他增强了他们的恐惧，增加了他的威胁。他说的是：\Quote{即使我应当做某事并惩罚和报复有罪的人，难道是我惩罚和报复吗？是住在我内的祂，基督自己。但如果你们不信这点，而渴望通过行为接受住在我内的祂的证明，你们立即就会知道；\BibleQuote{祂对你们不是软弱的，反而是有能力的。}为什么加上\Quote{对你们}，既然祂处处是大能的？因为如果祂有意惩罚不信者，祂能做到；或者魔鬼，或任何事物。那么加上的含义是什么？这句话要么是通过回忆他们已经接受的证明极大地羞辱他们；要么是宣告这一点：同时祂在你们这些应该被纠正的人中间显示祂的能力。正如他在另一处所说的：\BibleQuote{我何必审判那些外人呢？}\BibleRef{格前 5:12}\Quote{因为那些外人，}他说，\Quote{祂将在审判之日追究，但你们甚至现在，以便救你们脱离那惩罚。}但尽管如此，他的这份关怀虽然出于温柔的爱，注意他如何将其与恐惧和极大的愤怒结合，说：\BibleQuote{祂对你们不是软弱的，而是在你们中间大有能力的。}}\par

% structure: p0016 source=h3 target=subsubsection reason=relative-heading-rank; base=h2

\subsubsection{格后 13:4}

\BibleQuote{因为祂虽因软弱被钉十字架，却因天主的德能活着。}\par

\BibleQuote{祂虽因软弱而被钉死}是什么意思？\InnerQuote{因为祂选择，}他说，\InnerQuote{忍受一种似乎带有软弱概念的事，但这丝毫未侵入祂的德能。那仍是无敌的，而那似乎属于软弱的，丝毫未损害它，反而正是这事本身最显示了祂的德能，因为祂甚至忍受了这样的事，而祂的德能并未被损坏。}那么，不要让\Quote{软弱}这个说法困扰你；因为别处他也说：\BibleQuote{天主的愚妄比人更聪明，天主的软弱比人更强壮；}\BibleRef{格前 1:25} 尽管在天主内既无愚妄也无软弱；但他如此称呼十字架，是为了表达不信者对它的看法。至少听他亲自解释：\BibleQuote{原来十字架的道理，为那些丧亡的人是愚妄，但为我们得救的人，却是天主的德能。}\BibleRef{格前 1:18} 又说：\BibleQuote{我们却宣讲被钉的基督：这为犹太人固然是绊脚石，为外邦人是愚妄；但为那些蒙召的，不拘犹太人，希腊人，基督却是天主的德能和天主的智慧。}（同上，23-24节）又说：\BibleQuote{然而属血气的人不接受天主圣神的事，因为为他是愚妄。}\BibleRef{格前 2:14} 注意，他处处表达了不信者的看法，他们视十字架为愚妄和软弱。所以，实际上，这里他也不是指真正意义上的\Quote{软弱}，而是指不信者所怀疑的那种。因此他不是说，因为祂是软弱的而被钉死。断乎不可！因为祂本有权不让自己被钉死，祂自始至终显示了这一点：祂时而使人扑倒在地，时而使太阳暗晦，使无花果树枯萎，又使前来攻击祂的人眼睛失明，并行了无数其他奇事。那么他说的\Quote{藉着软弱}是什么意思呢？就是即使祂在经历了危险和出卖后被钉死（因为我们已说明危险和出卖被称为软弱），仍然丝毫未受损害。他说这话是为了将榜样引到自己的境况。因为既然格林多人看见他们被逼迫、被驱赶、被轻视，并不报复也不追究，为了教导他们，他们如此受苦并非因能力不足，也不是不能追究，他便将论证引向主，因为\BibleQuote{祂，}他说，\BibleQuote{也曾被钉死、被捆绑、遭受了无数事情，并不追究，反而继续忍受那些看似表明软弱的事，并以此显示出祂的德能，因为尽管祂不惩罚也不报复，却丝毫无损。例如，十字架没有切断祂的生命，也没有阻碍祂的复活，反而祂复活了并且活着。}当你听到十字架和生命时，要期待找到关于降生成人的教义，因为这里所说的一切都与此有关。如果他说\Quote{虽然天主的德能}，并非祂自己无力使自己的肉身复活；而是他提到父或子都无所谓。因为当他说\Quote{天主的德能}时，他是说藉着祂自己的德能。因为祂亲自复活并扶持它，请听祂说：\BibleQuote{你们拆毁这座圣殿，三天之内我要把它重建起来。}\BibleRef{若 2:19} 但如果祂所说的属于父的，也不要困扰；因为祂说：\BibleQuote{凡父所有的一切，都是我的。}\BibleRef{若 16:15} 又说：\BibleQuote{凡我所有的，都是祢的；凡祢所有的，都是我的。}\BibleRef{若 17:10} \BibleQuote{正如那被钉者毫无受损，}他说，\BibleQuote{同样，我们被逼迫、被攻击时也毫无受损；}因此他又加上：\par

\BibleQuote{因为我们虽然在祂内软弱，但我们将因天主的德能同祂一起生活。}\par

\Quote{我们在祂内软弱}是什么意思？\Quote{因为为了宣讲，}他说，\Quote{以及我们对祂的信德。但如果为了祂，我们承受了悲伤和不愉快的事，显然我们也必会承受愉快的事：}所以他又加上：\BibleQuote{但我们藉着天主的德能同祂一起得救。}\par

% structure: p0021 source=h2 target=subsection reason=relative-heading-rank; base=h2

\subsection{\BibleRef{格后 13:5-6}}

4. \BibleQuote{你们要试验自己，是否在信德中；要考验自己。难道你们自己不知道，除非你们是经不起考验的，基督耶稣就在你们内吗？但我希望你们会知道，我们并不是经不起考验的。}\par

因为藉着刚才所说的话，他已经表明：即使他不惩罚，也不是因为他内没有基督，而是因为他暗示了当初被钉十字架、不为自己报仇的那一位的含忍。接着，他又以另一种方式产生同样的效果，而且更加无可辩驳，用门徒来巩固他的论证。\Quote{因为我为什么要提到自己呢？} 他说，\Quote{我身为老师，肩负如此重担，受托管理整个世界，行了如此伟大的奇迹。因为如果你们愿意省察自己，你们身为门徒，就会看见基督也在你们内。但如果祂在你们内，那么更是在你们的老师内。因为如果你们有信德，基督也在你们内。} 因为那时相信的人都行了奇迹。因此他又加上：\BibleQuote{你们要试验自己，考验自己，看看你们是否在信德上。难道你们自己不知道，除非基督在你们内，你们就是不合格的吗？} \Quote{但如果祂在你们内，那不是更在你们的老师内吗？} 在我看来，他这里指的是与奇迹有关的\Quote{信德}。\Quote{因为如果你们有信德，} 他说，\Quote{基督就在你们内，除非你们成了不合格的人。} 你看，他如何再次使他们害怕，甚至过分地表明基督与他同在。因为在我看来，他这里是在暗示他们，甚至涉及他们的生活。既然单凭信德不足以引下圣神的德能，而他又曾说过\Quote{\InnerQuote{如果你们在信德中}你们就有基督在你们内}，并且发生许多有信德的人却缺乏那种德能的事；为了解决这个困难，他说：\BibleQuote{除非你们是不合格的，} 即除非你们在生活上是败坏的。\BibleQuote{但我希望你们知道，我们不是不合格的。} 按自然顺序，他应该说：\Quote{但如果你们成了不合格的，我们却不是。} 但他没有这样说，因为怕伤害他们；而是以模糊的方式暗示，既不说\Quote{你们是不合格的}，也不以问话的方式说\Quote{但你们是不合格的}，甚至连这种问话的方式也省略了，只是模糊地加了一句：\BibleQuote{但我希望你们知道，我们不是不合格的。} 这里又是极大的威胁，极大的恐慌。\Quote{因为既然你们愿意，} 他说，\Quote{藉此方式，通过自己的惩罚来接受证明，我们不难给你们这样的证明。} 但他确实没有这样表达，而是更有分量、更具威胁地说：\BibleQuote{但我希望你们知道，我们不是不合格的。} \Quote{因为你们本来应该，} 他说，\Quote{即使没有这个也知道我们是什么人，知道基督在我们内说话并工作；但既然你们愿意也通过行动来接受证明，你们就会知道我们不是不合格的。} 这样，当他将那威胁悬在他们头上，把惩罚带到他们的门前，使他们战栗，期待报复时，请看他又如何甜言蜜语，安抚他们的恐惧，表现出他不争强好胜的性情、对门徒的温柔关怀、原则坚定、高尚和没有虚荣。因为他在接下来的话中显示了所有这些美德，说道：\par

% structure: p0024 source=h2 target=subsection reason=relative-heading-rank; base=h2

\subsection{格后 13:7-9}

\BibleQuote{现在我祈求天主，使你们不作恶，不是要我们显得合格，而是要你们做那光荣的事，即使我们被视为不合格。因为我们不能反对真理，只能拥护真理。当我们软弱而你们坚强时，我们就喜乐。我们也为此祈求，使你们得以成全。}\par

5. 什么能与这种灵魂相比？他被轻视，被唾弃，被讥讽，被嘲笑，被视为卑贱、可鄙、吹牛、口出狂言而不见行动；虽然看到如此需要显示自己的权能，他不仅推迟，不仅退缩，甚至祈求自己不要陷入这种境地。因为他所说：\BibleQuote{现在我祈求天主，使你们不作恶，不是要我们显得合格，而是要你们做那光荣的事，即使我们被视为不合格。} 他是说什么呢？\Quote{我恳求天主。我祈求祂，} 他说，\Quote{使我找不到一个未改过的人，找不到一个未悔改的人？是的，不仅此，而且使没有人犯过罪。因为，} 他说，\Quote{你们没有作恶，但若偶然犯了罪，那么你们可能已改变行为，先我而改过，阻止一切愤怒。因为我所热切追求的不是要我们以这种方式显得合格，恰恰相反，是要我们不显得合格。因为如果你们继续，} 他说，\Quote{犯罪而不知悔改，我们就必须责打、惩罚、伤残你们的身体；（正如发生在撒斐辣和术士西满身上的事；）并且我们已经证明了我们的权能。但我们不是为此祈祷，而是相反，即不要我们以这种方式显得合格，不要我们以此来显示我们内在的权能——即通过责打和惩罚你们，视你们为犯罪和不治之症；而是什么呢？\BibleQuote{你们要做那光荣的事，} 我们为此祈祷，希望你们永远生活在德行中，不断改进；\BibleQuote{而我们被视为不合格，} 不显示我们惩罚的权能。} 而且他没有说\Quote{不合格的}，因为他不会\Quote{成为}不合格的，即使他不惩罚，正是为此他才是\Quote{合格的}；\Quote{但即使有人因我们不显示权能而怀疑我们，认为我们可鄙、被弃，我们也不在乎。我们宁可由那些人这样看待，也不愿用天主赐予我们的那些鞭笞和那种不悔改的心来显示权能。}\par

\BibleQuote{因为我们不能反对真理，只能拥护真理。}为了让他看起来不仅仅是在讨好他们（因为一个没有虚荣心的人可能会这样做），而是做了事情本性所要求的事，他又加上了这句：\BibleQuote{因为我们不能反对真理。}\Quote{因为如果我们发现你们在良好的声誉中，}他说，\Quote{借着悔改驱散了你们的罪过，并在天主面前拥有胆量；那么此后，即使我们非常愿意，也无法惩罚你们；即使我们尝试，天主也不会与我们合作。因为祂赐给我们权柄，正是为了让我们所作的审断是真实而正义的，不违反真理。}你们看到了吗？他如何尽一切可能使自己的话不带冒犯，并缓和了威胁的严厉？因为正如他热切地这样做了，他也渴望表明他的心是与他们紧密相连的；因此他又加上了这句：\BibleQuote{因为我们软弱而你们强壮时，我们就欢喜，我们也为此祈祷，使你们得以成全。}\Quote{因为最确实地，}他说，\Quote{我们不能作任何反对真理的事，也就是说，如果你们中悦天主，我们就不能惩罚你们；此外，因为我们不能，所以我们也不希望这样，甚至渴望相反。不仅如此，我们特别高兴的正是这件事，当我们发现你们没有给我们机会展示我们惩罚的权柄。因为即使做这样的事使人显得光荣、嘉许而强壮，我们仍然渴望相反，希望你们被嘉许、无可指摘，而我们永远不从中收取荣耀。}因此他说：\BibleQuote{因为我们软弱时，我们就欢喜。}什么是：\Quote{软弱？}\Quote{当我们可能被认为软弱时。}不是当我们软弱，而是当我们被认为软弱；因为他们在敌人眼中被这样认为，因为他们没有展示自己惩罚的权柄。\Quote{但是我们仍然欢喜，当你们的品行如此，以至于不给我们任何惩罚你们的借口。我们以这种方式被认为软弱是一种快乐，只要你们无可指摘；}因此他又加上了：\BibleQuote{而你们是强壮的，}也就是说，\Quote{是被嘉许的，是有德行的。我们不仅希望这一点，而且为此祈祷，希望你们无可指摘、成全，不给我们任何把柄。}\par

6. 这是父亲的慈爱：将门徒的救恩置于自己的好名声之上。这是一种摆脱虚荣的灵魂的表现；这最能使人从身体的束缚中释放出来，使人从地上升到天上，那便是纯洁无染于虚荣；因此，正如相反的情形会引人进入许多罪恶。因为一个没有摆脱虚荣的人，是不可能崇高、伟大和高贵的；他必然匍匐在地，造成许多损害，同时成为一个比任何野蛮人更残忍的污秽主母的奴隶。还有什么比她更凶猛的呢？当她最受奉承时，她最为野蛮。即使是野兽也不是这样，它们会因多加关注而被驯服。但虚荣却恰恰相反：被藐视时，她变得温顺；被尊敬时，她变得野蛮，并武装起来对抗尊敬她的人。犹太人尊敬了她，却受到了极其严厉的惩罚；门徒们藐视了她，却得到了冠冕。我何必提惩罚和冠冕呢？因为就被人看为荣耀这一点而言，最有效的莫过于唾弃虚荣。你甚至在这世上也会看到，那些尊敬虚荣的人受到了损害，而那些藐视它的人却获益了。因为那些藐视虚荣的门徒（我们可以再次用同样的例子，没有阻碍），他们选择了天主的事，比太阳还耀眼，即使在死后也为自己赢得了不朽的记忆；而那些向虚荣屈膝的犹太人却变得无城无家、无心无肝、堕落、流亡、卑贱、可鄙。因此，你若渴望得荣耀，就要拒绝荣耀；你若追求荣耀，就会失去荣耀。如果你愿意，让我们也在世俗事务上验证这一道理。在我们的玩笑中，我们嘲笑的是谁？不就是那些执著于虚荣的人吗？当然，这些人最缺乏它，因为他们有无数控告者，并被所有人藐视。我们钦佩的是谁呢？请告诉我，难道不是那些藐视它的人吗？当然，这些人正是被荣耀的人。因为正如富有的人不是需要很多东西的人，而是什么都不缺的人；同样，荣耀的人不是爱荣耀的人，而是藐视荣耀的人；因为这荣耀不过是荣耀的影子。没有人看到一幅画的饼，即使饿得极其厉害，也会去攻击那画。因此，你也不要追求这些影子，因为这是荣耀的影子，不是荣耀。为了让你知道它的本质是影子，请思考这一点：当这事物在人间有坏名声时，当所有人都认为它该被避免，甚至连那些渴望它的人也这样认为；当拥有它和渴望它的人都羞于以此称呼自己时。有人说：\InnerQuote{那么这欲望从何而来？这情欲是如何产生的？} 是由于灵魂的渺小（因为不仅应该指责它，还应该纠正它），是由于不完善的理智，是由于幼稚的判断。让我们不再做孩子，而成为成人；让我们在财富、享乐、奢侈、荣耀和权力等一切事物中追求现实，而不是影子；这样，这疾病就会止息，许多其他的也会止息。因为追求影子是疯子的行为。因此保罗也说：\BibleQuote{你们要按公道醒寤，不要犯罪。} \BibleRef{格前 15:34} 因为还有一种疯狂，比魔鬼所致的疯狂、比狂热所致的疯狂更严重。因为那种疯狂尚可原谅，但这种疯狂却无可推诿，因为灵魂本身被败坏，其正确判断丢失了；狂热确实是一种身体的疾病，但这种疯狂却存在于思想的工匠心中。因此，正如发烧中那些更加严重、甚至无法医治的，是那些侵袭坚固的身体、潜伏在神经的深处、隐藏在血管中的发烧；同样，这种疯狂也确实如此，因为它潜伏在思想本身的深处，败坏和摧毁它。因为轻视那些永存的事物，却极其热切地依附于那些消亡的事物，这难道不是明显而显著的疯狂，一种比任何疯狂更严重的疾病吗？因为，请告诉我，如果一个人去追逐风，或试图抓住风，我们难道不会说他疯了吗？又比如，如果一个人抓住影子而忽略现实，如果一个人厌恶自己的妻子却拥抱她的影子，或者厌恶自己的儿子却又爱他的影子，你还需要寻求其他更清晰的疯狂证据吗？那些贪婪地追逐现世事物的人也是如此。因为它们都是影子，无论你提到荣耀、权力、好名声、财富、奢侈，还是今生任何其他事物。因此，先知确实说过：\BibleQuote{世人只会虚幻经营，他们徒然焦虑。} \BibleRef{咏 38:6} 又说：\BibleQuote{我的岁月有如斜阳。} \BibleRef{咏 101:11} 在另一处，他称人间事物为烟和草上的花。但不仅是他的好事是影子，他的坏事也是影子，无论你提到死亡、贫穷、疾病，还是任何其他事。那么，哪些是永存的事物，无论是好的还是坏的？是永恒的王国和永远的地狱。因为\BibleQuote{虫子不死，火也不灭。} \BibleRef{谷 9:44} 并且\BibleQuote{这些人要进入永生，那些人要进入永罚。} \BibleRef{谷 25:46} 因此，为了逃避后者并享有前者，让我们放下影子，以全部的热忱抓住真实的事物，因为这样我们才能获得天主的国。愿我们全体都能因我们的主耶稣基督的恩宠和慈爱而获得它，愿光荣和权能归于他，至于无穷之世。阿们。\par

\SourceRecord{Translated by Talbot W. Chambers.；Nicene and Post-Nicene Fathers, First Series；Vol. 12.；Edited by Philip Schaff.；Buffalo, NY: Christian Literature Publishing Co.,；1889.；<http://www.newadvent.org/fathers/220229.htm>.}

% structure: p0001 source=h1 target=section reason=page-title

\section{\WorkTitle{格林多后书讲道集第三十篇}}

% structure: p0002 source=h2 target=subsection reason=relative-heading-rank; base=h2

\subsection{格林多后书 13:10}

\BibleQuote{为此，我身在异地而写这些事，为的是当我来到时，不必严厉对待，因为主赐给我的权柄是为了建立，而不是为了拆毁。}\par

他意识到自己说话比往常更加激烈，尤其是在书信的结尾。因为他先前说过，\Quote{如今我保禄亲自以基督的温和与良善劝勉你们；我当着你们的面是谦卑的，不在你们当中时却对你们勇敢。是的，我恳求你们，当我来到时，不要让我以那自信而勇敢地对待那些认为我们按血气行事的人。}\EditorialNote{参见格后10:1-2}，并且，\Quote{随时准备在你们的顺服得以成全时，惩罚一切不顺服。}\BibleRef{格后 10:6}，并且，\Quote{我担心当我来到时，会发现你们不是我愿意的那样，而你们也发现我不是你们愿意的那样。}\BibleRef{格后 12:20}，并且，\Quote{免得当我来到时，我的天主使我在你们面前谦卑，并且我为许多从前犯罪而尚未悔改他们曾犯的放荡和污秽的人哀哭。}\BibleRef{格后 12:21}，随后又说，\Quote{我从前说过，如今虽不在你们那里，却如同第二次来临一样预先告诉你们：如果我再来，我必不宽容；因为你们寻求基督在我内说话的证据。}\EditorialNote{参见格后13:2-3}。既然他说了这些以及更多恐吓、羞辱、责备、鞭策他们的话，他就为这一切开脱说，\Quote{为此，我身在异地而写这些事，为的是当我来到时，不必严厉对待。}因为我希望严厉停留在我的书信中，而非我的行为中。我希望我的威胁是强烈的，好让它们始终只是威胁，永远不会付诸行动。然而，即使在他这番道歉中，他也使自己的话更加可怕，表明惩罚他的不是他自己，而是天主；因为他补充说，\Quote{依据主赐给我的权柄。}并且，为了表明他不想使用自己的权柄来惩罚他们，他又补充说，\Quote{不是为了拆毁，而是为了建立。}正如我所说，他现在暗示了这一点，但他让他们自己得出结论：如果他们继续不思悔改，那么惩罚那些有这种心态的人，又是一种建立。事实的确如此，他知道这一点，并通过他的行为显示了这一点。\par

% structure: p0005 source=h2 target=subsection reason=relative-heading-rank; base=h2

\subsection{格林多后书 13:11}

\BibleQuote{此外，弟兄们，要喜乐，要成全，要受安慰，要同心，要和睦，如此，仁爱与平安的天主必与你们同在。}\par

\Quote{此外，弟兄们，要喜乐}是什么意思？你们使他们痛苦、恐惧、陷入焦虑，使他们战栗害怕，你怎么又命令他们喜乐？\Quote{正是为了这个原因，我命令他们喜乐。因为，他说，如果你们也尽到了你们的本分，就没有什么能阻止那喜乐。因为我的部分已经全部完成：我长久忍耐，我延缓，我忍住不剪除，我恳求，我劝告，我警告，我威胁，以便用各种方法将你们聚集到悔改的果实中。现在你们也应当尽到你们的本分，这样你们的喜乐就不会衰退。}\par

\Quote{你们要成全。}什么是\Quote{你们要成全}？\Quote{要完全，补足所欠缺的。}\par

\Quote{你们要受安慰。}因为，既然他们的考验众多，危险巨大，他说，\InnerQuote{你们要受安慰}，既彼此安慰，也由我们安慰，更因你们的改变而安慰。因为如果你们有了良心的喜乐并成为完全的，就没有什么缺乏你们的喜乐和安慰。因为没有什么能像纯洁的良心那样产生安慰，即使有无数的试炼围绕。\par

\Quote{你们要同心，要和睦。}这是他在前一封书信的开头也提出的请求。因为有可能同心，却并不和睦地生活，例如，当人们在教义上一致，但彼此交往时却不和。但\Person{保禄}要求两者兼备。\par

\Quote{如此，仁爱与平安的天主必与你们同在。}因为他不仅推荐和劝告，而且也祈祷。因为要么他为此祈祷，要么预言将要发生的事；或者更像是两者。\Quote{因为如果你们做这些事，他说，例如，如果你们\InnerQuote{同心}并\InnerQuote{和睦}，天主也必与你们同在，因为祂是\InnerQuote{仁爱与平安的天主}，祂喜爱这些事，祂为此喜乐。因此，从祂的爱中，你们也要有平安；因此，一切邪恶将被除去。这拯救了世界，这结束了长期的战争，这融合了天地，这使人们成为天使。那么让我们也效法这一点，因为爱是无数好事的母亲。藉此我们得救，藉此所有那些不可言说的好事都来到我们当中。}\par

2. 然后，为了引导他们进入这爱，他说：\par

% structure: p0013 source=h2 target=subsection reason=relative-heading-rank; base=h2

\subsection{格林多后书 13:12}

\Quote{你们要以圣洁的亲嘴彼此问安。}\par

\Quote{圣的？}不是虚伪的，不是奸诈的，不像犹达斯给基督的亲吻。因为亲吻因此被给予，好使它成为爱的燃料，点燃这种情操，使我们如此彼此相爱，如同兄弟爱兄弟，如同子女爱父母，如同父母爱子女；是的，甚至远超过这些。因为那些是由本性植入的情操，而这些是由属神的恩宠植入的。这样我们的灵魂彼此结合。因此当我们离别后归来时，我们彼此亲吻，我们的灵魂渴望彼此的交往。因为这个肢体最向我们显明灵魂的运作。但关于这圣吻，或许还可以说些别的。什么效果？我们是基督的圣殿；那么我们彼此亲吻时，我们亲吻了圣殿的门廊和入口。你们没看到有多少人甚至亲吻这圣殿的门廊吗？有些人弯腰，有些人用手抓住它，然后将手放到嘴上。通过这些门和通道，基督曾进入我们内，并且当我们领受时也进入我们内。你们这些参与奥秘的人，明白我所说的。因为我们的嘴唇受到荣耀，绝非寻常，当它们领受主的身体时。正是为此我们在这里亲吻。让那些说污秽话、出言辱骂的人听吧，让他们战栗，想到他们羞辱的是什么嘴；让那些以不洁方式亲吻的人听吧。听天主通过你的口宣告了什么，并保持它不受玷污。祂谈论了来世、复活、不朽，死亡不是死亡，以及其他无数奥秘。因为那将要受教的人来到司铎的口前，如同来到神谕前，聆听充满敬畏的事物。因为他从先祖那里就失去了生命，前来寻求重获它，并询问如何可能找到并收回它。然后天主向他宣告如何找到它，那张嘴变得比施恩座本身更可畏。因为施恩座从未发出这样的声音，而是谈论更小的事情，战争和此世的和平：但这张嘴谈论的全是关于天堂和来世，以及新而超越理解的事物。说完后，\par

% structure: p0016 source=h2 target=subsection reason=relative-heading-rank; base=h2

\subsection{\BibleRef{格后 13:13}}

\BibleQuote{你们要以圣吻彼此问候。}他接着说：\BibleQuote{诸位圣人问候你们。}\par

借此也给他们良好的希望。他以此代替亲吻加上了这个，通过问候使他们联结，因为话语也出自发出亲吻的同一张嘴。你看到他是如何将他们全部聚集在一起的吗？那些在身体上相距遥远的人，以及那些邻近的人；这些人通过亲吻，那些人通过书信。\par

% structure: p0019 source=h2 target=subsection reason=relative-heading-rank; base=h2

\subsection{\BibleRef{格后 13:14}}

3. \BibleQuote{我们主耶稣基督的恩宠，和天主的爱，}和圣父，\BibleQuote{和圣神的相通，与你们众人同在。}在用问候和亲吻使他们彼此联合之后，他又以祈祷结束他的讲论，非常谨慎地也将他们与天主联合。如今那些说，因为圣神在书信的开端未被提及，所以祂不是同一本体的人在哪里？因为，看，他现在将祂与圣父圣子并列。此外，可以注意到，当他写信给哥罗森人时说\BibleQuote{愿恩宠与平安由天主我们的父赐与你们，}他沉默不提圣子，并且没有像在他所有书信中那样加上\Quote{并由主耶稣基督}。那么圣子也不是同一本体吗，因为如此？不，这些推理是极其愚蠢的。因为正是这件事特别表明祂是同一本体，即保禄随意地使用（或不使用）这个表达。并且这里所说的并非臆测，听他是如何提到圣子和圣神，而对圣父闭口不言的。因为写信给格林多人时，他说：\BibleQuote{但你们已被洗净，但你们已被圣化，但你们已因我们主耶稣基督的名，并因我们天主的圣神而成义。}\BibleRef{格前 6:11}那么，告诉我，这些人难道不是因圣父之名受洗的吗？那么他们肯定既未被洗净也未被圣化。但他们难道没有给他们施洗吗？无疑正如他们确实施洗了。那么他为何不说\Quote{你们因父之名被洗净}？因为在他眼中，时而提这个人，时而提那个人，是无所谓的；你们可以在书信的许多地方观察到这个习惯。因为写信给罗马人时，他说：\BibleQuote{所以我以天主的仁慈恳求你们，}\BibleRef{罗 12:1}尽管那些仁慈是属于圣子的；并且：\BibleQuote{我以圣神的爱恳求你们，}\BibleRef{罗 15:30}尽管爱是属于圣父的。那么，他为何在\BibleQuote{仁慈}中没有提到圣子，在\BibleQuote{爱}中没有提到圣父？因为作为明显和公认的事情，他对此沉默。此外，他还会发现，也将恩赐本身颠倒放置。因为在这里说了\BibleQuote{基督的恩宠，和天主圣父的爱，以及圣神的相通；}他在另一处说\BibleQuote{圣子的相通，}和\BibleQuote{圣神的爱。}因为他\BibleQuote{我恳求你们，}他说，\BibleQuote{因圣神的爱。}\BibleRef{罗 15:30}并且在他致格林多人的书信中：\BibleQuote{天主是信实的，你们曾被祂召唤进入祂圣子的相通。}\BibleRef{格前 1:9}这样，天主圣三的事物是不可分的：既然相通属于圣神，它也被发现属于圣子；既然恩宠属于圣子，它也属于圣父和圣神；因为[我们读到]：\BibleQuote{愿恩宠由天主父赐与你们。}并且在另一处，列举了许多形式后，他加上：\BibleQuote{但这一切都是同一且相同的圣神所运作，随祂心意个别分给各人。}\BibleRef{格前 12:11}我说这些事，并非混淆位格，（切勿有这种想法！）而是知道祂们的个体性和区别，以及本体的合一。\par

4. 那么，让我们继续严格持守这些教义，并将天主的爱吸引到我们身上。因为此前祂爱我们时，我们还在恨祂，并和好那些曾是祂仇敌的我们；但从此以后，祂愿意爱我们如同爱祂。那么，让我们继续爱祂，好使我们也被祂所爱。因为如果当我们被有权势的人所爱时，我们就令所有人畏惧，何况是被天主所爱呢？若为这爱需要付出财富、身体，甚至生命本身，我们也不该吝惜。因为只在口头上说爱是不够的，我们也应以行动来证明；因为祂也不是仅以言语，也以行动显示爱。那么，你也要以你的行动显示这一点，做那些中悦祂的事，因为这样做你也会获得益处。因为祂不需要我们所能给予的任何东西，这也是真诚之爱的一个特别证明：当一个无所需求、毫无需要的人，为了被我们爱而做一切事。因此梅瑟也说：\BibleQuote{上主你的天主要求你什么？无非是爱祂，并乐意跟随祂？} \BibleRef{申 10:12} 所以当祂命令你爱祂时，祂恰恰最显示祂爱你。因为没有什么比爱祂更能保证我们的救恩。你看，祂的所有诫命都共同导向我们的安息、救恩和好名声。因为当祂说：\BibleQuote{怜悯人的人是有福的，心里洁净的人是有福的，温良的人是有福的，神贫的人是有福的，缔造和平的人是有福的；} \BibleRef{玛 5:3} 祂自己确实不从中得到任何好处，而是为了我们的装饰和调谐而命令这些。当祂说：\Quote{我饿了，} 并不是需要我们的服侍，而是激励你们要有同情心。因为即使没有你，祂也能养活穷人；但祂将这些命令加给你，是为了赐给你一个极大的宝藏。因为如果太阳——一个受造物——并不需要我们的眼睛；它自身永存于自己的光明中，即使无人仰望它，而当我们享受它的光芒时，我们才是受益者；何况是天主呢？但为让你以另一种方式明白这一点：你认为天主与我们之间的距离有多大？如同蚊蚋与我们之间的距离，还是更大？显然要大得多，甚至无限。那么，如果我们这些虚荣的受造物不需要蚊蚋的服务和尊敬，更何况神性本性（不需要我们从它得到什么），因为它是无情感的、一无所缺的。祂因我们而享受的限度，就是我们得到的益处之大，以及祂因我们的得救而喜悦。为此，祂也常常放弃自己的权利，而寻求你的利益。他说：\BibleQuote{因为若有人，} 有不信的妻子，且她乐意与他同住，就不要离弃她；\BibleRef{格前 7:12} 并说：\Quote{凡离弃妻子，除非为奸淫的缘故，就是使她犯奸淫。} 你看见何等不可言说的良善吗？他说：\Quote{如果妻子是个娼妓，我不强迫丈夫与她同住；如果她是个不信者，我也不禁止他。} 又说：\Quote{如果你对某人怀怨，我命令那得罪你的人放下我的祭品，来与你和好。} 因为祂说：\BibleQuote{你若在献祭时，想起你的兄弟有什么对不住你的事，就把你的祭品留在祭台前，先去与你的兄弟和好，然后再来献祭。} \BibleRef{玛 5:23} 那个耗尽一切的比喻说了什么？\BibleRef{玛 18:24} 它不正是显示这个吗？因为当他吞没了那一万塔冷通时，主人怜悯了他，放他走了；但当他向同伴索要一百银钱时，主人称他为恶人，并把他交给惩罚。祂如此看重你的安逸。那外邦人正要侵犯义人的妻子，祂说：\BibleQuote{我阻止了你，免得你得罪我。} \BibleRef{创 20:6} 保禄迫害宗徒，祂对他说：\Quote{你为什么迫害我？} 他人饥饿，祂自己说祂饿了，赤身露体，漂流在外，想要羞辱你，从而迫使你走上施舍的道路。\par

那么，既然反思了这爱——祂在万事中所显示并仍在显示的爱是何等伟大，祂既屈尊使自己为我们所认识（这是万善中最大的冠冕，理智的光照和德行的教导），又为最美好的生活方式制定了法律，且为我们做了万事，赐下了祂的圣子，应许了天国，邀请我们分享那些不可言说的美好，并为我们预备了最蒙福的生命——让我们言行一致，好显得配得上祂的爱，并获得将来的美好；愿我们众人都能到达那里，藉着我们的主耶稣基督对世人的恩宠和慈爱；愿光荣归于祂，与圣父和圣神，从现在到永远，世世无穷。阿们。\par

\SourceRecord{Translated by Talbot W. Chambers.；Nicene and Post-Nicene Fathers, First Series；Vol. 12.；Edited by Philip Schaff.；Buffalo, NY: Christian Literature Publishing Co.,；1889.；<http://www.newadvent.org/fathers/220230.htm>.}
